\chapt{किष्किन्धाकाण्डः}

\sect{प्रथमः सर्गः}

\twolineshloka
{स तां पुष्करिणीं गत्वा पद्मोत्पलझषाकुलाम्}
{रामः सौमित्रिसहितो विललापाकुलेन्द्रियः} %||4-1-1||

\twolineshloka
{तस्य दृष्ट्वैव तां हर्षादिन्द्रियाणि चकम्पिरे}
{स कामवशमापन्नः सौमित्रिमिदमब्रवीत्} %||4-1-2||

\twolineshloka
{सौमित्रे पश्य पम्पायाः काननं शुभदर्शनम्}
{यत्र राजन्ति शैलाभा द्रुमाः सशिखरा इव} %||4-1-3||

\twolineshloka
{मां तु शोकाभिसन्तप्तमाधयः पीडयन्ति वै}
{भरतस्य च दुःखेन वैदेह्या हरणेन च} %||4-1-4||

\twolineshloka
{अधिकं प्रविभात्येतन्नीलपीतं तु शाद्वलम्}
{द्रुमाणां विविधैः पुष्पैः परिस्तोमैरिवार्पितम्} %||4-1-5||

\twolineshloka
{सुखानिलोऽयं सौमित्रे कालः प्रचुरमन्मथः}
{गन्धवान्सुरभिर्मासो जातपुष्पफलद्रुमः} %||4-1-6||

\twolineshloka
{पश्य रूपाणि सौमित्रे वनानां पुष्पशालिनाम्}
{सृजतां पुष्पवर्षाणि वर्षं तोयमुचामिव} %||4-1-7||

\twolineshloka
{प्रस्तरेषु च रम्येषु विविधाः काननद्रुमाः}
{वायुवेगप्रचलिताः पुष्पैरवकिरन्ति गाम्} %||4-1-8||

\twolineshloka
{मारुतः सुखं संस्पर्शे वाति चन्दनशीतलः}
{षट्पदैरनुकूजद्भिर्वनेषु मधुगन्धिषु} %||4-1-9||

\twolineshloka
{गिरिप्रस्थेषु रम्येषु पुष्पवद्भिर्मनोरमैः}
{संसक्तशिखरा शैला विराजन्ति महाद्रुमैः} %||4-1-10||

\twolineshloka
{पुष्पिताग्रांश्च पश्येमान्कर्णिकारान्समन्ततः}
{हाटकप्रतिसञ्छन्नान्नरान्पीताम्बरानिव} %||4-1-11||

\twolineshloka
{अयं वसन्तः सौमित्रे नानाविहगनादितः}
{सीतया विप्रहीणस्य शोकसन्दीपनो मम} %||4-1-12||

\twolineshloka
{मां हि शोकसमाक्रान्तं सन्तापयति मन्मथः}
{हृष्टः प्रवदमानश्च समाह्वयति कोकिलः} %||4-1-13||

\twolineshloka
{एष दात्यूहको हृष्टो रम्ये मां वननिर्झरे}
{प्रणदन्मन्मथाविष्टं शोचयिष्यति लक्ष्मण} %||4-1-14||

\twolineshloka
{विमिश्रा विहगाः पुम्भिरात्मव्यूहाभिनन्दिताः}
{भृङ्गराजप्रमुदिताः सौमित्रे मधुरस्वराः} %||4-1-15||

\twolineshloka
{मां हि सा मृगशावाक्षी चिन्ताशोकबलात्कृतम्}
{सन्तापयति सौमित्रे क्रूरश्चैत्रवनानिलः} %||4-1-16||

\twolineshloka
{शिखिनीभिः परिवृता मयूरा गिरिसानुषु}
{मन्मथाभिपरीतस्य मम मन्मथवर्धनाः} %||4-1-17||

\twolineshloka
{पश्य लक्ष्णम नृत्यन्तं मयूरमुपनृत्यति}
{शिखिनी मन्मथार्तैषा भर्तारं गिरिसानुषु} %||4-1-18||

\twolineshloka
{मयूरस्य वने नूनं रक्षसा न हृता प्रिया}
{मम त्वयं विना वासः पुष्पमासे सुदुःसहः} %||4-1-19||

\twolineshloka
{पश्य लक्ष्मण पुष्पाणि निष्फलानि भवन्ति मे}
{पुष्पभारसमृद्धानां वनानां शिशिरात्यये} %||4-1-20||

\twolineshloka
{वदन्ति रावं मुदिताः शकुनाः सङ्घशः कलम्}
{आह्वयन्त इवान्योन्यं कामोन्मादकरा मम} %||4-1-21||

\twolineshloka
{नूनं परवशा सीता सापि शोचत्यहं यथा}
{श्यामा पद्मपलाशाक्षी मृदुभाषा च मे प्रिया} %||4-1-22||

\twolineshloka
{एष पुष्पवहो वायुः सुखस्पर्शो हिमावहः}
{तां विचिन्तयतः कान्तां पावकप्रतिमो मम} %||4-1-23||

\twolineshloka
{तां विनाथ विहङ्गोऽसौ पक्षी प्रणदितस्तदा}
{वायसः पादपगतः प्रहृष्टमभिनर्दति} %||4-1-24||

\twolineshloka
{एष वै तत्र वैदेह्या विहगः प्रतिहारकः}
{पक्षी मां तु विशालाक्ष्याः समीपमुपनेष्यति} %||4-1-25||

\twolineshloka
{पश्य लक्ष्मण संनादं वने मदविवर्धनम्}
{पुष्पिताग्रेषु वृक्षेषु द्विजानामुपकूजताम्} %||4-1-26||

\twolineshloka
{सौमित्रे पश्य पम्पायाश्चित्रासु वनराजिषु}
{नलिनानि प्रकाशन्ते जले तरुणसूर्यवत्} %||4-1-27||

\twolineshloka
{एषा प्रसन्नसलिला पद्मनीलोत्पलायता}
{हंसकारण्डवाकीर्णा पम्पा सौगन्धिकायुता} %||4-1-28||

\twolineshloka
{चक्रवाकयुता नित्यं चित्रप्रस्थवनान्तरा}
{मातङ्गमृगयूथैश्च शोभते सलिलार्थिभिः} %||4-1-29||

\twolineshloka
{पद्मकोशपलाशानि द्रष्टुं दृष्टिर्हि मन्यते}
{सीताया नेत्रकोशाभ्यां सदृशानीति लक्ष्मण} %||4-1-30||

\twolineshloka
{पद्मकेसरसंसृष्टो वृक्षान्तरविनिःसृतः}
{निःश्वास इव सीताया वाति वायुर्मनोहरः} %||4-1-31||

\twolineshloka
{सौमित्रे पश्य पम्पाया दक्षिणे गिरिसानुनि}
{पुष्पितां कर्णिकारस्य यष्टिं परमशोभनाम्} %||4-1-32||

\twolineshloka
{अधिकं शैलराजोऽयं धातुभिस्तु विभूषितः}
{विचित्रं सृजते रेणुं वायुवेगविघट्टितम्} %||4-1-33||

\twolineshloka
{गिरिप्रस्थास्तु सौमित्रे सर्वतः सम्प्रपुष्पितैः}
{निष्पत्रैः सर्वतो रम्यैः प्रदीपा इव कुंशुकैः} %||4-1-34||

\twolineshloka
{पम्पातीररुहाश्चेमे संसक्ता मधुगन्धिनः}
{मालतीमल्लिकाषण्डाः करवीराश्च पुष्पिताः} %||4-1-35||

\twolineshloka
{केतक्यः सिन्दुवाराश्च वासन्त्यश्च सुपुष्पिताः}
{माधव्यो गन्धपूर्णाश्च कुन्दगुल्माश्च सर्वशः} %||4-1-36||

\twolineshloka
{चिरिबिल्वा मधूकाश्च वञ्जुला बकुलास्तथा}
{चम्पकास्तिलकाश्चैव नागवृक्षाश्च पुष्पिताः} %||4-1-37||

\twolineshloka
{नीपाश्च वरणाश्चैव खर्जूराश्च सुपुष्पिताः}
{अङ्कोलाश्च कुरण्टाश्च चूर्णकाः पारिभद्रकाः} %||4-1-38||

\twolineshloka
{चूताः पाटलयश्चैव कोविदाराश्च पुष्पिताः}
{मुचुकुन्दार्जुनाश्चैव दृश्यन्ते गिरिसानुषु} %||4-1-39||

\threelineshloka
{केतकोद्दालकाश्चैव शिरीषाः शिंशपा धवाः}
{शाल्मल्यः किंशुकाश्चैव रक्ताः कुरबकास्तथा}
{तिनिशा नक्त मालाश्च चन्दनाः स्यन्दनास्तथा} %||4-1-40||

\twolineshloka
{विविधा विविधैः पुष्पैस्तैरेव नगसानुषु}
{विकीर्णैः पीतरक्ताभाः सौमित्रे प्रस्तराः कृताः} %||4-1-41||

\twolineshloka
{हिमान्ते पश्य सौमित्रे वृक्षाणां पुष्पसम्भवम्}
{पुष्पमासे हि तरवः सङ्घर्षादिव पुष्पिताः} %||4-1-42||

\threelineshloka
{पश्य शीतजलां चेमां सौमित्रे पुष्करायुताम्}
{चक्रवाकानुचरितां कारण्डवनिषेविताम्}
{प्लवैः क्रौञ्चैश्च सम्पूर्णां वराहमृगसेविताम्} %||4-1-43||

{अधिकं शोभते पम्पाविकूजद्भिर्विहङ्गमैः॥\devanumber{44}॥} %||4-1-44|| (Check)

\addtocounter{shlokacount}{1}

\twolineshloka
{दीपयन्तीव मे कामं विविधा मुदिता द्विजाः}
{श्यामां चन्द्रमुखीं स्मृत्वा प्रियां पद्मनिभेक्षणाम्} %||4-1-45||

\twolineshloka
{पय सानुषु चित्रेषु मृगीभिः सहितान्मृगान्}
{मां पुनर्मृगशावाक्ष्या वैदेह्या विरहीकृतम्} %||4-1-46||

\twolineshloka
{एवं स विलपंस्तत्र शोकोपहतचेतनः}
{अवेक्षत शिवां पम्पां रम्यवारिवहां शुभाम्} %||4-1-47||

\fourlineindentedshloka
{निरीक्षमाणः सहसा महात्मा}
{ सर्वं वनं निर्झरकन्दरं च}
{उद्विग्नचेताः सह लक्ष्मणेन}
{ विचार्य दुःखोपहतः प्रतस्थे} %||4-1-48||

\fourlineindentedshloka
{तावृष्यमूकं सहितौ प्रयातौ}
{ सुग्रीवशाखामृगसेवितं तम्}
{त्रस्तास्तु दृष्ट्वा हरयो बभूवुर्-}
{ महौजसौ राघवलक्ष्मणौ तौ} %||4-2-49||


{॥इत्यार्षे श्रीमद्रामायणे वाल्मीकीये आदिकाव्ये किष्किन्धाकाण्डे प्रथमः सर्गः॥}

\sect{द्वितीयः सर्गः}

\twolineshloka
{तौ तु दृष्ट्वा महात्मानौ भ्रातरौ रामलक्ष्मणौ}
{वरायुधधरौ वीरौ सुग्रीवः शङ्कितोऽभवत्} %||4-2-1||

\twolineshloka
{उद्विग्नहृदयः सर्वा दिशः समवलोकयन्}
{न व्यतिष्ठत कस्मिंश्चिद्देशे वानरपुङ्गवः} %||4-2-2||

\twolineshloka
{नैव चक्रे मनः स्थाने वीक्षमाणो महाबलौ}
{कपेः परमभीतस्य चित्तं व्यवससाद ह} %||4-2-3||

\twolineshloka
{चिन्तयित्वा स धर्मात्मा विमृश्य गुरुलाघवम्}
{सुग्रीवः परमोद्विग्नः सर्वैरनुचरैः सह} %||4-2-4||

\twolineshloka
{ततः स सचिवेभ्यस्तु सुग्रीवः प्लवगाधिपः}
{शशंस परमोद्विग्नः पश्यंस्तौ रामलक्ष्मणौ} %||4-2-5||

\twolineshloka
{एतौ वनमिदं दुर्गं वालिप्रणिहितौ ध्रुवम्}
{छद्मना चीरवसनौ प्रचरन्ताविहागतौ} %||4-2-6||

\twolineshloka
{ततः सुग्रीवसचिवा दृष्ट्वा परमधन्विनौ}
{जग्मुर्गिरितटात्तस्मादन्यच्छिखरमुत्तमम्} %||4-2-7||

\twolineshloka
{ते क्षिप्रमभिगम्याथ यूथपा यूथपर्षभम्}
{हरयो वानरश्रेष्ठं परिवार्योपतस्थिरे} %||4-2-8||

\twolineshloka
{एकमेकायनगताः प्लवमाना गिरेर्गिरिम्}
{प्रकम्पयन्तो वेगेन गिरीणां शिखराणि च} %||4-2-9||

\twolineshloka
{ततः शाखामृगाः सर्वे प्लवमाना महाबलाः}
{बभञ्जुश्च नगांस्तत्र पुष्पितान्दुर्गसंश्रितान्} %||4-2-10||

\twolineshloka
{आप्लवन्तो हरिवराः सर्वतस्तं महागिरिम्}
{मृगमार्जारशार्दूलांस्त्रासयन्तो ययुस्तदा} %||4-2-11||

\twolineshloka
{ततः सुग्रीवसचिवाः पर्वतेन्द्रं समाश्रिताः}
{सङ्गम्य कपिमुख्येन सर्वे प्राञ्जलयः स्थिताः} %||4-2-12||

\twolineshloka
{ततस्तं भयसन्त्रस्तं वालिकिल्बिषशङ्कितम्}
{उवाच हनुमान्वाक्यं सुग्रीवं वाक्यकोविदः} %||4-2-13||

\twolineshloka
{यस्मादुद्विग्नचेतास्त्वं प्रद्रुतो हरिपुङ्गव}
{तं क्रूरदर्शनं क्रूरं नेह पश्यामि वालिनम्} %||4-2-14||

\twolineshloka
{यस्मात्तव भयं सौम्य पूर्वजात्पापकर्मणः}
{स नेह वाली दुष्टात्मा न ते पश्याम्यहं भयम्} %||4-2-15||

\twolineshloka
{अहो शाखामृगत्वं ते व्यक्तमेव प्लवङ्गम}
{लघुचित्ततयात्मानं न स्थापयसि यो मतौ} %||4-2-16||

\twolineshloka
{बुद्धिविज्ञानसम्पन्न इङ्गितैः सर्वमाचर}
{न ह्यबुद्धिं गतो राजा सर्वभूतानि शास्ति हि} %||4-2-17||

\twolineshloka
{सुग्रीवस्तु शुभं वाक्यं श्रुत्वा सर्वं हनूमतः}
{ततः शुभतरं वाक्यं हनूमन्तमुवाच ह} %||4-2-18||

\twolineshloka
{दीर्घबाहू विशालाक्षौ शरचापासिधारिणौ}
{कस्य न स्याद्भयं दृष्ट्वा एतौ सुरसुतोपमौ} %||4-2-19||

\twolineshloka
{वालिप्रणिहितावेतौ शङ्केऽहं पुरुषोत्तमौ}
{राजानो बहुमित्राश्च विश्वासो नात्र हि क्षमः} %||4-2-20||

\twolineshloka
{अरयश्च मनुष्येण विज्ञेयाश्छन्नचारिणः}
{विश्वस्तानामविश्वस्ताश्छिद्रेषु प्रहरन्ति हि} %||4-2-21||

\twolineshloka
{कृत्येषु वाली मेधावी राजानो बहुदर्शनाः}
{भवन्ति परहन्तारस्ते ज्ञेयाः प्राकृतैर्नरैः} %||4-2-22||

\twolineshloka
{तौ त्वया प्राकृतेनैव गत्वा ज्ञेयौ प्लवङ्गम}
{शङ्कितानां प्रकारैश्च रूपव्याभाषणेन च} %||4-2-23||

\twolineshloka
{लक्षयस्व तयोर्भावं प्रहृष्टमनसौ यदि}
{विश्वासयन्प्रशंसाभिरिङ्गितैश्च पुनः पुनः} %||4-2-24||

\twolineshloka
{ममैवाभिमुखं स्थित्वा पृच्छ त्वं हरिपुङ्गव}
{प्रयोजनं प्रवेशस्य वनस्यास्य धनुर्धरौ} %||4-2-25||

\twolineshloka
{शुद्धात्मानौ यदि त्वेतौ जानीहि त्वं प्लवङ्गम}
{व्याभाषितैर्वा रूपैर्वा विज्ञेया दुष्टतानयोः} %||4-2-26||

\twolineshloka
{इत्येवं कपिराजेन सन्दिष्टो मारुतात्मजः}
{चकार गमने बुद्धिं यत्र तौ रामलक्ष्मणौ} %||4-2-27||

\fourlineindentedshloka
{तथेति सम्पूज्य वचस्तु तस्य}
{ कपेः सुभीतस्य दुरासदस्य}
{महानुभावो हनुमान्ययौ तदा}
{ स यत्र रामोऽतिबलश्च लक्ष्मणः} %||4-3-28||


{॥इत्यार्षे श्रीमद्रामायणे वाल्मीकीये आदिकाव्ये किष्किन्धाकाण्डे द्वितीयः सर्गः॥}

\sect{तृतीयः सर्गः}

\twolineshloka
{वचो विज्ञाय हनुमान्सुग्रीवस्य महात्मनः}
{पर्वतादृश्यमूकात्तु पुप्लुवे यत्र राघवौ} %||4-3-1||

\twolineshloka
{स तत्र गत्वा हनुमान्बलवान्वानरोत्तमः}
{उपचक्राम तौ वाग्भिर्मृद्वीभिः सत्यविक्रमः} %||4-3-2||

\twolineshloka
{स्वकं रूपं परित्यज्य भिक्षुरूपेण वानरः}
{आबभाषे च तौ वीरौ यथावत्प्रशशंस च} %||4-3-3||

\twolineshloka
{राजर्षिदेवप्रतिमौ तापसौ संशितव्रतौ}
{देशं कथमिमं प्राप्तौ भवन्तौ वरवर्णिनौ} %||4-3-4||

\twolineshloka
{त्रासयन्तौ मृगगणानन्यांश्च वनचारिणः}
{पम्पातीररुहान्वृक्षान्वीक्षमाणौ समन्ततः} %||4-3-5||

\twolineshloka
{इमां नदीं शुभजलां शोभयन्तौ तरस्विनौ}
{धैर्यवन्तौ सुवर्णाभौ कौ युवां चीरवाससौ} %||4-3-6||

\twolineshloka
{सिंहविप्रेक्षितौ वीरौ सिंहातिबलविक्रमौ}
{शक्रचापनिभे चापे प्रगृह्य विपुलैर्भुजैः} %||4-3-7||

\twolineshloka
{श्रीमन्तौ रूपसम्पन्नौ वृषभश्रेष्ठविक्रमौ}
{हस्तिहस्तोपमभुजौ द्युतिमन्तौ नरर्षभौ} %||4-3-8||

\twolineshloka
{प्रभया पर्वतेन्द्रोऽयं युवयोरवभासितः}
{राज्यार्हावमरप्रख्यौ कथं देशमिहागतौ} %||4-3-9||

\twolineshloka
{पद्मपत्रेक्षणौ वीरौ जटामण्डलधारिणौ}
{अन्योन्यसदृशौ वीरौ देवलोकादिवागतौ} %||4-3-10||

\twolineshloka
{यदृच्छयेव सम्प्राप्तौ चन्द्रसूर्यौ वसुन्धराम्}
{विशालवक्षसौ वीरौ मानुषौ देवरूपिणौ} %||4-3-11||

\threelineshloka
{सिंहस्कन्धौ महासत्त्वौ समदाविव गोवृषौ}
{आयताश्च सुवृत्ताश्च बाहवः परिघोत्तमाः}
{सर्वभूषणभूषार्हाः किमर्थं न विभूषितः} %||4-3-12||

\twolineshloka
{उभौ योग्यावहं मन्ये रक्षितुं पृथिवीमिमाम्}
{ससागरवनां कृत्स्नां विन्ध्यमेरुविभूषिताम्} %||4-3-13||

\twolineshloka
{इमे च धनुषी चित्रे श्लक्ष्णे चित्रानुलेपने}
{प्रकाशेते यथेन्द्रस्य वज्रे हेमविभूषिते} %||4-3-14||

\twolineshloka
{सम्पूर्णा निशितैर्बाणैर्तूणाश्च शुभदर्शनाः}
{जीवितान्तकरैर्घोरैर्ज्वलद्भिरिव पन्नगैः} %||4-3-15||

\twolineshloka
{महाप्रमाणौ विपुलौ तप्तहाटकभूषितौ}
{खड्गावेतौ विराजेते निर्मुक्तभुजगाविव} %||4-3-16||

{एवं मां परिभाषन्तं कस्माद्वै नाभिभाषथः॥\devanumber{17}॥} %||4-3-17|| (Check)

\addtocounter{shlokacount}{1}

\twolineshloka
{सुग्रीवो नाम धर्मात्मा कश्चिद्वानरयूथपः}
{वीरो विनिकृतो भ्रात्रा जगद्भ्रमति दुःखितः} %||4-3-18||

\twolineshloka
{प्राप्तोऽहं प्रेषितस्तेन सुग्रीवेण महात्मना}
{राज्ञा वानरमुख्यानां हनुमान्नाम वानरः} %||4-3-19||

\twolineshloka
{युवाभ्यां सह धर्मात्मा सुग्रीवः सख्यमिच्छति}
{तस्य मां सचिवं वित्तं वानरं पवनात्मजम्} %||4-3-20||

\twolineshloka
{भिक्षुरूपप्रतिच्छन्नं सुग्रीवप्रियकाम्यया}
{ऋश्यमूकादिह प्राप्तं कामगं कामरूपिणम्} %||4-3-21||

\twolineshloka
{एवमुक्त्वा तु हनुमांस्तौ वीरौ रामलक्ष्मणौ}
{वाक्यज्ञौ वाक्यकुशलः पुनर्नोवाच किञ्चन} %||4-3-22||

\twolineshloka
{एतच्छ्रुत्वा वचस्तस्य रामो लक्ष्मणमब्रवीत्}
{प्रहृष्टवदनः श्रीमान्भ्रातरं पार्श्वतः स्थितम्} %||4-3-23||

\twolineshloka
{सचिवोऽयं कपीन्द्रस्य सुग्रीवस्य महात्मनः}
{तमेव काङ्क्षमाणस्य ममान्तिकमुपागतः} %||4-3-24||

\twolineshloka
{तमभ्यभाष सौमित्रे सुग्रीवसचिवं कपिम्}
{वाक्यज्ञं मधुरैर्वाक्यैः स्नेहयुक्तमरिन्दमम्} %||4-4-25||


{॥इत्यार्षे श्रीमद्रामायणे वाल्मीकीये आदिकाव्ये किष्किन्धाकाण्डे तृतीयः सर्गः॥}

\sect{चतुर्थः सर्गः}

\twolineshloka
{ततः प्रहृष्टो हनुमान्कृत्यवानिति तद्वचः}
{श्रुत्वा मधुरसम्भाषं सुग्रीवं मनसा गतः} %||4-4-1||

\twolineshloka
{भव्यो राज्यागमस्तस्य सुग्रीवस्य महात्मनः}
{यदयं कृत्यवान्प्राप्तः कृत्यं चैतदुपागतम्} %||4-4-2||

\twolineshloka
{ततः परमसंहृष्टो हनूमान्प्लवगर्षभः}
{प्रत्युवाच ततो वाक्यं रामं वाक्यविशारदः} %||4-4-3||

\twolineshloka
{किमर्थं त्वं वनं घोरं पम्पाकाननमण्डितम्}
{आगतः सानुजो दुर्गं नानाव्यालमृगायुतम्} %||4-4-4||

\twolineshloka
{तस्य तद्वचनं श्रुत्वा लक्ष्मणो रामचोदितः}
{आचचक्षे महात्मानं रामं दशरथात्मजम्} %||4-4-5||

\twolineshloka
{राजा दशरथो नाम द्युतिमान्धर्मवत्सलः}
{तस्यायं पूर्वजः पुत्रो रामो नाम जनैः श्रुतः} %||4-4-6||

\twolineshloka
{शरण्यः सर्वभूतानां पितुर्निर्देशपारगः}
{वीरो दशरथस्यायं पुत्राणां गुणवत्तरः} %||4-4-7||

\threelineshloka
{राज्याद्भ्रष्टो वने वस्तुं मया सार्धमिहागतः}
{भार्यया च महातेजाः सीतयानुगतो वशी}
{दिनक्षये महातेजाः प्रभयेव दिवाकरः} %||4-4-8||

\twolineshloka
{अहमस्यावरो भ्राता गुणैर्दास्यमुपागतः}
{कृतज्ञस्य बहुज्ञस्य लक्ष्मणो नाम नामतः} %||4-4-9||

\twolineshloka
{सुखार्हस्य महार्हस्य सर्वभूतहितात्मनः}
{ऐश्वर्येण विहीनस्य वनवासाश्रितस्य च} %||4-4-10||

\twolineshloka
{रक्षसापहृता भार्या रहिते कामरूपिणा}
{तच्च न ज्ञायते रक्षः पत्नी येनास्य सा हृता} %||4-4-11||

\twolineshloka
{दनुर्नाम श्रियः पुत्रः शापाद्राक्षसतां गतः}
{आख्यातस्तेन सुग्रीवः समर्थो वानराधिपः} %||4-4-12||

\twolineshloka
{स ज्ञास्यति महावीर्यस्तव भार्यापहारिणम्}
{एवमुक्त्वा दनुः स्वर्गं भ्राजमानो गतः सुखम्} %||4-4-13||

\twolineshloka
{एतत्ते सर्वमाख्यातं याथातथ्येन पृच्छतः}
{अहं चैव हि रामश्च सुग्रीवं शरणं गतौ} %||4-4-14||

\twolineshloka
{एष दत्त्वा च वित्तानि प्राप्य चानुत्तमं यशः}
{लोकनाथः पुरा भूत्वा सुग्रीवं नाथमिच्छति} %||4-4-15||

\twolineshloka
{शोकाभिभूते रामे तु शोकार्ते शरणं गते}
{कर्तुमर्हति सुग्रीवः प्रसादं सह यूथपैः} %||4-4-16||

\twolineshloka
{एवं ब्रुवाणं सौमित्रिं करुणं साश्रुपातनम्}
{हनूमान्प्रत्युवाचेदं वाक्यं वाक्यविशारदः} %||4-4-17||

\twolineshloka
{ईदृशा बुद्धिसम्पन्ना जितक्रोधा जितेन्द्रियाः}
{द्रष्टव्या वानरेन्द्रेण दिष्ट्या दर्शनमागताः} %||4-4-18||

\twolineshloka
{स हि राज्याच्च विभ्रष्टः कृतवैरश्च वालिना}
{हृतदारो वने त्रस्तो भ्रात्रा विनिकृतो भृशम्} %||4-4-19||

\twolineshloka
{करिष्यति स साहाय्यं युवयोर्भास्करात्मजः}
{सुग्रीवः सह चास्माभिः सीतायाः परिमार्गणे} %||4-4-20||

\twolineshloka
{इत्येवमुक्त्वा हनुमाञ्श्लक्ष्णं मधुरया गिरा}
{बभाषे सोऽभिगच्छामः सुग्रीवमिति राघवम्} %||4-4-21||

\twolineshloka
{एवं ब्रुवाणं धर्मात्मा हनूमन्तं स लक्ष्मणः}
{प्रतिपूज्य यथान्यायमिदं प्रोवाच राघवम्} %||4-4-22||

\twolineshloka
{कपिः कथयते हृष्टो यथायं मारुतात्मजः}
{कृत्यवान्सोऽपि सम्प्राप्तः कृतकृत्योऽसि राघव} %||4-4-23||

\twolineshloka
{प्रसन्नमुखवर्णश्च व्यक्तं हृष्टश्च भाषते}
{नानृतं वक्ष्यते वीरो हनूमान्मारुतात्मजः} %||4-4-24||

\twolineshloka
{ततः स तु महाप्राज्ञो हनूमान्मारुतात्मजः}
{जगामादाय तौ वीरौ हरिराजाय राघवौ} %||4-4-25||

\fourlineindentedshloka
{स तु विपुल यशाः कपिप्रवीरः}
{ पवनसुतः कृतकृत्यवत्प्रहृष्टः}
{गिरिवरमुरुविक्रमः प्रयातः}
{ स शुभमतिः सह रामलक्ष्मणाभ्याम्} %||4-5-26||


{॥इत्यार्षे श्रीमद्रामायणे वाल्मीकीये आदिकाव्ये किष्किन्धाकाण्डे चतुर्थः सर्गः॥}

\sect{पञ्चमः सर्गः}

\twolineshloka
{ऋश्यमूकात्तु हनुमान्गत्वा तं मलयं गिरम्}
{आचचक्षे तदा वीरौ कपिराजाय राघवौ} %||4-5-1||

\twolineshloka
{अयं रामो महाप्राज्ञः सम्प्राप्तो दृढविक्रमः}
{लक्ष्मणेन सह भ्रात्रा रामोऽयं सत्यविक्रमः} %||4-5-2||

\twolineshloka
{इक्ष्वाकूणां कुले जातो रामो दशरथात्मजः}
{धर्मे निगदितश्चैव पितुर्निर्देशपालकः} %||4-5-3||

\twolineshloka
{तस्यास्य वसतोऽरण्ये नियतस्य महात्मनः}
{रक्षसापहृता भार्या स त्वां शरणमागतः} %||4-5-4||

\twolineshloka
{राजसूयाश्वमेधैश्च वह्निर्येनाभितर्पितः}
{दक्षिणाश्च तथोत्सृष्टा गावः शतसहस्रशः} %||4-5-5||

\twolineshloka
{तपसा सत्यवाक्येन वसुधा येन पालिता}
{स्त्रीहेतोस्तस्य पुत्रोऽयं रामस्त्वां शरणं गतः} %||4-5-6||

\twolineshloka
{भवता सख्यकामौ तौ भ्रातरौ रामलक्ष्मणौ}
{प्रतिगृह्यार्चयस्वेमौ पूजनीयतमावुभौ} %||4-5-7||

\twolineshloka
{श्रुत्वा हनुमतो वाक्यं सुग्रीवो हृष्टमानसः}
{भयं स राघवाद्घोरं प्रजहौ विगतज्वरः} %||4-5-8||

\twolineshloka
{स कृत्वा मानुषं रूपं सुग्रीवः प्लवगाधिपः}
{दर्शनीयतमो भूत्वा प्रीत्या प्रोवाच राघवम्} %||4-5-9||

\twolineshloka
{भवान्धर्मविनीतश्च विक्रान्तः सर्ववत्सलः}
{आख्याता वायुपुत्रेण तत्त्वतो मे भवद्गुणाः} %||4-5-10||

\twolineshloka
{तन्ममैवैष सत्कारो लाभश्चैवोत्तमः प्रभो}
{यत्त्वमिच्छसि सौहार्दं वानरेण मया सह} %||4-5-11||

\twolineshloka
{रोचते यदि वा सख्यं बाहुरेष प्रसारितः}
{गृह्यतां पाणिना पाणिर्मर्यादा वध्यतां ध्रुवा} %||4-5-12||

\threelineshloka
{एतत्तु वचनं श्रुत्वा सुग्रीवस्य सुभाषितम्}
{सम्प्रहृष्टमना हस्तं पीडयामास पाणिना}
{हृद्यं सौहृदमालम्ब्य पर्यष्वजत पीडितम्} %||4-5-13||

\twolineshloka
{ततो हनूमान्सन्त्यज्य भिक्षुरूपमरिन्दमः}
{काष्ठयोः स्वेन रूपेण जनयामास पावकम्} %||4-5-14||

\twolineshloka
{दीप्यमानं ततो वह्निं पुष्पैरभ्यर्च्य सत्कृतम्}
{तयोर्मध्ये तु सुप्रीतो निदधे सुसमाहितः} %||4-5-15||

\twolineshloka
{ततोऽग्निं दीप्यमानं तौ चक्रतुश्च प्रदक्षिणम्}
{सुग्रीवो राघवश्चैव वयस्यत्वमुपागतौ} %||4-5-16||

\twolineshloka
{ततः सुप्रीत मनसौ तावुभौ हरिराघवौ}
{अन्योन्यमभिवीक्षन्तौ न तृप्तिमुपजग्मतुः} %||4-5-17||

\twolineshloka
{ततः सर्वार्थविद्वांसं रामं दशरथात्मजम्}
{सुग्रीवः प्राह तेजस्वी वाक्यमेकमनास्तदा} %||4-6-18||


{॥इत्यार्षे श्रीमद्रामायणे वाल्मीकीये आदिकाव्ये किष्किन्धाकाण्डे पञ्चमः सर्गः॥}

\sect{षष्ठः सर्गः}

\twolineshloka
{अयमाख्याति मे राम सचिवो मन्त्रिसत्तमः}
{हनुमान्यन्निमित्तं त्वं निर्जनं वनमागतः} %||4-6-1||

\twolineshloka
{लक्ष्मणेन सह भ्रात्रा वसतश्च वने तव}
{रक्षसापहृता भार्या मैथिली जनकात्मजा} %||4-6-2||

\twolineshloka
{त्वया वियुक्ता रुदती लक्ष्मणेन च धीमता}
{अन्तरं प्रेप्सुना तेन हत्वा गृध्रं जटायुषम्} %||4-6-3||

\twolineshloka
{भार्यावियोगजं दुःखं नचिरात्त्वं विमोक्ष्यसे}
{अहं तामानयिष्यामि नष्टां वेदश्रुतिं यथा} %||4-6-4||

\twolineshloka
{रसातले वा वर्तन्तीं वर्तन्तीं वा नभस्तले}
{अहमानीय दास्यामि तव भार्यामरिन्दम} %||4-6-5||

\twolineshloka
{इदं तथ्यं मम वचस्त्वमवेहि च राघव}
{त्यज शोकं महाबाहो तां कान्तामानयामि ते} %||4-6-6||

\twolineshloka
{अनुमानात्तु जानामि मैथिली सा न संशयः}
{ह्रियमाणा मया दृष्टा रक्षसा क्रूरकर्मणा} %||4-6-7||

\twolineshloka
{क्रोशन्ती राम रामेति लक्ष्मणेति च विस्वरम्}
{स्फुरन्ती रावणस्याङ्के पन्नगेन्द्रवधूर्यथा} %||4-6-8||

\twolineshloka
{आत्मना पञ्चमं मां हि दृष्ट्वा शैलतटे स्थितम्}
{उत्तरीयं तया त्यक्तं शुभान्याभरणानि च} %||4-6-9||

\twolineshloka
{तान्यस्माभिर्गृहीतानि निहितानि च राघव}
{आनयिष्याम्यहं तानि प्रत्यभिज्ञातुमर्हसि} %||4-6-10||

\twolineshloka
{तमब्रवीत्ततो रामः सुग्रीवं प्रियवादिनम्}
{आनयस्व सखे शीघ्रं किमर्थं प्रविलम्बसे} %||4-6-11||

\twolineshloka
{एवमुक्तस्तु सुग्रीवः शैलस्य गहनां गुहाम्}
{प्रविवेश ततः शीघ्रं राघवप्रियकाम्यया} %||4-6-12||

\twolineshloka
{उत्तरीयं गृहीत्वा तु शुभान्याभरणानि च}
{इदं पश्येति रामाय दर्शयामास वानरः} %||4-6-13||

\twolineshloka
{ततो गृहीत्वा तद्वासः शुभान्याभरणानि च}
{अभवद्बाष्पसंरुद्धो नीहारेणेव चन्द्रमाः} %||4-6-14||

\twolineshloka
{सीतास्नेहप्रवृत्तेन स तु बाष्पेण दूषितः}
{हा प्रियेति रुदन्धैर्यमुत्सृज्य न्यपतत्क्षितौ} %||4-6-15||

\twolineshloka
{हृदि कृत्वा स बहुशस्तमलङ्कारमुत्तमम्}
{निशश्वास भृशं सर्पो बिलस्थ इव रोषितः} %||4-6-16||

\twolineshloka
{अविच्छिन्नाश्रुवेगस्तु सौमित्रिं वीक्ष्य पार्श्वतः}
{परिदेवयितुं दीनं रामः समुपचक्रमे} %||4-6-17||

\twolineshloka
{पश्य लक्ष्मण वैदेह्या सन्त्यक्तं ह्रियमाणया}
{उत्तरीयमिदं भूमौ शरीराद्भूषणानि च} %||4-6-18||

\twolineshloka
{शाद्वलिन्यां ध्रुवं भूम्यां सीतया ह्रियमाणया}
{उत्सृष्टं भूषणमिदं तथारूपं हि दृश्यते} %||4-6-19||

\twolineshloka
{ब्रूहि सुग्रीव कं देशं ह्रियन्ती लक्षिता त्वया}
{रक्षसा रौद्ररूपेण मम प्राणसमा प्रिया} %||4-6-20||

\twolineshloka
{क्व वा वसति तद्रक्षो महद्व्यसनदं मम}
{यन्निमित्तमहं सर्वान्नाशयिष्यामि राक्षसान्} %||4-6-21||

\twolineshloka
{हरता मैथिलीं येन मां च रोषयता भृशम्}
{आत्मनो जीवितान्ताय मृत्युद्वारमपावृतम्} %||4-6-22||

\fourlineindentedshloka
{मम दयिततमा हृता वनाद्-}
{ रजनिचरेण विमथ्य येन सा}
{कथय मम रिपुं तमद्य वै}
{ प्रवगपते यमसंनिधिं नयामि} %||4-7-23||


{॥इत्यार्षे श्रीमद्रामायणे वाल्मीकीये आदिकाव्ये किष्किन्धाकाण्डे षष्ठः सर्गः॥}

\sect{सप्तमः सर्गः}

\twolineshloka
{एवमुक्तस्तु सुग्रीवो रामेणार्तेन वानरः}
{अब्रवीत्प्राञ्जलिर्वाक्यं सबाष्पं बाष्पगद्गदः} %||4-7-1||

\twolineshloka
{न जाने निलयं तस्य सर्वथा पापरक्षसः}
{सामर्थ्यं विक्रमं वापि दौष्कुलेयस्य वा कुलम्} %||4-7-2||

\twolineshloka
{सत्यं तु प्रतिजानामि त्यज शोकमरिन्दम}
{करिष्यामि तथा यत्नं यथा प्राप्स्यसि मैथिलीम्} %||4-7-3||

\twolineshloka
{रावणं सगणं हत्वा परितोष्यात्मपौरुषम्}
{तथास्मि कर्ता नचिराद्यथा प्रीतो भविष्यसि} %||4-7-4||

\twolineshloka
{अलं वैक्लव्यमालम्ब्य धैर्यमात्मगतं स्मर}
{त्वद्विधानां न सदृशमीदृशं बुद्धिलाघवम्} %||4-7-5||

\twolineshloka
{मयापि व्यसनं प्राप्तं भार्या हरणजं महत्}
{न चाहमेवं शोचामि न च धैर्यं परित्यजे} %||4-7-6||

\twolineshloka
{नाहं तामनुशोचामि प्राकृतो वानरोऽपि सन्}
{महात्मा च विनीतश्चा किं पुनर्धृतिमान्भवान्} %||4-7-7||

\twolineshloka
{बाष्पमापतितं धैर्यान्निग्रहीतुं त्वमर्हसि}
{मर्यादां सत्त्वयुक्तानां धृतिं नोत्स्रष्टुमर्हसि} %||4-7-8||

\twolineshloka
{व्यसने वार्थ कृच्छ्रे वा भये वा जीवितान्तगे}
{विमृशन्वै स्वया बुद्ध्या धृतिमान्नावसीदति} %||4-7-9||

\twolineshloka
{बालिशस्तु नरो नित्यं वैक्लव्यं योऽनुवर्तते}
{स मज्जत्यवशः शोके भाराक्रान्तेव नौर्जले} %||4-7-10||

\twolineshloka
{एषोऽञ्जलिर्मया बद्धः प्रणयात्त्वां प्रसादये}
{पौरुषं श्रय शोकस्य नान्तरं दातुमर्हसि} %||4-7-11||

\twolineshloka
{ये शोकमनुवर्तन्ते न तेषां विद्यते सुखम्}
{तेजश्च क्षीयते तेषां न त्वं शोचितुमर्हसि} %||4-7-12||

\twolineshloka
{हितं वयस्य भावेन ब्रूहि नोपदिशामि ते}
{वयस्यतां पूजयन्मे न त्वं शोचितुमर्हसि} %||4-7-13||

\twolineshloka
{मधुरं सान्त्वितस्तेन सुग्रीवेण स राघवः}
{मुखमश्रुपरिक्लिन्नं वस्त्रान्तेन प्रमार्जयत्} %||4-7-14||

\twolineshloka
{प्रकृतिष्ठस्तु काकुत्स्थः सुग्रीववचनात्प्रभुः}
{सम्परिष्वज्य सुग्रीवमिदं वचनमब्रवीत्} %||4-7-15||

\twolineshloka
{कर्तव्यं यद्वयस्येन स्निग्धेन च हितेन च}
{अनुरूपं च युक्तं च कृतं सुग्रीव तत्त्वया} %||4-7-16||

\twolineshloka
{एष च प्रकृतिष्ठोऽहमनुनीतस्त्वया सखे}
{दुर्लभो हीदृशो बन्धुरस्मिन्काले विशेषतः} %||4-7-17||

\twolineshloka
{किं तु यत्नस्त्वया कार्यो मैथिल्याः परिमार्गणे}
{राक्षसस्य च रौद्रस्य रावणस्य दुरात्मनः} %||4-7-18||

\twolineshloka
{मया च यदनुष्ठेयं विस्रब्धेन तदुच्यताम्}
{वर्षास्विव च सुक्षेत्रे सर्वं सम्पद्यते तव} %||4-7-19||

\twolineshloka
{मया च यदिदं वाक्यमभिमानात्समीरितम्}
{तत्त्वया हरिशार्दूल तत्त्वमित्युपधार्यताम्} %||4-7-20||

\twolineshloka
{अनृतं नोक्तपूर्वं मे न च वक्ष्ये कदाचन}
{एतत्ते प्रतिजानामि सत्येनैव शपामि ते} %||4-7-21||

\twolineshloka
{ततः प्रहृष्टः सुग्रीवो वानरैः सचिवैः सह}
{राघवस्य वचः श्रुत्वा प्रतिज्ञातं विशेषतः} %||4-7-22||

\fourlineindentedshloka
{महानुभावस्य वचो निशम्य}
{ हरिर्नराणामृषभस्य तस्य}
{कृतं स मेने हरिवीर मुख्यः}
{ तदा स्वकार्यं हृदयेन विद्वान्} %||4-8-23||


{॥इत्यार्षे श्रीमद्रामायणे वाल्मीकीये आदिकाव्ये किष्किन्धाकाण्डे सप्तमः सर्गः॥}

\sect{अष्टमः सर्गः}

\twolineshloka
{परितुष्टस्तु सुग्रीवस्तेन वाक्येन वानरः}
{लक्ष्मणस्याग्रजं राममिदं वचनमब्रवीत्} %||4-8-1||

\twolineshloka
{सर्वथाहमनुग्राह्यो देवतानामसंशयः}
{उपपन्नगुणोपेतः सखा यस्य भवान्मम} %||4-8-2||

\twolineshloka
{शक्यं खलु भवेद्राम सहायेन त्वयानघ}
{सुरराज्यमपि प्राप्तुं स्वराज्यं किं पुनः प्रभो} %||4-8-3||

\twolineshloka
{सोऽहं सभाज्यो बन्धूनां सुहृदां चैव राघव}
{यस्याग्निसाक्षिकं मित्रं लब्धं राघववंशजम्} %||4-8-4||

\twolineshloka
{अहमप्यनुरूपस्ते वयस्यो ज्ञास्यसे शनैः}
{न तु वक्तुं समर्थोऽहं स्वयमात्मगतान्गुणान्} %||4-8-5||

\twolineshloka
{महात्मनां तु भूयिष्ठं त्वद्विधानां कृतात्मनाम्}
{निश्चला भवति प्रीतिर्धैर्यमात्मवतामिव} %||4-8-6||

\twolineshloka
{रजतं वा सुवर्णं वा वस्त्राण्याभरणानि वा}
{अविभक्तानि साधूनामवगच्छन्ति साधवः} %||4-8-7||

\twolineshloka
{आढ्यो वापि दरिद्रो वा दुःखितः सुखितोऽपि वा}
{निर्दोषो वा सदोषो वा वयस्यः परमा गतिः} %||4-8-8||

\twolineshloka
{धनत्यागः सुखत्यागो देहत्यागोऽपि वा पुनः}
{वयस्यार्थे प्रवर्तन्ते स्नेहं दृष्ट्वा तथाविधम्} %||4-8-9||

\twolineshloka
{तत्तथेत्यब्रवीद्रामः सुग्रीवं प्रियवादिनम्}
{लक्ष्मणस्याग्रतो लक्ष्म्या वासवस्येव धीमतः} %||4-8-10||

\twolineshloka
{ततो रामं स्थितं दृष्ट्वा लक्ष्मणं च महाबलम्}
{सुग्रीवः सर्वतश्चक्षुर्वने लोलमपातयत्} %||4-8-11||

\twolineshloka
{स ददर्श ततः सालमविदूरे हरीश्वरः}
{सुपुष्पमीषत्पत्राढ्यं भ्रमरैरुपशोभितम्} %||4-8-12||

\twolineshloka
{तस्यैकां पर्णबहुलां भङ्क्त्वा शाखां सुपुष्पिताम्}
{सालस्यास्तीर्य सुग्रीवो निषसाद सराघवः} %||4-8-13||

\twolineshloka
{तावासीनौ ततो दृष्ट्वा हनूमानपि लक्ष्मणम्}
{सालशाखां समुत्पाट्य विनीतमुपवेशयत्} %||4-8-14||

\twolineshloka
{ततः प्रहृष्टः सुग्रीवः श्लक्ष्णं मधुरया गिरा}
{उवाच प्रणयाद्रामं हर्षव्याकुलिताक्षरम्} %||4-8-15||

\twolineshloka
{अहं विनिकृतो भ्रात्रा चराम्येष भयार्दितः}
{ऋश्यमूकं गिरिवरं हृतभार्यः सुदुःखितः} %||4-8-16||

\twolineshloka
{सोऽहं त्रस्तो भये मग्नो वसाम्युद्भ्रान्तचेतनः}
{वालिना निकृतो भ्रात्रा कृतवैरश्च राघव} %||4-8-17||

\twolineshloka
{वालिनो मे भयार्तस्य सर्वलोकाभयङ्कर}
{ममापि त्वमनाथस्य प्रसादं कर्तुमर्हसि} %||4-8-18||

\twolineshloka
{एवमुक्तस्तु तेजस्वी धर्मज्ञो धर्मवत्सलः}
{प्रत्युवाच स काकुत्स्थः सुग्रीवं प्रहसन्निव} %||4-8-19||

\twolineshloka
{उपकारफलं मित्रमपकारोऽरिलक्षणम्}
{अद्यैव तं हनिष्यामि तव भार्यापहारिणम्} %||4-8-20||

\twolineshloka
{इमे हि मे महावेगाः पत्रिणस्तिग्मतेजसः}
{कार्तिकेयवनोद्भूताः शरा हेमविभूषिताः} %||4-8-21||

\twolineshloka
{कङ्कपत्रप्रतिच्छन्ना महेन्द्राशनिसंनिभाः}
{सुपर्वाणः सुतीक्ष्णाग्रा सरोषा भुजगा इव} %||4-8-22||

\twolineshloka
{भ्रातृसंज्ञममित्रं ते वालिनं कृतकिल्बिषम्}
{शरैर्विनिहतं पश्य विकीर्णमिव पर्वतम्} %||4-8-23||

\twolineshloka
{राघवस्य वचः श्रुत्वा सुग्रीवो वाहिनीपतिः}
{प्रहर्षमतुलं लेभे साधु साध्विति चाब्रवीत्} %||4-8-24||

\twolineshloka
{रामशोकाभिभूतोऽहं शोकार्तानां भवान्गतिः}
{वयस्य इति कृत्वा हि त्वय्यहं परिदेवये} %||4-8-25||

\twolineshloka
{त्वं हि पाणिप्रदानेन वयस्यो सोऽग्निसाक्षिकः}
{कृतः प्राणैर्बहुमतः सत्येनापि शपाम्यहम्} %||4-8-26||

\twolineshloka
{वयस्य इति कृत्वा च विस्रब्धं प्रवदाम्यहम्}
{दुःखमन्तर्गतं यन्मे मनो दहति नित्यशः} %||4-8-27||

\twolineshloka
{एतावदुक्त्वा वचनं बाष्पदूषितलोचनः}
{बाष्पोपहतया वाचा नोच्चैः शक्नोति भाषितुम्} %||4-8-28||

\twolineshloka
{बाष्पवेगं तु सहसा नदीवेगमिवागतम्}
{धारयामास धैर्येण सुग्रीवो रामसंनिधौ} %||4-8-29||

\twolineshloka
{संनिगृह्य तु तं बाष्पं प्रमृज्य नयने शुभे}
{विनिःश्वस्य च तेजस्वी राघवं पुनरब्रवीत्} %||4-8-30||

\twolineshloka
{पुराहं वलिना राम राज्यात्स्वादवरोपितः}
{परुषाणि च संश्राव्य निर्धूतोऽस्मि बलीयसा} %||4-8-31||

\twolineshloka
{हृता भार्या च मे तेन प्राणेभ्योऽपि गरीयसी}
{सुहृदश्च मदीया ये संयता बन्धनेषु ते} %||4-8-32||

\twolineshloka
{यत्नवांश्च सुदुष्टात्मा मद्विनाशाय राघव}
{बहुशस्तत्प्रयुक्ताश्च वानरा निहता मया} %||4-8-33||

\twolineshloka
{शङ्कया त्वेतया चाहं दृष्ट्वा त्वामपि राघव}
{नोपसर्पाम्यहं भीतो भये सर्वे हि बिभ्यति} %||4-8-34||

\twolineshloka
{केवलं हि सहाया मे हनुमत्प्रमुखास्त्विमे}
{अतोऽहं धारयाम्यद्य प्राणान्कृच्छ्र गतोऽपि सन्} %||4-8-35||

\twolineshloka
{एते हि कपयः स्निग्धा मां रक्षन्ति समन्ततः}
{सह गच्छन्ति गन्तव्ये नित्यं तिष्ठन्ति च स्थिते} %||4-8-36||

\twolineshloka
{सङ्क्षेपस्त्वेष मे राम किमुक्त्वा विस्तरं हि ते}
{स मे ज्येष्ठो रिपुर्भ्राता वाली विश्रुतपौरुषः} %||4-8-37||

\twolineshloka
{तद्विनाशाद्धि मे दुःखं प्रनष्टं स्यादनन्तरम्}
{सुखं मे जीवितं चैव तद्विनाशनिबन्धनम्} %||4-8-38||

\twolineshloka
{एष मे राम शोकान्तः शोकार्तेन निवेदितः}
{दुःखितोऽदुःखितो वापि सख्युर्नित्यं सखा गतिः} %||4-8-39||

\twolineshloka
{श्रुत्वैतच्च वचो रामः सुग्रीवमिदमब्रवीत्}
{किंनिमित्तमभूद्वैरं श्रोतुमिच्छामि तत्त्वतः} %||4-8-40||

\twolineshloka
{सुखं हि कारणं श्रुत्वा वैरस्य तव वानर}
{आनन्तर्यं विधास्यामि सम्प्रधार्य बलाबलम्} %||4-8-41||

\twolineshloka
{बलवान्हि ममामर्षः श्रुत्वा त्वामवमानितम्}
{वर्धते हृदयोत्कम्पी प्रावृड्वेग इवाम्भसः} %||4-8-42||

\twolineshloka
{हृष्टः कथय विस्रब्धो यावदारोप्यते धनुः}
{सृष्टश्च हि मया बाणो निरस्तश्च रिपुस्तव} %||4-8-43||

\twolineshloka
{एवमुक्तस्तु सुग्रीवः काकुत्स्थेन महात्मना}
{प्रहर्षमतुलं लेभे चतुर्भिः सह वानरैः} %||4-8-44||

\twolineshloka
{ततः प्रहृष्टवदनः सुग्रीवो लक्ष्मणाग्रजे}
{वैरस्य कारणं तत्त्वमाख्यातुमुपचक्रमे} %||4-9-45||


{॥इत्यार्षे श्रीमद्रामायणे वाल्मीकीये आदिकाव्ये किष्किन्धाकाण्डे अष्टमः सर्गः॥}

\sect{नवमः सर्गः}

\twolineshloka
{वाली नाम मम भ्राता ज्येष्ठः शत्रुनिषूदनः}
{पितुर्बहुमतो नित्यं मम चापि तथा पुरा} %||4-9-1||

\twolineshloka
{पितर्युपरतेऽस्माकं ज्येष्ठोऽयमिति मन्त्रिभिः}
{कपीनामीश्वरो राज्ये कृतः परमसम्मतः} %||4-9-2||

\twolineshloka
{राज्यं प्रशासतस्तस्य पितृपैतामहं महत्}
{अहं सर्वेषु कालेषु प्रणतः प्रेष्यवत्स्थितः} %||4-9-3||

\twolineshloka
{मायावी नाम तेजस्वी पूर्वजो दुन्दुभेः सुतः}
{तेन तस्य महद्वैरं स्त्रीकृतं विश्रुतं पुरा} %||4-9-4||

\twolineshloka
{स तु सुप्ते जने रात्रौ किष्किन्धाद्वारमागतः}
{नर्दति स्म सुसंरब्धो वालिनं चाह्वयद्रणे} %||4-9-5||

\twolineshloka
{प्रसुप्तस्तु मम भ्राता नर्दितं भैरवस्वनम्}
{श्रुत्वा न ममृषे वाली निष्पपात जवात्तदा} %||4-9-6||

\twolineshloka
{स तु वै निःसृतः क्रोधात्तं हन्तुमसुरोत्तमम्}
{वार्यमाणस्ततः स्त्रीभिर्मया च प्रणतात्मना} %||4-9-7||

\twolineshloka
{स तु निर्धूय सर्वान्नो निर्जगाम महाबलः}
{ततोऽहमपि सौहार्दान्निःसृतो वालिना सह} %||4-9-8||

\twolineshloka
{स तु मे भ्रातरं दृष्ट्वा मां च दूरादवस्थितम्}
{असुरो जातसन्त्रासः प्रदुद्राव तदा भृशम्} %||4-9-9||

\twolineshloka
{तस्मिन्द्रवति सन्त्रस्ते ह्यावां द्रुततरं गतौ}
{प्रकाशोऽपि कृतो मार्गश्चन्द्रेणोद्गच्छता तदा} %||4-9-10||

\twolineshloka
{स तृणैरावृतं दुर्गं धरण्या विवरं महत्}
{प्रविवेशासुरो वेगादावामासाद्य विष्ठितौ} %||4-9-11||

\twolineshloka
{तं प्रविष्टं रिपुं दृष्ट्वा बिलं रोषवशं गतः}
{मामुवाच तदा वाली वचनं क्षुभितेन्द्रियः} %||4-9-12||

\twolineshloka
{इह त्वं तिष्ठ सुग्रीव बिलद्वारि समाहितः}
{यावदत्र प्रविश्याहं निहन्मि समरे रिपुम्} %||4-9-13||

\twolineshloka
{मया त्वेतद्वचः श्रुत्वा याचितः स परन्तप}
{शापयित्वा च मां पद्भ्यां प्रविवेश बिलं तदा} %||4-9-14||

\twolineshloka
{तस्य प्रविष्टस्य बिलं साग्रः संवत्सरो गतः}
{स्थितस्य च मम द्वारि स कालो व्यत्यवर्तत} %||4-9-15||

\twolineshloka
{अहं तु नष्टं तं ज्ञात्वा स्नेहादागतसम्भ्रमः}
{भ्रातरं न हि पश्यामि पापशङ्कि च मे मनः} %||4-9-16||

\twolineshloka
{अथ दीर्घस्य कालस्य बिलात्तस्माद्विनिःसृतम्}
{सफेनं रुधिरं रक्तमहं दृष्ट्वा सुदुःखितः} %||4-9-17||

\twolineshloka
{नर्दतामसुराणां च ध्वनिर्मे श्रोत्रमागतः}
{निरस्तस्य च सङ्ग्रामे क्रोशतो निःस्वनो गुरोः} %||4-9-18||

\threelineshloka
{अहं त्ववगतो बुद्ध्या चिह्नैस्तैर्भ्रातरं हतम्}
{पिधाय च बिलद्वारं शिलया गिरिमात्रया}
{शोकार्तश्चोदकं कृत्वा किष्किन्धामागतः सखे} %||4-9-19||

\twolineshloka
{गूहमानस्य मे तत्त्वं यत्नतो मन्त्रिभिः श्रुतम्}
{ततोऽहं तैः समागम्य समेतैरभिषेचितः} %||4-9-20||

\twolineshloka
{राज्यं प्रशासतस्तस्य न्यायतो मम राघव}
{आजगाम रिपुं हत्वा वाली तमसुरोत्तमम्} %||4-9-21||

\twolineshloka
{अभिषिक्तं तु मां दृष्ट्वा क्रोधात्संरक्तलोचनः}
{मदीयान्मन्त्रिणो बद्ध्वा परुषं वाक्यमब्रवीत्} %||4-9-22||

\twolineshloka
{निग्रहेऽपि समर्थस्य तं पापं प्रति राघव}
{न प्रावर्तत मे बुद्धिर्भ्रातृगौरवयन्त्रिता} %||4-9-23||

\twolineshloka
{मानयंस्तं महात्मानं यथावच्चाभ्यवादयम्}
{उक्ताश्च नाशिषस्तेन सन्तुष्टेनान्तरात्मना} %||4-10-24||


{॥इत्यार्षे श्रीमद्रामायणे वाल्मीकीये आदिकाव्ये किष्किन्धाकाण्डे नवमः सर्गः॥}

\sect{दशमः सर्गः}

\twolineshloka
{ततः क्रोधसमाविष्टं संरब्धं तमुपागतम्}
{अहं प्रसादयां चक्रे भ्रातरं प्रियकाम्यया} %||4-10-1||

\twolineshloka
{दिष्ट्यासि कुशली प्राप्तो निहतश्च त्वया रिपुः}
{अनाथस्य हि मे नाथस्त्वमेकोऽनाथनन्दनः} %||4-10-2||

\twolineshloka
{इदं बहुशलाकं ते पूर्णचन्द्रमिवोदितम्}
{छत्रं सवालव्यजनं प्रतीच्छस्व मयोद्यतम्} %||4-10-3||

\twolineshloka
{त्वमेव राजा मानार्हः सदा चाहं यथापुरा}
{न्यासभूतमिदं राज्यं तव निर्यातयाम्यहम्} %||4-10-4||

\twolineshloka
{मा च रोषं कृथाः सौम्य मयि शत्रुनिबर्हण}
{याचे त्वां शिरसा राजन्मया बद्धोऽयमञ्जलिः} %||4-10-5||

\twolineshloka
{बलादस्मि समागम्य मन्त्रिभिः पुरवासिभिः}
{राजभावे नियुक्तोऽहं शून्यदेशजिगीषया} %||4-10-6||

\twolineshloka
{स्निग्धमेवं ब्रुवाणं मां स तु निर्भर्त्स्य वानरः}
{धिक्त्वामिति च मामुक्त्वा बहु तत्तदुवाच ह} %||4-10-7||

\twolineshloka
{प्रकृतीश्च समानीय मन्त्रिणश्चैव सम्मतान्}
{मामाह सुहृदां मध्ये वाक्यं परमगर्हितम्} %||4-10-8||

\twolineshloka
{विदितं वो यथा रात्रौ मायावी स महासुरः}
{मां समाह्वयत क्रूरो युद्धाकाङ्क्षी सुदुर्मतिः} %||4-10-9||

\twolineshloka
{तस्य तद्गर्जितं श्रुत्वा निःसृतोऽहं नृपालयात्}
{अनुयातश्च मां तूर्णमयं भ्राता सुदारुणः} %||4-10-10||

\threelineshloka
{स तु दृष्ट्वैव मां रात्रौ सद्वितीयं महाबलः}
{प्राद्रवद्भयसन्त्रस्तो वीक्ष्यावां तमनुद्रुतौ}
{अनुद्रुतस्तु वेगेन प्रविवेश महाबिलम्} %||4-10-11||

\twolineshloka
{तं प्रविष्टं विदित्वा तु सुघोरं सुमहद्बिलम्}
{अयमुक्तोऽथ मे भ्राता मया तु क्रूरदर्शनः} %||4-10-12||

\twolineshloka
{अहत्वा नास्ति मे शक्तिः प्रतिगन्तुमितः पुरीम्}
{बिलद्वारि प्रतीक्ष त्वं यावदेनं निहन्म्यहम्} %||4-10-13||

\twolineshloka
{स्थितोऽयमिति मत्वा तु प्रविष्टोऽहं दुरासदम्}
{तं च मे मार्गमाणस्य गतः संवत्सरस्तदा} %||4-10-14||

\twolineshloka
{स तु दृष्टो मया शत्रुरनिर्वेदाद्भयावहः}
{निहतश्च मया तत्र सोऽसुरो बन्धुभिः सह} %||4-10-15||

\twolineshloka
{तस्यास्यात्तु प्रवृत्तेन रुधिरौघेण तद्बिलम्}
{पूर्णमासीद्दुराक्रामं स्तनतस्तस्य भूतले} %||4-10-16||

\twolineshloka
{सूदयित्वा तु तं शत्रुं विक्रान्तं दुन्दुभेः सुतम्}
{निष्क्रामन्नेव पश्यामि बिलस्य पिहितं मुखम्} %||4-10-17||

\twolineshloka
{विक्रोशमानस्य तु मे सुग्रीवेति पुनः पुनः}
{यदा प्रतिवचो नास्ति ततोऽहं भृशदुःखितः} %||4-10-18||

\twolineshloka
{पादप्रहारैस्तु मया बहुशस्तद्विदारितम्}
{ततोऽहं तेन निष्क्रम्य यथा पुनरुपागतः} %||4-10-19||

\twolineshloka
{तत्रानेनास्मि संरुद्धो राज्यं मार्गयतात्मनः}
{सुग्रीवेण नृशंसेन विस्मृत्य भ्रातृसौहृदम्} %||4-10-20||

\twolineshloka
{एवमुक्त्वा तु मां तत्र वस्त्रेणैकेन वानरः}
{तदा निर्वासयामास वाली विगतसाध्वसः} %||4-10-21||

\twolineshloka
{तेनाहमपविद्धश्च हृतदारश्च राघव}
{तद्भयाच्च महीकृत्स्ना क्रान्तेयं सवनार्णवा} %||4-10-22||

\twolineshloka
{ऋश्यमूकं गिरिवरं भार्याहरणदुःखितः}
{प्रविष्टोऽस्मि दुराधर्षं वालिनः कारणान्तरे} %||4-10-23||

\twolineshloka
{एतत्ते सर्वमाख्यातं वैरानुकथनं महत्}
{अनागसा मया प्राप्तं व्यसनं पश्य राघव} %||4-10-24||

\twolineshloka
{वालिनस्तु भयार्तस्य सर्वलोकाभयङ्कर}
{कर्तुमर्हसि मे वीर प्रसादं तस्य निग्रहात्} %||4-10-25||

\twolineshloka
{एवमुक्तः स तेजस्वी धर्मज्ञो धर्मसंहितम्}
{वचनं वक्तुमारेभे सुग्रीवं प्रहसन्निव} %||4-10-26||

\twolineshloka
{अमोघाः सूर्यसङ्काशा ममेमे निशिताः शराः}
{तस्मिन्वालिनि दुर्वृत्ते पतिष्यन्ति रुषान्विताः} %||4-10-27||

\twolineshloka
{यावत्तं न हि पश्येयं तव भार्यापहारिणम्}
{तावत्स जीवेत्पापात्मा वाली चारित्रदूषकः} %||4-10-28||

\twolineshloka
{आत्मानुमानात्पश्यामि मग्नं त्वां शोकसागरे}
{त्वामहं तारयिष्यामि कामं प्राप्स्यसि पुष्कलम्} %||4-11-29||


{॥इत्यार्षे श्रीमद्रामायणे वाल्मीकीये आदिकाव्ये किष्किन्धाकाण्डे दशमः सर्गः॥}

\sect{एकादशः सर्गः}

\twolineshloka
{रामस्य वचनं श्रुत्वा हर्षपौरुषवर्धनम्}
{सुग्रीवः पूजयां चक्रे राघवं प्रशशंस च} %||4-11-1||

\twolineshloka
{असंशयं प्रज्वलितैस्तीक्ष्णैर्मर्मातिगैः शरैः}
{त्वं दहेः कुपितो लोकान्युगान्त इव भास्करः} %||4-11-2||

\twolineshloka
{वालिनः पौरुषं यत्तद्यच्च वीर्यं धृतिश्च या}
{तन्ममैकमनाः श्रुत्वा विधत्स्व यदनन्तरम्} %||4-11-3||

\twolineshloka
{समुद्रात्पश्चिमात्पूर्वं दक्षिणादपि चोत्तरम्}
{क्रामत्यनुदिते सूर्ये वाली व्यपगतक्लमः} %||4-11-4||

\twolineshloka
{अग्राण्यारुह्य शैलानां शिखराणि महान्त्यपि}
{ऊर्ध्वमुत्क्षिप्य तरसा प्रतिगृह्णाति वीर्यवान्} %||4-11-5||

\twolineshloka
{बहवः सारवन्तश्च वनेषु विविधा द्रुमाः}
{वालिना तरसा भग्ना बलं प्रथयतात्मनः} %||4-11-6||

\twolineshloka
{महिषो दुन्दुभिर्नाम कैलासशिखरप्रभः}
{बलं नागसहस्रस्य धारयामास वीर्यवान्} %||4-11-7||

\twolineshloka
{वीर्योत्सेकेन दुष्टात्मा वरदानाच्च मोहितः}
{जगाम स महाकायः समुद्रं सरितां पतिम्} %||4-11-8||

\twolineshloka
{ऊर्मिमन्तमतिक्रम्य सागरं रत्नसञ्चयम्}
{मम युद्धं प्रयच्छेति तमुवाच महार्णवम्} %||4-11-9||

\twolineshloka
{ततः समुद्रो धर्मात्मा समुत्थाय महाबलः}
{अब्रवीद्वचनं राजन्नसुरं कालचोदितम्} %||4-11-10||

\twolineshloka
{समर्थो नास्मि ते दातुं युद्धं युद्धविशारद}
{श्रूयतामभिधास्यामि यस्ते युद्धं प्रदास्यति} %||4-11-11||

\twolineshloka
{शैलराजो महारण्ये तपस्विशरणं परम्}
{शङ्करश्वशुरो नाम्ना हिमवानिति विश्रुतः} %||4-11-12||

\twolineshloka
{गुहा प्रस्रवणोपेतो बहुकन्दरनिर्झरः}
{स समर्थस्तव प्रीतिमतुलां कर्तुमाहवे} %||4-11-13||

\twolineshloka
{तं भीतमिति विज्ञाय समुद्रमसुरोत्तमः}
{हिमवद्वनमागच्छच्छरश्चापादिव च्युतः} %||4-11-14||

\twolineshloka
{ततस्तस्य गिरेः श्वेता गजेन्द्रविपुलाः शिलाः}
{चिक्षेप बहुधा भूमौ दुन्दुभिर्विननाद च} %||4-11-15||

\twolineshloka
{ततः श्वेताम्बुदाकारः सौम्यः प्रीतिकराकृतिः}
{हिमवानब्रवीद्वाक्यं स्व एव शिखरे स्थितः} %||4-11-16||

\twolineshloka
{क्लेष्टुमर्हसि मां न त्वं दुन्दुभे धर्मवत्सल}
{रणकर्मस्वकुशलस्तपस्विशरणं ह्यहम्} %||4-11-17||

\twolineshloka
{तस्य तद्वचनं श्रुत्वा गिरिराजस्य धीमतः}
{उवाच दुन्दुभिर्वाक्यं क्रोधात्संरक्तलोचनः} %||4-11-18||

\twolineshloka
{यदि युद्धेऽसमर्थस्त्वं मद्भयाद्वा निरुद्यमः}
{तमाचक्ष्व प्रदद्यान्मे योऽद्य युद्धं युयुत्सतः} %||4-11-19||

\twolineshloka
{हिमवानब्रवीद्वाक्यं श्रुत्वा वाक्यविशारदः}
{अनुक्तपूर्वं धर्मात्मा क्रोधात्तमसुरोत्तमम्} %||4-11-20||

\twolineshloka
{वाली नाम महाप्राज्ञः शक्रतुल्यपराक्रमः}
{अध्यास्ते वानरः श्रीमान्किष्किन्धामतुलप्रभाम्} %||4-11-21||

\twolineshloka
{स समर्थो महाप्राज्ञस्तव युद्धविशारदः}
{द्वन्द्वयुद्धं महद्दातुं नमुचेरिव वासवः} %||4-11-22||

\twolineshloka
{तं शीघ्रमभिगच्छ त्वं यदि युद्धमिहेच्छसि}
{स हि दुर्धर्षणो नित्यं शूरः समरकर्मणि} %||4-11-23||

\twolineshloka
{श्रुत्वा हिमवतो वाक्यं क्रोधाविष्टः स दुन्दुभिः}
{जगाम तां पुरीं तस्य किष्किन्धां वालिनस्तदा} %||4-11-24||

\twolineshloka
{धारयन्माहिषं रूपं तीक्ष्णशृङ्गो भयावहः}
{प्रावृषीव महामेघस्तोयपूर्णो नभस्तले} %||4-11-25||

\twolineshloka
{ततस्तु द्वारमागम्य किष्किन्धाया महाबलः}
{ननर्द कम्पयन्भूमिं दुन्दुभिर्दुन्दुभिर्यथा} %||4-11-26||

\twolineshloka
{समीपजान्द्रुमान्भञ्जन्वसुधां दारयन्खुरैः}
{विषाणेनोल्लेखन्दर्पात्तद्द्वारं द्विरदो यथा} %||4-11-27||

\twolineshloka
{अन्तःपुरगतो वाली श्रुत्वा शब्दममर्षणः}
{निष्पपात सह स्त्रीभिस्ताराभिरिव चन्द्रमाः} %||4-11-28||

\twolineshloka
{मितं व्यक्ताक्षरपदं तमुवाच स दुन्दुभिम्}
{हरीणामीश्वरो वाली सर्वेषां वनचारिणाम्} %||4-11-29||

\twolineshloka
{किमर्थं नगरद्वारमिदं रुद्ध्वा विनर्दसि}
{दुन्दुभे विदितो मेऽसि रक्ष प्राणान्महाबल} %||4-11-30||

\twolineshloka
{तस्य तद्वचनं श्रुत्वा वानरेन्द्रस्य धीमतः}
{उवाच दुन्दुभिर्वाक्यं क्रोधात्संरक्तलोचनः} %||4-11-31||

\twolineshloka
{न त्वं स्त्रीसंनिधौ वीर वचनं वक्तुमर्हसि}
{मम युद्धं प्रयच्छ त्वं ततो ज्ञास्यामि ते बलम्} %||4-11-32||

\twolineshloka
{अथ वा धारयिष्यामि क्रोधमद्य निशामिमाम्}
{गृह्यतामुदयः स्वैरं कामभोगेषु वानर} %||4-11-33||

\twolineshloka
{यो हि मत्तं प्रमत्तं वा सुप्तं वा रहितं भृशम्}
{हन्यात्स भ्रूणहा लोके त्वद्विधं मदमोहितम्} %||4-11-34||

\twolineshloka
{स प्रहस्याब्रवीन्मन्दं क्रोधात्तमसुरोत्तमम्}
{विसृज्य ताः स्त्रियः सर्वास्ताराप्रभृतिकास्तदा} %||4-11-35||

\twolineshloka
{मत्तोऽयमिति मा मंस्था यद्यभीतोऽसि संयुगे}
{मदोऽयं सम्प्रहारेऽस्मिन्वीरपानं समर्थ्यताम्} %||4-11-36||

\twolineshloka
{तमेवमुक्त्वा सङ्क्रुद्धो मालामुत्क्षिप्य काञ्चनीम्}
{पित्रा दत्तां महेन्द्रेण युद्धाय व्यवतिष्ठत} %||4-11-37||

\twolineshloka
{विषाणयोर्गृहीत्वा तं दुन्दुभिं गिरिसंनिभम्}
{वाली व्यापातयां चक्रे ननर्द च महास्वनम्} %||4-11-38||

\threelineshloka
{युद्धे प्राणहरे तस्मिन्निष्पिष्टो दुन्दुभिस्तदा}
{श्रोत्राभ्यामथ रक्तं तु तस्य सुस्राव पात्यतः}
{पपात च महाकायः क्षितौ पञ्चत्वमागतः} %||4-11-39||

\twolineshloka
{तं तोलयित्वा बाहुभ्यां गतसत्त्वमचेतनम्}
{चिक्षेप वेगवान्वाली वेगेनैकेन योजनम्} %||4-11-40||

\twolineshloka
{तस्य वेगप्रविद्धस्य वक्त्रात्क्षतजबिन्दवः}
{प्रपेतुर्मारुतोत्क्षिप्ता मतङ्गस्याश्रमं प्रति} %||4-11-41||

\threelineshloka
{तान्दृष्ट्वा पतितांस्तत्र मुनिः शोणितविप्रुषः}
{उत्ससर्ज महाशापं क्षेप्तारं वालिनं प्रति}
{इह तेनाप्रवेष्टव्यं प्रविष्टस्य बधो भवेत्} %||4-11-42||

{स महर्षिं समासाद्य याचते स्म कृताञ्जलिः॥\devanumber{43}॥} %||4-11-43|| (Check)

\addtocounter{shlokacount}{1}

\twolineshloka
{ततः शापभयाद्भीत ऋश्यमूकं महागिरिम्}
{प्रवेष्टुं नेच्छति हरिर्द्रष्टुं वापि नरेश्वर} %||4-11-44||

\twolineshloka
{तस्याप्रवेशं ज्ञात्वाहमिदं राम महावनम्}
{विचरामि सहामात्यो विषादेन विवर्जितः} %||4-11-45||

\twolineshloka
{एषोऽस्थिनिचयस्तस्य दुन्दुभेः सम्प्रकाशते}
{वीर्योत्सेकान्निरस्तस्य गिरिकूटनिभो महान्} %||4-11-46||

\twolineshloka
{इमे च विपुलाः सालाः सप्त शाखावलम्बिनः}
{यत्रैकं घटते वाली निष्पत्रयितुमोजसा} %||4-11-47||

\twolineshloka
{एतदस्यासमं वीर्यं मया राम प्रकाशितम्}
{कथं तं वालिनं हन्तुं समरे शक्ष्यसे नृप} %||4-11-48||

\twolineshloka
{यदि भिन्द्याद्भवान्सालानिमांस्त्वेकेषुणा ततः}
{जानीयां त्वां महाबाहो समर्थं वालिनो वधे} %||4-11-49||

\threelineshloka
{तस्य तद्वचनं श्रुत्वा सुग्रीवस्य महात्मनः}
{राघवो दुन्दुभेः कायं पादाङ्गुष्ठेन लीलया}
{तोलयित्वा महाबाहुश्चिक्षेप दशयोजनम्} %||4-11-50||

\twolineshloka
{क्षिप्तं दृष्ट्वा ततः कायं सुग्रीवः पुनरब्रवीत्}
{लक्ष्मणस्याग्रतो राममिदं वचनमर्थवत्} %||4-11-51||

\threelineshloka
{आर्द्रः समांसप्रत्यग्रः क्षिप्तः कायः पुरा सखे}
{लघुः सम्प्रति निर्मांसस्तृणभूतश्च राघव}
{नात्र शक्यं बलं ज्ञातुं तव वा तस्य वाधिकम्} %||4-12-52||


{॥इत्यार्षे श्रीमद्रामायणे वाल्मीकीये आदिकाव्ये किष्किन्धाकाण्डे एकादशः सर्गः॥}

\sect{द्वादशः सर्गः}

\twolineshloka
{एतच्च वचनं श्रुत्वा सुग्रीवेण सुभाषितम्}
{प्रत्ययार्थं महातेजा रामो जग्राह कार्मुकम्} %||4-12-1||

\twolineshloka
{स गृहीत्वा धनुर्घोरं शरमेकं च मानदः}
{सालानुद्दिश्य चिक्षेप ज्यास्वनैः पूरयन्दिशः} %||4-12-2||

\twolineshloka
{स विसृष्टो बलवता बाणः स्वर्णपरिष्कृतः}
{भित्त्वा सालान्गिरिप्रस्थे सप्त भूमिं विवेश ह} %||4-12-3||

\twolineshloka
{प्रविष्टस्तु मुहूर्तेन रसां भित्त्वा महाजवः}
{निष्पत्य च पुनस्तूर्णं स्वतूणीं प्रविवेश ह} %||4-12-4||

\twolineshloka
{तान्दृष्ट्वा सप्त निर्भिन्नान्सालान्वानरपुङ्गवः}
{रामस्य शरवेगेन विस्मयं परमं गतः} %||4-12-5||

\twolineshloka
{स मूर्ध्ना न्यपतद्भूमौ प्रलम्बीकृतभूषणः}
{सुग्रीवः परमप्रीतो राघवाय कृताञ्जलिः} %||4-12-6||

\twolineshloka
{इदं चोवाच धर्मज्ञं कर्मणा तेन हर्षितः}
{रामं सर्वास्त्रविदुषां श्रेष्ठं शूरमवस्थितम्} %||4-12-7||

\twolineshloka
{सेन्द्रानपि सुरान्सर्वांस्त्वं बाणैः पुरुषर्षभ}
{समर्थः समरे हन्तुं किं पुनर्वालिनं प्रभो} %||4-12-8||

\twolineshloka
{येन सप्त महासाला गिरिर्भूमिश्च दारिताः}
{बाणेनैकेन काकुत्स्थ स्थाता ते को रणाग्रतः} %||4-12-9||

\twolineshloka
{अद्य मे विगतः शोकः प्रीतिरद्य परा मम}
{सुहृदं त्वां समासाद्य महेन्द्रवरुणोपमम्} %||4-12-10||

\twolineshloka
{तमद्यैव प्रियार्थं मे वैरिणं भ्रातृरूपिणम्}
{वालिनं जहि काकुत्स्थ मया बद्धोऽयमञ्जलिः} %||4-12-11||

\twolineshloka
{ततो रामः परिष्वज्य सुग्रीवं प्रियदर्शनम्}
{प्रत्युवाच महाप्राज्ञो लक्ष्मणानुमतं वचः} %||4-12-12||

\twolineshloka
{अस्माद्गच्छाम किष्किन्धां क्षिप्रं गच्छ त्वमग्रतः}
{गत्वा चाह्वय सुग्रीव वालिनं भ्रातृगन्धिनम्} %||4-12-13||

\twolineshloka
{सर्वे ते त्वरितं गत्वा किष्किन्धां वालिनः पुरीम्}
{वृक्षैरात्मानमावृत्य व्यतिष्ठन्गहने वने} %||4-12-14||

\twolineshloka
{सुग्रीवो व्यनदद्घोरं वालिनो ह्वानकारणात्}
{गाढं परिहितो वेगान्नादैर्भिन्दन्निवाम्बरम्} %||4-12-15||

\twolineshloka
{तं श्रुत्वा निनदं भ्रातुः क्रुद्धो वाली महाबलः}
{निष्पपात सुसंरब्धो भास्करोऽस्ततटादिव} %||4-12-16||

\twolineshloka
{ततः सुतुमुलं युद्धं वालिसुग्रीवयोरभूत्}
{गगने ग्रहयोर्घोरं बुधाङ्गारकयोरिव} %||4-12-17||

\twolineshloka
{तलैरशनिकल्पैश्च वज्रकल्पैश्च मुष्टिभिः}
{जघ्नतुः समरेऽन्योन्यं भ्रातरौ क्रोधमूर्छितौ} %||4-12-18||

\twolineshloka
{ततो रामो धनुष्पाणिस्तावुभौ समुदीक्ष्य तु}
{अन्योन्यसदृशौ वीरावुभौ देवाविवाश्विनौ} %||4-12-19||

\twolineshloka
{यन्नावगच्छत्सुग्रीवं वालिनं वापि राघवः}
{ततो न कृतवान्बुद्धिं मोक्तुमन्तकरं शरम्} %||4-12-20||

\twolineshloka
{एतस्मिन्नन्तरे भग्नः सुग्रीवस्तेन वालिना}
{अपश्यन्राघवं नाथमृश्यमूकं प्रदुद्रुवे} %||4-12-21||

\twolineshloka
{क्लान्तो रुधिरसिक्ताङ्गः प्रहारैर्जर्जरीकृतः}
{वालिनाभिद्रुतः क्रोधात्प्रविवेश महावनम्} %||4-12-22||

\twolineshloka
{तं प्रविष्टं वनं दृष्ट्वा वाली शापभयात्ततः}
{मुक्तो ह्यसि त्वमित्युक्त्वा स निवृत्तो महाबलः} %||4-12-23||

\twolineshloka
{राघवोऽपि सह भ्रात्रा सह चैव हनूमता}
{तदेव वनमागच्छत्सुग्रीवो यत्र वानरः} %||4-12-24||

\twolineshloka
{तं समीक्ष्यागतं रामं सुग्रीवः सहलक्ष्मणम्}
{ह्रीमान्दीनमुवाचेदं वसुधामवलोकयन्} %||4-12-25||

\twolineshloka
{आह्वयस्वेति मामुक्त्वा दर्शयित्वा च विक्रमम्}
{वैरिणा घातयित्वा च किमिदानीं त्वया कृतम्} %||4-12-26||

\twolineshloka
{तामेव वेलां वक्तव्यं त्वया राघव तत्त्वतः}
{वालिनं न निहन्मीति ततो नाहमितो व्रजे} %||4-12-27||

\twolineshloka
{तस्य चैवं ब्रुवाणस्य सुग्रीवस्य महात्मनः}
{करुणं दीनया वाचा राघवः पुनरब्रवीत्} %||4-12-28||

\twolineshloka
{सुग्रीव श्रूयतां तातः क्रोधश्च व्यपनीयताम्}
{कारणं येन बाणोऽयं न मया स विसर्जितः} %||4-12-29||

\twolineshloka
{अलङ्कारेण वेषेण प्रमाणेन गतेन च}
{त्वं च सुग्रीव वाली च सदृशौ स्थः परस्परम्} %||4-12-30||

\twolineshloka
{स्वरेण वर्चसा चैव प्रेक्षितेन च वानर}
{विक्रमेण च वाक्यैश्च व्यक्तिं वां नोपलक्षये} %||4-12-31||

\twolineshloka
{ततोऽहं रूपसादृश्यान्मोहितो वानरोत्तम}
{नोत्सृजामि महावेगं शरं शत्रुनिबर्हणम्} %||4-12-32||

\twolineshloka
{एतन्मुहूर्ते तु मया पश्य वालिनमाहवे}
{निरस्तमिषुणैकेन वेष्टमानं महीतले} %||4-12-33||

\twolineshloka
{अभिज्ञानं कुरुष्व त्वमात्मनो वानरेश्वर}
{येन त्वामभिजानीयां द्वन्द्वयुद्धमुपागतम्} %||4-12-34||

\twolineshloka
{गजपुष्पीमिमां फुल्लामुत्पाट्य शुभलक्षणाम्}
{कुरु लक्ष्मण कण्ठेऽस्य सुग्रीवस्य महात्मनः} %||4-12-35||

\twolineshloka
{ततो गिरितटे जातामुत्पाट्य कुसुमायुताम्}
{लक्ष्मणो गजपुष्पीं तां तस्य कण्ठे व्यसर्जयत्} %||4-12-36||

\twolineshloka
{स तथा शुशुभे श्रीमाँल्लतया कण्ठसक्तया}
{मालयेव बलाकानां ससन्ध्य इव तोयदः} %||4-12-37||

\twolineshloka
{विभ्राजमानो वपुषा रामवाक्यसमाहितः}
{जगाम सह रामेण किष्किन्धां वालिपालिताम्} %||4-13-38||


{॥इत्यार्षे श्रीमद्रामायणे वाल्मीकीये आदिकाव्ये किष्किन्धाकाण्डे द्वादशः सर्गः॥}

\sect{त्रयोदशः सर्गः}

\twolineshloka
{ऋश्यमूकात्स धर्मात्मा किष्किन्धां लक्ष्मणाग्रजः}
{जगाम सहसुग्रीवो वालिविक्रमपालिताम्} %||4-13-1||

\twolineshloka
{समुद्यम्य महच्चापं रामः काञ्चनभूषितम्}
{शरांश्चादित्य सङ्काशान्गृहीत्वा रणसाधकान्} %||4-13-2||

\twolineshloka
{अग्रतस्तु ययौ तस्य राघवस्य महात्मनः}
{सुग्रीवः संहतग्रीवो लक्ष्मणश्च महाबलः} %||4-13-3||

\twolineshloka
{पृष्ठतो हनुमान्वीरो नलो नीलश्च वानरः}
{तारश्चैव महातेजा हरियूथप यूथपाः} %||4-13-4||

\twolineshloka
{ते वीक्षमाणा वृक्षांश्च पुष्पभारावलम्बिनः}
{प्रसन्नाम्बुवहाश्चैव सरितः सागरं गमाः} %||4-13-5||

\twolineshloka
{कन्दराणि च शैलांश्च निर्झराणि गुहास्तथा}
{शिखराणि च मुख्यानि दरीश्च प्रियदर्शनाः} %||4-13-6||

\twolineshloka
{वैदूर्यविमलैः पर्णैः पद्मैश्चाकाशकुड्मलैः}
{शोभितान्सजलान्मार्गे तटाकांश्च व्यलोकयन्} %||4-13-7||

\twolineshloka
{कारण्डैः सारसैर्हंसैर्वञ्जूलैर्जलकुक्कुटैः}
{चक्रवाकैस्तथा चान्यैः शकुनैः प्रतिनादितान्} %||4-13-8||

\twolineshloka
{मृदुशष्पाङ्कुराहारान्निर्भयान्वनगोचरान्}
{चरतः सर्वतोऽपश्यन्स्थलीषु हरिणान्स्थितान्} %||4-13-9||

\twolineshloka
{तटाकवैरिणश्चापि शुक्लदन्तविभूषितान्}
{घोरानेकचरान्वन्यान्द्विरदान्कूलघातिनः} %||4-13-10||

\twolineshloka
{वने वनचरांश्चान्यान्खेचरांश्च विहङ्गमान्}
{पश्यन्तस्त्वरिता जग्मुः सुग्रीववशवर्तिनः} %||4-13-11||

\twolineshloka
{तेषां तु गच्छतां तत्र त्वरितं रघुनन्दनः}
{द्रुमषण्डं वनं दृष्ट्वा रामः सुग्रीवमब्रवीत्} %||4-13-12||

\twolineshloka
{एष मेघ इवाकाशे वृक्षषण्डः प्रकाशते}
{मेघसङ्घातविपुलः पर्यन्तकदलीवृतः} %||4-13-13||

\twolineshloka
{किमेतज्ज्ञातुमिच्छामि सखे कौतूहलं मम}
{कौतूहलापनयनं कर्तुमिच्छाम्यहं त्वया} %||4-13-14||

\twolineshloka
{तस्य तद्वचनं श्रुत्वा राघवस्य महात्मनः}
{गच्छन्नेवाचचक्षेऽथ सुग्रीवस्तन्महद्वनम्} %||4-13-15||

\twolineshloka
{एतद्राघव विस्तीर्णमाश्रमं श्रमनाशनम्}
{उद्यानवनसम्पन्नं स्वादुमूलफलोदकम्} %||4-13-16||

\twolineshloka
{अत्र सप्तजना नाम मुनयः संशितव्रताः}
{सप्तैवासन्नधःशीर्षा नियतं जलशायिनः} %||4-13-17||

\twolineshloka
{सप्तरात्रकृताहारा वायुना वनवासिनः}
{दिवं वर्षशतैर्याताः सप्तभिः सकलेवराः} %||4-13-18||

\twolineshloka
{तेषामेवं प्रभावेन द्रुमप्राकारसंवृतम्}
{आश्रमं सुदुराधर्षमपि सेन्द्रैः सुरासुरैः} %||4-13-19||

\twolineshloka
{पक्षिणो वर्जयन्त्येतत्तथान्ये वनचारिणः}
{विशन्ति मोहाद्येऽप्यत्र निवर्तन्ते न ते पुनः} %||4-13-20||

\twolineshloka
{विभूषणरवाश्चात्र श्रूयन्ते सकलाक्षराः}
{तूर्यगीतस्वनाश्चापि गन्धो दिव्यश्च राघव} %||4-13-21||

\twolineshloka
{त्रेताग्नयोऽपि दीप्यन्ते धूमो ह्येष प्रदृश्यते}
{वेष्टयन्निव वृक्षाग्रान्कपोताङ्गारुणो घनः} %||4-13-22||

\twolineshloka
{कुरु प्रणामं धर्मात्मंस्तान्समुद्दिश्य राघवः}
{लक्ष्मणेन सह भ्रात्रा प्रयतः संयताञ्जलिः} %||4-13-23||

\twolineshloka
{प्रणमन्ति हि ये तेषामृषीणां भावितात्मनाम्}
{न तेषामशुभं किञ्चिच्छरीरे राम दृश्यते} %||4-13-24||

\twolineshloka
{ततो रामः सह भ्रात्रा लक्ष्मणेन कृताञ्जलिः}
{समुद्दिश्य महात्मानस्तानृषीनभ्यवादयत्} %||4-13-25||

\twolineshloka
{अभिवाद्य च धर्मात्मा रामो भ्राता च लक्ष्मणः}
{सुग्रीवो वानराश्चैव जग्मुः संहृष्टमानसाः} %||4-13-26||

\twolineshloka
{ते गत्वा दूरमध्वानं तस्मात्सप्तजनाश्रमात्}
{ददृशुस्तां दुराधर्षां किष्किन्धां वालिपालिताम्} %||4-14-27||


{॥इत्यार्षे श्रीमद्रामायणे वाल्मीकीये आदिकाव्ये किष्किन्धाकाण्डे त्रयोदशः सर्गः॥}

\sect{चतुर्दशः सर्गः}

\twolineshloka
{सर्वे ते त्वरितं गत्वा किष्किन्धां वालिपालिताम्}
{वृक्षैरात्मानमावृत्य व्यतिष्ठन्गहने वने} %||4-14-1||

\twolineshloka
{विचार्य सर्वतो दृष्टिं कानने काननप्रियः}
{सुग्रीवो विपुलग्रीवः क्रोधमाहारयद्भृशम्} %||4-14-2||

\twolineshloka
{ततः स निनदं घोरं कृत्वा युद्धाय चाह्वयत्}
{परिवारैः परिवृतो नादैर्भिन्दन्निवाम्बरम्} %||4-14-3||

\twolineshloka
{अथ बालार्कसदृशो दृप्तसिंहगतिस्तदा}
{दृष्ट्वा रामं क्रियादक्षं सुग्रीवो वाक्यमब्रवीत्} %||4-14-4||

\twolineshloka
{हरिवागुरया व्याप्तं तप्तकाञ्चनतोरणाम्}
{प्राप्ताः स्म ध्वजयन्त्राढ्यां किष्किन्धां वालिनः पुरीम्} %||4-14-5||

\twolineshloka
{प्रतिज्ञा या त्वया वीर कृता वालिवधे पुरा}
{सफलां तां कुरु क्षिप्रं लतां काल इवागतः} %||4-14-6||

\twolineshloka
{एवमुक्तस्तु धर्मात्मा सुग्रीवेण स राघवः}
{तमथोवाच सुग्रीवं वचनं शत्रुसूदनः} %||4-14-7||

\twolineshloka
{कृताभिज्ञान चिह्नस्त्वमनया गजसाह्वया}
{विपरीत इवाकाशे सूर्यो नक्षत्र मालया} %||4-14-8||

\twolineshloka
{अद्य वालिसमुत्थं ते भयं वैरं च वानर}
{एकेनाहं प्रमोक्ष्यामि बाणमोक्षेण संयुगे} %||4-14-9||

\twolineshloka
{मम दर्शय सुग्रीववैरिणं भ्रातृरूपिणम्}
{वाली विनिहतो यावद्वने पांसुषु वेष्टते} %||4-14-10||

\twolineshloka
{यदि दृष्टिपथं प्राप्तो जीवन्स विनिवर्तते}
{ततो दोषेण मा गच्छेत्सद्यो गर्हेच्च मा भवान्} %||4-14-11||

\twolineshloka
{प्रत्यक्षं सप्त ते साला मया बाणेन दारिताः}
{ततो वेत्सि बलेनाद्य बालिनं निहतं मया} %||4-14-12||

\twolineshloka
{अनृतं नोक्तपूर्वं मे वीर कृच्छ्रेऽपि तिष्ठता}
{धर्मलोभपरीतेन न च वक्ष्ये कथञ्चन} %||4-14-13||

\twolineshloka
{सफलां च करिष्यामि प्रतिज्ञां जहि सम्भ्रमम्}
{प्रसूतं कलमं क्षेत्रे वर्षेणेव शतक्रतुः} %||4-14-14||

\twolineshloka
{तदाह्वाननिमित्तं त्वं वालिनो हेममालिनः}
{सुग्रीव कुरु तं शब्दं निष्पतेद्येन वानरः} %||4-14-15||

\twolineshloka
{जितकाशी जयश्लाघी त्वया चाधर्षितः पुरात्}
{निष्पतिष्यत्यसङ्गेन वाली स प्रियसंयुगः} %||4-14-16||

\twolineshloka
{रिपूणां धर्षणं शूरा मर्षयन्ति न संयुगे}
{जानन्तस्तु स्वकं वीर्यं स्त्रीसमक्षं विशेषतः} %||4-14-17||

\twolineshloka
{स तु रामवचः श्रुत्वा सुग्रीवो हेमपिङ्गलः}
{ननर्द क्रूरनादेन विनिर्भिन्दन्निवाम्बरम्} %||4-14-18||

\twolineshloka
{तस्य शब्देन वित्रस्ता गावो यान्ति हतप्रभाः}
{राजदोषपरामृष्टाः कुलस्त्रिय इवाकुलाः} %||4-14-19||

\twolineshloka
{द्रवन्ति च मृगाः शीघ्रं भग्ना इव रणे हयाः}
{पतन्ति च खगा भूमौ क्षीणपुण्या इव ग्रहाः} %||4-14-20||

\fourlineindentedshloka
{ततः स जीमूतगणप्रणादो}
{ नादं व्यमुञ्चत्त्वरया प्रतीतः}
{सूर्यात्मजः शौर्यविवृद्धतेजाः}
{ सरित्पतिर्वानिलचञ्चलोर्मिः} %||4-15-21||


{॥इत्यार्षे श्रीमद्रामायणे वाल्मीकीये आदिकाव्ये किष्किन्धाकाण्डे चतुर्दशः सर्गः॥}

\sect{पञ्चदशः सर्गः}

\twolineshloka
{अथ तस्य निनादं तं सुग्रीवस्य महात्मनः}
{शुश्रावान्तःपुरगतो वाली भ्रातुरमर्षणः} %||4-15-1||

\twolineshloka
{श्रुत्वा तु तस्य निनदं सर्वभूतप्रकम्पनम्}
{मदश्चैकपदे नष्टः क्रोधश्चापतितो महान्} %||4-15-2||

\twolineshloka
{स तु रोषपरीताङ्गो वाली सन्ध्यातपप्रभः}
{उपरक्त इवादित्यः सद्यो निष्प्रभतां गतः} %||4-15-3||

\twolineshloka
{वाली दंष्ट्रा करालस्तु क्रोधाद्दीप्ताग्निसंनिभः}
{भात्युत्पतितपद्माभः समृणाल इव ह्रदः} %||4-15-4||

\twolineshloka
{शब्दं दुर्मर्षणं श्रुत्वा निष्पपात ततो हरिः}
{वेगेन चरणन्यासैर्दारयन्निव मेदिनीम्} %||4-15-5||

\twolineshloka
{तं तु तारा परिष्वज्य स्नेहाद्दर्शितसौहृदा}
{उवाच त्रस्तसम्भ्रान्ता हितोदर्कमिदं वचः} %||4-15-6||

\twolineshloka
{साधु क्रोधमिमं वीर नदी वेगमिवागतम्}
{शयनादुत्थितः काल्यं त्यज भुक्तामिव स्रजम्} %||4-15-7||

\twolineshloka
{सहसा तव निष्क्रामो मम तावन्न रोचते}
{श्रूयतामभिधास्यामि यन्निमित्तं निवार्यसे} %||4-15-8||

\twolineshloka
{पूर्वमापतितः क्रोधात्स त्वामाह्वयते युधि}
{निष्पत्य च निरस्तस्ते हन्यमानो दिशो गतः} %||4-15-9||

\twolineshloka
{त्वया तस्य निरस्तस्य पीडितस्य विशेषतः}
{इहैत्य पुनराह्वानं शङ्कां जनयतीव मे} %||4-15-10||

\twolineshloka
{दर्पश्च व्यवसायश्च यादृशस्तस्य नर्दतः}
{निनादस्य च संरम्भो नैतदल्पं हि कारणम्} %||4-15-11||

\twolineshloka
{नासहायमहं मन्ये सुग्रीवं तमिहागतम्}
{अवष्टब्धसहायश्च यमाश्रित्यैष गर्जति} %||4-15-12||

\twolineshloka
{प्रकृत्या निपुणश्चैव बुद्धिमांश्चैव वानरः}
{अपरीक्षितवीर्येण सुग्रीवः सह नैष्यति} %||4-15-13||

\twolineshloka
{पूर्वमेव मया वीर श्रुतं कथयतो वचः}
{अङ्गदस्य कुमारस्य वक्ष्यामि त्वा हितं वचः} %||4-15-14||

\twolineshloka
{तव भ्रातुर्हि विख्यातः सहायो रणकर्कशः}
{रामः परबलामर्दी युगान्ताग्निरिवोत्थितः} %||4-15-15||

\twolineshloka
{निवासवृक्षः साधूनामापन्नानां परा गतिः}
{आर्तानां संश्रयश्चैव यशसश्चैकभाजनम्} %||4-15-16||

\twolineshloka
{ज्ञानविज्ञानसम्पन्नो निदेशो निरतः पितुः}
{धातूनामिव शैलेन्द्रो गुणानामाकरो महान्} %||4-15-17||

\twolineshloka
{तत्क्षमं न विरोधस्ते सह तेन महात्मना}
{दुर्जयेनाप्रमेयेन रामेण रणकर्मसु} %||4-15-18||

\twolineshloka
{शूर वक्ष्यामि ते किञ्चिन्न चेच्छाम्यभ्यसूयितुम्}
{श्रूयतां क्रियतां चैव तव वक्ष्यामि यद्धितम्} %||4-15-19||

\twolineshloka
{यौवराज्येन सुग्रीवं तूर्णं साध्वभिषेचय}
{विग्रहं मा कृथा वीर भ्रात्रा राजन्बलीयसा} %||4-15-20||

\twolineshloka
{अहं हि ते क्षमं मन्ये तव रामेण सौहृदम्}
{सुग्रीवेण च सम्प्रीतिं वैरमुत्सृज्य दूरतः} %||4-15-21||

\twolineshloka
{लालनीयो हि ते भ्राता यवीयानेष वानरः}
{तत्र वा सन्निहस्थो वा सर्वथा बन्धुरेव ते} %||4-15-22||

\twolineshloka
{यदि ते मत्प्रियं कार्यं यदि चावैषि मां हिताम्}
{याच्यमानः प्रयत्नेन साधु वाक्यं कुरुष्व मे} %||4-16-23||


{॥इत्यार्षे श्रीमद्रामायणे वाल्मीकीये आदिकाव्ये किष्किन्धाकाण्डे पञ्चदशः सर्गः॥}

\sect{षोडशः सर्गः}

\twolineshloka
{तामेवं ब्रुवतीं तारां ताराधिपनिभाननाम्}
{वाली निर्भर्त्सयामास वचनं चेदमब्रवीत्} %||4-16-1||

\twolineshloka
{गर्जतोऽस्य च संरम्भं भ्रातुः शत्रोर्विशेषतः}
{मर्षयिष्याम्यहं केन कारणेन वरानने} %||4-16-2||

\twolineshloka
{अधर्षितानां शूराणां समरेष्वनिवर्तिनाम्}
{धर्षणामर्षणं भीरु मरणादतिरिच्यते} %||4-16-3||

\twolineshloka
{सोढुं न च समर्थोऽहं युद्धकामस्य संयुगे}
{सुग्रीवस्य च संरम्भं हीनग्रीवस्य गर्जतः} %||4-16-4||

\twolineshloka
{न च कार्यो विषादस्ते राघवं प्रति मत्कृते}
{धर्मज्ञश्च कृतज्ञश्च कथं पापं करिष्यति} %||4-16-5||

\twolineshloka
{निवर्तस्व सह स्त्रीभिः कथं भूयोऽनुगच्छसि}
{सौहृदं दर्शितं तारे मयि भक्तिः कृता त्वया} %||4-16-6||

\twolineshloka
{प्रतियोत्स्याम्यहं गत्वा सुग्रीवं जहि सम्भ्रमम्}
{दर्पं चास्य विनेष्यामि न च प्राणैर्विमोक्ष्यते} %||4-16-7||

\twolineshloka
{शापितासि मम प्राणैर्निवर्तस्व जयेन च}
{अहं जित्वा निवर्तिष्ये तमलं भ्रातरं रणे} %||4-16-8||

\twolineshloka
{तं तु तारा परिष्वज्य वालिनं प्रियवादिनी}
{चकार रुदती मन्दं दक्षिणा सा प्रदक्षिणम्} %||4-16-9||

\twolineshloka
{ततः स्वस्त्ययनं कृत्वा मन्त्रवद्विजयैषिणी}
{अन्तःपुरं सह स्त्रीभिः प्रविष्टा शोकमोहिता} %||4-16-10||

\twolineshloka
{प्रविष्टायां तु तारायां सह स्त्रीभिः स्वमालयम्}
{नगरान्निर्ययौ क्रुद्धो महासर्प इव श्वसन्} %||4-16-11||

\twolineshloka
{स निःश्वस्य महावेगो वाली परमरोषणः}
{सर्वतश्चारयन्दृष्टिं शत्रुदर्शनकाङ्क्षया} %||4-16-12||

\twolineshloka
{स ददर्श ततः श्रीमान्सुग्रीवं हेमपिङ्गलम्}
{सुसंवीतमवष्टब्धं दीप्यमानमिवानलम्} %||4-16-13||

\twolineshloka
{स तं दृष्ट्वा महावीर्यं सुग्रीवं पर्यवस्थितम्}
{गाढं परिदधे वासो वाली परमरोषणः} %||4-16-14||

\twolineshloka
{स वाली गाढसंवीतो मुष्टिमुद्यम्य वीर्यवान्}
{सुग्रीवमेवाभिमुखो ययौ योद्धुं कृतक्षणः} %||4-16-15||

\twolineshloka
{श्लिष्टमुष्टिं समुद्यम्य संरब्धतरमागतः}
{सुग्रीवोऽपि समुद्दिश्य वालिनं हेममालिनम्} %||4-16-16||

\twolineshloka
{तं वाली क्रोधताम्राक्षः सुग्रीवं रणपण्डितम्}
{आपतन्तं महावेगमिदं वचनमब्रवीत्} %||4-16-17||

\twolineshloka
{एष मुष्टिर्मया बद्धो गाढः सुनिहिताङ्गुलिः}
{मया वेगविमुक्तस्ते प्राणानादाय यास्यति} %||4-16-18||

\twolineshloka
{एवमुक्तस्तु सुग्रीवः क्रुद्धो वालिनमब्रवीत्}
{तवैव च हरन्प्राणान्मुष्टिः पततु मूर्धनि} %||4-16-19||

\twolineshloka
{ताडितस्तेन सङ्क्रुद्धः समभिक्रम्य वेगतः}
{अभवच्छोणितोद्गारी सोत्पीड इव पर्वतः} %||4-16-20||

\twolineshloka
{सुग्रीवेण तु निःसङ्गं सालमुत्पाट्य तेजसा}
{गात्रेष्वभिहतो वाली वज्रेणेव महागिरिः} %||4-16-21||

\twolineshloka
{स तु वाली प्रचरितः सालताडनविह्वलः}
{गुरुभारसमाक्रान्ता सागरे नौरिवाभवत्} %||4-16-22||

\twolineshloka
{तौ भीमबलविक्रान्तौ सुपर्णसमवेगिनौ}
{प्रवृद्धौ घोरवपुषौ चन्द्रसूर्याविवाम्बरे} %||4-16-23||

\twolineshloka
{वालिना भग्नदर्पस्तु सुग्रीवो मन्दविक्रमः}
{वालिनं प्रति सामर्षो दर्शयामास लाघवम्} %||4-16-24||

\twolineshloka
{ततो धनुषि सन्धाय शरमाशीविषोपमम्}
{राघवेण महाबाणो वालिवक्षसि पातितः} %||4-16-25||

{वेगेनाभिहतो वाली निपपात महीतले॥\devanumber{26}॥} %||4-16-26|| (Check)

\addtocounter{shlokacount}{1}

\fourlineindentedshloka
{अथोक्षितः शोणिततोयविस्रवैः}
{ सुपुष्पिताशोक इवानिलोद्धतः}
{विचेतनो वासवसूनुराहवे}
{ प्रभ्रंशितेन्द्रध्वजवत्क्षितिं गतः} %||4-17-27||


{॥इत्यार्षे श्रीमद्रामायणे वाल्मीकीये आदिकाव्ये किष्किन्धाकाण्डे षोडशः सर्गः॥}

\sect{सप्तदशः सर्गः}

\twolineshloka
{ततः शरेणाभिहतो रामेण रणकर्कशः}
{पपात सहसा वाली निकृत्त इव पादपः} %||4-17-1||

\twolineshloka
{स भूमौ न्यस्तसर्वाङ्गस्तप्तकाञ्चनभूषणः}
{अपतद्देवराजस्य मुक्तरश्मिरिव ध्वजः} %||4-17-2||

\twolineshloka
{तस्मिन्निपतिते भूमौ हर्यृषाणां गणेश्वरे}
{नष्टचन्द्रमिव व्योम न व्यराजत भूतलम्} %||4-17-3||

\twolineshloka
{भूमौ निपतितस्यापि तस्य देहं महात्मनः}
{न श्रीर्जहाति न प्राणा न तेजो न पराक्रमः} %||4-17-4||

\twolineshloka
{शक्रदत्ता वरा माला काञ्चनी रत्नभूषिता}
{दधार हरिमुख्यस्य प्राणांस्तेजः श्रियं च सा} %||4-17-5||

\twolineshloka
{स तया मालया वीरो हैमया हरियूथपः}
{सन्ध्यानुगतपर्यन्तः पयोधर इवाभवत्} %||4-17-6||

\twolineshloka
{तस्य माला च देहश्च मर्मघाती च यः शरः}
{त्रिधेव रचिता लक्ष्मीः पतितस्यापि शोभते} %||4-17-7||

\twolineshloka
{तदस्त्रं तस्य वीरस्य स्वर्गमार्गप्रभावनम्}
{रामबाणासनक्षिप्तमावहत्परमां गतिम्} %||4-17-8||

\twolineshloka
{तं तथा पतितं सङ्ख्ये गतार्चिषमिवानलम्}
{ययातिमिव पुण्यान्ते देवलोकात्परिच्युतम्} %||4-17-9||

\twolineshloka
{आदित्यमिव कालेन युगान्ते भुवि पातितम्}
{महेन्द्रमिव दुर्धर्षं महेन्द्रमिव दुःसहम्} %||4-17-10||

\threelineshloka
{महेन्द्रपुत्रं पतितं वालिनं हेममालिनम्}
{सिंहोरस्कं महाबाहुं दीप्तास्यं हरिलोचनम्}
{लक्ष्मणानुगतो रामो ददर्शोपससर्प च} %||4-17-11||

\twolineshloka
{स दृष्ट्वा राघवं वाली लक्ष्मणं च महाबलम्}
{अब्रवीत्प्रश्रितं वाक्यं परुषं धर्मसंहितम्} %||4-17-12||

\twolineshloka
{पराङ्मुखवधं कृत्वा को नु प्राप्तस्त्वया गुणः}
{यदहं युद्धसंरब्धस्त्वत्कृते निधनं गतः} %||4-17-13||

\twolineshloka
{कुलीनः सत्त्वसम्पन्नस्तेजस्वी चरितव्रतः}
{रामः करुणवेदी च प्रजानां च हिते रतः} %||4-17-14||

\twolineshloka
{सानुक्रोशो महोत्साहः समयज्ञो दृढव्रतः}
{इति ते सर्वभूतानि कथयन्ति यशो भुवि} %||4-17-15||

\twolineshloka
{तान्गुणान्सम्प्रधार्याहमग्र्यं चाभिजनं तव}
{तारया प्रतिषिद्धः सन्सुग्रीवेण समागतः} %||4-17-16||

\twolineshloka
{न मामन्येन संरब्धं प्रमत्तं वेद्धुमर्हसि}
{इति मे बुद्धिरुत्पन्ना बभूवादर्शने तव} %||4-17-17||

\twolineshloka
{न त्वां विनिहतात्मानं धर्मध्वजमधार्मिकम्}
{जाने पापसमाचारं तृणैः कूपमिवावृतम्} %||4-17-18||

\twolineshloka
{सतां वेषधरं पापं प्रच्छन्नमिव पावकम्}
{नाहं त्वामभिजानानि धर्मच्छद्माभिसंवृतम्} %||4-17-19||

\twolineshloka
{विषये वा पुरे वा ते यदा नापकरोम्यहम्}
{न च त्वां प्रतिजानेऽहं कस्मात्त्वं हंस्यकिल्बिषम्} %||4-17-20||

\twolineshloka
{फलमूलाशनं नित्यं वानरं वनगोचरम्}
{मामिहाप्रतियुध्यन्तमन्येन च समागतम्} %||4-17-21||

\twolineshloka
{त्वं नराधिपतेः पुत्रः प्रतीतः प्रियदर्शनः}
{लिङ्गमप्यस्ति ते राजन्दृश्यते धर्मसंहितम्} %||4-17-22||

\twolineshloka
{कः क्षत्रियकुले जातः श्रुतवान्नष्टसंशयः}
{धर्मलिङ्ग प्रतिच्छन्नः क्रूरं कर्म समाचरेत्} %||4-17-23||

\twolineshloka
{राम राजकुले जातो धर्मवानिति विश्रुतः}
{अभव्यो भव्यरूपेण किमर्थं परिधावसि} %||4-17-24||

\twolineshloka
{साम दानं क्षमा धर्मः सत्यं धृतिपराक्रमौ}
{पार्थिवानां गुणा राजन्दण्डश्चाप्यपकारिषु} %||4-17-25||

\twolineshloka
{वयं वनचरा राम मृगा मूलफलाशनाः}
{एषा प्रकृतिरस्माकं पुरुषस्त्वं नरेश्वरः} %||4-17-26||

\twolineshloka
{भूमिर्हिरण्यं रूप्यं च निग्रहे कारणानि च}
{तत्र कस्ते वने लोभो मदीयेषु फलेषु वा} %||4-17-27||

\twolineshloka
{नयश्च विनयश्चोभौ निग्रहानुग्रहावपि}
{राजवृत्तिरसङ्कीर्णा न नृपाः कामवृत्तयः} %||4-17-28||

\twolineshloka
{त्वं तु कामप्रधानश्च कोपनश्चानवस्थितः}
{राजवृत्तैश्च सङ्कीर्णः शरासनपरायणः} %||4-17-29||

\twolineshloka
{न तेऽस्त्यपचितिर्धर्मे नार्थे बुद्धिरवस्थिता}
{इन्द्रियैः कामवृत्तः सन्कृष्यसे मनुजेश्वर} %||4-17-30||

\twolineshloka
{हत्वा बाणेन काकुत्स्थ मामिहानपराधिनम्}
{किं वक्ष्यसि सतां मध्ये कर्म कृत्वा जुगुप्सितम्} %||4-17-31||

\twolineshloka
{राजहा ब्रह्महा गोघ्नश्चोरः प्राणिवधे रतः}
{नास्तिकः परिवेत्ता च सर्वे निरयगामिनः} %||4-17-32||

\twolineshloka
{अधार्यं चर्म मे सद्भी रोमाण्यस्थि च वर्जितम्}
{अभक्ष्याणि च मांसानि त्वद्विधैर्धर्मचारिभिः} %||4-17-33||

\twolineshloka
{पञ्च पञ्चनखा भक्ष्या ब्रह्मक्षत्रेण राघव}
{शल्यकः श्वाविधो गोधा शशः कूर्मश्च पञ्चमः} %||4-17-34||

\twolineshloka
{चर्म चास्थि च मे राजन्न स्पृशन्ति मनीषिणः}
{अभक्ष्याणि च मांसानि सोऽहं पञ्चनखो हतः} %||4-17-35||

\twolineshloka
{त्वया नाथेन काकुत्स्थ न सनाथा वसुन्धरा}
{प्रमदा शीलसम्पन्ना धूर्तेन पतिता यथा} %||4-17-36||

\twolineshloka
{शठो नैकृतिकः क्षुद्रो मिथ्या प्रश्रितमानसः}
{कथं दशरथेन त्वं जातः पापो महात्मना} %||4-17-37||

\twolineshloka
{छिन्नचारित्र्यकक्ष्येण सतां धर्मातिवर्तिना}
{त्यक्तधर्माङ्कुशेनाहं निहतो रामहस्तिना} %||4-17-38||

\twolineshloka
{दृश्यमानस्तु युध्येथा मया युधि नृपात्मज}
{अद्य वैवस्वतं देवं पश्येस्त्वं निहतो मया} %||4-17-39||

\twolineshloka
{त्वयादृश्येन तु रणे निहतोऽहं दुरासदः}
{प्रसुप्तः पन्नगेनेव नरः पानवशं गतः} %||4-17-40||

\twolineshloka
{सुग्रीवप्रियकामेन यदहं निहतस्त्वया}
{कण्ठे बद्ध्वा प्रदद्यां तेऽनिहतं रावणं रणे} %||4-17-41||

\twolineshloka
{न्यस्तां सागरतोये वा पाताले वापि मैथिलीम्}
{जानयेयं तवादेशाच्छ्वेतामश्वतरीमिव} %||4-17-42||

\twolineshloka
{युक्तं यत्प्रप्नुयाद्राज्यं सुग्रीवः स्वर्गते मयि}
{अयुक्तं यदधर्मेण त्वयाहं निहतो रणे} %||4-17-43||

\twolineshloka
{काममेवंविधं लोकः कालेन विनियुज्यते}
{क्षमं चेद्भवता प्राप्तमुत्तरं साधु चिन्त्यताम्} %||4-17-44||

\fourlineindentedshloka
{इत्येवमुक्त्वा परिशुष्कवक्त्रः}
{ शराभिघाताद्व्यथितो महात्मा}
{समीक्ष्य रामं रविसंनिकाशम्}
{ तूष्णीं बभूवामरराजसूनुः} %||4-18-45||


{॥इत्यार्षे श्रीमद्रामायणे वाल्मीकीये आदिकाव्ये किष्किन्धाकाण्डे सप्तदशः सर्गः॥}

\sect{अष्टादशः सर्गः}

\twolineshloka
{इत्युक्तः प्रश्रितं वाक्यं धर्मार्थसहितं हितम्}
{परुषं वालिना रामो निहतेन विचेतसा} %||4-18-1||

\twolineshloka
{तं निष्प्रभमिवादित्यं मुक्ततोयमिवाम्बुदम्}
{उक्तवाक्यं हरिश्रेष्ठमुपशान्तमिवानलम्} %||4-18-2||

\twolineshloka
{धर्मार्थगुणसम्पन्नं हरीश्वरमनुत्तमम्}
{अधिक्षिप्तस्तदा रामः पश्चाद्वालिनमब्रवीत्} %||4-18-3||

\twolineshloka
{धर्ममर्थं च कामं च समयं चापि लौकिकम्}
{अविज्ञाय कथं बाल्यान्मामिहाद्य विगर्हसे} %||4-18-4||

\twolineshloka
{अपृष्ट्वा बुद्धिसम्पन्नान्वृद्धानाचार्यसम्मतान्}
{सौम्य वानरचापल्यात्त्वं मां वक्तुमिहेच्छसि} %||4-18-5||

\twolineshloka
{इक्ष्वाकूणामियं भूमिः सशैलवनकानना}
{मृगपक्षिमनुष्याणां निग्रहानुग्रहावपि} %||4-18-6||

\twolineshloka
{तां पालयति धर्मात्मा भरतः सत्यवागृजुः}
{धर्मकामार्थतत्त्वज्ञो निग्रहानुग्रहे रतः} %||4-18-7||

\twolineshloka
{नयश्च विनयश्चोभौ यस्मिन्सत्यं च सुस्थितम्}
{विक्रमश्च यथा दृष्टः स राजा देशकालवित्} %||4-18-8||

\twolineshloka
{तस्य धर्मकृतादेशा वयमन्ये च पार्थिवः}
{चरामो वसुधां कृत्स्नां धर्मसन्तानमिच्छवः} %||4-18-9||

\twolineshloka
{तस्मिन्नृपतिशार्दूल भरते धर्मवत्सले}
{पालयत्यखिलां भूमिं कश्चरेद्धर्मनिग्रहम्} %||4-18-10||

\twolineshloka
{ते वयं मार्गविभ्रष्टं स्वधर्मे परमे स्थिताः}
{भरताज्ञां पुरस्कृत्य निगृह्णीमो यथाविधि} %||4-18-11||

\twolineshloka
{त्वं तु सङ्क्लिष्टधर्मा च कर्मणा च विगर्हितः}
{कामतन्त्रप्रधानश्च न स्थितो राजवर्त्मनि} %||4-18-12||

\twolineshloka
{ज्येष्ठो भ्राता पिता चैव यश्च विद्यां प्रयच्छति}
{त्रयस्ते पितरो ज्ञेया धर्मे च पथि वर्तिनः} %||4-18-13||

\twolineshloka
{यवीयानात्मनः पुत्रः शिष्यश्चापि गुणोदितः}
{पुत्रवत्ते त्रयश्चिन्त्या धर्मश्चेदत्र कारणम्} %||4-18-14||

\twolineshloka
{सूक्ष्मः परमदुर्ज्ञेयः सतां धर्मः प्लवङ्गम}
{हृदिस्थः सर्वभूतानामात्मा वेद शुभाशुभम्} %||4-18-15||

\twolineshloka
{चपलश्चपलैः सार्धं वानरैरकृतात्मभिः}
{जात्यन्ध इव जात्यन्धैर्मन्त्रयन्द्रक्ष्यसे नु किम्} %||4-18-16||

\twolineshloka
{अहं तु व्यक्ततामस्य वचनस्य ब्रवीमि ते}
{न हि मां केवलं रोषात्त्वं विगर्हितुमर्हसि} %||4-18-17||

\twolineshloka
{तदेतत्कारणं पश्य यदर्थं त्वं मया हतः}
{भ्रातुर्वर्तसि भार्यायां त्यक्त्वा धर्मं सनातनम्} %||4-18-18||

\twolineshloka
{अस्य त्वं धरमाणस्य सुग्रीवस्य महात्मनः}
{रुमायां वर्तसे कामात्स्नुषायां पापकर्मकृत्} %||4-18-19||

\twolineshloka
{तद्व्यतीतस्य ते धर्मात्कामवृत्तस्य वानर}
{भ्रातृभार्याभिमर्शेऽस्मिन्दण्डोऽयं प्रतिपादितः} %||4-18-20||

\twolineshloka
{न हि धर्मविरुद्धस्य लोकवृत्तादपेयुषः}
{दण्डादन्यत्र पश्यामि निग्रहं हरियूथप} %||4-18-21||

\twolineshloka
{औरसीं भगिनीं वापि भार्यां वाप्यनुजस्य यः}
{प्रचरेत नरः कामात्तस्य दण्डो वधः स्मृतः} %||4-18-22||

\twolineshloka
{भरतस्तु महीपालो वयं त्वादेशवर्तिनः}
{त्वं च धर्मादतिक्रान्तः कथं शक्यमुपेक्षितुम्} %||4-18-23||

\twolineshloka
{गुरुधर्मव्यतिक्रान्तं प्राज्ञो धर्मेण पालयन्}
{भरतः कामवृत्तानां निग्रहे पर्यवस्थितः} %||4-18-24||

\twolineshloka
{वयं तु भरतादेशं विधिं कृत्वा हरीश्वर}
{त्वद्विधान्भिन्नमर्यादान्नियन्तुं पर्यवस्थिताः} %||4-18-25||

\twolineshloka
{सुग्रीवेण च मे सख्यं लक्ष्मणेन यथा तथा}
{दारराज्यनिमित्तं च निःश्रेयसि रतः स मे} %||4-18-26||

\twolineshloka
{प्रतिज्ञा च मया दत्ता तदा वानरसंनिधौ}
{प्रतिज्ञा च कथं शक्या मद्विधेनानवेक्षितुम्} %||4-18-27||

\twolineshloka
{तदेभिः कारणैः सर्वैर्महद्भिर्धर्मसंहितैः}
{शासनं तव यद्युक्तं तद्भवाननुमन्यताम्} %||4-18-28||

\twolineshloka
{सर्वथा धर्म इत्येव द्रष्टव्यस्तव निग्रहः}
{वयस्यस्योपकर्तव्यं धर्ममेवानुपश्यता} %||4-18-29||

\twolineshloka
{राजभिर्धृतदण्डास्तु कृत्वा पापानि मानवाः}
{निर्मलाः स्वर्गमायान्ति सन्तः सुकृतिनो यथा} %||4-18-30||

\twolineshloka
{आर्येण मम मान्धात्रा व्यसनं घोरमीप्सितम्}
{श्रमणेन कृते पापे यथा पापं कृतं त्वया} %||4-18-31||

\twolineshloka
{अन्यैरपि कृतं पापं प्रमत्तैर्वसुधाधिपैः}
{प्रायश्चित्तं च कुर्वन्ति तेन तच्छाम्यते रजः} %||4-18-32||

\twolineshloka
{तदलं परितापेन धर्मतः परिकल्पितः}
{वधो वानरशार्दूल न वयं स्ववशे स्थिताः} %||4-18-33||

\threelineshloka
{वागुराभिश्च पाशैश्च कूटैश्च विविधैर्नराः}
{प्रतिच्छन्नाश्च दृश्याश्च गृह्णन्ति सुबहून्मृगान्}
{प्रधावितान्वा वित्रस्तान्विस्रब्धानतिविष्ठितान्} %||4-18-34||

\twolineshloka
{प्रमत्तानप्रमत्तान्वा नरा मांसार्थिनो भृशम्}
{विध्यन्ति विमुखांश्चापि न च दोषोऽत्र विद्यते} %||4-18-35||

\threelineshloka
{यान्ति राजर्षयश्चात्र मृगयां धर्मकोविदाः}
{तस्मात्त्वं निहतो युद्धे मया बाणेन वानर}
{अयुध्यन्प्रतियुध्यन्वा यस्माच्छाखामृगो ह्यसि} %||4-18-36||

\twolineshloka
{दुर्लभस्य च धर्मस्य जीवितस्य शुभस्य च}
{राजानो वानरश्रेष्ठ प्रदातारो न संशयः} %||4-18-37||

\twolineshloka
{तान्न हिंस्यान्न चाक्रोशेन्नाक्षिपेन्नाप्रियं वदेत्}
{देवा मानुषरूपेण चरन्त्येते महीतले} %||4-18-38||

\twolineshloka
{त्वं तु धर्ममविज्ञाय केवलं रोषमास्थितः}
{प्रदूषयसि मां धर्मे पितृपैतामहे स्थितम्} %||4-18-39||

\twolineshloka
{एवमुक्तस्तु रामेण वाली प्रव्यथितो भृशम्}
{प्रत्युवाच ततो रामं प्राञ्जलिर्वानरेश्वरः} %||4-18-40||

\twolineshloka
{यत्त्वमात्थ नरश्रेष्ठ तदेवं नात्र संशयः}
{प्रतिवक्तुं प्रकृष्टे हि नापकृष्टस्तु शक्नुयात्} %||4-18-41||

\twolineshloka
{यदयुक्तं मया पूर्वं प्रमादाद्वाक्यमप्रियम्}
{तत्रापि खलु मे दोषं कर्तुं नार्हसि राघव} %||4-18-42||

\twolineshloka
{त्वं हि दृष्टार्थतत्त्वज्ञः प्रजानां च हिते रतः}
{कार्यकारणसिद्धौ ते प्रसन्ना बुद्धिरव्यया} %||4-18-43||

\twolineshloka
{मामप्यवगतं धर्माद्व्यतिक्रान्तपुरस्कृतम्}
{धर्मसंहितया वाचा धर्मज्ञ परिपालय} %||4-18-44||

\twolineshloka
{बाष्पसंरुद्धकण्ठस्तु वाली सार्तरवः शनैः}
{उवाच रामं सम्प्रेक्ष्य पङ्कलग्न इव द्विपः} %||4-18-45||

\twolineshloka
{न त्वात्मानमहं शोचे न तारां नापि बान्धवान्}
{यथा पुत्रं गुणश्रेष्ठमङ्गदं कनकाङ्गदम्} %||4-18-46||

\twolineshloka
{स ममादर्शनाद्दीनो बाल्यात्प्रभृति लालितः}
{तटाक इव पीताम्बुरुपशोषं गमिष्यति} %||4-18-47||

\twolineshloka
{सुग्रीवे चाङ्गदे चैव विधत्स्व मतिमुत्तमाम्}
{त्वं हि शास्ता च गोप्ता च कार्याकार्यविधौ स्थितः} %||4-18-48||

\twolineshloka
{या ते नरपते वृत्तिर्भरते लक्ष्मणे च या}
{सुग्रीवे चाङ्गदे राजंस्तां चिन्तयितुमर्हसि} %||4-18-49||

\twolineshloka
{मद्दोषकृतदोषां तां यथा तारां तपस्विनीम्}
{सुग्रीवो नावमन्येत तथावस्थातुमर्हसि} %||4-18-50||

\twolineshloka
{त्वया ह्यनुगृहीतेन शक्यं राज्यमुपासितुम्}
{त्वद्वशे वर्तमानेन तव चित्तानुवर्तिना} %||4-18-51||

{स तमाश्वासयद्रामो वालिनं व्यक्तदर्शनम्॥\devanumber{52}॥} %||4-18-52|| (Check)

\addtocounter{shlokacount}{1}

\twolineshloka
{न वयं भवता चिन्त्या नाप्यात्मा हरिसत्तम}
{वयं भवद्विशेषेण धर्मतः कृतनिश्चयाः} %||4-18-53||

\twolineshloka
{दण्ड्ये यः पातयेद्दण्डं दण्ड्यो यश्चापि दण्ड्यते}
{कार्यकारणसिद्धार्थावुभौ तौ नावसीदतः} %||4-18-54||

\twolineshloka
{तद्भवान्दण्डसंयोगादस्माद्विगतकल्मषः}
{गतः स्वां प्रकृतिं धर्म्यां धर्मदृष्ट्तेन वर्त्मना} %||4-18-55||

\fourlineindentedshloka
{स तस्य वाक्यं मधुरं महात्मनः}
{ समाहितं धर्मपथानुवर्तिनः}
{निशम्य रामस्य रणावमर्दिनो}
{ वचः सुयुक्तं निजगाद वानरः} %||4-18-56||

\fourlineindentedshloka
{शराभितप्तेन विचेतसा मया}
{ प्रदूषितस्त्वं यदजानता प्रभो}
{इदं महेन्द्रोपमभीमविक्रम}
{ प्रसादितस्त्वं क्षम मे महीश्वर} %||4-19-57||


{॥इत्यार्षे श्रीमद्रामायणे वाल्मीकीये आदिकाव्ये किष्किन्धाकाण्डे अष्टादशः सर्गः॥}

\sect{एकोनविंशः सर्गः}

\twolineshloka
{स वानरमहाराजः शयानः शरविक्षतः}
{प्रत्युक्तो हेतुमद्वाक्यैर्नोत्तरं प्रत्यपद्यत} %||4-19-1||

\twolineshloka
{अश्मभिः परिभिन्नाङ्गः पादपैराहतो भृशम्}
{रामबाणेन चाक्रान्तो जीवितान्ते मुमोह सः} %||4-19-2||

\twolineshloka
{तं भार्याबाणमोक्षेण रामदत्तेन संयुगे}
{हतं प्लवगशार्दूलं तारा शुश्राव वालिनम्} %||4-19-3||

\twolineshloka
{सा सपुत्राप्रियं श्रुत्वा वधं भर्तुः सुदारुणम्}
{निष्पपात भृशं त्रस्ता विविधाद्गिरिगह्वरात्} %||4-19-4||

\twolineshloka
{ये त्वङ्गदपरीवारा वानरा हि महाबलाः}
{ते सकार्मुकमालोक्य रामं त्रस्ताः प्रदुद्रुवुः} %||4-19-5||

\twolineshloka
{सा ददर्श ततस्त्रस्तान्हरीनापततो द्रुतम्}
{यूथादिव परिभ्रष्टान्मृगान्निहतयूथपान्} %||4-19-6||

\twolineshloka
{तानुवाच समासाद्य दुःखितान्दुःखिता सती}
{राम वित्रासितान्सर्वाननुबद्धानिवेषुभिः} %||4-19-7||

\twolineshloka
{वानरा राजसिंहस्य यस्य यूयं पुरःसराः}
{तं विहाय सुवित्रस्ताः कस्माद्द्रवत दुर्गताः} %||4-19-8||

\twolineshloka
{राज्यहेतोः स चेद्भ्राता भ्राता रौद्रेण पातितः}
{रामेण प्रसृतैर्दूरान्मार्गणैर्दूर पातिभिः} %||4-19-9||

\twolineshloka
{कपिपत्न्या वचः श्रुत्वा कपयः कामरूपिणः}
{प्राप्तकालमविश्लिष्टमूचुर्वचनमङ्गनाम्} %||4-19-10||

\twolineshloka
{जीव पुत्रे निवर्तस्य पुत्रं रक्षस्व चान्दगम्}
{अन्तको राम रूपेण हत्वा नयति वालिनम्} %||4-19-11||

\twolineshloka
{क्षिप्तान्वृक्षान्समाविध्य विपुलाश्च शिलास्तथा}
{वाली वज्रसमैर्बाणैर्वज्रेणेव निपातितः} %||4-19-12||

\twolineshloka
{अभिद्रुतमिदं सर्वं विद्रुतं प्रसृतं बलम्}
{अस्मिन्प्लवगशार्दूले हते शक्रसमप्रभे} %||4-19-13||

\twolineshloka
{रक्ष्यतां नगरं शूरैरङ्गदश्चाभिषिच्यताम्}
{पदस्थं वालिनः पुत्रं भजिष्यन्ति प्लवङ्गमाः} %||4-19-14||

\twolineshloka
{अथ वा रुचिरं स्थानमिह ते रुचिरानने}
{आविशन्ति हि दुर्गाणि क्षिप्रमद्यैव वानराः} %||4-19-15||

\twolineshloka
{अभार्याः सह भार्याश्च सन्त्यत्र वनचारिणः}
{लुब्धेभ्यो विप्रयुक्तेभ्यः स्वेभ्यो नस्तुमुलं भयम्} %||4-19-16||

\twolineshloka
{अल्पान्तरगतानां तु श्रुत्वा वचनमङ्गना}
{आत्मनः प्रतिरूपं सा बभाषे चारुहासिनी} %||4-19-17||

\twolineshloka
{पुत्रेण मम किं कार्यं किं राज्येन किमात्मना}
{कपिसिंहे महाभागे तस्मिन्भर्तरि नश्यति} %||4-19-18||

\twolineshloka
{पादमूलं गमिष्यामि तस्यैवाहं महात्मनः}
{योऽसौ रामप्रयुक्तेन शरेण विनिपातितः} %||4-19-19||

\twolineshloka
{एवमुक्त्वा प्रदुद्राव रुदती शोककर्शिता}
{शिरश्चोरश्च बाहुभ्यां दुःखेन समभिघ्नती} %||4-19-20||

\twolineshloka
{आव्रजन्ती ददर्शाथ पतिं निपतितं भुवि}
{हन्तारं दानवेन्द्राणां समरेष्वनिवर्तिनाम्} %||4-19-21||

\twolineshloka
{क्षेप्तारं पर्वतेन्द्राणां वज्राणामिव वासवम्}
{महावातसमाविष्टं महामेघौघनिःस्वनम्} %||4-19-22||

\twolineshloka
{शक्रतुल्यपराक्रान्तं वृष्ट्वेवोपरतं घनम्}
{नर्दन्तं नर्दतां भीमं शूरं शूरेण पातितम्} %||4-19-23||

\twolineshloka
{शार्दूलेनामिषस्यार्थे मृगराजं यथा हतम्}
{अर्चितं सर्वलोकस्य सपताकं सवेदिकम्} %||4-19-24||

\twolineshloka
{नागहेतोः सुपर्णेन चैत्यमुन्मथितं यथा}
{अवष्टभ्यावतिष्ठन्तं ददर्श धनुरूर्जितम्} %||4-19-25||

\twolineshloka
{रामं रामानुजं चैव भर्तुश्चैवानुजं शुभा}
{तानतीत्य समासाद्य भर्तारं निहतं रणे} %||4-19-26||

\twolineshloka
{समीक्ष्य व्यथिता भूमौ सम्भ्रान्ता निपपात ह}
{सुप्तेव पुनरुत्थाय आर्यपुत्रेति क्रोशती} %||4-19-27||

\twolineshloka
{रुरोद सा पतिं दृष्ट्वा सन्दितं मृत्युदामभिः}
{तामवेक्ष्य तु सुग्रीवः क्रोशन्तीं कुररीमिव} %||4-19-28||

{विषादमगमत्कष्टं दृष्ट्वा चाङ्गदमागतम्॥\devanumber{29}॥} %||4-20-29|| (Check)

\addtocounter{shlokacount}{1}


{॥इत्यार्षे श्रीमद्रामायणे वाल्मीकीये आदिकाव्ये किष्किन्धाकाण्डे एकोनविंशः सर्गः॥}

\sect{विंशः सर्गः}

\twolineshloka
{रामचापविसृष्टेन शरेणान्तकरेण तम्}
{दृष्ट्वा विनिहतं भूमौ तारा ताराधिपानना} %||4-20-1||

\twolineshloka
{सा समासाद्य भर्तारं पर्यष्वजत भामिनी}
{इषुणाभिहतं दृष्ट्वा वालिनं कुञ्जरोपमम्} %||4-20-2||

\twolineshloka
{वानरेन्द्रं महेन्द्राभं शोकसन्तप्तमानसा}
{तारा तरुमिवोन्मूलं पर्यदेवयदातुरा} %||4-20-3||

\twolineshloka
{रणे दारुणविक्रान्त प्रवीर प्लवतां वर}
{किं दीनामपुरोभागामद्य त्वं नाभिभाषसे} %||4-20-4||

\twolineshloka
{उत्तिष्ठ हरिशार्दूल भजस्व शयनोत्तमम्}
{नैवंविधाः शेरते हि भूमौ नृपतिसत्तमाः} %||4-20-5||

\twolineshloka
{अतीव खलु ते कान्ता वसुधा वसुधाधिप}
{गतासुरपि यां गात्रैर्मां विहाय निषेवसे} %||4-20-6||

\twolineshloka
{व्यक्तमन्या त्वया वीर धर्मतः सम्प्रवर्तता}
{किष्किन्धेव पुरी रम्या स्वर्गमार्गे विनिर्मिता} %||4-20-7||

\twolineshloka
{यान्यस्माभिस्त्वया सार्धं वनेषु मधुगन्धिषु}
{विहृतानि त्वया काले तेषामुपरमः कृतः} %||4-20-8||

\twolineshloka
{निरानन्दा निराशाहं निमग्ना शोकसागरे}
{त्वयि पञ्चत्वमापन्ने महायूथपयूथपे} %||4-20-9||

\twolineshloka
{हृदयं सुस्थिरं मह्यं दृष्ट्वा विनिहतं भुवि}
{यन्न शोकाभिसन्तप्तं स्फुटतेऽद्य सहस्रधा} %||4-20-10||

\twolineshloka
{सुग्रीवस्य त्वया भार्या हृता स च विवासितः}
{यत्तत्तस्य त्वया व्युष्टिः प्राप्तेयं प्लवगाधिप} %||4-20-11||

\twolineshloka
{निःश्रेयसपरा मोहात्त्वया चाहं विगर्हिता}
{यैषाब्रुवं हितं वाक्यं वानरेन्द्रहितैषिणी} %||4-20-12||

\twolineshloka
{कालो निःसंशयो नूनं जीवितान्तकरस्तव}
{बलाद्येनावपन्नोऽसि सुग्रीवस्यावशो वशम्} %||4-20-13||

\twolineshloka
{वैधव्यं शोकसन्तापं कृपणं कृपणा सती}
{अदुःखोपचिता पूर्वं वर्तयिष्याम्यनाथवत्} %||4-20-14||

\twolineshloka
{लालितश्चाङ्गदो वीरः सुकुमारः सुखोचितः}
{वत्स्यते कामवस्थां मे पितृव्ये क्रोधमूर्छिते} %||4-20-15||

\twolineshloka
{कुरुष्व पितरं पुत्र सुदृष्टं धर्मवत्सलम्}
{दुर्लभं दर्शनं त्वस्य तव वत्स भविष्यति} %||4-20-16||

\twolineshloka
{समाश्वासय पुत्रं त्वं सन्देशं सन्दिशस्व च}
{मूर्ध्नि चैनं समाघ्राय प्रवासं प्रस्थितो ह्यसि} %||4-20-17||

\twolineshloka
{रामेण हि महत्कर्म कृतं त्वामभिनिघ्नता}
{आनृण्यं तु गतं तस्य सुग्रीवस्य प्रतिश्रवे} %||4-20-18||

\twolineshloka
{सकामो भव सुग्रीव रुमां त्वं प्रतिपत्स्यसे}
{भुङ्क्ष्व राज्यमनुद्विग्नः शस्तो भ्राता रिपुस्तव} %||4-20-19||

\twolineshloka
{किं मामेवं विलपतीं प्रेंणा त्वं नाभिभाषसे}
{इमाः पश्य वरा बह्वीर्भार्यास्ते वानरेश्वर} %||4-20-20||

\twolineshloka
{तस्या विलपितं श्रुत्वा वानर्यः सर्वतश्च ताः}
{परिगृह्याङ्गदं दीनं दुःखार्ताः परिचुक्रुशुः} %||4-20-21||

\fourlineindentedshloka
{किमङ्गदं साङ्गद वीर बाहो}
{ विहाय यास्यद्य चिरप्रवासम्}
{न युक्तमेवं गुणसंनिकृष्टम्}
{ विहाय पुत्रं प्रियपुत्र गन्तुम्} %||4-20-22||

\fourlineindentedshloka
{किमप्रियं ते प्रियचारुवेष}
{ कृतं मया नाथ सुतेन वा ते}
{सहायिनीमद्य विहाय वीर}
{ यमक्षयं गच्छसि दुर्विनीतम्} %||4-20-23||

\fourlineindentedshloka
{यद्यप्रियं किञ्चिदसम्प्रधार्य}
{ कृतं मया स्यात्तव दीर्घबाहो}
{क्षमस्व मे तद्धरिवंश नाथ}
{ व्रजामि मूर्ध्ना तव वीर पादौ} %||4-20-24||

\fourlineindentedshloka
{तथा तु तारा करुणं रुदन्ती}
{ भर्तुः समीपे सह वानरीभिः}
{व्यवस्यत प्रायमनिन्द्यवर्णा}
{ उपोपवेष्टुं भुवि यत्र वाली} %||4-21-25||


{॥इत्यार्षे श्रीमद्रामायणे वाल्मीकीये आदिकाव्ये किष्किन्धाकाण्डे विंशः सर्गः॥}

\sect{एकविंशः सर्गः}

\twolineshloka
{ततो निपतितां तारां च्युतां तारामिवाम्बरात्}
{शनैराश्वासयामास हनूमान्हरियूथपः} %||4-21-1||

\twolineshloka
{गुणदोषकृतं जन्तुः स्वकर्मफलहेतुकम्}
{अव्यग्रस्तदवाप्नोति सर्वं प्रेत्य शुभाशुभम्} %||4-21-2||

\twolineshloka
{शोच्या शोचसि कं शोच्यं दीनं दीनानुकम्पसे}
{कश्च कस्यानुशोच्योऽस्ति देहेऽस्मिन्बुद्बुदोपमे} %||4-21-3||

\twolineshloka
{अङ्गदस्तु कुमारोऽयं द्रष्टव्यो जीवपुत्रया}
{आयत्या च विधेयानि समर्थान्यस्य चिन्तय} %||4-21-4||

\twolineshloka
{जानास्यनियतामेवं भूतानामागतिं गतिम्}
{तस्माच्छुभं हि कर्तव्यं पण्डिते नैहलौकिकम्} %||4-21-5||

\twolineshloka
{यस्मिन्हरिसहस्राणि प्रयुतान्यर्बुदानि च}
{वर्तयन्ति कृतांशानि सोऽयं दिष्टान्तमागतः} %||4-21-6||

\twolineshloka
{यदयं न्यायदृष्टार्थः सामदानक्षमापरः}
{गतो धर्मजितां भूमिं नैनं शोचितुमर्हसि} %||4-21-7||

\twolineshloka
{सर्वे च हरिशार्दूल पुत्रश्चायं तवाङ्गदः}
{हर्यृष्कपतिराज्यं च त्वत्सनाथमनिन्दिते} %||4-21-8||

\twolineshloka
{ताविमौ शोकसन्तप्तौ शनैः प्रेरय भामिनि}
{त्वया परिगृहीतोऽयमङ्गदः शास्तु मेदिनीम्} %||4-21-9||

\twolineshloka
{सन्ततिश्च यथादृष्टा कृत्यं यच्चापि साम्प्रतम्}
{राज्ञस्तत्क्रियतां सर्वमेष कालस्य निश्चयः} %||4-21-10||

\twolineshloka
{संस्कार्यो हरिराजस्तु अङ्गदश्चाभिषिच्यताम्}
{सिंहासनगतं पुत्रं पश्यन्ती शान्तिमेष्यसि} %||4-21-11||

\twolineshloka
{सा तस्य वचनं श्रुत्वा भर्तृव्यसनपीडिता}
{अब्रवीदुत्तरं तारा हनूमन्तमवस्थितम्} %||4-21-12||

\twolineshloka
{अङ्गद प्रतिरूपाणां पुत्राणामेकतः शतम्}
{हतस्याप्यस्य वीरस्य गात्रसंश्लेषणं वरम्} %||4-21-13||

\twolineshloka
{न चाहं हरिराजस्य प्रभवाम्यङ्गदस्य वा}
{पितृव्यस्तस्य सुग्रीवः सर्वकार्येष्वनन्तरः} %||4-21-14||

\twolineshloka
{न ह्येषा बुद्धिरास्थेया हनूमन्नङ्गदं प्रति}
{पिता हि बन्धुः पुत्रस्य न माता हरिसत्तम} %||4-21-15||

\fourlineindentedshloka
{न हि मम हरिराजसंश्रयात्}
{ क्षमतरमस्ति परत्र चेह वा}
{अभिमुखहतवीरसेवितम्}
{ शयनमिदं मम सेवितुं क्षमम्} %||4-22-16||


{॥इत्यार्षे श्रीमद्रामायणे वाल्मीकीये आदिकाव्ये किष्किन्धाकाण्डे एकविंशः सर्गः॥}

\sect{द्वाविंशः सर्गः}

\twolineshloka
{वीक्षमाणस्तु मन्दासुः सर्वतो मन्दमुच्छ्वसन्}
{आदावेव तु सुग्रीवं ददर्श त्वात्मजाग्रतः} %||4-22-1||

\twolineshloka
{तं प्राप्तविजयं वाली सुग्रीवं प्लवगेश्वरम्}
{आभाष्य व्यक्तया वाचा सस्नेहमिदमब्रवीत्} %||4-22-2||

\twolineshloka
{सुग्रीवदोषेण न मां गन्तुमर्हसि किल्बिषात्}
{कृष्यमाणं भविष्येण बुद्धिमोहेन मां बलात्} %||4-22-3||

\twolineshloka
{युगपद्विहितं तात न मन्ये सुखमावयोः}
{सौहार्दं भ्रातृयुक्तं हि तदिदं जातमन्यथा} %||4-22-4||

\twolineshloka
{प्रतिपद्य त्वमद्यैव राज्यमेषां वनौकसाम्}
{मामप्यद्यैव गच्छन्तं विद्धि वैवस्वतक्षयम्} %||4-22-5||

\twolineshloka
{जीवितं च हि राज्यं च श्रियं च विपुलामिमाम्}
{प्रजहाम्येष वै तूर्णं महच्चागर्हितं यशः} %||4-22-6||

\twolineshloka
{अस्यां त्वहमवस्थायां वीर वक्ष्यामि यद्वचः}
{यद्यप्यसुकरं राजन्कर्तुमेव तदर्हसि} %||4-22-7||

\twolineshloka
{सुखार्हं सुखसंवृद्धं बालमेनमबालिशम्}
{बाष्पपूर्णमुखं पश्य भूमौ पतितमङ्गदम्} %||4-22-8||

\twolineshloka
{मम प्राणैः प्रियतरं पुत्रं पुत्रमिवौरसम्}
{मया हीनमहीनार्थं सर्वतः परिपालय} %||4-22-9||

\twolineshloka
{त्वमप्यस्य हि दाता च परित्राता च सर्वतः}
{भयेष्वभयदश्चैव यथाहं प्लवगेश्वर} %||4-22-10||

\twolineshloka
{एष तारात्मजः श्रीमांस्त्वया तुल्यपराक्रमः}
{रक्षसां तु वधे तेषामग्रतस्ते भविष्यति} %||4-22-11||

\twolineshloka
{अनुरूपाणि कर्माणि विक्रम्य बलवान्रणे}
{करिष्यत्येष तारेयस्तरस्वी तरुणोऽङ्गदः} %||4-22-12||

\twolineshloka
{सुषेणदुहिता चेयमर्थसूक्ष्मविनिश्चये}
{औत्पातिके च विविधे सर्वतः परिनिष्ठिता} %||4-22-13||

\twolineshloka
{यदेषा साध्विति ब्रूयात्कार्यं तन्मुक्तसंशयम्}
{न हि तारामतं किञ्चिदन्यथा परिवर्तते} %||4-22-14||

\twolineshloka
{राघवस्य च ते कार्यं कर्तव्यमविशङ्कया}
{स्यादधर्मो ह्यकरणे त्वां च हिंस्याद्विमानितः} %||4-22-15||

\twolineshloka
{इमां च मालामाधत्स्व दिव्यां सुग्रीवकाञ्चनीम्}
{उदारा श्रीः स्थिता ह्यस्यां सम्प्रजह्यान्मृते मयि} %||4-22-16||

\twolineshloka
{इत्येवमुक्तः सुग्रीवो वालिना भ्रातृसौहृदात्}
{हर्षं त्यक्त्वा पुनर्दीनो ग्रहग्रस्त इवोडुराट्} %||4-22-17||

\twolineshloka
{तद्वालिवचनाच्छान्तः कुर्वन्युक्तमतन्द्रितः}
{जग्राह सोऽभ्यनुज्ञातो मालां तां चैव काञ्चनीम्} %||4-22-18||

\twolineshloka
{तां मालां काञ्चनीं दत्त्वा वाली दृष्ट्वात्मजं स्थितम्}
{संसिद्धः प्रेत्य भावाय स्नेहादङ्गदमब्रवीत्} %||4-22-19||

\twolineshloka
{देशकालौ भजस्वाद्य क्षममाणः प्रियाप्रिये}
{सुखदुःखसहः काले सुग्रीववशगो भव} %||4-22-20||

\twolineshloka
{यथा हि त्वं महाबाहो लालितः सततं मया}
{न तथा वर्तमानं त्वां सुग्रीवो बहु मंस्यते} %||4-22-21||

\twolineshloka
{मास्यामित्रैर्गतं गच्छेर्मा शत्रुभिररिन्दम}
{भर्तुरर्थपरो दान्तः सुग्रीववशगो भव} %||4-22-22||

\twolineshloka
{न चातिप्रणयः कार्यः कर्तव्योऽप्रणयश्च ते}
{उभयं हि महादोषं तस्मादन्तरदृग्भव} %||4-22-23||

\twolineshloka
{इत्युक्त्वाथ विवृत्ताक्षः शरसम्पीडितो भृशम्}
{विवृतैर्दशनैर्भीमैर्बभूवोत्क्रान्तजीवितः} %||4-22-24||

\fourlineindentedshloka
{हते तु वीरे प्लवगाधिपे तदा}
{ प्लवङ्गमास्तत्र न शर्म लेभिरे}
{वनेचराः सिंहयुते महावने}
{ यथा हि गावो निहते गवां पतौ} %||4-22-25||

\fourlineindentedshloka
{ततस्तु तारा व्यसनार्णव प्लुता}
{ मृतस्या भर्तुर्वदनं समीक्ष्य सा}
{जगाम भूमिं परिरभ्य वालिनम्}
{ महाद्रुमं छिन्नमिवाश्रिता लता} %||4-23-26||


{॥इत्यार्षे श्रीमद्रामायणे वाल्मीकीये आदिकाव्ये किष्किन्धाकाण्डे द्वाविंशः सर्गः॥}

\sect{त्रयोविंशः सर्गः}

\twolineshloka
{ततः समुपजिघ्रन्ती कपिराजस्य तन्मुखम्}
{पतिं लोकाच्च्युतं तारा मृतं वचनमब्रवीत्} %||4-23-1||

\twolineshloka
{शेषे त्वं विषमे दुःखमकृत्वा वचनं मम}
{उपलोपचिते वीर सुदुःखे वसुधातले} %||4-23-2||

\twolineshloka
{मत्तः प्रियतरा नूनं वानरेन्द्र मही तव}
{शेषे हि तां परिष्वज्य मां च न प्रतिभाषसे} %||4-23-3||

\twolineshloka
{सुग्रीव एव विक्रान्तो वीर साहसिक प्रिय}
{ऋक्षवानरमुख्यास्त्वां बलिनं पर्युपासते} %||4-23-4||

\twolineshloka
{एषां विलपितं कृच्छ्रमङ्गदस्य च शोचतः}
{मम चेमां गिरं श्रुत्वा किं त्वं न प्रतिबुध्यसे} %||4-23-5||

\twolineshloka
{इदं तच्छूरशयनं यत्र शेषे हतो युधि}
{शायिता निहता यत्र त्वयैव रिपवः पुरा} %||4-23-6||

\twolineshloka
{विशुद्धसत्त्वाभिजन प्रिययुद्ध मम प्रिय}
{मामनाथां विहायैकां गतस्त्वमसि मानद} %||4-23-7||

\twolineshloka
{शूराय न प्रदातव्या कन्या खलु विपश्चिता}
{शूरभार्यां हतां पश्य सद्यो मां विधवां कृताम्} %||4-23-8||

\twolineshloka
{अवभग्नश्च मे मानो भग्ना मे शाश्वती गतिः}
{अगाधे च निमग्नास्मि विपुले शोकसागरे} %||4-23-9||

\twolineshloka
{अश्मसारमयं नूनमिदं मे हृदयं दृढम्}
{भर्तारं निहतं दृष्ट्वा यन्नाद्य शतधा गतम्} %||4-23-10||

\twolineshloka
{सुहृच्चैव हि भर्ता च प्रकृत्या च मम प्रियः}
{आहवे च पराक्रान्तः शूरः पञ्चत्वमागतः} %||4-23-11||

\twolineshloka
{पतिहीना तु या नारी कामं भवतु पुत्रिणी}
{धनधान्यैः सुपूर्णापि विधवेत्युच्यते बुधैः} %||4-23-12||

\twolineshloka
{स्वगात्रप्रभवे वीर शेषे रुधिरमण्डले}
{कृमिरागपरिस्तोमे त्वमेवं शयने यथा} %||4-23-13||

\twolineshloka
{रेणुशोणितसंवीतं गात्रं तव समन्ततः}
{परिरब्धुं न शक्नोमि भुजाभ्यां प्लवगर्षभ} %||4-23-14||

\twolineshloka
{कृतकृत्योऽद्य सुग्रीवो वैरेऽस्मिन्नतिदारुणे}
{यस्य रामविमुक्तेन हृतमेकेषुणा भयम्} %||4-23-15||

\twolineshloka
{शरेण हृदि लग्नेन गात्रसंस्पर्शने तव}
{वार्यामि त्वां निरीक्षन्ती त्वयि पञ्चत्वमागते} %||4-23-16||

\twolineshloka
{उद्बबर्ह शरं नीलस्तस्य गात्रगतं तदा}
{गिरिगह्वरसंलीनं दीप्तमाशीविषं यथा} %||4-23-17||

\twolineshloka
{तस्य निष्कृष्यमाणस्य बाणस्य च बभौ द्युतिः}
{अस्तमस्तकसंरुद्धो रश्मिर्दिनकरादिव} %||4-23-18||

\twolineshloka
{पेतुः क्षतजधारास्तु व्रणेभ्यस्तस्य सर्वशः}
{ताम्रगैरिकसम्पृक्ता धारा इव धराधरात्} %||4-23-19||

\twolineshloka
{अवकीर्णं विमार्जन्ती भर्तारं रणरेणुना}
{अस्रैर्नयनजैः शूरं सिषेचास्त्रसमाहतम्} %||4-23-20||

\twolineshloka
{रुधिरोक्षितसर्वाङ्गं दृष्ट्वा विनिहतं पतिम्}
{उवाच तारा पिङ्गाक्षं पुत्रमङ्गदमङ्गना} %||4-23-21||

\twolineshloka
{अवस्थां पश्चिमां पश्य पितुः पुत्र सुदारुणाम्}
{सम्प्रसक्तस्य वैरस्य गतोऽन्तः पापकर्मणा} %||4-23-22||

\twolineshloka
{बालसूर्योदयतनुं प्रयान्तं यमसादनम्}
{अभिवादय राजानं पितरं पुत्र मानदम्} %||4-23-23||

\twolineshloka
{एवमुक्तः समुत्थाय जग्राह चरणौ पितुः}
{भुजाभ्यां पीनवृताभ्यामङ्गदोऽहमिति ब्रुवन्} %||4-23-24||

\twolineshloka
{अभिवादयमानं त्वामङ्गदं त्वं यथापुरा}
{दीर्घायुर्भव पुत्रेति किमर्थं नाभिभाषसे} %||4-23-25||

\twolineshloka
{अहं पुत्रसहाया त्वामुपासे गतचेतनम्}
{सिंहेन निहतं सद्यो गौः सवत्सेव गोवृषम्} %||4-23-26||

\twolineshloka
{इष्ट्वा सङ्ग्रामयज्ञेन नानाप्रहरणाम्भसा}
{अस्मिन्नवभृथे स्नातः कथं पत्न्या मया विना} %||4-23-27||

\twolineshloka
{या दत्ता देवराजेन तव तुष्टेन संयुगे}
{शातकुम्भमयीं मालां तां ते पश्यामि नेह किम्} %||4-23-28||

\twolineshloka
{राजश्रीर्न जहाति त्वां गतासुमपि मानद}
{सूर्यस्यावर्तमानस्य शैलराजमिव प्रभा} %||4-23-29||

\fourlineindentedshloka
{न मे वचः पथ्यमिदं त्वया कृतम्}
{ न चास्मि शक्ता हि निवारणे तव}
{हता सपुत्रास्मि हतेन संयुगे}
{ सह त्वया श्रीर्विजहाति मामिह} %||4-24-30||


{॥इत्यार्षे श्रीमद्रामायणे वाल्मीकीये आदिकाव्ये किष्किन्धाकाण्डे त्रयोविंशः सर्गः॥}

\sect{चतुर्विंशः सर्गः}

\twolineshloka
{गतासुं वालिनं दृष्ट्वा राघवस्तदनन्तरम्}
{अब्रवीत्प्रश्रितं वाक्यं सुग्रीवं शत्रुतापनः} %||4-24-1||

\twolineshloka
{न शोकपरितापेन श्रेयसा युज्यते मृतः}
{यदत्रानन्तरं कार्यं तत्समाधातुमर्हथ} %||4-24-2||

\twolineshloka
{लोकवृत्तमनुष्ठेयं कृतं वो बाष्पमोक्षणम्}
{न कालादुत्तरं किञ्चित्कर्म शक्यमुपासितुम्} %||4-24-3||

\twolineshloka
{नियतः कारणं लोके नियतिः कर्मसाधनम्}
{नियतिः सर्वभूतानां नियोगेष्विह कारणम्} %||4-24-4||

\twolineshloka
{न कर्ता कस्यचित्कश्चिन्नियोगे चापि नेश्वरः}
{स्वभावे वर्तते लोकस्तस्य कालः परायणम्} %||4-24-5||

\twolineshloka
{न कालः कालमत्येति न कालः परिहीयते}
{स्वभावं वा समासाद्य न कश्चिदतिवर्तते} %||4-24-6||

\twolineshloka
{न कालस्यास्ति बन्धुत्वं न हेतुर्न पराक्रमः}
{न मित्रज्ञातिसम्बन्धः कारणं नात्मनो वशः} %||4-24-7||

\twolineshloka
{किं तु काल परीणामो द्रष्टव्यः साधु पश्यता}
{धर्मश्चार्थश्च कामश्च कालक्रमसमाहिताः} %||4-24-8||

\twolineshloka
{इतः स्वां प्रकृतिं वाली गतः प्राप्तः क्रियाफलम्}
{धर्मार्थकामसंयोगैः पवित्रं प्लवगेश्वर} %||4-24-9||

\twolineshloka
{स्वधर्मस्य च संयोगाज्जितस्तेन महात्मना}
{स्वर्गः परिगृहीतश्च प्राणानपरिरक्षता} %||4-24-10||

\twolineshloka
{एषा वै नियतिः श्रेष्ठा यां गतो हरियूथपः}
{तदलं परितापेन प्राप्तकालमुपास्यताम्} %||4-24-11||

\twolineshloka
{वचनान्ते तु रामस्य लक्ष्मणः परवीरहा}
{अवदत्प्रश्रितं वाक्यं सुग्रीवं गतचेतसम्} %||4-24-12||

\twolineshloka
{कुरु त्वमस्य सुग्रीव प्रेतकार्यमनन्तरम्}
{ताराङ्गदाभ्यां सहितो वालिनो दहनं प्रति} %||4-24-13||

\twolineshloka
{समाज्ञापय काष्ठानि शुष्काणि च बहूनि च}
{चन्दनानि च दिव्यानि वालिसंस्कारकारणात्} %||4-24-14||

\twolineshloka
{समाश्वासय चैनं त्वमङ्गदं दीनचेतसम्}
{मा भूर्बालिशबुद्धिस्त्वं त्वदधीनमिदं पुरम्} %||4-24-15||

\twolineshloka
{अङ्गदस्त्वानयेन्माल्यं वस्त्राणि विविधानि च}
{घृतं तैलमथो गन्धान्यच्चात्र समनन्तरम्} %||4-24-16||

\twolineshloka
{त्वं तार शिबिकां शीघ्रमादायागच्छ सम्भ्रमात्}
{त्वरा गुणवती युक्ता ह्यस्मिन्काले विशेषतः} %||4-24-17||

\twolineshloka
{सज्जीभवन्तु प्लवगाः शिबिकावाहनोचिताः}
{समर्था बलिनश्चैव निर्हरिष्यन्ति वालिनम्} %||4-24-18||

\twolineshloka
{एवमुक्त्वा तु सुग्रीवं सुमित्रानन्दवर्धनः}
{तस्थौ भ्रातृसमीपस्थो लक्ष्मणः परवीरहा} %||4-24-19||

\twolineshloka
{लक्ष्मणस्य वचः श्रुत्वा तारः सम्भ्रान्तमानसः}
{प्रविवेश गुहां शीघ्रं शिबिकासक्तमानसः} %||4-24-20||

\twolineshloka
{आदाय शिबिकां तारः स तु पर्यापयत्पुनः}
{वानरैरुह्यमानां तां शूरैरुद्वहनोचितैः} %||4-24-21||

\twolineshloka
{ततो वालिनमुद्यम्य सुग्रीवः शिबिकां तदा}
{आरोपयत विक्रोशन्नङ्गदेन सहैव तु} %||4-24-22||

\twolineshloka
{आरोप्य शिबिकां चैव वालिनं गतजीवितम्}
{अलङ्कारैश्च विविधैर्माल्यैर्वस्त्रैश्च भूषितम्} %||4-24-23||

\twolineshloka
{आज्ञापयत्तदा राजा सुग्रीवः प्लवगेश्वरः}
{और्ध्वदेहिकमार्यस्य क्रियतामनुरूपतः} %||4-24-24||

\twolineshloka
{विश्राणयन्तो रत्नानि विविधानि बहूनि च}
{अग्रतः प्लवगा यान्तु शिबिका तदनन्तरम्} %||4-24-25||

\twolineshloka
{राज्ञामृद्धिविशेषा हि दृश्यन्ते भुवि यादृशाः}
{तादृशं वालिनः क्षिप्रं प्राकुर्वन्नौर्ध्वदेहिकम्} %||4-24-26||

\twolineshloka
{अङ्गदमप्रिगृह्याशु तारप्रभृतयस्तथा}
{क्रोशन्तः प्रययुः सर्वे वानरा हतबान्धवाः} %||4-24-27||

\twolineshloka
{ताराप्रभृतयः सर्वा वानर्यो हतयूथपाः}
{अनुजग्मुर्हि भर्तारं क्रोशन्त्यः करुणस्वनाः} %||4-24-28||

\twolineshloka
{तासां रुदितशब्देन वानरीणां वनान्तरे}
{वनानि गिरयः सर्वे विक्रोशन्तीव सर्वतः} %||4-24-29||

\twolineshloka
{पुलिने गिरिनद्यास्तु विविक्ते जलसंवृते}
{चितां चक्रुः सुबहवो वानरा वनचारिणः} %||4-24-30||

\twolineshloka
{अवरोप्य ततः स्कन्धाच्छिबिकां वहनोचिताः}
{तस्थुरेकान्तमाश्रित्य सर्वे शोकसमन्विताः} %||4-24-31||

\twolineshloka
{ततस्तारा पतिं दृष्ट्वा शिबिकातलशायिनम्}
{आरोप्याङ्के शिरस्तस्य विललाप सुदुःखिता} %||4-24-32||

\threelineshloka
{जनं च पश्यसीमं त्वं कस्माच्छोकाभिपीडितम्}
{प्रहृष्टमिव ते वक्त्रं गतासोरपि मानद}
{अस्तार्कसमवर्णं च लक्ष्यते जीवतो यथा} %||4-24-33||

\twolineshloka
{एष त्वां रामरूपेण कालः कर्षति वानर}
{येन स्म विधवाः सर्वाः कृता एकेषुणा रणे} %||4-24-34||

\twolineshloka
{इमास्तास्तव राजेन्द्रवानर्यो वल्लभाः सदा}
{पादैर्विकृष्टमध्वानमागताः किं न बुध्यसे} %||4-24-35||

\twolineshloka
{तवेष्टा ननु नामैता भार्याश्चन्द्रनिभाननाः}
{इदानीं नेक्षसे कस्मात्सुग्रीवं प्लवगेश्वरम्} %||4-24-36||

\twolineshloka
{एते हि सचिवा राजंस्तारप्रभृतयस्तव}
{पुरवासिजनश्चायं परिवार्यासतेऽनघ} %||4-24-37||

\twolineshloka
{विसर्जयैनान्प्रवलान्यथोचितमरिन्दम}
{ततः क्रीडामहे सर्वा वनेषु मदिरोत्कटाः} %||4-24-38||

\twolineshloka
{एवं विलपतीं तारां पतिशोकपरिप्लुताम्}
{उत्थापयन्ति स्म तदा वानर्यः शोककर्शिताः} %||4-24-39||

\twolineshloka
{सुग्रीवेण ततः सार्धमङ्गदः पितरं रुदन्}
{चितामारोपयामास शोकेनाभिहतेन्द्रियः} %||4-24-40||

\twolineshloka
{ततोऽग्निं विधिवद्दत्त्वा सोऽपसव्यं चकार ह}
{पितरं दीर्घमध्वानं प्रस्थितं व्याकुलेन्द्रियः} %||4-24-41||

\twolineshloka
{संस्कृत्य वालिनं ते तु विधिपूर्वं प्लवङ्गमाः}
{आजग्मुरुदकं कर्तुं नदीं शीतजलां शुभाम्} %||4-24-42||

\twolineshloka
{ततस्ते सहितास्तत्र अङ्गदं स्थाप्य चाग्रतः}
{सुग्रीवतारासहिताः सिषिचुर्वालिने जलम्} %||4-24-43||

\twolineshloka
{सुग्रीवेणैव दीनेन दीनो भूत्वा महाबलः}
{समानशोकः काकुत्स्थः प्रेतकार्याण्यकारयत्} %||4-25-44||


{॥इत्यार्षे श्रीमद्रामायणे वाल्मीकीये आदिकाव्ये किष्किन्धाकाण्डे चतुर्विंशः सर्गः॥}

\sect{पञ्चविंशः सर्गः}

\twolineshloka
{ततः शोकाभिसन्तप्तं सुग्रीवं क्लिन्नवासनम्}
{शाखामृगमहामात्राः परिवार्योपतस्थिरे} %||4-25-1||

\twolineshloka
{अभिगम्य महाबाहुं राममक्लिष्टकारिणम्}
{स्थिताः प्राञ्जलयः सर्वे पितामहमिवर्षयः} %||4-25-2||

\twolineshloka
{ततः काञ्चनशैलाभस्तरुणार्कनिभाननः}
{अब्रवीत्प्राञ्जलिर्वाक्यं हनुमान्मारुतात्मजः} %||4-25-3||

\twolineshloka
{भवत्प्रसादात्सुग्रीवः पितृपैतामहं महत्}
{वानराणां सुदुष्प्रापं प्राप्तो राज्यमिदं प्रभो} %||4-25-4||

\twolineshloka
{भवता समनुज्ञातः प्रविश्य नगरं शुभम्}
{संविधास्यति कार्याणि सर्वाणि ससुहृज्जनः} %||4-25-5||

\twolineshloka
{स्नातोऽयं विविधैर्गन्धैरौषधैश्च यथाविधि}
{अर्चयिष्यति रत्नैश्च माल्यैश्च त्वां विशेषतः} %||4-25-6||

\twolineshloka
{इमां गिरिगुहां रम्यामभिगन्तुमितोऽर्हसि}
{कुरुष्व स्वामि सम्बन्धं वानरान्सम्प्रहर्षयन्} %||4-25-7||

\twolineshloka
{एवमुक्तो हनुमता राघवः परवीरहा}
{प्रत्युवाच हनूमन्तं बुद्धिमान्वाक्यकोविदः} %||4-25-8||

\twolineshloka
{चतुर्दशसमाः सौम्य ग्रामं वा यदि वा पुरम्}
{न प्रवेक्ष्यामि हनुमन्पितुर्निर्देशपालकः} %||4-25-9||

\twolineshloka
{सुसमृद्धां गुहां दिव्यां सुग्रीवो वानरर्षभः}
{प्रविष्टो विधिवद्वीरः क्षिप्रं राज्येऽभिषिच्यताम्} %||4-25-10||

\twolineshloka
{एवमुक्त्वा हनूमन्तं रामः सुग्रीवमब्रवीत्}
{इममप्यङ्गदं वीर यौवराज्येऽभिषेचय} %||4-25-11||

\twolineshloka
{पूर्वोऽयं वार्षिको मासः श्रावणः सलिलागमः}
{प्रवृत्ताः सौम्य चत्वारो मासा वार्षिकसंज्ञिताः} %||4-25-12||

\twolineshloka
{नायमुद्योगसमयः प्रविश त्वं पुरीं शुभाम्}
{अस्मिन्वत्स्याम्यहं सौम्य पर्वते सहलक्ष्मणः} %||4-25-13||

\twolineshloka
{इयं गिरिगुहा रम्या विशाला युक्तमारुता}
{प्रभूतसलिला सौम्य प्रभूतकमलोत्पला} %||4-25-14||

\threelineshloka
{कार्तिके समनुप्राप्ते त्वं रावणवधे यत}
{एष नः समयः सौम्य प्रविश त्वं स्वमालयम्}
{अभिषिञ्चस्व राज्ये च सुहृदः सम्प्रहर्षय} %||4-25-15||

\twolineshloka
{इति रामाभ्यनुज्ञातः सुग्रीवो वानरर्षभः}
{प्रविवेश पुरीं रम्यां किष्किन्धां वालिपालिताम्} %||4-25-16||

\twolineshloka
{तं वानरसहस्राणि प्रविष्टं वानरेश्वरम्}
{अभिवाद्य प्रहृष्टानि सर्वतः पर्यवारयन्} %||4-25-17||

\twolineshloka
{ततः प्रकृतयः सर्वा दृष्ट्वा हरिगणेश्वरम्}
{प्रणम्य मूर्ध्ना पतिता वसुधायां समाहिताः} %||4-25-18||

\twolineshloka
{सुग्रीवः प्रकृतीः सर्वाः सम्भाष्योत्थाप्य वीर्यवान्}
{भ्रातुरन्तःपुरं सौम्यं प्रविवेश महाबलः} %||4-25-19||

\twolineshloka
{प्रविश्य त्वभिनिष्क्रान्तं सुग्रीवं वानरर्षभम्}
{अभ्यषिञ्चन्त सुहृदः सहस्राक्षमिवामराः} %||4-25-20||

\twolineshloka
{तस्य पाण्डुरमाजह्रुश्छत्रं हेमपरिष्कृतम्}
{शुक्ले च बालव्यजने हेमदण्डे यशस्करे} %||4-25-21||

\twolineshloka
{तथा सर्वाणि रत्नानि सर्वबीजौषधानि च}
{सक्षीराणां च वृक्षाणां प्ररोहान्कुसुमानि च} %||4-25-22||

\twolineshloka
{शुक्लानि चैव वस्त्राणि श्वेतं चैवानुलेपनम्}
{सुगन्धीनि च माल्यानि स्थलजान्यम्बुजानि च} %||4-25-23||

\twolineshloka
{चन्दनानि च दिव्यानि गन्धांश्च विविधान्बहून्}
{अक्षतं जातरूपं च प्रियङ्गुमधुसर्पिषी} %||4-25-24||

\threelineshloka
{दधिचर्म च वैयाघ्रं वाराही चाप्युपानहौ}
{समालम्भनमादाय रोचनां समनःशिलाम्}
{आजग्मुस्तत्र मुदिता वराः कन्यास्तु षोडश} %||4-25-25||

\twolineshloka
{ततस्ते वानरश्रेष्ठं यथाकालं यथाविधि}
{रत्नैर्वस्त्रैश्च भक्ष्यैश्च तोषयित्वा द्विजर्षभान्} %||4-25-26||

\twolineshloka
{ततः कुशपरिस्तीर्णं समिद्धं जातवेदसम्}
{मन्त्रपूतेन हविषा हुत्वा मन्त्रविदो जनाः} %||4-25-27||

\twolineshloka
{ततो हेमप्रतिष्ठाने वरास्तरणसंवृते}
{प्रासादशिखरे रम्ये चित्रमाल्योपशोभिते} %||4-25-28||

\twolineshloka
{प्राङ्मुखं विविधैर्मन्त्रैः स्थापयित्वा वरासने}
{नदीनदेभ्यः संहृत्य तीर्थेभ्यश्च समन्ततः} %||4-25-29||

\twolineshloka
{आहृत्य च समुद्रेभ्यः सर्वेभ्यो वानरर्षभाः}
{अपः कनककुम्भेषु निधाय विमलाः शुभाः} %||4-25-30||

\twolineshloka
{शुभैर्वृषभशृङ्गैश्च कलशैश्चापि काञ्चनैः}
{शास्त्रदृष्टेन विधिना महर्षिविहितेन च} %||4-25-31||

\twolineshloka
{गजो गवाक्षो गवयः शरभो गन्धमादनः}
{मैन्दश्च द्विविदश्चैव हनूमाञ्जाम्बवान्नलः} %||4-25-32||

\twolineshloka
{अभ्यषिञ्चन्त सुग्रीवं प्रसन्नेन सुगन्धिना}
{सलिलेन सहस्राक्षं वसवो वासवं यथा} %||4-25-33||

\twolineshloka
{अभिषिक्ते तु सुग्रीवे सर्वे वानरपुङ्गवाः}
{प्रचुक्रुशुर्महात्मानो हृष्टास्तत्र सहस्रशः} %||4-25-34||

\twolineshloka
{रामस्य तु वचः कुर्वन्सुग्रीवो हरिपुङ्गवः}
{अङ्गदं सम्परिष्वज्य यौवराज्येऽभिषेचयत्} %||4-25-35||

\twolineshloka
{अङ्गदे चाभिषिक्ते तु सानुक्रोशाः प्लवङ्गमाः}
{साधु साध्विति सुग्रीवं महात्मानोऽभ्यपूजयन्} %||4-25-36||

\twolineshloka
{हृष्टपुष्टजनाकीर्णा पताकाध्वजशोभिता}
{बभूव नगरी रम्या क्षिकिन्धा गिरिगह्वरे} %||4-25-37||

\fourlineindentedshloka
{निवेद्य रामाय तदा महात्मने}
{ महाभिषेकं कपिवाहिनीपतिः}
{रुमां च भार्यां प्रतिलभ्य वीर्यवा}
{नवाप राज्यं त्रिदशाधिपो यथा} %||4-26-38||


{॥इत्यार्षे श्रीमद्रामायणे वाल्मीकीये आदिकाव्ये किष्किन्धाकाण्डे पञ्चविंशः सर्गः॥}

\sect{षड्विंशः सर्गः}

\twolineshloka
{अभिषिक्ते तु सुग्रीवे प्रविष्टे वानरे गुहाम्}
{आजगाम सह भ्रात्रा रामः प्रस्रवणं गिरिम्} %||4-26-1||

\twolineshloka
{शार्दूलमृगसङ्घुष्टं सिंहैर्भीमरवैर्वृतम्}
{नानागुल्मलतागूढं बहुपादपसङ्कुलम्} %||4-26-2||

\twolineshloka
{ऋक्षवानरगोपुच्छैर्मार्जारैश्च निषेवितम्}
{मेघराशिनिभं शैलं नित्यं शुचिजलाश्रयम्} %||4-26-3||

\twolineshloka
{तस्य शैलस्य शिखरे महतीमायतां गुहाम्}
{प्रत्यगृह्णत वासार्थं रामः सौमित्रिणा सह} %||4-26-4||

\twolineshloka
{अवसत्तत्र धर्मात्मा राघवः सहलक्ष्मणः}
{बहुदृश्यदरीकुञ्जे तस्मिन्प्रस्रवणे गिरौ} %||4-26-5||

\threelineshloka
{सुसुखेऽपि बहुद्रव्ये तस्मिन्हि धरणीधरे}
{वसतस्तस्य रामस्य रतिरल्पापि नाभवत्}
{हृतां हि भार्यां स्मरतः प्राणेभ्योऽपि गरीयसीम्} %||4-26-6||

\twolineshloka
{उदयाभ्युदितं दृष्ट्वा शशाङ्कं च विशेषतः}
{आविवेश न तं निद्रा निशासु शयनं गतम्} %||4-26-7||

\threelineshloka
{तत्समुत्थेन शोकेन बाष्पोपहतचेतसम्}
{तं शोचमानं काकुत्स्थं नित्यं शोकपरायणम्}
{तुल्यदुःखोऽब्रवीद्भ्राता लक्ष्मणोऽनुनयन्वचः} %||4-26-8||

\twolineshloka
{अलं वीर व्यथां गत्वा न त्वं शोचितुमर्हसि}
{शोचतो ह्यवसीदन्ति सर्वार्था विदितं हि ते} %||4-26-9||

\twolineshloka
{भवान्क्रियापरो लोके भवान्देवपरायणः}
{आस्तिको धर्मशीलश्च व्यवसायी च राघव} %||4-26-10||

\twolineshloka
{न ह्यव्यवसितः शत्रुं राक्षसं तं विशेषतः}
{समर्थस्त्वं रणे हन्तुं विक्रमैर्जिह्मकारिणम्} %||4-26-11||

\twolineshloka
{समुन्मूलय शोकं त्वं व्यवसायं स्थिरं कुरु}
{ततः सपरिवारं तं निर्मूलं कुरु राक्षसम्} %||4-26-12||

\twolineshloka
{पृथिवीमपि काकुत्स्थ ससागरवनाचलाम्}
{परिवर्तयितुं शक्तः किमङ्ग पुन रावणम्} %||4-26-13||

\twolineshloka
{अहं तु खलु ते वीर्यं प्रसुप्तं प्रतिबोधये}
{दीप्तैराहुतिभिः काले भस्मच्छन्नमिवानलम्} %||4-26-14||

\twolineshloka
{लक्ष्मणस्य तु तद्वाक्यं प्रतिपूज्य हितं शुभम्}
{राघवः सुहृदं स्निग्धमिदं वचनमब्रवीत्} %||4-26-15||

\twolineshloka
{वाच्यं यदनुरक्तेन स्निग्धेन च हितेन च}
{सत्यविक्रम युक्तेन तदुक्तं लक्ष्मण त्वया} %||4-26-16||

\twolineshloka
{एष शोकः परित्यक्तः सर्वकार्यावसादकः}
{विक्रमेष्वप्रतिहतं तेजः प्रोत्साहयाम्यहम्} %||4-26-17||

\twolineshloka
{शरत्कालं प्रतीक्षेऽहमियं प्रावृडुपस्थिता}
{ततः सराष्ट्रं सगणं राक्षसं तं निहन्म्यहम्} %||4-26-18||

\twolineshloka
{तस्य तद्वचनं श्रुत्वा हृष्टो रामस्य लक्ष्मणः}
{पुनरेवाब्रवीद्वाक्यं सौमित्रिर्मित्रनन्दनः} %||4-26-19||

\twolineshloka
{एतत्ते सदृशं वाक्यमुक्तं शत्रुनिबर्हण}
{इदानीमसि काकुत्स्थ प्रकृतिं स्वामुपागतः} %||4-26-20||

\twolineshloka
{विज्ञाय ह्यात्मनो वीर्यं तथ्यं भवितुमर्हसि}
{एतत्सदृशमुक्तं ते श्रुतस्याभिजनस्य च} %||4-26-21||

\twolineshloka
{तस्मात्पुरुषशार्दूल चिन्तयञ्शत्रुनिग्रहम्}
{वर्षारात्रमनुप्राप्तमतिक्रामय राघव} %||4-26-22||

\fourlineindentedshloka
{नियम्य कोपं प्रतिपाल्यतां शरत्}
{ क्षमस्व मासांश्चतुरो मया सह}
{वसाचलेऽस्मिन्मृगराजसेविते}
{ संवर्धयञ्शत्रुवधे समुद्यतः} %||4-27-23||


{॥इत्यार्षे श्रीमद्रामायणे वाल्मीकीये आदिकाव्ये किष्किन्धाकाण्डे षड्विंशः सर्गः॥}

\sect{सप्तविंशः सर्गः}

\twolineshloka
{स तदा वालिनं हत्वा सुग्रीवमभिषिच्य च}
{वसन्माल्यवतः पृष्टे रामो लक्ष्मणमब्रवीत्} %||4-27-1||

\twolineshloka
{अयं स कालः सम्प्राप्तः समयोऽद्य जलागमः}
{सम्पश्य त्वं नभो मेघैः संवृतं गिरिसंनिभैः} %||4-27-2||

\twolineshloka
{नव मास धृतं गर्भं भास्कारस्य गभस्तिभिः}
{पीत्वा रसं समुद्राणां द्यौः प्रसूते रसायनम्} %||4-27-3||

\twolineshloka
{शक्यमम्बरमारुह्य मेघसोपानपङ्क्तिभिः}
{कुटजार्जुनमालाभिरलङ्कर्तुं दिवाकरम्} %||4-27-4||

\twolineshloka
{सन्ध्यारागोत्थितैस्ताम्रैरन्तेष्वधिकपाण्डुरैः}
{स्निग्धैरभ्रपटच्छदैर्बद्धव्रणमिवाम्बरम्} %||4-27-5||

\twolineshloka
{मन्दमारुतनिःश्वासं सन्ध्याचन्दनरञ्जितम्}
{आपाण्डुजलदं भाति कामातुरमिवाम्बरम्} %||4-27-6||

\twolineshloka
{एषा धर्मपरिक्लिष्टा नववारिपरिप्लुता}
{सीतेव शोकसन्तप्ता मही बाष्पं विमुञ्चति} %||4-27-7||

\twolineshloka
{मेघोदरविनिर्मुक्ताः कह्लारसुखशीतलाः}
{शक्यमञ्जलिभिः पातुं वाताः केतकिगन्धिनः} %||4-27-8||

\twolineshloka
{एष फुल्लार्जुनः शैलः केतकैरधिवासितः}
{सुग्रीव इव शान्तारिर्धाराभिरभिषिच्यते} %||4-27-9||

\twolineshloka
{मेघकृष्णाजिनधरा धारायज्ञोपवीतिनः}
{मारुतापूरितगुहाः प्राधीता इव पर्वताः} %||4-27-10||

\twolineshloka
{कशाभिरिव हैमीभिर्विद्युद्भिरिव ताडितम्}
{अन्तःस्तनितनिर्घोषं सवेदनमिवाम्बरम्} %||4-27-11||

\twolineshloka
{नीलमेघाश्रिता विद्युत्स्फुरन्ती प्रतिभाति मे}
{स्फुरन्ती रावणस्याङ्के वैदेहीव तपस्विनी} %||4-27-12||

\twolineshloka
{इमास्ता मन्मथवतां हिताः प्रतिहता दिशः}
{अनुलिप्ता इव घनैर्नष्टग्रहनिशाकराः} %||4-27-13||

\threelineshloka
{क्वचिद्बाष्पाभिसंरुद्धान्वर्षागमसमुत्सुकान्}
{कुटजान्पश्य सौमित्रे पुष्टितान्गिरिसानुषु}
{मम शोकाभिभूतस्य कामसन्दीपनान्स्थितान्} %||4-27-14||

\fourlineindentedshloka
{रजः प्रशान्तं सहिमोऽद्य वायुर्-}
{ निदाघदोषप्रसराः प्रशान्ताः}
{स्थिता हि यात्रा वसुधाधिपानाम्}
{ प्रवासिनो यान्ति नराः स्वदेशान्} %||4-27-15||

\fourlineindentedshloka
{सम्प्रस्थिता मानसवासलुब्धाः}
{ प्रियान्विताः सम्प्रति चक्रवाकः}
{अभीक्ष्णवर्षोदकविक्षतेषु}
{ यानानि मार्गेषु न सम्पतन्ति} %||4-27-16||

\fourlineindentedshloka
{क्वचित्प्रकाशं क्वचिदप्रकाशम्}
{ नभः प्रकीर्णाम्बुधरं विभाति}
{क्वचित्क्वचित्पर्वतसंनिरुद्धम्}
{ रूपं यथा शान्तमहार्णवस्य} %||4-27-17||

\fourlineindentedshloka
{व्यामिश्रितं सर्जकदम्बपुष्पैर्-}
{ नवं जलं पर्वतधातुताम्रम्}
{मयूरकेकाभिरनुप्रयातम्}
{ शैलापगाः शीघ्रतरं वहन्ति} %||4-27-18||

\fourlineindentedshloka
{रसाकुलं षट्पदसंनिकाशम्}
{ प्रभुज्यते जम्बुफलं प्रकामम्}
{अनेकवर्णं पवनावधूतम्}
{ भूमौ पतत्याम्रफलं विपक्वम्} %||4-27-19||

\fourlineindentedshloka
{विद्युत्पताकाः सबलाक मालाः}
{ शैलेन्द्रकूटाकृतिसंनिकाशाः}
{गर्जन्ति मेघाः समुदीर्णनादा}
{ मत्तगजेन्द्रा इव संयुगस्थः} %||4-27-20||

\fourlineindentedshloka
{मेघाभिकामी परिसम्पतन्ती}
{ सम्मोदिता भाति बलाकपङ्क्तिः}
{वातावधूता वरपौण्डरीकी}
{ लम्बेव माला रचिताम्बरस्य} %||4-27-21||

\fourlineindentedshloka
{निद्रा शनैः केशवमभ्युपैति}
{ द्रुतं नदी सागरमभ्युपैति}
{हृष्टा बलाका घनमभ्युपैति}
{ कान्ता सकामा प्रियमभ्युपैति} %||4-27-22||

\fourlineindentedshloka
{जाता वनान्ताः शिखिसुप्रनृत्ता}
{ जाताः कदम्बाः सकदम्बशाखाः}
{जाता वृषा गोषु समानकामा}
{ जाता मही सस्यवनाभिरामा} %||4-27-23||

\fourlineindentedshloka
{वहन्ति वर्षन्ति नदन्ति भान्ति}
{ ध्यायन्ति नृत्यन्ति समाश्वसन्ति}
{नद्यो घना मत्तगजा वनान्ताः}
{ प्रियाविनीहाः शिखिनः प्लवङ्गाः} %||4-27-24||

\fourlineindentedshloka
{प्रहर्षिताः केतकपुष्पगन्धम्}
{ आघ्राय हृष्टा वननिर्झरेषु}
{प्रपात शब्दाकुलिता गजेन्द्राः}
{ सार्धं मयूरैः समदा नदन्ति} %||4-27-25||

\fourlineindentedshloka
{धारानिपातैरभिहन्यमानाः}
{ कदम्बशाखासु विलम्बमानाः}
{क्षणार्जितं पुष्परसावगाढम्}
{ शनैर्मदं षट्चरणास्त्यजन्ति} %||4-27-26||

\fourlineindentedshloka
{अङ्गारचूर्णोत्करसंनिकाशैः}
{ फलैः सुपर्याप्त रसैः समृद्धैः}
{जम्बूद्रुमाणां प्रविभान्ति शाखा}
{ निलीयमाना इव षट्पदौघैः} %||4-27-27||

\fourlineindentedshloka
{तडित्पताकाभिरलङ्कृतानाम्}
{ उदीर्णगम्भीरमहारवाणाम्}
{विभान्ति रूपाणि बलाहकानाम्}
{ रणोद्यतानामिव वारणानाम्} %||4-27-28||

\fourlineindentedshloka
{मार्गानुगः शैलवनानुसारी}
{ सम्प्रस्थितो मेघरवं निशम्य}
{युद्धाभिकामः प्रतिनागशङ्की}
{ मत्तो गजेन्द्रः प्रतिसंनिवृत्तः} %||4-27-29||

\fourlineindentedshloka
{मुक्तासकाशं सलिलं पतद्वै}
{ सुनिर्मलं पत्रपुटेषु लग्नम्}
{हृष्टा विवर्णच्छदना विहङ्गाः}
{ सुरेन्द्रदत्तं तृषिताः पिबन्ति} %||4-27-30||

\fourlineindentedshloka
{नीलेषु नीला नववारिपूर्णा}
{ मेघेषु मेघाः प्रविभान्ति सक्ताः}
{दवाग्निदग्धेषु दवाग्निदग्धाः}
{ शैलेषु शैला इव बद्धमूलाः} %||4-27-31||

\fourlineindentedshloka
{मत्ता गजेन्द्रा मुदिता गवेन्द्रा}
{ वनेषु विश्रान्ततरा मृगेन्द्राः}
{रम्या नगेन्द्रा निभृता नगेन्द्राः}
{ प्रक्रीडितो वारिधरैः सुरेन्द्रः} %||4-27-32||

\twolineshloka
{वृत्ता यात्रा नरेन्द्राणां सेना प्रतिनिवर्तते}
{वैराणि चैव मार्गाश्च सलिलेन समीकृताः} %||4-27-33||

\twolineshloka
{मासि प्रौष्ठपदे ब्रह्म ब्राह्मणानां विवक्षताम्}
{अयमध्यायसमयः सामगानामुपस्थितः} %||4-27-34||

\twolineshloka
{निवृत्तकर्मायतनो नूनं सञ्चितसञ्चयः}
{आषाढीमभ्युपगतो भरतः कोषकाधिपः} %||4-27-35||

\twolineshloka
{नूनमापूर्यमाणायाः सरय्वा वधते रयः}
{मां समीक्ष्य समायान्तमयोध्याया इव स्वनः} %||4-27-36||

\twolineshloka
{इमाः स्फीतगुणा वर्षाः सुग्रीवः सुखमश्नुते}
{विजितारिः सदारश्च राज्ये महति च स्थितः} %||4-27-37||

\twolineshloka
{अहं तु हृतदारश्च राज्याच्च महतश्च्युतः}
{नदीकूलमिव क्लिन्नमवसीदामि लक्ष्मण} %||4-27-38||

\twolineshloka
{शोकश्च मम विस्तीर्णो वर्षाश्च भृशदुर्गमाः}
{रावणश्च महाञ्शत्रुरपारं प्रतिभाति मे} %||4-27-39||

\twolineshloka
{अयात्रां चैव दृष्ट्वेमां मार्गांश्च भृशदुर्गमान्}
{प्रणते चैव सुग्रीवे न मया किञ्चिदीरितम्} %||4-27-40||

\twolineshloka
{अपि चातिपरिक्लिष्टं चिराद्दारैः समागतम्}
{आत्मकार्यगरीयस्त्वाद्वक्तुं नेच्छामि वानरम्} %||4-27-41||

\twolineshloka
{स्वयमेव हि विश्रम्य ज्ञात्वा कालमुपागतम्}
{उपकारं च सुग्रीवो वेत्स्यते नात्र संशयः} %||4-27-42||

\twolineshloka
{तस्मात्कालप्रतीक्षोऽहं स्थितोऽस्मि शुभलक्षण}
{सुग्रीवस्य नदीनां च प्रसादमनुपालयन्} %||4-27-43||

\twolineshloka
{उपकारेण वीरो हि प्रतिकारेण युज्यते}
{अकृतज्ञोऽप्रतिकृतो हन्ति सत्त्ववतां मनः} %||4-27-44||

\fourlineindentedshloka
{अथैवमुक्तः प्रणिधाय लक्ष्मणः}
{ कृताञ्जलिस्तत्प्रतिपूज्य भाषितम्}
{उवाच रामं स्वभिराम दर्शनम्}
{ प्रदर्शयन्दर्शनमात्मनः शुभम्} %||4-27-45||

\fourlineindentedshloka
{यथोक्तमेतत्तव सर्वमीप्सितम्}
{ नरेन्द्र कर्ता नचिराद्धरीश्वरः}
{शरत्प्रतीक्षः क्षमतामिमं भवा}
{ञ्जलप्रपातं रिपुनिग्रहे धृतः} %||4-28-46||


{॥इत्यार्षे श्रीमद्रामायणे वाल्मीकीये आदिकाव्ये किष्किन्धाकाण्डे सप्तविंशः सर्गः॥}

\sect{अष्टाविंशः सर्गः}

\twolineshloka
{समीक्ष्य विमलं व्योम गतविद्युद्बलाहकम्}
{सारसारवसङ्घुष्टं रम्यज्योत्स्नानुलेपनम्} %||4-28-1||

\twolineshloka
{समृद्धार्थं च सुग्रीवं मन्दधर्मार्थसङ्ग्रहम्}
{अत्यर्थमसतां मार्गमेकान्तगतमानसम्} %||4-28-2||

\twolineshloka
{निवृत्तकार्यं सिद्धार्थं प्रमदाभिरतं सदा}
{प्राप्तवन्तमभिप्रेतान्सर्वानेव मनोरथान्} %||4-28-3||

\twolineshloka
{स्वां च पात्नीमभिप्रेतां तारां चापि समीप्सिताम्}
{विहरन्तमहोरात्रं कृतार्थं विगतज्वलम्} %||4-28-4||

\twolineshloka
{क्रीडन्तमिव देवेशं नन्दनेऽप्सरसां गणैः}
{मन्त्रिषु न्यस्तकार्यं च मन्त्रिणामनवेक्षकम्} %||4-28-5||

\twolineshloka
{उत्सन्नराज्यसन्देशं कामवृत्तमवस्थितम्}
{निश्चितार्थोऽर्थतत्त्वज्ञः कालधर्मविशेषवित्} %||4-28-6||

\twolineshloka
{प्रसाद्य वाक्यैर्मधुरैर्हेतुमद्भिर्मनोरमैः}
{वाक्यविद्वाक्यतत्त्वज्ञं हरीशं मारुतात्मजः} %||4-28-7||

\threelineshloka
{हितं तथ्यं च पथ्यं च सामधर्मार्थनीतिमत्}
{प्रणयप्रीतिसंयुक्तं विश्वासकृतनिश्चयम्}
{हरीश्वरमुपागम्य हनुमान्वाक्यमब्रवीत्} %||4-28-8||

\twolineshloka
{राज्यं प्राप्तं यशश्चैव कौली श्रीरभिवर्थिता}
{मित्राणां सङ्ग्रहः शेषस्तद्भवान्कर्तुमर्हति} %||4-28-9||

\twolineshloka
{यो हि मित्रेषु कालज्ञः सततं साधु वर्तते}
{तस्य राज्यं च कीर्तिश्च प्रतापश्चाभिवर्धते} %||4-28-10||

\twolineshloka
{यस्य कोशश्च दण्डश्च मित्राण्यात्मा च भूमिप}
{समवेतानि सर्वाणि स राज्यं महदश्नुते} %||4-28-11||

\twolineshloka
{तद्भवान्वृत्तसम्पन्नः स्थितः पथि निरत्यये}
{मित्रार्थमभिनीतार्थं यथावत्कर्तुमर्हति} %||4-28-12||

\twolineshloka
{यस्तु कालव्यतीतेषु मित्रकार्येषु वर्तते}
{स कृत्वा महतोऽप्यर्थान्न मित्रार्थेन युज्यते} %||4-28-13||

\twolineshloka
{क्रियतां राघवस्यैतद्वैदेह्याः परिमार्गणम्}
{तदिदं वीर कार्यं ते कालातीतमरिन्दम} %||4-28-14||

\twolineshloka
{न च कालमतीतं ते निवेदयति कालवित्}
{त्वरमाणोऽपि सन्प्राज्ञस्तव राजन्वशानुगः} %||4-28-15||

\twolineshloka
{कुलस्य केतुः स्फीतस्य दीर्घबन्धुश्च राघवः}
{अप्रमेयप्रभावश्च स्वयं चाप्रतिमो गुणैः} %||4-28-16||

\twolineshloka
{तस्य त्वं कुरु वै कार्यं पूर्वं तेन कृतं तव}
{हरीश्वर हरिश्रेष्ठानाज्ञापयितुमर्हसि} %||4-28-17||

\twolineshloka
{न हि तावद्भवेत्कालो व्यतीतश्चोदनादृते}
{चोदितस्य हि कार्यस्य भवेत्कालव्यतिक्रमः} %||4-28-18||

\twolineshloka
{अकर्तुरपि कार्यस्य भवान्कर्ता हरीश्वर}
{किं पुनः प्रतिकर्तुस्ते राज्येन च धनेन च} %||4-28-19||

\twolineshloka
{शक्तिमानसि विक्रान्तो वानरर्ष्क गणेश्वर}
{कर्तुं दाशरथेः प्रीतिमाज्ञायां किं नु सज्जसे} %||4-28-20||

\twolineshloka
{कामं खलु शरैर्शक्तः सुरासुरमहोरगान्}
{वशे दाशरथिः कर्तुं त्वत्प्रतिज्ञां तु काङ्क्षते} %||4-28-21||

\twolineshloka
{प्राणत्यागाविशङ्केन कृतं तेन तव प्रियम्}
{तस्य मार्गाम वैदेहीं पृथिव्यामपि चाम्बरे} %||4-28-22||

\twolineshloka
{न देवा न च गन्धर्वा नासुरा न मरुद्गणाः}
{न च यक्षा भयं तस्य कुर्युः किमुत राक्षसाः} %||4-28-23||

\twolineshloka
{तदेवं शक्तियुक्तस्य पूर्वं प्रियकृतस्तथा}
{रामस्यार्हसि पिङ्गेश कर्तुं सर्वात्मना प्रियम्} %||4-28-24||

\twolineshloka
{नाधस्तादवनौ नाप्सु गतिर्नोपरि चाम्बरे}
{कस्यचित्सज्जतेऽस्माकं कपीश्वर तवाज्ञया} %||4-28-25||

\twolineshloka
{तदाज्ञापय कः किं ते कृते वसतु कुत्रचित्}
{हरयो ह्यप्रधृष्यास्ते सन्ति कोट्यग्रतोऽनघ} %||4-28-26||

\twolineshloka
{तस्य तद्वचनं श्रुत्वा काले साधुनिवेदितम्}
{सुग्रीवः सत्त्वसम्पन्नश्चकार मतिमुत्तमाम्} %||4-28-27||

\twolineshloka
{स सन्दिदेशाभिमतं नीलं नित्यकृतोद्यमम्}
{दिक्षु सर्वासु सर्वेषां सैन्यानामुपसङ्ग्रहे} %||4-28-28||

\twolineshloka
{यथा सेना समग्रा मे यूथपालाश्च सर्वशः}
{समागच्छन्त्यसङ्गेन सेनाग्राणि तथा कुरु} %||4-28-29||

\threelineshloka
{ये त्वन्तपालाः प्लवगाः शीघ्रगा व्यवसायिनः}
{समानयन्तु ते सैन्यं त्वरिताः शासनान्मम}
{स्वयं चानन्तरं सैन्यं भवानेवानुपश्यतु} %||4-28-30||

\twolineshloka
{त्रिपञ्चरात्रादूर्ध्वं यः प्राप्नुयान्नेह वानरः}
{तस्य प्राणान्तिको दण्डो नात्र कार्या विचारणा} %||4-28-31||

\fourlineindentedshloka
{हरींश्च वृद्धानुपयातु साङ्गदो}
{ भवान्ममाज्ञामधिकृत्य निश्चिताम्}
{इति व्यवस्थां हरिपुङ्गवेश्वरो}
{ विधाय वेश्म प्रविवेश वीर्यवान्} %||4-29-32||


{॥इत्यार्षे श्रीमद्रामायणे वाल्मीकीये आदिकाव्ये किष्किन्धाकाण्डे अष्टाविंशः सर्गः॥}

\sect{एकोनत्रिंशः सर्गः}

\twolineshloka
{गुहां प्रविष्टे सुग्रीवे विमुक्ते गगने घनैः}
{वर्षरात्रोषितो रामः कामशोकाभिपीडितः} %||4-29-1||

\twolineshloka
{पाण्डुरं गगनं दृष्ट्वा विमलं चन्द्रमण्डलम्}
{शारदीं रजनीं चैव दृष्ट्वा ज्योत्स्नानुलेपनाम्} %||4-29-2||

\twolineshloka
{कामवृत्तं च सुग्रीवं नष्टां च जनकात्मजाम्}
{बुद्ध्वा कालमतीतं च मुमोह परमातुरः} %||4-29-3||

\twolineshloka
{स तु संज्ञामुपागम्य मुहूर्तान्मतिमान्पुनः}
{मनःस्थामपि वैदेहीं चिन्तयामास राघवः} %||4-29-4||

\twolineshloka
{आसीनः पर्वतस्याग्रे हेमधातुविभूषिते}
{शारदं गगनं दृष्ट्व जगाम मनसा प्रियाम्} %||4-29-5||

\twolineshloka
{दृष्ट्वा च विमलं व्योम गतविद्युद्बलाहकम्}
{सारसारवसङ्घुष्टं विललापार्तया गिरा} %||4-29-6||

\twolineshloka
{सारसारवसंनादैः सारसारवनादिनी}
{याश्रमे रमते बाला साद्य मे रमते कथम्} %||4-29-7||

\twolineshloka
{पुष्पितांश्चासनान्दृष्ट्वा काञ्चनानिव निर्मलान्}
{कथं स रमते बाला पश्यन्ती मामपश्यती} %||4-29-8||

\twolineshloka
{या पुरा कलहंसानां स्वरेण कलभाषिणी}
{बुध्यते चारुसर्वाङ्गी साद्य मे बुध्यते कथम्} %||4-29-9||

\twolineshloka
{निःस्वनं चक्रवाकानां निशम्य सहचारिणाम्}
{पुण्डरीकविशालाक्षी कथमेषा भविष्यति} %||4-29-10||

\twolineshloka
{सरांसि सरितो वापीः काननानि वनानि च}
{तां विना मृगशावाक्षीं चरन्नाद्य सुखं लभे} %||4-29-11||

\twolineshloka
{अपि तां मद्वियोगाच्च सौकुमार्याच्च भामिनीम्}
{न दूरं पीडयेत्कामः शरद्गुणनिरन्तरः} %||4-29-12||

\twolineshloka
{एवमादि नरश्रेष्ठो विललाप नृपात्मजः}
{विहङ्ग इव सारङ्गः सलिलं त्रिदशेश्वरात्} %||4-29-13||

\twolineshloka
{ततश्चञ्चूर्य रम्येषु फलार्थी गिरिसानुषु}
{ददर्श पर्युपावृत्तो लक्ष्मीवाँल्लक्ष्मणोऽग्रजम्} %||4-29-14||

\fourlineindentedshloka
{तं चिन्तया दुःसहया परीतम्}
{ विसंज्ञमेकं विजने मनस्वी}
{भ्रातुर्विषादात्परितापदीनः}
{ समीक्ष्य सौमित्रिरुवाच रामम्} %||4-29-15||

\fourlineindentedshloka
{किमार्य कामस्य वशङ्गतेन}
{ किमात्मपौरुष्यपराभवेन}
{अयं सदा संहृइयते समाधिः}
{ किमत्र योगेन निवर्तितेन} %||4-29-16||

\fourlineindentedshloka
{क्रियाभियोगं मनसः प्रसादम्}
{ समाधियोगानुगतं च कालम्}
{सहायसामर्थ्यमदीनसत्त्व}
{ स्वकर्महेतुं च कुरुष्व हेतुम्} %||4-29-17||

\fourlineindentedshloka
{न जानकी मानववंशनाथ}
{ त्वया सनाथा सुलभा परेण}
{न चाग्निचूडां ज्वलितामुपेत्य}
{ न दह्यते वीरवरार्ह कश्चित्} %||4-29-18||

\fourlineindentedshloka
{सलक्ष्मणं लक्ष्मणमप्रधृष्यम्}
{ स्वभावजं वाक्यमुवाच रामः}
{हितं च पथ्यं च नयप्रसक्तम्}
{ ससामधर्मार्थसमाहितं च} %||4-29-19||

\fourlineindentedshloka
{निःसंशयं कार्यमवेक्षितव्यम्}
{ क्रियाविशेषो ह्यनुवर्तितव्यः}
{ननु प्रवृत्तस्य दुरासदस्य}
{ कुमारकार्यस्य फलं न चिन्त्यम्} %||4-29-20||

\twolineshloka
{अथ पद्मपलाशाक्षीं मैथिलीमनुचिन्तयन्}
{उवाच लक्ष्मणं रामो मुखेन परिशुष्यता} %||4-29-21||

\twolineshloka
{तर्पयित्वा सहस्राक्षः सलिलेन वसुन्धराम्}
{निर्वर्तयित्वा सस्यानि कृतकर्मा व्यवस्थितः} %||4-29-22||

\twolineshloka
{स्निग्धगम्भीरनिर्घोषाः शैलद्रुमपुरोगमाः}
{विसृज्य सलिलं मेघाः परिश्रान्ता नृपात्मज} %||4-29-23||

\twolineshloka
{नीलोत्पलदलश्यामः श्यामीकृत्वा दिशो दश}
{विमदा इव मातङ्गाः शान्तवेगाः पयोधराः} %||4-29-24||

\twolineshloka
{जलगर्भा महावेगाः कुटजार्जुनगन्धिनः}
{चरित्वा विरताः सौम्य वृष्टिवाताः समुद्यताः} %||4-29-25||

\twolineshloka
{घनानां वारणानां च मयूराणां च लक्ष्मण}
{नादः प्रस्रवणानां च प्रशान्तः सहसानघ} %||4-29-26||

\twolineshloka
{अभिवृष्टा महामेघैर्निर्मलाश्चित्रसानवः}
{अनुलिप्ता इवाभान्ति गिरयश्चन्द्ररश्मिभिः} %||4-29-27||

\twolineshloka
{दर्शयन्ति शरन्नद्यः पुलिनानि शनैः शनैः}
{नवसङ्गमसव्रीडा जघनानीव योषितः} %||4-29-28||

\twolineshloka
{प्रसन्नसलिलाः सौम्य कुररीभिर्विनादिताः}
{चक्रवाकगणाकीर्णा विभान्ति सलिलाशयाः} %||4-29-29||

\twolineshloka
{अन्योन्यबद्धवैराणां जिगीषूणां नृपात्मज}
{उद्योगसमयः सौम्य पार्थिवानामुपस्थितः} %||4-29-30||

\twolineshloka
{इयं सा प्रथमा यात्रा पार्थिवानां नृपात्मज}
{न च पश्यामि सुग्रीवमुद्योगं वा तथाविधम्} %||4-29-31||

\twolineshloka
{चत्वारो वार्षिका मासा गता वर्षशतोपमाः}
{मम शोकाभितप्तस्य सौम्य सीतामपश्यतः} %||4-29-32||

\twolineshloka
{प्रियाविहीने दुःखार्ते हृतराज्ये विवासिते}
{कृपां न कुरुते राजा सुग्रीवो मयि लक्ष्मण} %||4-29-33||

\twolineshloka
{अनाथो हृतराज्योऽयं रावणेन च धर्षितः}
{दीनो दूरगृहः कामी मां चैव शरणं गतः} %||4-29-34||

\twolineshloka
{इत्येतैः कारणैः सौम्य सुग्रीवस्य दुरात्मनः}
{अहं वानरराजस्य परिभूतः परन्तप} %||4-29-35||

\twolineshloka
{स कालं परिसङ्ख्याय सीतायाः परिमार्गणे}
{कृतार्थः समयं कृत्वा दुर्मतिर्नावबुध्यते} %||4-29-36||

\twolineshloka
{त्वं प्रविश्य च किष्किन्धां ब्रूहि वानरपुङ्गवम्}
{मूर्खं ग्राम्य सुखे सक्तं सुग्रीवं वचनान्मम} %||4-29-37||

\twolineshloka
{अर्थिनामुपपन्नानां पूर्वं चाप्युपकारिणाम्}
{आशां संश्रुत्य यो हन्ति स लोके पुरुषाधमः} %||4-29-38||

\twolineshloka
{शुभं वा यदि वा पापं यो हि वाक्यमुदीरितम्}
{सत्येन परिगृह्णाति स वीरः पुरुषोत्तमः} %||4-29-39||

\twolineshloka
{कृतार्था ह्यकृतार्थानां मित्राणां न भवन्ति ये}
{तान्मृतानपि क्रव्यादः कृतघ्नान्नोपभुञ्जते} %||4-29-40||

\twolineshloka
{नूनं काञ्चनपृष्ठस्य विकृष्टस्य मया रणे}
{द्रष्टुमिच्छन्ति चापस्य रूपं विद्युद्गणोपमम्} %||4-29-41||

\twolineshloka
{घोरं ज्यातलनिर्घोषं क्रुद्धस्य मम संयुगे}
{निर्घोषमिव वज्रस्य पुनः संश्रोतुमिच्छति} %||4-29-42||

\twolineshloka
{काममेवं गतेऽप्यस्य परिज्ञाते पराक्रमे}
{त्वत्सहायस्य मे वीर न चिन्ता स्यान्नृपात्मज} %||4-29-43||

\twolineshloka
{यदर्थमयमारम्भः कृतः परपुरंजय}
{समयं नाभिजानाति कृतार्थः प्लवगेश्वरः} %||4-29-44||

\twolineshloka
{वर्षासमयकालं तु प्रतिज्ञाय हरीश्वरः}
{व्यतीतांश्चतुरो मासान्विहरन्नावबुध्यते} %||4-29-45||

\twolineshloka
{सामात्यपरिषत्क्रीडन्पानमेवोपसेवते}
{शोकदीनेषु नास्मासु सुग्रीवः कुरुते दयाम्} %||4-29-46||

\twolineshloka
{उच्यतां गच्छ सुग्रीवस्त्वया वत्स महाबल}
{मम रोषस्य यद्रूपं ब्रूयाश्चैनमिदं वचः} %||4-29-47||

\twolineshloka
{न च सङ्कुचितः पन्था येन वाली हतो गतः}
{समये तिष्ठ सुग्रीवमा वालिपथमन्वगाः} %||4-29-48||

\twolineshloka
{एक एव रणे वाली शरेण निहतो मया}
{त्वां तु सत्यादतिक्रान्तं हनिष्यामि सबान्धवम्} %||4-29-49||

\twolineshloka
{तदेवं विहिते कार्ये यद्धितं पुरुषर्षभ}
{तत्तद्ब्रूहि नरश्रेष्ठ त्वर कालव्यतिक्रमः} %||4-29-50||

\fourlineindentedshloka
{कुरुष्व सत्यं मयि वानरेश्वर}
{ प्रतिश्रुतं धर्ममवेक्ष्य शाश्वतम्}
{मा वालिनं प्रेत्य गतो यमक्षयम्}
{ त्वमद्य पश्येर्मम चोदितैः शरैः} %||4-29-51||

\fourlineindentedshloka
{स पूर्वजं तीव्रविवृद्धकोपम्}
{ लालप्यमानं प्रसमीक्ष्य दीनम्}
{चकार तीव्रां मतिमुग्रतेजा}
{ हरीश्वरमानववंशनाथः} %||4-30-52||


{॥इत्यार्षे श्रीमद्रामायणे वाल्मीकीये आदिकाव्ये किष्किन्धाकाण्डे एकोनत्रिंशः सर्गः॥}

\sect{त्रिंशः सर्गः}

\fourlineindentedshloka
{स कामिनं दीनमदीनसत्त्वः}
{ शोकाभिपन्नं समुदीर्णकोपम्}
{नरेन्द्रसूनुर्नरदेवपुत्रम्}
{ रामानुजः पूर्वजमित्युवाच} %||4-30-1||

\fourlineindentedshloka
{न वानरः स्थास्यति साधुवृत्ते}
{ न मंस्यते कार्यफलानुषङ्गान्}
{न भक्ष्यते वानरराज्यलक्ष्मीम्}
{ तथा हि नाभिक्रमतेऽस्य बुद्धिः} %||4-30-2||

\fourlineindentedshloka
{मतिक्षयाद्ग्राम्यसुखेषु सक्तः}
{ तव प्रसादाप्रतिकारबुद्धिः}
{हतोऽग्रजं पश्यतु वालिनं स}
{ न राज्यमेवं विगुणस्य देयम्} %||4-30-3||

\fourlineindentedshloka
{न धारये कोपमुदीर्णवेगम्}
{ निहन्मि सुग्रीवमसत्यमद्य}
{हरिप्रवीरैः सह वालिपुत्रो}
{ नरेन्द्रपत्न्या विचयं करोतु} %||4-30-4||

\fourlineindentedshloka
{तमात्तबाणासनमुत्पतन्तम्}
{ निवेदितार्थं रणचण्डकोपम्}
{उवच रामः परवीरहन्ता}
{ स्ववेक्षितं सानुनयं च वाक्यम्} %||4-30-5||

\twolineshloka
{न हि वै त्वद्विधो लोके पापमेवं समाचरेत्}
{पापमार्येण यो हन्ति स वीरः पुरुषोत्तमः} %||4-30-6||

\twolineshloka
{नेदमद्य त्वया ग्राह्यं साधुवृत्तेन लक्ष्मण}
{तां प्रीतिमनुवर्तस्व पूर्ववृत्तं च सङ्गतम्} %||4-30-7||

\twolineshloka
{सामोपहितया वाचा रूक्षाणि परिवर्जयन्}
{वक्तुमर्हसि सुग्रीवं व्यतीतं कालपर्यये} %||4-30-8||

\twolineshloka
{सोऽ ग्रजेनानुशिष्टार्थो यथावत्पुरुषर्षभः}
{प्रविवेश पुरीं वीरो लक्ष्मणः परवीरहा} %||4-30-9||

\twolineshloka
{ततः शुभमतिः प्राज्ञो भ्रातुः प्रियहिते रतः}
{लक्ष्मणः प्रतिसंरब्धो जगाम भवनं कपेः} %||4-30-10||

\twolineshloka
{शक्रबाणासनप्रख्यं धनुः कालान्तकोपमः}
{प्रगृह्य गिरिशृङ्गाभं मन्दरः सानुमानिव} %||4-30-11||

\twolineshloka
{यथोक्तकारी वचनमुत्तरं चैव सोत्तरम्}
{बृहस्पतिसमो बुद्ध्या मत्त्वा रामानुजस्तदा} %||4-30-12||

\twolineshloka
{कामक्रोधसमुत्थेन भ्रातुः कोपाग्निना वृतः}
{प्रभञ्जन इवाप्रीतः प्रययौ लक्ष्मणस्तदा} %||4-30-13||

\twolineshloka
{सालतालाश्वकर्णांश्च तरसा पातयन्बहून्}
{पर्यस्यन्गिरिकूटानि द्रुमानन्यांश्च वेगतः} %||4-30-14||

\twolineshloka
{शिलाश्च शकलीकुर्वन्पद्भ्यां गज इवाशुगः}
{दूरमेकपदं त्यक्त्वा ययौ कार्यवशाद्द्रुतम्} %||4-30-15||

\twolineshloka
{तामपश्यद्बलाकीर्णां हरिराजमहापुरीम्}
{दुर्गामिक्ष्वाकुशार्दूलः किष्किन्धां गिरिसङ्कटे} %||4-30-16||

\twolineshloka
{रोषात्प्रस्फुरमाणौष्ठः सुग्रीवं प्रति कल्ष्मणः}
{ददर्श वानरान्भीमान्किष्किन्धाया बहिश्चरान्} %||4-30-17||

\twolineshloka
{शैलशृङ्गाणि शतशः प्रवृद्धांश्च महीरुहान्}
{जगृहुः कुञ्जरप्रख्या वानराः पर्वतान्तरे} %||4-30-18||

\twolineshloka
{तान्गृहीतप्रहरणान्हरीन्दृष्ट्वा तु लक्ष्मणः}
{बभूव द्विगुणं क्रुद्धो बह्विन्धन इवानलः} %||4-30-19||

\twolineshloka
{तं ते भयपरीताङ्गाः क्रुद्धं दृष्ट्वा प्लवङ्गमाः}
{कालमृत्युयुगान्ताभं शतशो विद्रुता दिशः} %||4-30-20||

\twolineshloka
{ततः सुग्रीवभवनं प्रविश्य हरिपुङ्गवाः}
{क्रोधमागमनं चैव लक्ष्मणस्य न्यवेदयन्} %||4-30-21||

\twolineshloka
{तारया सहितः कामी सक्तः कपिवृषो रहः}
{न तेषां कपिवीराणां शुश्राव वचनं तदा} %||4-30-22||

\twolineshloka
{ततः सचिवसन्दिष्टा हरयो रोमहर्षणाः}
{गिरिकुञ्जरमेघाभा नगर्या निर्ययुस्तदा} %||4-30-23||

\twolineshloka
{नखदंष्ट्रायुधा घोराः सर्वे विकृतदर्शनाः}
{सर्वे शार्दूलदर्पाश्च सर्वे च विकृताननाः} %||4-30-24||

\twolineshloka
{दशनागबलाः केचित्केचिद्दशगुणोत्तराः}
{केचिन्नागसहस्रस्य बभूवुस्तुल्यविक्रमाः} %||4-30-25||

\twolineshloka
{कृत्स्नां हि कपिभिर्व्याप्तां द्रुमहस्तैर्महाबलैः}
{अपश्यल्लक्ष्मणः क्रुद्धः किष्किन्धां तां दुरासदम्} %||4-30-26||

\twolineshloka
{ततस्ते हरयः सर्वे प्राकारपरिखान्तरात्}
{निष्क्रम्योदग्रसत्त्वास्तु तस्थुराविष्कृतं तदा} %||4-30-27||

\twolineshloka
{सुग्रीवस्य प्रमादं च पूर्वजं चार्तमात्मवान्}
{बुद्ध्वा कोपवशं वीरः पुनरेव जगाम सः} %||4-30-28||

\twolineshloka
{स दीर्घोष्णमहोच्छ्वासः कोपसंरक्तलोचनः}
{बभूव नरशार्दूलसधूम इव पावकः} %||4-30-29||

\twolineshloka
{बाणशल्यस्फुरज्जिह्वः सायकासनभोगवान्}
{स्वतेजोविषसङ्घातः पञ्चास्य इव पन्नगः} %||4-30-30||

\twolineshloka
{तं दीप्तमिव कालाग्निं नागेन्द्रमिव कोपितम्}
{समासाद्याङ्गदस्त्रासाद्विषादमगमद्भृशम्} %||4-30-31||

\twolineshloka
{सोऽङ्गदं रोषताम्राक्षः सन्दिदेश महायशाः}
{सुग्रीवः कथ्यतां वत्स ममागमनमित्युत} %||4-30-32||

\twolineshloka
{एष रामानुजः प्राप्तस्त्वत्सकाशमरिन्दमः}
{भ्रातुर्व्यसनसन्तप्तो द्वारि तिष्ठति लक्ष्मणः} %||4-30-33||

\twolineshloka
{लक्ष्मणस्य वचः श्रुत्वा शोकाविष्टोऽङ्गदोऽब्रवीत्}
{पितुः समीपमागम्य सौमित्रिरयमागतः} %||4-30-34||

\twolineshloka
{ते महौघनिभं दृष्ट्वा वज्राशनिसमस्वनम्}
{सिंहनादं समं चक्रुर्लक्ष्मणस्य समीपतः} %||4-30-35||

\twolineshloka
{तेन शब्देन महता प्रत्यबुध्यत वानरः}
{मदविह्वलताम्राक्षो व्याकुलस्रग्विभूषणः} %||4-30-36||

\twolineshloka
{अथाङ्गदवचः श्रुत्वा तेनैव च समागतौ}
{मन्त्रिणो वानरेन्द्रस्य सम्मतोदारदर्शिनौ} %||4-30-37||

\twolineshloka
{प्लक्षश्चैव प्रभावश्च मन्त्रिणावर्थधर्मयोः}
{वक्तुमुच्चावचं प्राप्तं लक्ष्मणं तौ शशंसतुः} %||4-30-38||

\twolineshloka
{प्रसादयित्वा सुग्रीवं वचनैः सामनिश्चितैः}
{आसीनं पर्युपासीनौ यथा शक्रं मरुत्पतिम्} %||4-30-39||

\twolineshloka
{सत्यसन्धौ महाभागौ भ्रातरौ रामलक्ष्मणौ}
{वयस्य भावं सम्प्राप्तौ राज्यार्हौ राज्यदायिनौ} %||4-30-40||

\twolineshloka
{तयोरेको धनुष्पाणिर्द्वारि तिष्ठति लक्ष्मणः}
{यस्य भीताः प्रवेपन्ते नादान्मुञ्चन्ति वानराः} %||4-30-41||

\twolineshloka
{स एष राघवभ्राता लक्ष्मणो वाक्यसारथिः}
{व्यवसाय रथः प्राप्तस्तस्य रामस्य शासनात्} %||4-30-42||

\twolineshloka
{तस्य मूर्ध्ना प्रणम्य त्वं सपुत्रः सह बन्धुभिः}
{राजंस्तिष्ठ स्वसमये भव सत्यप्रतिश्रवः} %||4-31-43||


{॥इत्यार्षे श्रीमद्रामायणे वाल्मीकीये आदिकाव्ये किष्किन्धाकाण्डे त्रिंशः सर्गः॥}

\sect{एकत्रिंशः सर्गः}

\twolineshloka
{अङ्गदस्य वचः श्रुत्वा सुग्रीवः सचिवैः सह}
{लक्ष्मणं कुपितं श्रुत्वा मुमोचासनमात्मवान्} %||4-31-1||

\twolineshloka
{सचिवानब्रवीद्वाक्यं निश्चित्य गुरुलाघवम्}
{मन्त्रज्ञान्मन्त्रकुशलो मन्त्रेषु परिनिष्ठितः} %||4-31-2||

\twolineshloka
{न मे दुर्व्याहृतं किञ्चिन्नापि मे दुरनुष्ठितम्}
{लक्ष्मणो राघवभ्राता क्रुद्धः किमिति चिन्तये} %||4-31-3||

\twolineshloka
{असुहृद्भिर्ममामित्रैर्नित्यमन्तरदर्शिभिः}
{मम दोषानसम्भूताञ्श्रावितो राघवानुजः} %||4-31-4||

\twolineshloka
{अत्र तावद्यथाबुद्धि सर्वैरेव यथाविधि}
{भवद्भिर्निश्चयस्तस्य विज्ञेयो निपुणं शनैः} %||4-31-5||

\twolineshloka
{न खल्वस्ति मम त्रासो लक्ष्मणान्नापि राघवात्}
{मित्रं त्वस्थान कुपितं जनयत्येव सम्भ्रमम्} %||4-31-6||

\twolineshloka
{सर्वथा सुकरं मित्रं दुष्करं परिपालनम्}
{अनित्यत्वात्तु चित्तानां प्रीतिरल्पेऽपि भिद्यते} %||4-31-7||

\twolineshloka
{अतोनिमित्तं त्रस्तोऽहं रामेण तु महात्मना}
{यन्ममोपकृतं शक्यं प्रतिकर्तुं न तन्मया} %||4-31-8||

\twolineshloka
{सुग्रीवेणैवमुक्तस्तु हनुमान्हरिपुङ्गवः}
{उवाच स्वेन तर्केण मध्ये वानरमन्त्रिणाम्} %||4-31-9||

\twolineshloka
{सर्वथा नैतदाश्चर्यं यत्त्वं हरिगणेश्वर}
{न विस्मरसि सुस्निग्धमुपकारकृतं शुभम्} %||4-31-10||

\twolineshloka
{राघवेण तु शूरेण भयमुत्सृज्य दूरतः}
{त्वत्प्रियार्थं हतो वाली शक्रतुल्यपराक्रमः} %||4-31-11||

\twolineshloka
{सर्वथा प्रणयात्क्रुद्धो राघवो नात्र संशयः}
{भ्रातरं स प्रहितवाँल्लक्ष्मणं लक्ष्मिवर्धनम्} %||4-31-12||

\twolineshloka
{त्वं प्रमत्तो न जानीषे कालं कलविदां वर}
{फुल्लसप्तच्छदश्यामा प्रवृत्ता तु शरच्छिवा} %||4-31-13||

\twolineshloka
{निर्मल ग्रहनक्षत्रा द्यौः प्रनष्टबलाहका}
{प्रसन्नाश्च दिशः सर्वाः सरितश्च सरांसि च} %||4-31-14||

\twolineshloka
{प्राप्तमुद्योगकालं तु नावैषि हरिपुङ्गव}
{त्वं प्रमत्त इति व्यक्तं लक्ष्मणोऽयमिहागतः} %||4-31-15||

\twolineshloka
{आर्तस्य हृतदारस्य परुषं पुरुषान्तरात्}
{वचनं मर्षणीयं ते राघवस्य महात्मनः} %||4-31-16||

\twolineshloka
{कृतापराधस्य हि ते नान्यत्पश्याम्यहं क्षमम्}
{अन्तरेणाञ्जलिं बद्ध्वा लक्ष्मणस्य प्रसादनात्} %||4-31-17||

\twolineshloka
{नियुक्तैर्मन्त्रिभिर्वाच्यो अवश्यं पार्थिवो हितम्}
{अत एव भयं त्यक्त्वा ब्रवीम्यवधृतं वचः} %||4-31-18||

\twolineshloka
{अभिक्रुद्धः समर्थो हि चापमुद्यम्य राघवः}
{सदेवासुरगन्धर्वं वशे स्थापयितुं जगत्} %||4-31-19||

\twolineshloka
{न स क्षमः कोपयितुं यः प्रसाद्य पुनर्भवेत्}
{पूर्वोपकारं स्मरता कृतज्ञेन विशेषतः} %||4-31-20||

\twolineshloka
{तस्य मूर्ध्ना प्रणम्य त्वं सपुत्रः ससुहृज्जनः}
{राजंस्तिष्ठ स्वसमये भर्तुर्भार्येव तद्वशे} %||4-31-21||

\fourlineindentedshloka
{न रामरामानुजशासनं त्वया}
{ कपीन्द्रयुक्तं मनसाप्यपोहितुम्}
{मनो हि ते ज्ञास्यति मानुषं बलम्}
{ सराघवस्यास्य सुरेन्द्रवर्चसः} %||4-32-22||


{॥इत्यार्षे श्रीमद्रामायणे वाल्मीकीये आदिकाव्ये किष्किन्धाकाण्डे एकत्रिंशः सर्गः॥}

\sect{द्वात्रिंशः सर्गः}

\twolineshloka
{अथ प्रतिसमादिष्टो लक्ष्मणः परवीरहा}
{प्रविवेश गुहां घोरां किष्किन्धां रामशासनात्} %||4-32-1||

\twolineshloka
{द्वारस्था हरयस्तत्र महाकाया महाबलाः}
{बभूवुर्लक्ष्मणं दृष्ट्वा सर्वे प्राञ्जलयः स्थिताः} %||4-32-2||

\twolineshloka
{निःश्वसन्तं तु तं दृष्ट्वा क्रुद्धं दशरथात्मजम्}
{बभूवुर्हरयस्त्रस्ता न चैनं पर्यवारयन्} %||4-32-3||

\twolineshloka
{स तं रत्नमयीं श्रीमान्दिव्यां पुष्पितकाननाम्}
{रम्यां रत्नसमाकीर्णां ददर्श महतीं गुहाम्} %||4-32-4||

\twolineshloka
{हर्म्यप्रासादसम्बाधां नानापण्योपशोभिताम्}
{सर्वकामफलैर्वृक्षैः पुष्पितैरुपशोभिताम्} %||4-32-5||

\twolineshloka
{देवगन्धर्वपुत्रैश्च वानरैः कामरूपिभिः}
{दिव्य माल्याम्बरधारैः शोभितां प्रियदर्शनैः} %||4-32-6||

\twolineshloka
{चन्दनागरुपद्मानां गन्धैः सुरभिगन्धिनाम्}
{मैरेयाणां मधूनां च सम्मोदितमहापथाम्} %||4-32-7||

\twolineshloka
{विन्ध्यमेरुगिरिप्रस्थैः प्रासादैर्नैकभूमिभिः}
{ददर्श गिरिनद्यश्च विमलास्तत्र राघवः} %||4-32-8||

\twolineshloka
{अङ्गदस्य गृहं रम्यं मैन्दस्य द्विविदस्य च}
{गवयस्य गवाक्षस्य गजस्य शरभस्य च} %||4-32-9||

\twolineshloka
{विद्युन्मालेश्च सम्पातेः सूर्याक्षस्य हनूमतः}
{वीरबाहोः सुबाहोश्च नलस्य च महात्मनः} %||4-32-10||

\twolineshloka
{कुमुदस्य सुषेणस्य तारजाम्बवतोस्तथा}
{दधिवक्त्रस्य नीलस्य सुपाटलसुनेत्रयोः} %||4-32-11||

\twolineshloka
{एतेषां कपिमुख्यानां राजमार्गे महात्मनाम्}
{ददर्श गृहमुख्यानि महासाराणि लक्ष्मणः} %||4-32-12||

\twolineshloka
{पाण्डुराभ्रप्रकाशानि दिव्यमाल्ययुतानि च}
{प्रभूतधनधान्यानि स्त्रीरत्नैः शोभितानि च} %||4-32-13||

\twolineshloka
{पाण्डुरेण तु शैलेन परिक्षिप्तं दुरासदम्}
{वानरेन्द्रगृहं रम्यं महेन्द्रसदनोपमम्} %||4-32-14||

\twolineshloka
{शुल्कैः प्रासादशिखरैः कैलासशिखरोपमैः}
{सर्वकामफलैर्वृक्षैः पुष्टितैरुपशोभितम्} %||4-32-15||

\twolineshloka
{महेन्द्रदत्तैः श्रीमद्भिर्नीलजीमूतसंनिभैः}
{दिव्यपुष्पफलैर्वृक्षैः शीतच्छायैर्मनोरमैः} %||4-32-16||

\twolineshloka
{हरिभिः संवृतद्वारं बलिभिः शस्त्रपाणिभिः}
{दिव्यमाल्यावृतं शुभ्रं तप्तकाञ्चनतोरणम्} %||4-32-17||

\twolineshloka
{सुग्रीवस्य गृहं रम्यं प्रविवेश महाबलः}
{अवार्यमाणः सौमित्रिर्महाभ्रमिव भास्करः} %||4-32-18||

\twolineshloka
{स सप्त कक्ष्या धर्मात्मा यानासनसमावृताः}
{प्रविश्य सुमहद्गुप्तं ददर्शान्तःपुरं महत्} %||4-32-19||

\twolineshloka
{हैमराजतपर्यङ्कैर्बहुभिश्च वरासनैः}
{महार्हास्तरणोपेतैस्तत्र तत्रोपशोभितम्} %||4-32-20||

\twolineshloka
{प्रविशन्नेव सततं शुश्राव मधुरस्वरम्}
{तन्त्रीगीतसमाकीर्णं समगीतपदाक्षरम्} %||4-32-21||

\twolineshloka
{बह्वीश्च विविधाकारा रूपयौवनगर्विताः}
{स्त्रियः सुग्रीवभवने ददर्श स महाबलः} %||4-32-22||

\twolineshloka
{दृष्ट्वाभिजनसम्पन्नाश्चित्रमाल्यकृतस्रजः}
{वरमाल्यकृतव्यग्रा भूषणोत्तमभूषिताः} %||4-32-23||

\twolineshloka
{नातृप्तान्नाति च व्यग्रान्नानुदात्तपरिच्छदान्}
{सुग्रीवानुचरांश्चापि लक्षयामास लक्ष्मणः} %||4-32-24||

\twolineshloka
{ततः सुग्रीवमासीनं काञ्चने परमासने}
{महार्हास्तरणोपेते ददर्शादित्यसंनिभम्} %||4-32-25||

\threelineshloka
{दिव्याभरणचित्राङ्गं दिव्यरूपं यशस्विनम्}
{दिव्यमाल्याम्बरधरं महेन्द्रमिव दुर्जयम्}
{दिव्याभरणमाल्याभिः प्रमदाभिः समावृतम्} %||4-32-26||

\fourlineindentedshloka
{रुमां तु वीरः परिरभ्य गाढम्}
{ वरासनस्थो वरहेमवर्णः}
{ददर्श सौमित्रिमदीनसत्त्वम्}
{ विशालनेत्रः सुविशालनेत्रम्} %||4-33-27||


{॥इत्यार्षे श्रीमद्रामायणे वाल्मीकीये आदिकाव्ये किष्किन्धाकाण्डे द्वात्रिंशः सर्गः॥}

\sect{त्रयस्त्रिंशः सर्गः}

\twolineshloka
{तमप्रतिहतं क्रुद्धं प्रविष्टं पुरुषर्षभम्}
{सुग्रीवो लक्ष्मणं दृष्ट्वा बभूव व्यथितेन्द्रियः} %||4-33-1||

\twolineshloka
{क्रुद्धं निःश्वसमानं तं प्रदीप्तमिव तेजसा}
{भ्रातुर्व्यसनसन्तप्तं दृष्ट्वा दशरथात्मजम्} %||4-33-2||

\twolineshloka
{उत्पपात हरिश्रेष्ठो हित्वा सौवर्णमासनम्}
{महान्महेन्द्रस्य यथा स्वलङ्कृत इव ध्वजः} %||4-33-3||

\twolineshloka
{उत्पतन्तमनूत्पेतू रुमाप्रभृतयः स्त्रियः}
{सुग्रीवं गगने पूर्णं चन्द्रं तारागणा इव} %||4-33-4||

\twolineshloka
{संरक्तनयनः श्रीमान्विचचाल कृताञ्जलिः}
{बभूवावस्थितस्तत्र कल्पवृक्षो महानिव} %||4-33-5||

\twolineshloka
{रुमा द्वितीयं सुग्रीवं नारीमध्यगतं स्थितम्}
{अब्रवील्लक्ष्मणः क्रुद्धः सतारं शशिनं यथा} %||4-33-6||

\twolineshloka
{सत्त्वाभिजनसम्पन्नः सानुक्रोशो जितेन्द्रियः}
{कृतज्ञः सत्यवादी च राजा लोके महीयते} %||4-33-7||

\twolineshloka
{यस्तु राजा स्थितोऽधर्मे मित्राणामुपकारिणाम्}
{मिथ्याप्रतिज्ञां कुरुते को नृशंसतरस्ततः} %||4-33-8||

\twolineshloka
{शतमश्वानृते हन्ति सहस्रं तु गवानृते}
{आत्मानं स्वजनं हन्ति पुरुषः पुरुषानृते} %||4-33-9||

\twolineshloka
{पूर्वं कृतार्थो मित्राणां न तत्प्रतिकरोति यः}
{कृतघ्नः सर्वभूतानां स वध्यः प्लवगेश्वर} %||4-33-10||

\twolineshloka
{गीतोऽयं ब्रह्मणा श्लोकः सर्वलोकनमस्कृतः}
{दृष्ट्वा कृतघ्नं क्रुद्धेन तं निबोध प्लवङ्गम} %||4-33-11||

\twolineshloka
{ब्रह्मघ्ने च सुरापे च चोरे भग्नव्रते तथा}
{निष्कृतिर्विहिता सद्भिः कृतघ्ने नास्ति निष्कृतिः} %||4-33-12||

\twolineshloka
{अनार्यस्त्वं कृतघ्नश्च मिथ्यावादी च वानर}
{पूर्वं कृतार्थो रामस्य न तत्प्रतिकरोषि यत्} %||4-33-13||

\twolineshloka
{ननु नाम कृतार्थेन त्वया रामस्य वानर}
{सीताया मार्गणे यत्नः कर्तव्यः कृतमिच्छता} %||4-33-14||

\twolineshloka
{स त्वं ग्राम्येषु भोगेषु सक्तो मिथ्या प्रतिश्रवः}
{न त्वां रामो विजानीते सर्पं मण्डूकराविणम्} %||4-33-15||

\twolineshloka
{महाभागेन रामेण पापः करुणवेदिना}
{हरीणां प्रापितो राज्यं त्वं दुरात्मा महात्मना} %||4-33-16||

\twolineshloka
{कृतं चेन्नाभिजानीषे रामस्याक्लिष्टकर्मणः}
{सद्यस्त्वं निशितैर्बाणैर्हतो द्रक्ष्यसि वालिनम्} %||4-33-17||

\twolineshloka
{न च सङ्कुचितः पन्था येन वाली हतो गतः}
{समये तिष्ठ सुग्रीव मा वालिपथमन्वगाः} %||4-33-18||

\fourlineindentedshloka
{न नूनमिक्ष्वाकुवरस्य कार्मुका}
{च्च्युताञ्शरान्पश्यसि वज्रसंनिभान्}
{ततः सुखं नाम निषेवसे सुखी}
{ न रामकार्यं मनसाप्यवेक्षसे} %||4-34-19||


{॥इत्यार्षे श्रीमद्रामायणे वाल्मीकीये आदिकाव्ये किष्किन्धाकाण्डे त्रयस्त्रिंशः सर्गः॥}

\sect{चतुस्त्रिंशः सर्गः}

\twolineshloka
{तथा ब्रुवाणं सौमित्रिं प्रदीप्तमिव तेजसा}
{अब्रवील्लक्ष्मणं तारा ताराधिपनिभानना} %||4-34-1||

\twolineshloka
{नैवं लक्ष्मण वक्तव्यो नायं परुषमर्हति}
{हरीणामीश्वरः श्रोतुं तव वक्त्राद्विशेषतः} %||4-34-2||

\twolineshloka
{नैवाकृतज्ञः सुग्रीवो न शठो नापि दारुणः}
{नैवानृतकथो वीर न जिह्मश्च कपीश्वरः} %||4-34-3||

\twolineshloka
{उपकारं कृतं वीरो नाप्ययं विस्मृतः कपिः}
{रामेण वीर सुग्रीवो यदन्यैर्दुष्करं रणे} %||4-34-4||

\twolineshloka
{रामप्रसादात्कीर्तिं च कपिराज्यं च शाश्वतम्}
{प्राप्तवानिह सुग्रीवो रुमां मां च परन्तप} %||4-34-5||

\twolineshloka
{सुदुःखं शायितः पूर्वं प्राप्येदं सुखमुत्तमम्}
{प्राप्तकालं न जानीते विश्वामित्रो यथा मुनिः} %||4-34-6||

\twolineshloka
{घृताच्यां किल संसक्तो दशवर्षाणि लक्ष्मण}
{अहोऽमन्यत धर्मात्मा विश्वामित्रो महामुनिः} %||4-34-7||

\twolineshloka
{स हि प्राप्तं न जानीते कालं कालविदां वरः}
{विश्वामित्रो महातेजाः किं पुनर्यः पृथग्जनः} %||4-34-8||

\twolineshloka
{देहधर्मं गतस्यास्य परिश्रान्तस्य लक्ष्मण}
{अवितृप्तस्य कामेषु रामः क्षन्तुमिहार्हति} %||4-34-9||

\twolineshloka
{न च रोषवशं तात गन्तुमर्हसि लक्ष्मण}
{निश्चयार्थमविज्ञाय सहसा प्राकृतो यथा} %||4-34-10||

\twolineshloka
{सत्त्वयुक्ता हि पुरुषास्त्वद्विधाः पुरुषर्षभ}
{अविमृश्य न रोषस्य सहसा यान्ति वश्यताम्} %||4-34-11||

\twolineshloka
{प्रसादये त्वां धर्मज्ञ सुग्रीवार्थे समाहिता}
{महान्रोषसमुत्पन्नः संरम्भस्त्यज्यतामयम्} %||4-34-12||

\twolineshloka
{रुमां मां कपिराज्यं च धनधान्यवसूनि च}
{रामप्रियार्थं सुग्रीवस्त्यजेदिति मतिर्मम} %||4-34-13||

\twolineshloka
{समानेष्व्यति सुग्रीवः सीतया सह राघवम्}
{शशाङ्कमिव रोहिष्या निहत्वा रावणं रणे} %||4-34-14||

\twolineshloka
{शतकोटिसहस्राणि लङ्कायां किल रक्षसाम्}
{अयुतानि च षट्त्रिंशत्सहस्राणि शतानि च} %||4-34-15||

\twolineshloka
{अहत्वा तांश्च दुर्धर्षान्राक्षसान्कामरूपिणः}
{न शक्यो रावणो हन्तुं येन सा मैथिली हृता} %||4-34-16||

\twolineshloka
{ते न शक्या रणे हन्तुमसहायेन लक्ष्मण}
{रावणः क्रूरकर्मा च सुग्रीवेण विशेषतः} %||4-34-17||

\twolineshloka
{एवमाख्यातवान्वाली स ह्यभिज्ञो हरीश्वरः}
{आगमस्तु न मे व्यक्तः श्रवात्तस्य ब्रवीम्यहम्} %||4-34-18||

\twolineshloka
{त्वत्सहायनिमित्तं वै प्रेषिता हरिपुङ्गवाः}
{आनेतुं वानरान्युद्धे सुबहून्हरियूथपान्} %||4-34-19||

\twolineshloka
{तांश्च प्रतीक्षमाणोऽयं विक्रान्तान्सुमहाबलान्}
{राघवस्यार्थसिद्ध्यर्थं न निर्याति हरीश्वरः} %||4-34-20||

\twolineshloka
{कृता तु संस्था सौमित्रे सुग्रीवेण यथापुरा}
{अद्य तैर्वानरैर्सर्वैरागन्तव्यं महाबलैः} %||4-34-21||

\threelineshloka
{ऋक्षकोटिसहस्राणि गोलाङ्गूलशतानि च}
{अद्य त्वामुपयास्यन्ति जहि कोपमरिन्दम}
{कोट्योऽनेकास्तु काकुत्स्थ कपीनां दीप्ततेजसाम्} %||4-34-22||

\fourlineindentedshloka
{तव हि मुखमिदं निरीक्ष्य कोपात्}
{ क्षतजनिभे नयने निरीक्षमाणाः}
{हरिवरवनिता न यान्ति शान्तिम्}
{ प्रथमभयस्य हि शङ्किताः स्म सर्वाः} %||4-35-23||


{॥इत्यार्षे श्रीमद्रामायणे वाल्मीकीये आदिकाव्ये किष्किन्धाकाण्डे चतुस्त्रिंशः सर्गः॥}

\sect{पञ्चत्रिंशः सर्गः}

\twolineshloka
{इत्युक्तस्तारया वाक्यं प्रश्रितं धर्मसंहितम्}
{मृदुस्वभावः सौमित्रिः प्रतिजग्राह तद्वचः} %||4-35-1||

\twolineshloka
{तस्मिन्प्रतिगृहीते तु वाक्ये हरिगणेश्वरः}
{लक्ष्मणात्सुमहत्त्रासं वस्त्रं क्लिन्नमिवात्यजत्} %||4-35-2||

\twolineshloka
{ततः कण्ठगतं माल्यं चित्रं बहुगुणं महत्}
{चिच्छेद विमदश्चासीत्सुग्रीवो वानरेश्वरः} %||4-35-3||

\twolineshloka
{स लक्ष्मणं भीमबलं सर्ववानरसत्तमः}
{अब्रवीत्प्रश्रितं वाक्यं सुग्रीवः सम्प्रहर्षयन्} %||4-35-4||

\twolineshloka
{प्रनष्टा श्रीश्च कीर्तिश्च कपिराज्यं च शाश्वतम्}
{रामप्रसादात्सौमित्रे पुनः प्राप्तमिदं मया} %||4-35-5||

\twolineshloka
{कः शक्तस्तस्य देवस्य ख्यातस्य स्वेन कर्मणा}
{तादृशं विक्रमं वीर प्रतिकर्तुमरिन्दम} %||4-35-6||

\twolineshloka
{सीतां प्राप्स्यति धर्मात्मा वधिष्यति च रावणम्}
{सहायमात्रेण मया राघवः स्वेन तेजसा} %||4-35-7||

\twolineshloka
{सहायकृत्यं हि तस्य येन सप्त महाद्रुमाः}
{शैलश्च वसुधा चैव बाणेनैकेन दारिताः} %||4-35-8||

\twolineshloka
{धनुर्विस्फारमाणस्य यस्य शब्देन लक्ष्मण}
{सशैला कम्पिता भूमिः सहायैस्तस्य किं नु वै} %||4-35-9||

\twolineshloka
{अनुयात्रां नरेन्द्रस्य करिष्येऽहं नरर्षभ}
{गच्छतो रावणं हन्तुं वैरिणं सपुरःसरम्} %||4-35-10||

\twolineshloka
{यदि किञ्चिदतिक्रान्तं विश्वासात्प्रणयेन वा}
{प्रेष्यस्य क्षमितव्यं मे न कश्चिन्नापराध्यति} %||4-35-11||

\twolineshloka
{इति तस्य ब्रुवाणस्य सुग्रीवस्य महात्मनः}
{अभवल्लक्ष्मणः प्रीतः प्रेंणा चेदमुवाच ह} %||4-35-12||

\twolineshloka
{सर्वथा हि मम भ्राता सनाथो वानरेश्वर}
{त्वया नाथेन सुग्रीव प्रश्रितेन विशेषतः} %||4-35-13||

\twolineshloka
{यस्ते प्रभावः सुग्रीव यच्च ते शौचमुत्तमम्}
{अर्हस्तं कपिराज्यस्य श्रियं भोक्तुमनुत्तमाम्} %||4-35-14||

\twolineshloka
{सहायेन च सुग्रीव त्वया रामः प्रतापवान्}
{वधिष्यति रणे शत्रूनचिरान्नात्र संशयः} %||4-35-15||

\twolineshloka
{धर्मज्ञस्य कृतज्ञस्य सङ्ग्रामेष्वनिवर्तिनः}
{उपपन्नं च युक्तं च सुग्रीव तव भाषितम्} %||4-35-16||

\twolineshloka
{दोषज्ञः सति सामर्थ्ये कोऽन्यो भाषितुमर्हति}
{वर्जयित्वा मम ज्येष्ठं त्वां च वानरसत्तम} %||4-35-17||

\twolineshloka
{सदृशश्चासि रामस्य विक्रमेण बलेन च}
{सहायो दैवतैर्दत्तश्चिराय हरिपुङ्गव} %||4-35-18||

\twolineshloka
{किं तु शीघ्रमितो वीर निष्क्राम त्वं मया सह}
{सान्त्वयस्व वयस्यं च भार्याहरणदुःखितम्} %||4-35-19||

\twolineshloka
{यच्च शोकाभिभूतस्य श्रुत्वा रामस्य भाषितम्}
{मया त्वं परुषाण्युक्तस्तच्च त्वं क्षन्तुमर्हसि} %||4-36-20||


{॥इत्यार्षे श्रीमद्रामायणे वाल्मीकीये आदिकाव्ये किष्किन्धाकाण्डे पञ्चत्रिंशः सर्गः॥}

\sect{षट्त्रिंशः सर्गः}

\twolineshloka
{एवमुक्तस्तु सुग्रीवो लक्ष्मणेन महात्मना}
{हनुमन्तं स्थितं पार्श्वे सचिवं वाक्यमब्रवीत्} %||4-36-1||

\twolineshloka
{महेन्द्रहिमवद्विन्ध्यकैलासशिखरेषु च}
{मन्दरे पाण्डुशिखरे पञ्चशैलेषु ये स्थिताः} %||4-36-2||

\twolineshloka
{तरुणादित्यवर्णेषु भ्राजमानेषु सर्वशः}
{पर्वतेषु समुद्रान्ते पश्चिमस्यां तु ये दिशि} %||4-36-3||

\twolineshloka
{आदित्यभवने चैव गिरौ सन्ध्याभ्रसंनिभे}
{पद्मतालवनं भीमं संश्रिता हरिपुङ्गवाः} %||4-36-4||

\twolineshloka
{अञ्जनाम्बुदसङ्काशाः कुञ्जरप्रतिमौजसः}
{अञ्जने परते चैव ये वसन्ति प्लवङ्गमाः} %||4-36-5||

\twolineshloka
{मनःशिला गुहावासा वानराः कनकप्रभाः}
{मेरुपार्श्वगताश्चैव ये च धूम्रगिरिं श्रिताः} %||4-36-6||

\twolineshloka
{तरुणादित्यवर्णाश्च पर्वते ये महारुणे}
{पिबन्तो मधुमैरेयं भीमवेगाः प्लवङ्गमाः} %||4-36-7||

\twolineshloka
{वनेषु च सुरम्येषु सुगन्धिषु महत्सु च}
{तापसानां च रम्येषु वनान्तेषु समन्ततः} %||4-36-8||

\twolineshloka
{तांस्तांस्त्वमानय क्षिप्रं पृथिव्यां सर्ववानरान्}
{सामदानादिभिः कल्पैराशु प्रेषय वानरान्} %||4-36-9||

\twolineshloka
{प्रेषिताः प्रथमं ये च मया दूता महाजवाः}
{त्वरणार्थं तु भूयस्त्वं हरीन्सम्प्रेषयापरान्} %||4-36-10||

\twolineshloka
{ये प्रसक्ताश्च कामेषु दीर्घसूत्राश्च वानराः}
{इहानयस्व तान्सर्वाञ्शीघ्रं तु मम शासनात्} %||4-36-11||

\twolineshloka
{अहोभिर्दशभिर्ये च नागच्छन्ति ममाज्ञया}
{हन्तव्यास्ते दुरात्मानो राजशासनदूषकाः} %||4-36-12||

\twolineshloka
{शतान्यथ सहस्राणि कोट्यश्च मम शासनात्}
{प्रयान्तु कपिसिंहानां दिशो मम मते स्थिताः} %||4-36-13||

\twolineshloka
{मेघपर्वतसङ्काशाश्छादयन्त इवाम्बरम्}
{घोररूपाः कपिश्रेष्ठा यान्तु मच्छासनादितः} %||4-36-14||

\twolineshloka
{ते गतिज्ञा गतिं गत्वा पृथिव्यां सर्ववानराः}
{आनयन्तु हरीन्सर्वांस्त्वरिताः शासनान्मम} %||4-36-15||

\twolineshloka
{तस्य वानरराजस्य श्रुत्वा वायुसुतो वचः}
{दिक्षु सर्वासु विक्रान्तान्प्रेषयामास वानरान्} %||4-36-16||

\twolineshloka
{ते पदं विष्णुविक्रान्तं पतत्रिज्योतिरध्वगाः}
{प्रयाताः प्रहिता राज्ञा हरयस्तत्क्षणेन वै} %||4-36-17||

\twolineshloka
{ते समुद्रेषु गिरिषु वनेषु च सरित्सु च}
{वानरा वानरान्सर्वान्रामहेतोरचोदयन्} %||4-36-18||

\twolineshloka
{मृत्युकालोपमस्याज्ञां राजराजस्य वानराः}
{सुग्रीवस्याययुः श्रुत्वा सुग्रीवभयदर्शिनः} %||4-36-19||

\twolineshloka
{ततस्तेऽञ्जनसङ्काशा गिरेस्तस्मान्महाजवाः}
{तिस्रः कोट्यः प्लवङ्गानां निर्ययुर्यत्र राघवः} %||4-36-20||

\twolineshloka
{अस्तं गच्छति यत्रार्कस्तस्मिन्गिरिवरे रताः}
{तप्तहेमसमाभासास्तस्मात्कोट्यो दशच्युताः} %||4-36-21||

\twolineshloka
{कैलास शिखरेभ्यश्च सिंहकेसरवर्चसाम्}
{ततः कोटिसहस्राणि वानराणामुपागमन्} %||4-36-22||

\twolineshloka
{फलमूलेन जीवन्तो हिमवन्तमुपाश्रिताः}
{तेषां कोटिसहस्राणां सहस्रं समवर्तत} %||4-36-23||

\twolineshloka
{अङ्गारक समानानां भीमानां भीमकर्मणाम्}
{विन्ध्याद्वानरकोटीनां सहस्राण्यपतन्द्रुतम्} %||4-36-24||

\twolineshloka
{क्षीरोदवेलानिलयास्तमालवनवासिनः}
{नारिकेलाशनाश्चैव तेषां सङ्ख्या न विद्यते} %||4-36-25||

\twolineshloka
{वनेभ्यो गह्वरेभ्यश्च सरिद्भ्यश्च महाजवाः}
{आगच्छद्वानरी सेना पिबन्तीव दिवाकरम्} %||4-36-26||

\twolineshloka
{ये तु त्वरयितुं याता वानराः सर्ववानरान्}
{ते वीरा हिमवच्छैलं ददृशुस्तं महाद्रुमम्} %||4-36-27||

\twolineshloka
{तस्मिन्गिरिवरे रम्ये यज्ञो महेश्वरः पुरा}
{सर्वदेवमनस्तोषो बभौ दिव्यो मनोहरः} %||4-36-28||

\twolineshloka
{अन्नविष्यन्दजातानि मूलानि च फलानि च}
{अमृतस्वादुकल्पानि ददृशुस्तत्र वानराः} %||4-36-29||

\twolineshloka
{तदन्न सम्भवं दिव्यं फलं मूलं मनोहरम्}
{यः कश्चित्सकृदश्नाति मासं भवति तर्पितः} %||4-36-30||

\twolineshloka
{तानि मूलानि दिव्यानि फलानि च फलाशनाः}
{औषधानि च दिव्यानि जगृहुर्हरियूथपाः} %||4-36-31||

\twolineshloka
{तस्माच्च यज्ञायतनात्पुष्पाणि सुरभीणि च}
{आनिन्युर्वानरा गत्वा सुग्रीवप्रियकारणात्} %||4-36-32||

\twolineshloka
{ते तु सर्वे हरिवराः पृथिव्यां सर्ववानरान्}
{सञ्चोदयित्वा त्वरितं यूथानां जग्मुरग्रतः} %||4-36-33||

\twolineshloka
{ते तु तेन मुहूर्तेन यूथपाः शीघ्रकारिणः}
{किष्किन्धां त्वरया प्राप्ताः सुग्रीवो यत्र वानरः} %||4-36-34||

\twolineshloka
{ते गृहीत्वौषधीः सर्वाः फलं मूलं च वानराः}
{तं प्रतिग्राहयामासुर्वचनं चेदमब्रुवन्} %||4-36-35||

\twolineshloka
{सर्वे परिगताः शैलाः समुद्राश्च वनानि च}
{पृथिव्यां वानराः सर्वे शासनादुपयान्ति ते} %||4-36-36||

\twolineshloka
{एवं श्रुत्वा ततो हृष्टः सुग्रीवः प्लवगाधिपः}
{प्रतिजग्राह च प्रीतस्तेषां सर्वमुपायनम्} %||4-37-37||


{॥इत्यार्षे श्रीमद्रामायणे वाल्मीकीये आदिकाव्ये किष्किन्धाकाण्डे षट्त्रिंशः सर्गः॥}

\sect{सप्तत्रिंशः सर्गः}

\twolineshloka
{प्रतिगृह्य च तत्सर्वमुपानयमुपाहृतम्}
{वानरान्सान्त्वयित्वा च सर्वानेव व्यसर्जयत्} %||4-37-1||

\twolineshloka
{विसर्जयित्वा स हरीञ्शूरांस्तान्कृतकर्मणः}
{मेने कृतार्थमात्मानं राघवं च महाबलम्} %||4-37-2||

\threelineshloka
{स लक्ष्मणो भीमबलं सर्ववानरसत्तमम्}
{अब्रवीत्प्रश्रितं वाक्यं सुग्रीवं सम्प्रहर्षयन्}
{किष्किन्धाया विनिष्क्राम यदि ते सौम्य रोचते} %||4-37-3||

\twolineshloka
{तस्य तद्वचनं श्रुत्वा लक्ष्मणस्य सुभाषितम्}
{सुग्रीवः परमप्रीतो वाक्यमेतदुवाच ह} %||4-37-4||

\twolineshloka
{एवं भवतु गच्छामः स्थेयं त्वच्छासने मया}
{तमेवमुक्त्वा सुग्रीवो लक्ष्मणं शुभलक्ष्मणम्} %||4-37-5||

\twolineshloka
{विसर्जयामास तदा तारामन्याश्च योषितः}
{एतेत्युच्चैर्हरिवरान्सुग्रीवः समुदाहरत्} %||4-37-6||

\twolineshloka
{तस्य तद्वचनं श्रुत्वा हरयः शीघ्रमाययुः}
{बद्धाञ्जलिपुटाः सर्वे ये स्युः स्त्रीदर्शनक्षमाः} %||4-37-7||

\twolineshloka
{तानुवाच ततः प्राप्तान्राजार्कसदृशप्रभः}
{उपस्थापयत क्षिप्रं शिबिकां मम वानराः} %||4-37-8||

\twolineshloka
{श्रुत्वा तु वचनं तस्य हरयः शीघ्रविक्रमाः}
{समुपस्थापयामासुः शिबिकां प्रियदर्शनाम्} %||4-37-9||

\twolineshloka
{तामुपस्थापितां दृष्ट्वा शिबिकां वानराधिपः}
{लक्ष्मणारुह्यतां शीघ्रमिति सौमित्रिमब्रवीत्} %||4-37-10||

\twolineshloka
{इत्युक्त्वा काञ्चनं यानं सुग्रीवः सूर्यसंनिभम्}
{बृहद्भिर्हरिभिर्युक्तमारुरोह सलक्ष्मणः} %||4-37-11||

\twolineshloka
{पाण्डुरेणातपत्रेण ध्रियमाणेन मूर्धनि}
{शुक्लैश्च बालव्यजनैर्धूयमानैः समन्ततः} %||4-37-12||

\twolineshloka
{शङ्खभेरीनिनादैश्च बन्दिभिश्चाभिवन्दितः}
{निर्ययौ प्राप्य सुग्रीवो राज्यश्रियमनुत्तमाम्} %||4-37-13||

\twolineshloka
{स वानरशतैस्तीष्क्णैर्बहुभिः शस्त्रपाणिभिः}
{परिकीर्णो ययौ तत्र यत्र रामो व्यवस्थितः} %||4-37-14||

\twolineshloka
{स तं देशमनुप्राप्य श्रेष्ठं रामनिषेवितम्}
{अवातरन्महातेजाः शिबिकायाः सलक्ष्मणः} %||4-37-15||

\twolineshloka
{आसाद्य च ततो रामं कृताञ्जलिपुटोऽभवत्}
{कृताञ्जलौ स्थिते तस्मिन्वानराश्चभवंस्तथा} %||4-37-16||

\twolineshloka
{तटाकमिव तद्दृष्ट्वा रामः कुड्मलपङ्कजम्}
{वानराणां महत्सैन्यं सुग्रीवे प्रीतिमानभूत्} %||4-37-17||

\twolineshloka
{पादयोः पतितं मूर्ध्ना तमुत्थाप्य हरीश्वरम्}
{प्रेंणा च बहुमानाच्च राघवः परिषस्वजे} %||4-37-18||

\twolineshloka
{परिष्वज्य च धर्मात्मा निषीदेति ततोऽब्रवीत्}
{तं निषण्णं ततो दृष्ट्वा क्षितौ रामोऽब्रवीद्वचः} %||4-37-19||

\twolineshloka
{धर्ममर्थं च कामं च काले यस्तु निषेवते}
{विभज्य सततं वीर स राजा हरिसत्तम} %||4-37-20||

\twolineshloka
{हित्वा धर्मं तथार्थं च कामं यस्तु निषेवते}
{स वृक्षाग्रे यथा सुप्तः पतितः प्रतिबुध्यते} %||4-37-21||

\twolineshloka
{अमित्राणां वधे युक्तो मित्राणां सङ्ग्रहे रतः}
{त्रिवर्गफलभोक्ता तु राजा धर्मेण युज्यते} %||4-37-22||

\twolineshloka
{उद्योगसमयस्त्वेष प्राप्तः शत्रुविनाशन}
{सञ्चिन्त्यतां हि पिङ्गेश हरिभिः सह मन्त्रिभिः} %||4-37-23||

{एवमुक्तस्तु सुग्रीवो रामं वचनमब्रवीत्॥\devanumber{24}॥} %||4-37-24|| (Check)

\addtocounter{shlokacount}{1}

\twolineshloka
{प्रनष्टा श्रीश्च कीर्तिश्च कपिराज्यं च शाश्वतम्}
{त्वत्प्रसादान्महाबाहो पुनः प्राप्तमिदं मया} %||4-37-25||

\twolineshloka
{तव देवप्रसदाच्च भ्रातुश्च जयतां वर}
{कृतं न प्रतिकुर्याद्यः पुरुषाणां स दूषकः} %||4-37-26||

\twolineshloka
{एते वानरमुख्याश्च शतशः शत्रुसूदन}
{प्राप्ताश्चादाय बलिनः पृथिव्यां सर्ववानरान्} %||4-37-27||

\twolineshloka
{ऋक्षाश्चावहिताः शूरा गोलाङ्गूलाश्च राघव}
{कान्तार वनदुर्गाणामभिज्ञा घोरदर्शनाः} %||4-37-28||

\twolineshloka
{देवगन्धर्वपुत्राश्च वानराः कामरूपिणः}
{स्वैः स्वैः परिवृताः सैन्यैर्वर्तन्ते पथि राघव} %||4-37-29||

\twolineshloka
{शतैः शतसहस्रैश्च कोटिभिश्च प्लवङ्गमाः}
{अयुतैश्चावृता वीरा शङ्कुभिश्च परन्तप} %||4-37-30||

\twolineshloka
{अर्बुदैरर्बुदशतैर्मध्यैश्चान्तैश्च वानराः}
{समुद्रैश्च परार्धैश्च हरयो हरियूथपाः} %||4-37-31||

\twolineshloka
{आगमिष्यन्ति ते राजन्महेन्द्रसमविक्रमाः}
{मेरुमन्दरसङ्काशा विन्ध्यमेरुकृतालयाः} %||4-37-32||

\twolineshloka
{ते त्वामभिगमिष्यन्ति राक्षसं ये सबान्धवम्}
{निहत्य रावणं सङ्ख्ये ह्यानयिष्यन्ति मैथिलीम्} %||4-37-33||

\fourlineindentedshloka
{ततस्तमुद्योगमवेक्ष्य बुद्धिमा}
{न्हरिप्रवीरस्य निदेशवर्तिनः}
{बभूव हर्षाद्वसुधाधिपात्मजः}
{ प्रबुद्धनीलोत्पलतुल्यदर्शनः} %||4-38-34||


{॥इत्यार्षे श्रीमद्रामायणे वाल्मीकीये आदिकाव्ये किष्किन्धाकाण्डे सप्तत्रिंशः सर्गः॥}

\sect{अष्टात्रिंशः सर्गः}

\twolineshloka
{इति ब्रुवाणं सुग्रीवं रामो धर्मभृतां वरः}
{बाहुभ्यां सम्परिष्वज्य प्रत्युवाच कृताञ्जलिम्} %||4-38-1||

\twolineshloka
{यदिन्द्रो वर्षते वर्षं न तच्चित्रं भवेद्भुवि}
{आदित्यो वा सहस्रांशुः कुर्याद्वितिमिरं नभः} %||4-38-2||

\twolineshloka
{चन्द्रमा रश्मिभिः कुर्यात्पृथिवीं सौम्य निर्मलाम्}
{त्वद्विधो वापि मित्राणां प्रतिकुर्यात्परन्तप} %||4-38-3||

\twolineshloka
{एवं त्वयि न तच्चित्रं भवेद्यत्सौम्य शोभनम्}
{जानाम्यहं त्वां सुग्रीव सततं प्रियवादिनम्} %||4-38-4||

\twolineshloka
{त्वत्सनाथः सखे सङ्ख्ये जेतास्मि सकलानरीन्}
{त्वमेव मे सुहृन्मित्रं साहाय्यं कर्तुमर्हसि} %||4-38-5||

\twolineshloka
{जहारात्मविनाशाय वैदेहीं राक्षसाधमः}
{वञ्चयित्वा तु पौलोमीमनुह्लादो यथा शचीम्} %||4-38-6||

\twolineshloka
{नचिरात्तं हनिष्यामि रावणं निशितैः शरैः}
{पौलोम्याः पितरं दृप्तं शतक्रतुरिवारिहा} %||4-38-7||

\twolineshloka
{एतस्मिन्नन्तरे चैव रजः समभिवर्तत}
{उष्णां तीव्रां सहस्रांशोश्छादयद्गगने प्रभाम्} %||4-38-8||

\twolineshloka
{दिशः पर्याकुलाश्चासन्रजसा तेन मूर्छिताः}
{चचाल च मही सर्वा सशैलवनकानना} %||4-38-9||

\twolineshloka
{ततो नगेन्द्रसङ्काशैस्तीक्ष्ण दंष्ट्रैर्महाबलैः}
{कृत्स्ना सञ्छादिता भूमिरसङ्ख्येयैः प्लवङ्गमैः} %||4-38-10||

\twolineshloka
{निमेषान्तरमात्रेण ततस्तैर्हरियूथपैः}
{कोटीशतपरीवारैः कामरूपिभिरावृता} %||4-38-11||

\twolineshloka
{नादेयैः पार्वतीयैश्च सामुद्रैश्च महाबलैः}
{हरिभिर्मेघनिर्ह्रादैरन्यैश्च वनचारिभिः} %||4-38-12||

\twolineshloka
{तरुणादित्यवर्णैश्च शशिगौरैश्च वानरैः}
{पद्मकेसरवर्णैश्च श्वेतैर्मेरुकृतालयैः} %||4-38-13||

\twolineshloka
{कोटीसहस्रैर्दशभिः श्रीमान्परिवृतस्तदा}
{वीरः शतबलिर्नाम वानरः प्रत्यदृश्यत} %||4-38-14||

\twolineshloka
{ततः काञ्चनशैलाभस्ताराया वीर्यवान्पिता}
{अनेकैर्दशसाहस्रैः कोटिभिः प्रत्यदृश्यत} %||4-38-15||

\twolineshloka
{पद्मकेसरसङ्काशस्तरुणार्कनिभाननः}
{बुद्धिमान्वानरश्रेष्ठः सर्ववानरसत्तमः} %||4-38-16||

\twolineshloka
{अनीकैर्बहुसाहस्रैर्वानराणां समन्वितः}
{पिता हनुमतः श्रीमान्केसरी प्रत्यदृश्यत} %||4-38-17||

\twolineshloka
{गोलाङ्गूलमहाराजो गवाक्षो भीमविक्रमः}
{वृतः कोटिसहस्रेण वानराणामदृश्यत} %||4-38-18||

\twolineshloka
{ऋक्षाणां भीमवेगानां धूम्रः शत्रुनिबर्हणः}
{वृतः कोटिसहस्राभ्यां द्वाभ्यां समभिवर्तत} %||4-38-19||

\twolineshloka
{महाचलनिभैर्घोरैः पनसो नाम यूथपः}
{आजगाम महावीर्यस्तिसृभिः कोटिभिर्वृतः} %||4-38-20||

\twolineshloka
{नीलाञ्जनचयाकारो नीलो नामाथ यूथपः}
{अदृश्यत महाकायः कोटिभिर्दशभिर्वृतः} %||4-38-21||

\twolineshloka
{दरीमुखश्च बलवान्यूथपोऽभ्याययौ तदा}
{वृतः कोटिसहस्रेण सुग्रीवं समुपस्थितः} %||4-38-22||

\twolineshloka
{मैन्दश्च द्विविदश्चोभावश्विपुत्रौ महावलौ}
{कोटिकोटिसहस्रेण वानराणामदृश्यताम्} %||4-38-23||

\twolineshloka
{ततः कोटिसहस्राणां सहस्रेण शतेन च}
{पृष्ठतोऽनुगतः प्राप्तो हरिभिर्गन्धमादनः} %||4-38-24||

\twolineshloka
{ततः पद्मसहस्रेण वृतः शङ्कुशतेन च}
{युवराजोऽङ्गदः प्राप्तः पितृतुल्यपराक्रमः} %||4-38-25||

\twolineshloka
{ततस्ताराद्युतिस्तारो हरिर्भीमपराक्रमः}
{पञ्चभिर्हरिकोटीभिर्दूरतः प्रत्यदृश्यत} %||4-38-26||

\twolineshloka
{इन्द्रजानुः कपिर्वीरो यूथपः प्रत्यदृश्यत}
{एकादशानां कोटीनामीश्वरस्तैश्च संवृतः} %||4-38-27||

\twolineshloka
{ततो रम्भस्त्वनुप्राप्तस्तरुणादित्यसंनिभः}
{अयुतेन वृतश्चैव सहस्रेण शतेन च} %||4-38-28||

\twolineshloka
{ततो यूथपतिर्वीरो दुर्मुखो नाम वानरः}
{प्रत्यदृश्यत कोटिभ्यां द्वाभ्यां परिवृतो बली} %||4-38-29||

\twolineshloka
{कैलासशिखराकारैर्वानरैर्भीमविक्रमैः}
{वृतः कोटिसहस्रेण हनुमान्प्रत्यदृश्यत} %||4-38-30||

\twolineshloka
{नलश्चापि महावीर्यः संवृतो द्रुमवासिभिः}
{कोटीशतेन सम्प्राप्तः सहस्रेण शतेन च} %||4-38-31||

\twolineshloka
{शरभः कुमुदो वह्निर्वानरो रम्भ एव च}
{एते चान्ये च बहवो वानराः कामरूपिणः} %||4-38-32||

\threelineshloka
{आवृत्य पृथिवीं सर्वां पर्वतांश्च वनानि च}
{आप्लवन्तः प्लवन्तश्च गर्जन्तश्च प्लवङ्गमाः}
{अभ्यवर्तन्त सुग्रीवं सूर्यमभ्रगणा इव} %||4-38-33||

\twolineshloka
{कुर्वाणा बहुशब्दांश्च प्रहृष्टा बलशालिनः}
{शिरोभिर्वानरेन्द्राय सुग्रीवाय न्यवेदयन्} %||4-38-34||

\twolineshloka
{अपरे वानरश्रेष्ठाः सङ्गम्य च यथोचितम्}
{सुग्रीवेण समागम्य स्थिताः प्राञ्जलयस्तदा} %||4-38-35||

\twolineshloka
{सुग्रीवस्त्वरितो रामे सर्वांस्तान्वानरर्षभान्}
{निवेदयित्वा धर्मज्ञः स्थितः प्राञ्जलिरब्रवीत्} %||4-38-36||

\fourlineindentedshloka
{यथा सुखं पर्वतनिर्झरेषु}
{ वनेषु सर्वेषु च वानरेन्द्राः}
{निवेशयित्वा विधिवद्बलानि}
{ बलं बलज्ञः प्रतिपत्तुमीष्टे} %||4-39-37||


{॥इत्यार्षे श्रीमद्रामायणे वाल्मीकीये आदिकाव्ये किष्किन्धाकाण्डे अष्टात्रिंशः सर्गः॥}

\sect{एकोनचत्वारिंशः सर्गः}

\twolineshloka
{अथ राजा समृद्धार्थः सुग्रीवः प्लवगेश्वरः}
{उवाच नरशार्दूलं रामं परबलार्दनम्} %||4-39-1||

\twolineshloka
{आगता विनिविष्टाश्च बलिनः कामरूपिणः}
{वानरेन्द्रा महेन्द्राभा ये मद्विषयवासिनः} %||4-39-2||

\twolineshloka
{त इमे बहुसाहस्रैर्हरिभिर्भीमविक्रमैः}
{आगता वानरा घोरा दैत्यदानवसंनिभाः} %||4-39-3||

\twolineshloka
{ख्यातकर्मापदानाश्च बलवन्तो जितक्लमाः}
{पराक्रमेषु विख्याता व्यवसायेषु चोत्तमाः} %||4-39-4||

\twolineshloka
{पृथिव्यम्बुचरा राम नानानगनिवासिनः}
{कोट्यग्रश इमे प्राप्ता वानरास्तव किङ्कराः} %||4-39-5||

\twolineshloka
{निदेशवर्तिनः सर्वे सर्वे गुरुहिते रताः}
{अभिप्रेतमनुष्ठातुं तव शक्ष्यन्त्यरिन्दम} %||4-39-6||

\twolineshloka
{यन्मन्यसे नरव्याघ्र प्राप्तकालं तदुच्यताम्}
{तत्सैन्यं त्वद्वशे युक्तमाज्ञापयितुमर्हसि} %||4-39-7||

\twolineshloka
{काममेषामिदं कार्यं विदितं मम तत्त्वतः}
{तथापि तु यथा तत्त्वमाज्ञापयितुमर्हसि} %||4-39-8||

\twolineshloka
{तथा ब्रुवाणं सुग्रीवं रामो दशरथात्मजः}
{बाहुभ्यां सम्परिष्वज्य इदं वचनमब्रवीत्} %||4-39-9||

\twolineshloka
{ज्ञायतां सौम्य वैदेही यदि जीवति वा न वा}
{स च देशो महाप्राज्ञ यस्मिन्वसति रावणः} %||4-39-10||

\twolineshloka
{अधिगम्य च वैदेहीं निलयं रावणस्य च}
{प्राप्तकालं विधास्यामि तस्मिन्काले सह त्वया} %||4-39-11||

\twolineshloka
{नाहमस्मिन्प्रभुः कार्ये वानरेश न लक्ष्मणः}
{त्वमस्य हेतुः कार्यस्य प्रभुश्च प्लवगेश्वर} %||4-39-12||

\twolineshloka
{त्वमेवाज्ञापय विभो मम कार्यविनिश्चयम्}
{त्वं हि जानासि यत्कार्यं मम वीर न संशयः} %||4-39-13||

\twolineshloka
{सुहृद्द्वितीयो विक्रान्तः प्राज्ञः कालविशेषवित्}
{भवानस्मद्धिते युक्तः सुकृतार्थोऽर्थवित्तमः} %||4-39-14||

\threelineshloka
{एवमुक्तस्तु सुग्रीवो विनतं नाम यूथपम्}
{अब्रवीद्राम साम्निध्ये लक्ष्मणस्य च धीमतः}
{शैलाभं मेघनिर्घोषमूर्जितं प्लवगेश्वरम्} %||4-39-15||

\twolineshloka
{सोमसूर्यात्मजैः सार्धं वानरैर्वानरोत्तम}
{देशकालनयैर्युक्तः कार्याकार्यविनिश्चये} %||4-39-16||

\twolineshloka
{वृतः शतसहस्रेण वानराणां तरस्विनाम्}
{अधिगच्छ दिशं पूर्वां सशैलवनकाननाम्} %||4-39-17||

\twolineshloka
{तत्र सीतां च वैदेहीं निलयं रावणस्य च}
{मार्गध्वं गिरिदुर्गेषु वनेषु च नदीषु च} %||4-39-18||

\twolineshloka
{नदीं भागीरथीं रम्यां सरयूं कौशिकीं तथा}
{कालिन्दीं यमुनां रम्यां यामुनं च महागिरिम्} %||4-39-19||

\twolineshloka
{सरस्वतीं च सिन्धुं च शोणं मणिनिभोदकम्}
{महीं कालमहीं चैव शैलकाननशोभिताम्} %||4-39-20||

\twolineshloka
{ब्रह्ममालान्विदेहांश्च मालवान्काशिकोसलान्}
{मागधांश्च महाग्रामान्पुण्ड्रान्वङ्गांस्तथैव च} %||4-39-21||

\twolineshloka
{पत्तनं कोशकाराणां भूमिं च रजताकराम्}
{सर्वमेतद्विचेतव्यं मृगयद्भिर्ततस्ततः} %||4-39-22||

\twolineshloka
{रामस्य दयितां भार्यां सीतां दशरतः स्नुषाम्}
{समुद्रमवगाढांश्च पर्वतान्पत्तनानि च} %||4-39-23||

\twolineshloka
{मन्दरस्य च ये कोटिं संश्रिताः केचिदायताम्}
{कर्णप्रावरणाश्चैव तथा चाप्योष्ठकर्णकाः} %||4-39-24||

\twolineshloka
{घोरा लोहमुखाश्चैव जवनाश्चैकपादकाः}
{अक्षया बलवन्तश्च पुरुषाः पुरुषादकाः} %||4-39-25||

\twolineshloka
{किराताः कर्णचूडाश्च हेमाङ्गाः प्रियदर्शनाः}
{आममीनाशनास्तत्र किराता द्वीपवासिनः} %||4-39-26||

\twolineshloka
{अन्तर्जलचरा घोरा नरव्याघ्रा इति श्रुताः}
{एतेषामालयाः सर्वे विचेयाः काननौकसः} %||4-39-27||

\twolineshloka
{गिरिभिर्ये च गम्यन्ते प्लवनेन प्लवेन च}
{रत्नवन्तं यवद्वीपं सप्तराज्योपशोभितम्} %||4-39-28||

\twolineshloka
{सुवर्णरूप्यकं चैव सुवर्णाकरमण्डितम्}
{यवद्वीपमतिक्रम्य शिशिरो नाम पर्वतः} %||4-39-29||

\twolineshloka
{दिवं स्पृशति शृङ्गेण देवदानवसेवितः}
{एतेषां गिरिदुर्गेषु प्रतापेषु वनेषु च} %||4-39-30||

\twolineshloka
{रावणः सह वैदेह्या मार्गितव्यस्ततस्ततः}
{ततः समुद्रद्वीपांश्च सुभीमान्द्रष्टुमर्हथ} %||4-39-31||

\twolineshloka
{तत्रासुरा महाकायाश्छायां गृह्णन्ति नित्यशः}
{ब्रह्मणा समनुज्ञाता दीर्घकालं बुभुक्षिताः} %||4-39-32||

\twolineshloka
{तं कालमेघप्रतिमं महोरगनिषेवितम्}
{अभिगम्य महानादं तीर्थेनैव महोदधिम्} %||4-39-33||

\twolineshloka
{ततो रक्तजलं भीमं लोहितं नाम सागरम्}
{गता द्रक्ष्यथ तां चैव बृहतीं कूटशाल्मलीम्} %||4-39-34||

\twolineshloka
{गृहं च वैनतेयस्य नानारत्नविभूषितम्}
{तत्र कैलाससङ्काशं विहितं विश्वकर्मणा} %||4-39-35||

\twolineshloka
{तत्र शैलनिभा भीमा मन्देहा नाम राक्षसाः}
{शैलशृङ्गेषु लम्बन्ते नानारूपा भयावहाः} %||4-39-36||

\twolineshloka
{ते पतन्ति जले नित्यं सूर्यस्योदयनं प्रति}
{अभितप्ताश्च सूर्येण लम्बन्ते स्म पुनः पुनः} %||4-39-37||

\twolineshloka
{ततः पाण्डुरमेघाभं क्षीरौदं नाम सागरम्}
{गता द्रक्ष्यथ दुर्धर्षा मुखा हारमिवोर्मिभिः} %||4-39-38||

\twolineshloka
{तस्य मध्ये महाश्वेत ऋषभो नाम पर्वतः}
{दिव्यगन्धैः कुसुमितै रजतैश्च नगैर्वृतः} %||4-39-39||

\twolineshloka
{सरश्च राजतैः पद्मैर्ज्वलितैर्हेमकेसरैः}
{नाम्ना सुदर्शनं नाम राजहंसैः समाकुलम्} %||4-39-40||

\twolineshloka
{विबुधाश्चारणा यक्षाः किंनराः साप्सरोगणाः}
{हृष्टाः समभिगच्छन्ति नलिनीं तां रिरंसवः} %||4-39-41||

\twolineshloka
{क्षीरोदं समतिक्रम्य ततो द्रक्ष्यथ वानराः}
{जलोदं सागरश्रेष्ठं सर्वभूतभयावहम्} %||4-39-42||

\twolineshloka
{तत्र तत्कोपजं तेजः कृतं हयमुखं महत्}
{अस्याहुस्तन्महावेगमोदनं सचराचरम्} %||4-39-43||

\twolineshloka
{तत्र विक्रोशतां नादो भूतानां सागरौकसाम्}
{श्रूयते चासमर्थानां दृष्ट्वा तद्वडवामुखम्} %||4-39-44||

\twolineshloka
{स्वादूदस्योत्तरे देशे योजनानि त्रयोदश}
{जातरूपशिलो नाम महान्कनकपर्वतः} %||4-39-45||

\twolineshloka
{आसीनं पर्वतस्याग्रे सर्वभूतनमस्कृतम्}
{सहस्रशिरसं देवमनन्तं नीलवाससम्} %||4-39-46||

\twolineshloka
{त्रिशिराः काञ्चनः केतुस्तालस्तस्य महात्मनः}
{स्थापितः पर्वतस्याग्रे विराजति सवेदिकः} %||4-39-47||

\twolineshloka
{पूर्वस्यां दिशि निर्माणं कृतं तत्त्रिदशेश्वरैः}
{ततः परं हेममयः श्रीमानुदयपर्वतः} %||4-39-48||

\twolineshloka
{तस्य कोटिर्दिवं स्पृष्ट्वा शतयोजनमायता}
{जातरूपमयी दिव्या विराजति सवेदिका} %||4-39-49||

\twolineshloka
{सालैस्तालैस्तमालैश्च कर्णिकारैश्च पुष्पितैः}
{जातरूपमयैर्दिव्यैः शोभते सूर्यसंनिभैः} %||4-39-50||

\twolineshloka
{तत्र योजनविस्तारमुच्छ्रितं दशयोजनम्}
{शृङ्गं सौमनसं नाम जातरूपमयं ध्रुवम्} %||4-39-51||

\twolineshloka
{तत्र पूर्वं पदं कृत्वा पुरा विष्णुस्त्रिविक्रमे}
{द्वितीयं शिखरं मेरोश्चकार पुरुषोत्तमः} %||4-39-52||

\twolineshloka
{उत्तरेण परिक्रम्य जम्बूद्वीपं दिवाकरः}
{दृश्यो भवति भूयिष्ठं शिखरं तन्महोच्छ्रयम्} %||4-39-53||

\twolineshloka
{तत्र वैखानसा नाम वालखिल्या महर्षयः}
{प्रकाशमाना दृश्यन्ते सूर्यवर्णास्तपस्विनः} %||4-39-54||

\twolineshloka
{अयं सुदर्शनो द्वीपः पुरो यस्य प्रकाशते}
{यस्मिंस्तेजश्च चक्षुश्च सर्वप्रानभृतामपि} %||4-39-55||

\twolineshloka
{शैलस्य तस्य कुञ्जेषु कन्दरेषु वनेषु च}
{रावणः सह वैदेह्या मार्गितव्यस्ततस्ततः} %||4-39-56||

\twolineshloka
{काञ्चनस्य च शैलस्य सूर्यस्य च महात्मनः}
{आविष्टा तेजसा सन्ध्या पूर्वा रक्ता प्रकाशते} %||4-39-57||

\twolineshloka
{ततः परमगम्या स्याद्दिक्पूर्वा त्रिदशावृता}
{रहिता चन्द्रसूर्याभ्यामदृश्या तिमिरावृता} %||4-39-58||

\twolineshloka
{शैलेषु तेषु सर्वेषु कन्दरेषु वनेषु च}
{य च नोक्ता मया देशा विचेया तेषु जानकी} %||4-39-59||

\twolineshloka
{एतावद्वानरैः शक्यं गन्तुं वानरपुङ्गवाः}
{अभास्करममर्यादं न जानीमस्ततः परम्} %||4-39-60||

\twolineshloka
{अधिगम्य तु वैदेहीं निलयं रावणस्य च}
{मासे पूर्णे निवर्तध्वमुदयं प्राप्य पर्वतम्} %||4-39-61||

\twolineshloka
{ऊर्ध्वं मासान्न वस्तव्यं वसन्वध्यो भवेन्मम}
{सिद्धार्थाः संनिवर्तध्वमधिगम्य च मैथिलीम्} %||4-39-62||

\fourlineindentedshloka
{महेन्द्रकान्तां वनषण्ड मण्डिताम्}
{ दिशं चरित्वा निपुणेन वानराः}
{अवाप्य सीतां रघुवंशजप्रियाम्}
{ ततो निवृत्ताः सुखितो भविष्यथ} %||4-40-63||


{॥इत्यार्षे श्रीमद्रामायणे वाल्मीकीये आदिकाव्ये किष्किन्धाकाण्डे एकोनचत्वारिंशः सर्गः॥}

\sect{चत्वारिंशः सर्गः}

\twolineshloka
{ततः प्रस्थाप्य सुग्रीवस्तन्महद्वानरं बलम्}
{दक्षिणां प्रेषयामास वानरानभिलक्षितान्} %||4-40-1||

\twolineshloka
{नीलमग्निसुतं चैव हनुमन्तं च वानरम्}
{पितामहसुतं चैव जाम्बवन्तं महाकपिम्} %||4-40-2||

\twolineshloka
{सुहोत्रं च शरीरं च शरगुल्मं तथैव च}
{गजं गवाक्षं गवयं सुषेणमृषभं तथा} %||4-40-3||

\twolineshloka
{मैन्दं च द्विविदं चैव विजयं गन्धमादनम्}
{उल्कामुखमसङ्गं च हुताशन सुतावुभौ} %||4-40-4||

\twolineshloka
{अङ्गदप्रमुखान्वीरान्वीरः कपिगणेश्वरः}
{वेगविक्रमसम्पन्नान्सन्दिदेश विशेषवित्} %||4-40-5||

\twolineshloka
{तेषामग्रेषरं चैव महद्बलमसङ्गगम्}
{विधाय हरिवीराणामादिशद्दक्षिणां दिशम्} %||4-40-6||

\twolineshloka
{ये केचन समुद्देशास्तस्यां दिशि सुदुर्गमाः}
{कपीशः कपिमुख्यानां स तेषां तानुदाहरत्} %||4-40-7||

\twolineshloka
{सहस्रशिरसं विन्ध्यं नानाद्रुमलतावृतम्}
{नर्मदां च नदीं दुर्गां महोरगनिषेविताम्} %||4-40-8||

\twolineshloka
{ततो गोदावरीं रम्यां कृष्णावेणीं महानदीम्}
{वरदां च महाभागां महोरगनिषेविताम्} %||4-40-9||

\twolineshloka
{मेखलानुत्कलांश्चैव दशार्णनगराण्यपि}
{अवन्तीमभ्रवन्तीं च सर्वमेवानुपश्यत} %||4-40-10||

\twolineshloka
{विदर्भानृषिकांश्चैव रम्यान्माहिषकानपि}
{तथा बङ्गान्कलिङ्गांश्च कौशिकांश्च समन्ततः} %||4-40-11||

\twolineshloka
{अन्वीक्ष्य दण्डकारण्यं सपर्वतनदीगुहम्}
{नदीं गोदावरीं चैव सर्वमेवानुपश्यत} %||4-40-12||

\twolineshloka
{तथैवान्ध्रांश्च पुण्ड्रांश्च चोलान्पाण्ड्यान्सकेरलान्}
{अयोमुखश्च गन्तव्यः पर्वतो धातुमण्डितः} %||4-40-13||

\twolineshloka
{विचित्रशिखरः श्रीमांश्चित्रपुष्पितकाननः}
{सचन्दनवनोद्देशो मार्गितव्यो महागिरिः} %||4-40-14||

\twolineshloka
{ततस्तामापगां दिव्यां प्रसन्नसलिलां शिवाम्}
{तत्र द्रक्ष्यथ कावेरीं विहृतामप्सरोगणैः} %||4-40-15||

\twolineshloka
{तस्यासीनं नगस्याग्रे मलयस्य महौजसम्}
{द्रक्ष्यथादित्यसङ्काशमगस्त्यमृषिसत्तमम्} %||4-40-16||

\twolineshloka
{ततस्तेनाभ्यनुज्ञाताः प्रसन्नेन महात्मना}
{ताम्रपर्णीं ग्राहजुष्टां तरिष्यथ महानदीम्} %||4-40-17||

\twolineshloka
{सा चन्दनवनैर्दिव्यैः प्रच्छन्ना द्वीप शालिनी}
{कान्तेव युवतिः कान्तं समुद्रमवगाहते} %||4-40-18||

\twolineshloka
{ततो हेममयं दिव्यं मुक्तामणिविभूषितम्}
{युक्तं कवाटं पाण्ड्यानां गता द्रक्ष्यथ वानराः} %||4-40-19||

\twolineshloka
{ततः समुद्रमासाद्य सम्प्रधार्यार्थनिश्चयम्}
{अगस्त्येनान्तरे तत्र सागरे विनिवेशितः} %||4-40-20||

\twolineshloka
{चित्रनानानगः श्रीमान्महेन्द्रः पर्वतोत्तमः}
{जातरूपमयः श्रीमानवगाढो महार्णवम्} %||4-40-21||

\twolineshloka
{नानाविधैर्नगैः फुल्लैर्लताभिश्चोपशोभितम्}
{देवर्षियक्षप्रवरैरप्सरोभिश्च सेवितम्} %||4-40-22||

\twolineshloka
{सिद्धचारणसङ्घैश्च प्रकीर्णं सुमनोहरम्}
{तमुपैति सहस्राक्षः सदा पर्वसु पर्वसु} %||4-40-23||

\threelineshloka
{द्वीपस्तस्यापरे पारे शतयोजनमायतः}
{अगम्यो मानुषैर्दीप्तस्तं मार्गध्वं समन्ततः}
{तत्र सर्वात्मना सीता मार्गितव्या विशेषतः} %||4-40-24||

\twolineshloka
{स हि देशस्तु वध्यस्य रावणस्य दुरात्मनः}
{राक्षसाधिपतेर्वासः सहस्राक्षसमद्युतेः} %||4-40-25||

\twolineshloka
{दक्षिणस्य समुद्रस्य मध्ये तस्य तु राक्षसी}
{अङ्गारकेति विख्याता छायामाक्षिप्य भोजिनी} %||4-40-26||

\twolineshloka
{तमतिक्रम्य लक्ष्मीवान्समुद्रे शतयोजने}
{गिरिः पुष्पितको नाम सिद्धचारणसेवितः} %||4-40-27||

\twolineshloka
{चन्द्रसूर्यांशुसङ्काशः सागराम्बुसमावृतः}
{भ्राजते विपुलैः शृङ्गैरम्बरं विलिखन्निव} %||4-40-28||

\twolineshloka
{तस्यैकं काञ्चनं शृङ्गं सेवते यं दिवाकरः}
{श्वेतं राजतमेकं च सेवते यं निशाकरः} %||4-40-29||

\twolineshloka
{न तं कृतघ्नाः पश्यन्ति न नृशंसा न नास्तिकाः}
{प्रणम्य शिरसा शैलं तं विमार्गत वानराः} %||4-40-30||

\twolineshloka
{तमतिक्रम्य दुर्धर्षाः सूर्यवान्नाम पर्वतः}
{अध्वना दुर्विगाहेन योजनानि चतुर्दश} %||4-40-31||

\twolineshloka
{ततस्तमप्यतिक्रम्य वैद्युतो नाम पर्वतः}
{सर्वकामफलैर्वृक्षैः सर्वकालमनोहरैः} %||4-40-32||

\twolineshloka
{तत्र भुक्त्वा वरार्हाणि मूलानि च फलानि च}
{मधूनि पीत्वा मुख्यानि परं गच्छत वानराः} %||4-40-33||

\twolineshloka
{तत्र नेत्रमनःकान्तः कुञ्जरो नाम पर्वतः}
{अगस्त्यभवनं यत्र निर्मितं विश्वकर्मणा} %||4-40-34||

\twolineshloka
{तत्र योजनविस्तारमुच्छ्रितं दशयोजनम्}
{शरणं काञ्चनं दिव्यं नानारत्नविभूषितम्} %||4-40-35||

\threelineshloka
{तत्र भोगवती नाम सर्पाणामालयः पुरी}
{विशालरथ्या दुर्धर्षा सर्वतः परिरक्षिता}
{रक्षिता पन्नगैर्घोरैस्तीक्ष्णदंष्ट्रैर्महाविषैः} %||4-40-36||

\twolineshloka
{सर्पराजो महाघोरो यस्यां वसति वासुकिः}
{निर्याय मार्गितव्या च सा च भोगवती पुरी} %||4-40-37||

\twolineshloka
{तं च देशमतिक्रम्य महानृषभसंस्थितः}
{सर्वरत्नमयः श्रीमानृषभो नाम पर्वतः} %||4-40-38||

\twolineshloka
{गोशीर्षकं पद्मकं च हरिश्यामं च चन्दनम्}
{दिव्यमुत्पद्यते यत्र तच्चैवाग्निसमप्रभम्} %||4-40-39||

\twolineshloka
{न तु तच्चन्दनं दृष्ट्वा स्प्रष्टव्यं च कदाचन}
{रोहिता नाम गन्धर्वा घोरा रक्षन्ति तद्वनम्} %||4-40-40||

\twolineshloka
{तत्र गन्धर्वपतयः पञ्चसूर्यसमप्रभाः}
{शैलूषो ग्रामणीर्भिक्षुः शुभ्रो बभ्रुस्तथैव च} %||4-40-41||

\threelineshloka
{अन्ते पृथिव्या दुर्धर्षास्तत्र स्वर्गजितः स्थिताः}
{ततः परं न वः सेव्यः पितृलोकः सुदारुणः}
{राजधानी यमस्यैषा कष्टेन तमसावृता} %||4-40-42||

\twolineshloka
{एतावदेव युष्माभिर्वीरा वानरपुङ्गवाः}
{शक्यं विचेतुं गन्तुं वा नातो गतिमतां गतिः} %||4-40-43||

\twolineshloka
{सर्वमेतत्समालोक्य यच्चान्यदपि दृश्यते}
{गतिं विदित्वा वैदेह्याः संनिवर्तितमर्हथ} %||4-40-44||

\twolineshloka
{यस्तु मासान्निवृत्तोऽग्रे दृष्टा सीतेति वक्ष्यति}
{मत्तुल्यविभवो भोगैः सुखं स विहरिष्यति} %||4-40-45||

\twolineshloka
{ततः प्रियतरो नास्ति मम प्राणाद्विशेषतः}
{कृतापराधो बहुशो मम बन्धुर्भविष्यति} %||4-40-46||

\fourlineindentedshloka
{अमितबलपराक्रमा भवन्तो}
{ विपुलगुणेषु कुलेषु च प्रसूताः}
{मनुजपतिसुतां यथा लभध्वम्}
{ तदधिगुणं पुरुषार्थमारभध्वम्} %||4-41-47||


{॥इत्यार्षे श्रीमद्रामायणे वाल्मीकीये आदिकाव्ये किष्किन्धाकाण्डे चत्वारिंशः सर्गः॥}

\sect{एकचत्वारिंशः सर्गः}

\twolineshloka
{ततः प्रस्थाप्य सुग्रीवस्तान्हरीन्दक्षिणां दिशम्}
{बुद्धिविक्रमसम्पन्नान्वायुवेगसमाञ्जवे} %||4-41-1||

\twolineshloka
{अथाहूय महातेजाः सुषेणं नाम यूथपम्}
{तारायाः पितरं राजा श्वशुरभीमविक्रमम्} %||4-41-2||

\twolineshloka
{अब्रवीत्प्राञ्जलिर्वाक्यमभिगम्य प्रणम्य च}
{साहाय्यं कुरु रामस्य कृत्येऽस्मिन्समुपस्थिते} %||4-41-3||

\twolineshloka
{वृतः शतसहस्रेण वानराणां तरस्विनाम्}
{अभिगच्छ दिशं सौम्य पश्चिमां वारुणीं प्रभो} %||4-41-4||

\twolineshloka
{सुराष्ट्रान्सह बाह्लीकाञ्शूराभीरांस्तथैव च}
{स्फीताञ्जनपदान्रम्यान्विपुलानि पुराणि च} %||4-41-5||

\twolineshloka
{पुंनागगहनं कुक्षिं बहुलोद्दालकाकुलम्}
{तथा केतकषण्डांश्च मार्गध्वं हरियूथपाः} %||4-41-6||

\twolineshloka
{प्रत्यक्स्रोतोगमाश्चैव नद्यः शीतजलाः शिवाः}
{तापसानामरण्यानि कान्तारा गिरयश्च ये} %||4-41-7||

\threelineshloka
{गिरिजालावृतां दुर्गां मार्गित्वा पश्चिमां दिशम्}
{ततः पश्चिममासाद्य समुद्रं द्रष्टुमर्हथ}
{तिमि नक्रायुत जलमक्षोभ्यमथ वानरः} %||4-41-8||

\twolineshloka
{ततः केतकषण्डेषु तमालगहनेषु च}
{कपयो विहरिष्यन्ति नारिकेलवनेषु च} %||4-41-9||

\twolineshloka
{तत्र सीतां च मार्गध्वं निलयं रावणस्य च}
{मरीचिपत्तनं चैव रम्यं चैव जटीपुरम्} %||4-41-10||

\twolineshloka
{अवन्तीमङ्गलोपां च तथा चालक्षितं वनम्}
{राष्ट्राणि च विशालानि पत्तनानि ततस्ततः} %||4-41-11||

\twolineshloka
{सिन्धुसागरयोश्चैव सङ्गमे तत्र पर्वतः}
{महान्हेमगिरिर्नाम शतशृङ्गो महाद्रुमः} %||4-41-12||

\twolineshloka
{तस्य प्रस्थेषु रम्येषु सिंहाः पक्षगमाः स्थिताः}
{तिमिमत्स्यगजांश्चैव नीडान्यारोपयन्ति ते} %||4-41-13||

\threelineshloka
{तानि नीडानि सिंहानां गिरिशृङ्गगताश्च ये}
{दृप्तास्तृप्ताश्च मातङ्गास्तोयदस्वननिःस्वनाः}
{विचरन्ति विशालेऽस्मिंस्तोयपूर्णे समन्ततः} %||4-41-14||

\twolineshloka
{तस्य शृङ्गं दिवस्पर्शं काञ्चनं चित्रपादपम्}
{सर्वमाशु विचेतव्यं कपिभिः कामरूपिभिः} %||4-41-15||

\twolineshloka
{कोटिं तत्र समुद्रे तु काञ्चनीं शतयोजनम्}
{दुर्दर्शां परियात्रस्य गता द्रक्ष्यथ वानराः} %||4-41-16||

\twolineshloka
{कोट्यस्तत्र चतुर्विंशद्गन्धर्वाणां तरस्विनाम्}
{वसन्त्यग्निनिकाशानां घोराणां कामरूपिणाम्} %||4-41-17||

\twolineshloka
{नात्यासादयितव्यास्ते वानरैर्भीमविक्रमैः}
{नादेयं च फलं तस्माद्देशात्किञ्चित्प्लवङ्गमैः} %||4-41-18||

\twolineshloka
{दुरासदा हि ते वीराः सत्त्ववन्तो महाबलाः}
{फलमूलानि ते तत्र रक्षन्ते भीमविक्रमाः} %||4-41-19||

\twolineshloka
{तत्र यत्नश्च कर्तव्यो मार्गितव्या च जानकी}
{न हि तेभ्यो भयं किञ्चित्कपित्वमनुवर्तताम्} %||4-41-20||

\twolineshloka
{चतुर्भागे समुद्रस्य चक्रवान्नाम पर्वतः}
{तत्र चक्रं सहस्रारं निर्मितं विश्वकर्मणा} %||4-41-21||

\twolineshloka
{तत्र पञ्चजनं हत्वा हयग्रीवं च दानवम्}
{आजहार ततश्चक्रं शङ्खं च पुरुषोत्तमः} %||4-41-22||

\twolineshloka
{तस्य सानुषु चित्रेषु विशालासु गुहासु च}
{रावणः सह वैदेह्या मार्गितव्यस्ततस्ततः} %||4-41-23||

\twolineshloka
{योजनानि चतुःषष्टिर्वराहो नाम पर्वतः}
{सुवर्णशृङ्गः सुश्रीमानगाधे वरुणालये} %||4-41-24||

\twolineshloka
{तत्र प्राग्ज्योतिषं नाम जातरूपमयं पुरम्}
{यस्मिन्वस्ति दुष्टात्मा नरको नाम गुहासु च} %||4-41-25||

\twolineshloka
{तस्य सानुषु चित्रेषु विशालासु गुहासु च}
{रावणः सह वैदेह्या मार्गितव्यस्ततस्ततः} %||4-41-26||

\twolineshloka
{तमतिक्रम्य शैलेन्द्रं काञ्चनान्तरनिर्दरः}
{पर्वतः सर्वसौवर्णो धारा प्रस्रवणायुतः} %||4-41-27||

\twolineshloka
{तं गजाश्च वराहाश्च सिंहा व्याघ्राश्च सर्वतः}
{अभिगर्जन्ति सततं तेन शब्देन दर्पिताः} %||4-41-28||

\twolineshloka
{तस्मिन्हरिहयः श्रीमान्महेन्द्रः पाकशासनः}
{अभिषिक्तः सुरै राजा मेघवान्नाम पर्वतः} %||4-41-29||

\twolineshloka
{तमतिक्रम्य शैलेन्द्रं महेन्द्रपरिपालितम्}
{षष्टिं गिरिसहस्राणि काञ्चनानि गमिष्यथ} %||4-41-30||

\twolineshloka
{तरुणादित्यवर्णानि भ्राजमानानि सर्वतः}
{जातरूपमयैर्वृक्षैः शोभितानि सुपुष्पितैः} %||4-41-31||

\twolineshloka
{तेषां मध्ये स्थितो राजा मेरुरुत्तमपर्वतः}
{आदित्येन प्रसन्नेन शैलो दत्तवरः पुरा} %||4-41-32||

\twolineshloka
{तेनैवमुक्तः शैलेन्द्रः सर्व एव त्वदाश्रयाः}
{मत्प्रसादाद्भविष्यन्ति दिवारात्रौ च काञ्चनाः} %||4-41-33||

\twolineshloka
{त्वयि ये चापि वत्स्यन्ति देवगन्धर्वदानवाः}
{ते भविष्यन्ति रक्ताश्च प्रभया काञ्चनप्रभाः} %||4-41-34||

\twolineshloka
{आदित्या वसवो रुद्रा मरुतश्च दिवौकसः}
{आगम्य पश्चिमां सन्ध्यां मेरुमुत्तमपर्वतम्} %||4-41-35||

\twolineshloka
{आदित्यमुपतिष्ठन्ति तैश्च सूर्योऽभिपूजितः}
{अदृश्यः सर्वभूतानामस्तं गच्छति पर्वतम्} %||4-41-36||

\twolineshloka
{योजनानां सहस्राणि दशतानि दिवाकरः}
{मुहूर्तार्धेन तं शीघ्रमभियाति शिलोच्चयम्} %||4-41-37||

\twolineshloka
{शृङ्गे तस्य महद्दिव्यं भवनं सूर्यसंनिभम्}
{प्रासादगुणसम्बाधं विहितं विश्वकर्मणा} %||4-41-38||

\twolineshloka
{शोभितं तरुभिश्चित्रैर्नानापक्षिसमाकुलैः}
{निकेतं पाशहस्तस्य वरुणस्य महात्मनः} %||4-41-39||

\twolineshloka
{अन्तरा मेरुमस्तं च तालो दशशिरा महान्}
{जातरूपमयः श्रीमान्भ्राजते चित्रवेदिकः} %||4-41-40||

\twolineshloka
{तेषु सर्वेषु दुर्गेषु सरःसु च सरित्सु च}
{रावणः सह वैदेह्या मार्गितव्यस्ततस्ततः} %||4-41-41||

\twolineshloka
{यत्र तिष्ठति धर्मात्मा तपसा स्वेन भावितः}
{मेरुसावर्णिरित्येव ख्यातो वै ब्रह्मणा समः} %||4-41-42||

\twolineshloka
{प्रष्टव्यो मेरुसावर्णिर्महर्षिः सूर्यसंनिभः}
{प्रणम्य शिरसा भूमौ प्रवृत्तिं मैथिलीं प्रति} %||4-41-43||

\twolineshloka
{एतावज्जीवलोकस्य भास्करो रजनीक्षये}
{कृत्वा वितिमिरं सर्वमस्तं गच्छति पर्वतम्} %||4-41-44||

\twolineshloka
{एतावद्वानरैः शक्यं गन्तुं वानरपुङ्गवाः}
{अभास्करममर्यादं न जानीमस्ततः परम्} %||4-41-45||

\twolineshloka
{अधिगम्य तु वैदेहीं निलयं रावणस्य च}
{अस्तं पर्वतमासाद्य पूर्णे मासे निवर्तत} %||4-41-46||

\twolineshloka
{ऊर्ध्वं मासान्न वस्तव्यं वसन्वध्यो भवेन्मम}
{सहैव शूरो युष्माभिः श्वशुरो मे गमिष्यति} %||4-41-47||

\twolineshloka
{श्रोतव्यं सर्वमेतस्य भवद्भिर्दिष्ट कारिभिः}
{गुरुरेष महाबाहुः श्वशुरो मे महाबलः} %||4-41-48||

\twolineshloka
{भवन्तश्चापि विक्रान्ताः प्रमाणं सर्वकर्मसु}
{प्रमाणमेनं संस्थाप्य पश्यध्वं पश्चिमां दिशम्} %||4-41-49||

\twolineshloka
{दृष्टायां तु नरेन्द्रस्या पत्न्याममिततेजसः}
{कृतकृत्या भविष्यामः कृतस्य प्रतिकर्मणा} %||4-41-50||

\twolineshloka
{अतोऽन्यदपि यत्किञ्चित्कार्यस्यास्य हितं भवेत्}
{सम्प्रधार्य भवद्भिश्च देशकालार्थसंहितम्} %||4-41-51||

\fourlineindentedshloka
{ततः सुषेण प्रमुखाः प्लवङ्गमाः}
{ सुग्रीववाक्यं निपुणं निशम्य}
{आमन्त्र्य सर्वे प्लवगाधिपं ते}
{ जग्मुर्दिशं तां वरुणाभिगुप्ताम्} %||4-42-52||


{॥इत्यार्षे श्रीमद्रामायणे वाल्मीकीये आदिकाव्ये किष्किन्धाकाण्डे एकचत्वारिंशः सर्गः॥}

\sect{द्विचत्वारिंशः सर्गः}

\twolineshloka
{ततः सन्दिश्य सुग्रीवः श्वशुरं पश्चिमां दिशम्}
{वीरं शतबलिं नाम वानरं वानरर्षभः} %||4-42-1||

\twolineshloka
{उवाच राजा मन्त्रज्ञः सर्ववानरसम्मतम्}
{वाक्यमात्महितं चैव रामस्य च हितं तथा} %||4-42-2||

\twolineshloka
{वृतः शतसहस्रेण त्वद्विधानां वनौकसाम्}
{वैवस्वत सुतैः सार्धं प्रतिष्ठस्व स्वमन्त्रिभिः} %||4-42-3||

\twolineshloka
{दिशं ह्युदीचीं विक्रान्तां हिमशैलावतंसकाम्}
{सर्वतः परिमार्गध्वं रामपत्नीमनिन्दिताम्} %||4-42-4||

\twolineshloka
{अस्मिन्कार्ये विनिवृत्ते कृते दाशरथेः प्रिये}
{ऋणान्मुक्ता भविष्यामः कृतार्थार्थविदां वराः} %||4-42-5||

\twolineshloka
{कृतं हि प्रियमस्माकं राघवेण महात्मना}
{तस्य चेत्प्रतिकारोऽस्ति सफलं जीवितं भवेत्} %||4-42-6||

\twolineshloka
{एतां बुद्धिं समास्थाय दृश्यते जानकी यथा}
{तथा भवद्भिः कर्तव्यमस्मत्प्रियहितैषिभिः} %||4-42-7||

\twolineshloka
{अयं हि सर्वभूतानां मान्यस्तु नरसत्तमः}
{अस्मासु चागतप्रीती रामः परपुरंजयः} %||4-42-8||

\twolineshloka
{इमानि वनदुर्गाणि नद्यः शैलान्तराणि च}
{भवन्तः परिमार्गंस्तु बुद्धिविक्रमसम्पदा} %||4-42-9||

\twolineshloka
{तत्र म्लेच्छान्पुलिन्दांश्च शूरसेनांस्तथैव च}
{प्रस्थालान्भरतांश्चैव कुरूंश्च सह मद्रकैः} %||4-42-10||

\twolineshloka
{काम्बोजान्यवनांश्चैव शकानारट्टकानपि}
{बाह्लीकानृषिकांश्चैव पौरवानथ टङ्कणान्} %||4-42-11||

\twolineshloka
{चीनान्परमचीनांश्च नीहारांश्च पुनः पुनः}
{अन्विष्य दरदांश्चैव हिमवन्तं विचिन्वथ} %||4-42-12||

\twolineshloka
{लोध्रपद्मकषण्डेषु देवदारुवनेषु च}
{रावणः सह वैदेह्य मार्गितव्यस्ततस्ततः} %||4-42-13||

\twolineshloka
{ततः सोमाश्रमं गत्वा देवगन्धर्वसेवितम्}
{कालं नाम महासानुं पर्वतं तं गमिष्यथ} %||4-42-14||

\twolineshloka
{महत्सु तस्य शृङ्गेषु निर्दरेषु गुहासु च}
{विचिनुध्वं महाभागां रामपत्नीं यशस्विनीम्} %||4-42-15||

\twolineshloka
{तमतिक्रम्य शैलेन्द्रं हेमवर्गं महागिरिम्}
{ततः सुदर्शनं नाम पर्वतं गन्तुमर्हथ} %||4-42-16||

\twolineshloka
{तस्य काननषण्डेषु निर्दरेषु गुहासु च}
{रावणः सह वैदेह्या मार्गितव्यस्ततस्ततः} %||4-42-17||

\twolineshloka
{तमतिक्रम्य चाकाशं सर्वतः शतयोजनम्}
{अपर्वतनदी वृक्षं सर्वसत्त्वविवर्जितम्} %||4-42-18||

\twolineshloka
{तं तु शीघ्रमतिक्रम्य कान्तारं रोमहर्षणम्}
{कैलासं पाण्डुरं शैलं प्राप्य हृष्टा भविष्यथ} %||4-42-19||

\twolineshloka
{तत्र पाण्डुरमेघाभं जाम्बूनदपरिष्कृतम्}
{कुबेरभवनं दिव्यं निर्मितं विश्वकर्मणा} %||4-42-20||

\twolineshloka
{विशाला नलिनी यत्र प्रभूतकमलोत्पला}
{हंसकारण्डवाकीर्णा अप्सरोगणसेविता} %||4-42-21||

\twolineshloka
{तत्र वैश्रवणो राजा सर्वभूतनमस्कृतः}
{धनदो रमते श्रीमान्गुह्यकैः सह यक्षराट्} %||4-42-22||

\twolineshloka
{तस्य चन्द्रनिकशेषु पर्वतेषु गुहासु च}
{रावणः सह वैदेह्या मार्गितव्यस्ततस्ततः} %||4-42-23||

\twolineshloka
{क्रौञ्चं तु गिरिमासाद्य बिलं तस्य सुदुर्गमम्}
{अप्रमत्तैः प्रवेष्टव्यं दुष्प्रवेशं हि तत्स्मृतम्} %||4-42-24||

\twolineshloka
{वसन्ति हि महात्मानस्तत्र सूर्यसमप्रभाः}
{देवैरप्यर्चिताः सम्यग्देवरूपा महर्षयः} %||4-42-25||

\twolineshloka
{क्रौञ्चस्य तु गुहाश्चान्याः सानूनि शिखराणि च}
{निर्दराश्च नितम्बाश्च विचेतव्यास्ततस्ततः} %||4-42-26||

\twolineshloka
{क्रौञ्चस्य शिखरं चापि निरीक्ष्य च ततस्ततः}
{अवृक्षं कामशैलं च मानसं विहगालयम्} %||4-42-27||

\twolineshloka
{न गतिस्तत्र भूतानां देवदानवरक्षसाम्}
{स च सर्वैर्विचेतव्यः ससानुप्रस्थभूधरः} %||4-42-28||

\twolineshloka
{क्रौञ्चं गिरिमतिक्रम्य मैनाको नाम पर्वतः}
{मयस्य भवनं तत्र दानवस्य स्वयं कृतम्} %||4-42-29||

\twolineshloka
{मैनाकस्तु विचेतव्यः ससानुप्रस्थकन्दरः}
{स्त्रीणामश्वमुखीनां च निकेतास्तत्र तत्र तु} %||4-42-30||

\twolineshloka
{तं देशं समतिक्रम्य आश्रमं सिद्धसेवितम्}
{सिद्धा वैखानसास्तत्र वालखिल्याश्च तापसाः} %||4-42-31||

\twolineshloka
{वन्द्यास्ते तु तपःसिद्धास्तापसा वीतकल्मषाः}
{प्रष्टव्याश्चापि सीतायाः प्रवृत्तं विनयान्वितैः} %||4-42-32||

\twolineshloka
{हेमपुष्करसञ्छन्नं तत्र वैखानसं सरः}
{तरुणादित्यसङ्काशैर्हंसैर्विचरितं शुभैः} %||4-42-33||

\twolineshloka
{औपवाह्यः कुबेरस्य सर्वभौम इति स्मृतः}
{गजः पर्येति तं देशं सदा सह करेणुभिः} %||4-42-34||

\twolineshloka
{तत्सारः समतिक्रम्य नष्टचन्द्रदिवाकरम्}
{अनक्षत्रगणं व्योम निष्पयोदमनादिमत्} %||4-42-35||

\twolineshloka
{गभस्तिभिरिवार्कस्य स तु देशः प्रकाशते}
{विश्राम्यद्भिस्तपः सिद्धैर्देवकल्पैः स्वयम्प्रभैः} %||4-42-36||

\twolineshloka
{तं तु देशमतिक्रम्य शैलोदा नाम निम्नगा}
{उभयोस्तीरयोर्यस्याः कीचका नाम वेणवः} %||4-42-37||

\twolineshloka
{ते नयन्ति परं तीरं सिद्धान्प्रत्यानयन्ति च}
{उत्तराः कुरवस्तत्र कृतपुण्यप्रतिश्रियाः} %||4-42-38||

\twolineshloka
{ततः काञ्चनपद्माभिः पद्मिनीभिः कृतोदकाः}
{नीलवैदूर्यपत्राढ्या नद्यस्तत्र सहस्रशः} %||4-42-39||

\twolineshloka
{रक्तोत्पलवनैश्चात्र मण्डिताश्च हिरण्मयैः}
{तरुणादित्यसदृशैर्भान्ति तत्र जलाशयाः} %||4-42-40||

\twolineshloka
{महार्हमणिपत्रैश्च काञ्चनप्रभ केसरैः}
{नीलोत्पलवनैश्चित्रैः स देशः सर्वतोवृतः} %||4-42-41||

\twolineshloka
{निस्तुलाभिश्च मुक्ताभिर्मणिभिश्च महाधनैः}
{उद्भूतपुलिनास्तत्र जातरूपैश्च निम्नगाः} %||4-42-42||

\twolineshloka
{सर्वरत्नमयैश्चित्रैरवगाढा नगोत्तमैः}
{जातरूपमयैश्चापि हुताशनसमप्रभैः} %||4-42-43||

\twolineshloka
{नित्यपुष्पफलाश्चात्र नगाः पत्ररथाकुलाः}
{दिव्यगन्धरसस्पर्शाः सर्वकामान्स्रवन्ति च} %||4-42-44||

\twolineshloka
{नानाकाराणि वासांसि फलन्त्यन्ये नगोत्तमाः}
{मुक्तावैदूर्यचित्राणि भूषणानि तथैव च} %||4-42-45||

\twolineshloka
{स्त्रीणां यान्यनुरूपाणि पुरुषाणां तथैव च}
{सर्वर्तुसुखसेव्यानि फलन्त्यन्ये नगोत्तमाः} %||4-42-46||

\twolineshloka
{महार्हाणि विचित्राणि हैमान्यन्ये नगोत्तमाः}
{शयनानि प्रसूयन्ते चित्रास्तारणवन्ति च} %||4-42-47||

\twolineshloka
{मनःकान्तानि माल्यानि फलन्त्यत्रापरे द्रुमाः}
{पानानि च महार्हाणि भक्ष्याणि विविधानि च} %||4-42-48||

\threelineshloka
{स्त्रियश्च गुणसम्पन्ना रूपयौवनलक्षिताः}
{गन्धर्वाः किंनरा सिद्धा नागा विद्याधरास्तथा}
{रमन्ते सहितास्तत्र नारीभिर्भास्करप्रभाः} %||4-42-49||

\twolineshloka
{सर्वे सुकृतकर्माणः सर्वे रतिपरायणाः}
{सर्वे कामार्थसहिता वसन्ति सह योषितः} %||4-42-50||

\twolineshloka
{गीतवादित्रनिर्घोषः सोत्कृष्टहसितस्वनः}
{श्रूयते सततं तत्र सर्वभूतमनोहरः} %||4-42-51||

\twolineshloka
{तत्र नामुदितः कश्चिन्नास्ति कश्चिदसत्प्रियः}
{अहन्यहनि वर्धन्ते गुणास्तत्र मनोरमाः} %||4-42-52||

\twolineshloka
{समतिक्रम्य तं देशमुत्तरस्तोयसां निधिः}
{तत्र सोमगिरिर्नाम मध्ये हेममयो महान्} %||4-42-53||

\twolineshloka
{इन्द्रलोकगता ये च ब्रह्मलोकगताश्च ये}
{देवास्तं समवेक्षन्ते गिरिराजं दिवं गतम्} %||4-42-54||

\twolineshloka
{स तु देशो विसूर्योऽपि तस्य भासा प्रकाशते}
{सूर्यलक्ष्म्याभिविज्ञेयस्तपसेव विवस्वता} %||4-42-55||

\twolineshloka
{भगवानपि विश्वात्मा शम्भुरेकादशात्मकः}
{ब्रह्मा वसति देवेशो ब्रह्मर्षिपरिवारितः} %||4-42-56||

\twolineshloka
{न कथञ्चन गन्तव्यं कुरूणामुत्तरेण वः}
{अन्येषामपि भूतानां नातिक्रामति वै गतिः} %||4-42-57||

\twolineshloka
{सा हि सोमगिरिर्नाम देवानामपि दुर्गमः}
{तमालोक्य ततः क्षिप्रमुपावर्तितुमर्हथ} %||4-42-58||

\twolineshloka
{एतावद्वानरैः शक्यं गन्तुं वानरपुङ्गवाः}
{अभास्करममर्यादं न जानीमस्ततः परम्} %||4-42-59||

\twolineshloka
{सर्वमेतद्विचेतव्यं यन्मया परिकीर्तितम्}
{यदन्यदपि नोक्तं च तत्रापि क्रियतां मतिः} %||4-42-60||

\fourlineindentedshloka
{ततः कृतं दाशरथेर्महत्प्रियम्}
{ महत्तरं चापि ततो मम प्रियम्}
{कृतं भविष्यत्यनिलानलोपमा}
{ विदेहजा दर्शनजेन कर्मणा} %||4-42-61||

\fourlineindentedshloka
{ततः कृतार्थाः सहिताः सबान्धवा}
{ मयार्चिताः सर्वगुणैर्मनोरमैः}
{चरिष्यथोर्वीं प्रतिशान्तशत्रवः}
{ सहप्रिया भूतधराः प्लवङ्गमाः} %||4-43-62||


{॥इत्यार्षे श्रीमद्रामायणे वाल्मीकीये आदिकाव्ये किष्किन्धाकाण्डे द्विचत्वारिंशः सर्गः॥}

\sect{त्रिचत्वारिंशः सर्गः}

\twolineshloka
{विशेषेण तु सुग्रीवो हनुमत्यर्थमुक्तवान्}
{स हि तस्मिन्हरिश्रेष्ठे निश्चितार्थोऽर्थसाधने} %||4-43-1||

\twolineshloka
{न भूमौ नान्तरिक्षे वा नाम्बरे नामरालये}
{नाप्सु वा गतिसङ्गं ते पश्यामि हरिपुङ्गव} %||4-43-2||

\twolineshloka
{सासुराः सहगन्धर्वाः सनागनरदेवताः}
{विदिताः सर्वलोकास्ते ससागरधराधराः} %||4-43-3||

\twolineshloka
{गतिर्वेगश्च तेजश्च लाघवं च महाकपे}
{पितुस्ते सदृशं वीर मारुतस्य महौजसः} %||4-43-4||

\twolineshloka
{तेजसा वापि ते भूतं समं भुवि न विद्यते}
{तद्यथा लभ्यते सीता तत्त्वमेवोपपादय} %||4-43-5||

\twolineshloka
{त्वय्येव हनुमन्नस्ति बलं बुद्धिः पराक्रमः}
{देशकालानुवृत्तश्च नयश्च नयपण्डित} %||4-43-6||

\twolineshloka
{ततः कार्यसमासङ्गमवगम्य हनूमति}
{विदित्वा हनुमन्तं च चिन्तयामास राघवः} %||4-43-7||

\twolineshloka
{सर्वथा निश्चितार्थोऽयं हनूमति हरीश्वरः}
{निश्चितार्थतरश्चापि हनूमान्कार्यसाधने} %||4-43-8||

\twolineshloka
{तदेवं प्रस्थितस्यास्य परिज्ञातस्य कर्मभिः}
{भर्त्रा परिगृहीतस्य ध्रुवः कार्यफलोदयः} %||4-43-9||

\twolineshloka
{तं समीक्ष्य महातेजा व्यवसायोत्तरं हरिम्}
{कृतार्थ इव संवृत्तः प्रहृष्टेन्द्रियमानसः} %||4-43-10||

\twolineshloka
{ददौ तस्य ततः प्रीतः स्वनामाङ्कोपशोभितम्}
{अङ्गुलीयमभिज्ञानं राजपुत्र्याः परन्तपः} %||4-43-11||

\twolineshloka
{अनेन त्वां हरिश्रेष्ठ चिह्नेन जनकात्मजा}
{मत्सकाशादनुप्राप्तमनुद्विग्नानुपश्यति} %||4-43-12||

\twolineshloka
{व्यवसायश्च ते वीर सत्त्वयुक्तश्च विक्रमः}
{सुग्रीवस्य च सन्देशः सिद्धिं कथयतीव मे} %||4-43-13||

\twolineshloka
{स तद्गृह्य हरिश्रेष्ठः स्थाप्य मूर्ध्नि कृताञ्जलिः}
{वन्दित्वा चरणौ चैव प्रस्थितः प्लवगोत्तमः} %||4-43-14||

\fourlineindentedshloka
{स तत्प्रकर्षन्हरिणां बलं महद्-}
{ बभूव वीरः पवनात्मजः कपि}
{गताम्बुदे व्योम्नि विशुद्धमण्डलः}
{ शशीव नक्षत्रगणोपशोभितः} %||4-43-15||

\fourlineindentedshloka
{अतिबलबलमाश्रितस्तवाहम्}
{ हरिवरविक्रमविक्रमैरनल्पैः}
{पवनसुत यथाभिगम्यते सा}
{ जनकसुता हनुमंस्तथा कुरुष्व} %||4-44-16||


{॥इत्यार्षे श्रीमद्रामायणे वाल्मीकीये आदिकाव्ये किष्किन्धाकाण्डे त्रिचत्वारिंशः सर्गः॥}

\sect{चतुश्चत्वारिंशः सर्गः}

\twolineshloka
{तदुग्रशासनं भर्तुर्विज्ञाय हरिपुङ्गवाः}
{शलभा इव सञ्छाद्य मेदिनीं सम्प्रतस्थिरे} %||4-44-1||

\twolineshloka
{रामः प्रस्रवणे तस्मिन्न्यवसत्सहलक्ष्मणः}
{प्रतीक्षमाणस्तं मासं यः सीताधिगमे कृतः} %||4-44-2||

\twolineshloka
{उत्तरां तु दिशं रम्यां गिरिराजसमावृताम्}
{प्रतस्थे सहसा वीरो हरिः शतबलिस्तदा} %||4-44-3||

{पूर्वां दिशं प्रति ययौ विनतो हरियूथपः॥\devanumber{4}॥} %||4-44-4|| (Check)

\addtocounter{shlokacount}{1}

\twolineshloka
{ताराङ्गदादि सहितः प्लवगः पवनात्मजः}
{अगस्त्यचरितामाशां दक्षिणां हरियूथपः} %||4-44-5||

\twolineshloka
{पश्चिमां तु दिशं घोरां सुषेणः प्लवगेश्वरः}
{प्रतस्थे हरिशार्दूलो भृशं वरुणपालिताम्} %||4-44-6||

\twolineshloka
{ततः सर्वा दिशो राजा चोदयित्वा यथा तथम्}
{कपिसेना पतीन्मुख्यान्मुमोद सुखितः सुखम्} %||4-44-7||

\twolineshloka
{एवं सञ्चोदिताः सर्वे राज्ञा वानरयूथपाः}
{स्वां स्वां दिशमभिप्रेत्य त्वरिताः सम्प्रतस्थिरे} %||4-44-8||

\threelineshloka
{नदन्तश्चोन्नदन्तश्च गर्जन्तश्च प्लवङ्गमाः}
{क्ष्वेलन्तो धावमानाश्च ययुः प्लवगसत्तमाः}
{आनयिष्यामहे सीतां हनिष्यामश्च रावणम्} %||4-44-9||

\twolineshloka
{अहमेको हनिष्यामि प्राप्तं रावणमाहवे}
{ततश्चोन्मथ्य सहसा हरिष्ये जनकात्मजाम्} %||4-44-10||

\twolineshloka
{वेपमानं श्रमेणाद्य भवद्भिः स्थीयतामिति}
{एक एवाहरिष्यामि पातालादपि जानकीम्} %||4-44-11||

\twolineshloka
{विधमिष्याम्यहं वृक्षान्दारयिष्याम्यहं गिरीन्}
{धरणीं दारयिष्यामि क्षोभयिष्यामि सागरान्} %||4-44-12||

\twolineshloka
{अहं योजनसङ्ख्यायाः प्लविता नात्र संशयः}
{शतं योजनसङ्ख्यायाः शतं समधिकं ह्यहम्} %||4-44-13||

\twolineshloka
{भूतले सागरे वापि शैलेषु च वनेषु च}
{पातालस्यापि वा मध्ये न ममाच्छिद्यते गतिः} %||4-44-14||

\twolineshloka
{इत्येकैकं तदा तत्र वानरा बलदर्पिताः}
{ऊचुश्च वचनं तस्मिन्हरिराजस्य संनिधौ} %||4-45-15||


{॥इत्यार्षे श्रीमद्रामायणे वाल्मीकीये आदिकाव्ये किष्किन्धाकाण्डे चतुश्चत्वारिंशः सर्गः॥}

\sect{पञ्चचत्वारिंशः सर्गः}

\twolineshloka
{गतेषु वानरेन्द्रेषु रामः सुग्रीवमब्रवीत्}
{कथं भवान्विनाजीते सर्वं वै मण्डलं भुवः} %||4-45-1||

\twolineshloka
{सुग्रीवस्तु ततो राममुवाच प्रणतात्मवान्}
{श्रूयतां सर्वमाख्यास्ये विस्तरेण नरर्षभ} %||4-45-2||

\twolineshloka
{यदा तु दुन्दुभिं नाम दानवं महिषाकृतिम्}
{परिकालयते वाली मलयं प्रति पर्वतम्} %||4-45-3||

\twolineshloka
{तदा विवेश महिषो मलयस्य गुहां प्रति}
{विवेश वाली तत्रापि मलयं तज्जिघांसया} %||4-45-4||

\twolineshloka
{ततोऽहं तत्र निक्षिप्तो गुहाद्वारिविनीतवत्}
{न च निष्क्रमते वाली तदा संवत्सरे गते} %||4-45-5||

\twolineshloka
{ततः क्षतजवेगेन आपुपूरे तदा बिलम्}
{तदहं विस्मितो दृष्ट्वा भ्रातृशोकविषार्दितः} %||4-45-6||

\threelineshloka
{अथाहं कृतबुद्धिस्तु सुव्यक्तं निहतो गुरुः}
{शिलापर्वतसङ्काशा बिलद्वारि मया कृता}
{अशक्नुवन्निष्क्रमितुं महिषो विनशेदिति} %||4-45-7||

\threelineshloka
{ततोऽहमागां किष्किन्धां निराशस्तस्य जीविते}
{राज्यं च सुमहत्प्राप्तं तारा च रुमया सह}
{मित्रैश्च सहितस्तत्र वसामि विगतज्वरः} %||4-45-8||

\twolineshloka
{आजगाम ततो वाली हत्वा तं दानवर्षभम्}
{ततोऽहमददां राज्यं गौरवाद्भययन्त्रितः} %||4-45-9||

\twolineshloka
{स मां जिघांसुर्दुष्टात्मा वाली प्रव्यथितेन्द्रियः}
{परिलाकयते क्रोधाद्धावन्तं सचिवैः सह} %||4-45-10||

\twolineshloka
{ततोऽहं वालिना तेन सानुबन्धः प्रधावितः}
{नदीश्च विविधाः पश्यन्वनानि नगराणि च} %||4-45-11||

\twolineshloka
{आदर्शतलसङ्काशा ततो वै पृथिवी मया}
{अलातचक्रप्रतिमा दृष्टा गोष्पदवत्तदा} %||4-45-12||

\threelineshloka
{ततः पूर्वमहं गत्वा दक्षिणामहमाश्रितः}
{दिशं च पश्चिमां भूयो गतोऽस्मि भयशङ्कितः}
{उत्तरां तु दिशं यान्तं हनुमान्मामथाब्रवीत्} %||4-45-13||

\twolineshloka
{इदानीं मे स्मृतं राजन्यथा वाली हरीश्वरः}
{मतङ्गेन तदा शप्तो ह्यस्मिन्नाश्रममण्डले} %||4-45-14||

\twolineshloka
{प्रविशेद्यदि वा वाली मूर्धास्य शतधा भवेत्}
{तत्र वासः सुखोऽस्माकं निरुद्विग्नो भविष्यति} %||4-45-15||

\twolineshloka
{ततः पर्वतमासाद्य ऋश्यमूकं नृपात्मज}
{न विवेश तदा वाली मतङ्गस्य भयात्तदा} %||4-45-16||

\twolineshloka
{एवं मया तदा राजन्प्रत्यक्षमुपलक्षितम्}
{पृथिवीमण्डलं कृत्स्नं गुहामस्म्यागतस्ततः} %||4-46-17||


{॥इत्यार्षे श्रीमद्रामायणे वाल्मीकीये आदिकाव्ये किष्किन्धाकाण्डे पञ्चचत्वारिंशः सर्गः॥}

\sect{षट्चत्वारिंशः सर्गः}

\twolineshloka
{दर्शनार्थं तु वैदेह्याः सर्वतः कपियूथपाः}
{व्यादिष्टाः कपिराजेन यथोक्तं जग्मुरञ्जसा} %||4-46-1||

\twolineshloka
{सरांसि सरितः कक्षानाकाशं नगराणि च}
{नदीदुर्गांस्तथा शैलान्विचिन्वन्ति समन्ततः} %||4-46-2||

\twolineshloka
{सुग्रीवेण समाख्यातान्सर्वे वानरयूथपाः}
{प्रदेशान्प्रविचिन्वन्ति सशैलवनकाननान्} %||4-46-3||

\twolineshloka
{विचिन्त्य दिवसं सर्वे सीताधिगमने धृताः}
{समायान्ति स्म मेदिन्यां निशाकालेशु वानराः} %||4-46-4||

\twolineshloka
{सर्वर्तुकांश्च देशेषु वानराः सफलान्द्रुमान्}
{आसाद्य रजनीं शय्यां चक्रुः सर्वेष्वहःसु ते} %||4-46-5||

\twolineshloka
{तदहः प्रथमं कृत्वा मासे प्रस्रवणं गताः}
{कपिराजेन सङ्गम्य निराशाः कपियूथपाः} %||4-46-6||

\twolineshloka
{विचित्य तु दिशं पूर्वां यथोक्तां सचिवैः सह}
{अदृष्ट्वा विनतः सीतामाजगाम महाबलः} %||4-46-7||

\twolineshloka
{उत्तरां तु दिशं सर्वां विचित्य स महाकपिः}
{आगतः सह सैन्येन वीरः शतबलिस्तदा} %||4-46-8||

\twolineshloka
{सुषेणः पश्चिमामाशां विचित्य सह वानरैः}
{समेत्य मासे सम्पूर्णे सुग्रीवमुपचक्रमे} %||4-46-9||

\twolineshloka
{तं प्रस्रवणपृष्ठस्थं समासाद्याभिवाद्य च}
{आसीनं सह रामेण सुग्रीवमिदमब्रुवन्} %||4-46-10||

\twolineshloka
{विचिताः पर्वताः सर्वे वनानि नगराणि च}
{निम्नगाः सागरान्ताश्च सर्वे जनपदास्तथा} %||4-46-11||

\twolineshloka
{गुहाश्च विचिताः सर्वा यास्त्वया परिकीर्तिताः}
{विचिताश्च महागुल्मा लताविततसन्तताः} %||4-46-12||

\threelineshloka
{गहनेषु च देशेषु दुर्गेषु विषमेषु च}
{सत्त्वान्यतिप्रमाणानि विचितानि हतानि च}
{ये चैव गहना देशा विचितास्ते पुनः पुनः} %||4-46-13||

\fourlineindentedshloka
{उदारसत्त्वाभिजनो महात्मा}
{ स मैथिलीं द्रक्ष्यति वानरेन्द्रः}
{दिशं तु यामेव गता तु सीता}
{ तामास्थितो वायुसुतो हनूमान्} %||4-47-14||


{॥इत्यार्षे श्रीमद्रामायणे वाल्मीकीये आदिकाव्ये किष्किन्धाकाण्डे षट्चत्वारिंशः सर्गः॥}

\sect{सप्तचत्वारिंशः सर्गः}

\twolineshloka
{सहताराङ्गदाभ्यां तु गत्वा स हनुमान्कपिः}
{सुग्रीवेण यथोद्दिष्टं तं देशमुपचक्रमे} %||4-47-1||

\twolineshloka
{स तु दूरमुपागम्य सर्वैस्तैः कपिसत्तमैः}
{विचिनोति स्म विन्ध्यस्य गुहाश्च गहनानि च} %||4-47-2||

\twolineshloka
{पर्वताग्रान्नदीदुर्गान्सरांसि विपुलान्द्रुमान्}
{वृक्षषण्डांश्च विविधान्पर्वतान्घनपादपान्} %||4-47-3||

\twolineshloka
{अन्वेषमाणास्ते सर्वे वानराः सर्वतो दिशम्}
{न सीतां ददृशुर्वीरा मैथिलीं जनकात्मजाम्} %||4-47-4||

\threelineshloka
{ते भक्षयन्तो मूलानि फलानि विविधानि च}
{अन्वेषमाणा दुर्धर्षा न्यवसंस्तत्र तत्र ह}
{स तु देशो दुरन्वेषो गुहागहनवान्महान्} %||4-47-5||

\twolineshloka
{त्यक्त्वा तु तं तदा देशं सर्वे वै हरियूथपाः}
{देशमन्यं दुराधर्षं विविशुश्चाकुतोभयाः} %||4-47-6||

\twolineshloka
{यत्र वन्ध्यफला वृक्षा विपुष्पाः पर्णवर्जिताः}
{निस्तोयाः सरितो यत्र मूलं यत्र सुदुर्लभम्} %||4-47-7||

\twolineshloka
{न सन्ति महिषा यत्र न मृगा न च हस्तिनः}
{शार्दूलाः पक्षिणो वापि ये चान्ये वनगोचराः} %||4-47-8||

\twolineshloka
{स्निग्धपत्राः स्थले यत्र पद्मिन्यः फुल्लपङ्कजाः}
{प्रेक्षणीयाः सुगन्धाश्च भ्रमरैश्चापि वर्जिताः} %||4-47-9||

\twolineshloka
{कण्डुर्नाम महाभागः सत्यवादी तपोधनः}
{महर्षिः परमामर्षी नियमैर्दुष्प्रधर्षणः} %||4-47-10||

\twolineshloka
{तस्य तस्मिन्वने पुत्रो बालको दशवार्षिकः}
{प्रनष्टो जीवितान्ताय क्रुद्धस्तत्र महामुनिः} %||4-47-11||

\twolineshloka
{तेन धर्मात्मना शप्तं कृत्स्नं तत्र महद्वनम्}
{अशरण्यं दुराधर्षं मृगपक्षिविवर्जितम्} %||4-47-12||

\twolineshloka
{तस्य ते काननान्तांस्तु गिरीणां कन्दराणि च}
{प्रभवानि नदीनाञ्च विचिन्वन्ति समाहिताः} %||4-47-13||

\twolineshloka
{तत्र चापि महात्मानो नापश्यञ्जनकात्मजाम्}
{हर्तारं रावणं वापि सुग्रीवप्रियकारिणः} %||4-47-14||

\twolineshloka
{ते प्रविश्य तु तं भीमं लतागुल्मसमावृतम्}
{ददृशुः क्रूरकर्माणमसुरं सुरनिर्भयम्} %||4-47-15||

\twolineshloka
{तं दृष्ट्वा वनरा घोरं स्थितं शैलमिवापरम्}
{गाढं परिहिताः सर्वे दृष्ट्वा तं पर्वतोपमम्} %||4-47-16||

\twolineshloka
{सोऽपि तान्वानरान्सर्वान्नष्टाः स्थेत्यब्रवीद्बली}
{अभ्यधावत सङ्क्रुद्धो मुष्टिमुद्यम्य संहितम्} %||4-47-17||

{तमापतन्तं सहसा वालिपुत्रोऽङ्गदस्तदा॥\devanumber{18}॥} %||4-47-18|| (Check)

\addtocounter{shlokacount}{1}

{स वालिपुत्राभिहतो वक्त्राच्छोणितमुद्वमन्॥\devanumber{19}॥} %||4-47-19|| (Check)

\addtocounter{shlokacount}{1}

\threelineshloka
{असुरो न्यपतद्भूमौ पर्यस्त इव पर्वतः}
{ते तु तस्मिन्निरुच्छ्वासे वानरा जितकाशिनः}
{व्यचिन्वन्प्रायशस्तत्र सर्वं तद्गिरिगह्वरम्} %||4-47-20||

\twolineshloka
{विचितं तु ततः कृत्वा सर्वे ते काननं पुनः}
{अन्यदेवापरं घोरं विविशुर्गिरिगह्वरम्} %||4-47-21||

\twolineshloka
{ते विचिन्त्य पुनः खिन्ना विनिष्पत्य समागताः}
{एकान्ते वृक्षमूले तु निषेदुर्दीनमानसाः} %||4-48-22||


{॥इत्यार्षे श्रीमद्रामायणे वाल्मीकीये आदिकाव्ये किष्किन्धाकाण्डे सप्तचत्वारिंशः सर्गः॥}

\sect{अष्टचत्वारिंशः सर्गः}

\twolineshloka
{अथाङ्गदस्तदा सर्वान्वानरानिदमब्रवीत्}
{परिश्रान्तो महाप्राज्ञः समाश्वास्य शनैर्वचः} %||4-48-1||

\twolineshloka
{वनानि गिरयो नद्यो दुर्गाणि गहनानि च}
{दर्यो गिरिगुहाश्चैव विचिता नः समन्ततः} %||4-48-2||

\twolineshloka
{तत्र तत्र सहास्माभिर्जानकी न च दृश्यते}
{तद्वा रक्षो हृता येन सीता सुरसुतोपमा} %||4-48-3||

\twolineshloka
{कालश्च नो महान्यातः सुग्रीवश्चोग्रशासनः}
{तस्माद्भवन्तः सहिता विचिन्वन्तु समन्ततः} %||4-48-4||

\twolineshloka
{विहाय तन्द्रीं शोकं च निद्रां चैव समुत्थिताम्}
{विचिनुध्वं यथा सीतां पश्यामो जनकात्मजाम्} %||4-48-5||

\twolineshloka
{अनिर्वेदं च दाक्ष्यं च मनसश्चापराजयम्}
{कार्यसिद्धिकराण्याहुस्तस्मादेतद्ब्रवीम्यहम्} %||4-48-6||

\twolineshloka
{अद्यापीदं वनं दुर्गं विचिन्वन्तु वनौकसः}
{खेदं त्यक्त्वा पुनः सर्वं वनमेतद्विचीयताम्} %||4-48-7||

\twolineshloka
{अवश्यं क्रियमाणस्य दृश्यते कर्मणः फलम्}
{अलं निर्वेदमागम्य न हि नो मलिनं क्षमम्} %||4-48-8||

\twolineshloka
{सुग्रीवः क्रोधनो राजा तीक्ष्णदण्डश्च वानराः}
{भेतव्यं तस्य सततं रामस्य च महात्मनः} %||4-48-9||

\twolineshloka
{हितार्थमेतदुक्तं वः क्रियतां यदि रोचते}
{उच्यतां वा क्षमं यन्नः सर्वेषामेव वानराः} %||4-48-10||

\twolineshloka
{अङ्गदस्य वचः श्रुत्वा वचनं गन्धमादनः}
{उवाचाव्यक्तया वाचा पिपासा श्रमखिन्नया} %||4-48-11||

\twolineshloka
{सदृशं खलु वो वाक्यमङ्गदो यदुवाच ह}
{हितं चैवानुकूलं च क्रियतामस्य भाषितम्} %||4-48-12||

\twolineshloka
{पुनर्मार्गामहे शैलान्कन्दरांश्च दरीस्तथा}
{काननानि च शून्यानि गिरिप्रस्रवणानि च} %||4-48-13||

\twolineshloka
{यथोद्दिष्ठानि सर्वाणि सुग्रीवेण महात्मना}
{विचिन्वन्तु वनं सर्वे गिरिदुर्गाणि सर्वशः} %||4-48-14||

\twolineshloka
{ततः समुत्थाय पुनर्वानरास्ते महाबलाः}
{विन्ध्यकाननसङ्कीर्णां विचेरुर्दक्षिणां दिशम्} %||4-48-15||

\twolineshloka
{ते शारदाभ्रप्रतिमं श्रीमद्रजतपर्वतम्}
{शृङ्गवन्तं दरीवन्तमधिरुह्य च वानराः} %||4-48-16||

\twolineshloka
{तत्र लोध्रवनं रम्यं सप्तपर्णवनानि च}
{विचिन्वन्तो हरिवराः सीतादर्शनकाङ्क्षिणः} %||4-48-17||

\twolineshloka
{तस्याग्रमधिरूढास्ते श्रान्ता विपुलविक्रमाः}
{न पश्यन्ति स्म वैदेहीं रामस्य महिषीं प्रियाम्} %||4-48-18||

\twolineshloka
{ते तु दृष्टिगतं कृत्वा तं शैलं बहुकन्दरम्}
{अवारोहन्त हरयो वीक्षमाणाः समन्ततः} %||4-48-19||

\twolineshloka
{अवरुह्य ततो भूमिं श्रान्ता विगतचेतसः}
{स्थित्वा मुहूर्तं तत्राथ वृक्षमूलमुपाश्रिताः} %||4-48-20||

\twolineshloka
{ते मुहूर्तं समाश्वस्ताः किञ्चिद्भग्नपरिश्रमाः}
{पुनरेवोद्यताः कृत्स्नां मार्गितुं दक्षिणां दिशम्} %||4-48-21||

\twolineshloka
{हनुमत्प्रमुखास्ते तु प्रस्थिताः प्लवगर्षभाः}
{विन्ध्यमेवादितस्तावद्विचेरुस्ते समन्ततः} %||4-49-22||


{॥इत्यार्षे श्रीमद्रामायणे वाल्मीकीये आदिकाव्ये किष्किन्धाकाण्डे अष्टचत्वारिंशः सर्गः॥}

\sect{एकोनपञ्चाशत्तमः सर्गः}

\twolineshloka
{सह ताराङ्गदाभ्यां तु सङ्गम्य हनुमान्कपिः}
{विचिनोति स्म विन्ध्यस्य गुहाश्च गहनानि च} %||4-49-1||

\twolineshloka
{सिंहशार्दूलजुष्टाश्च गुहाश्च परितस्तथा}
{विषमेषु नगेन्द्रस्य महाप्रस्रवणेषु च} %||4-49-2||

{तेषां तत्रैव वसतां स कालो व्यत्यवर्तत॥\devanumber{3}॥} %||4-49-3|| (Check)

\addtocounter{shlokacount}{1}

\twolineshloka
{स हि देशो दुरन्वेषो गुहा गहनवान्महान्}
{तत्र वायुसुतः सर्वं विचिनोति स्म पर्वतम्} %||4-49-4||

\twolineshloka
{परस्परेण रहिता अन्योन्यस्याविदूरतः}
{गजो गवाक्षो गवयः शरभो गन्धमादनः} %||4-49-5||

\twolineshloka
{मैन्दश्च द्विविदश्चैव हनुमाञ्जाम्बवानपि}
{अङ्गदो युवराजश्च तारश्च वनगोचरः} %||4-49-6||

\threelineshloka
{गिरिजालावृतान्देशान्मार्गित्वा दक्षिणां दिशम्}
{क्षुत्पिपासा परीताश्च श्रान्ताश्च सलिलार्थिनः}
{अवकीर्णं लतावृक्षैर्ददृशुस्ते महाबिलम्} %||4-49-7||

\twolineshloka
{ततः क्रौञ्चाश्च हंसाश्च सारसाश्चापि निष्क्रमन्}
{जलार्द्राश्चक्रवाकाश्च रक्ताङ्गाः पद्मरेणुभिः} %||4-49-8||

\twolineshloka
{ततस्तद्बिलमासाद्य सुगन्धि दुरतिक्रमम्}
{विस्मयव्यग्रमनसो बभूवुर्वानरर्षभाः} %||4-49-9||

\twolineshloka
{संजातपरिशङ्कास्ते तद्बिलं प्लवगोत्तमाः}
{अभ्यपद्यन्त संहृष्टास्तेजोवन्तो महाबलाः} %||4-49-10||

\twolineshloka
{ततः पर्वतकूटाभो हनुमान्मारुतात्मजः}
{अब्रवीद्वानरान्सर्वान्कान्तार वनकोविदः} %||4-49-11||

\twolineshloka
{गिरिजालावृतान्देशान्मार्गित्वा दक्षिणां दिशम्}
{वयं सर्वे परिश्रान्ता न च पश्यामि मैथिलीम्} %||4-49-12||

\twolineshloka
{अस्माच्चापि बिलाद्धंसाः क्रौञ्चाश्च सह सारसैः}
{जलार्द्राश्चक्रवाकाश्च निष्पतन्ति स्म सर्वशः} %||4-49-13||

\twolineshloka
{नूनं सलिलवानत्र कूपो वा यदि वा ह्रदः}
{तथा चेमे बिलद्वारे स्निग्धास्तिष्ठन्ति पादपाः} %||4-49-14||

\twolineshloka
{इत्युक्तास्तद्बिलं सर्वे विविशुस्तिमिरावृतम्}
{अचन्द्रसूर्यं हरयो ददृशू रोमहर्षणम्} %||4-49-15||

\twolineshloka
{ततस्तस्मिन्बिले दुर्गे नानापादपसङ्कुले}
{अन्योन्यं सम्परिष्वज्य जग्मुर्योजनमन्तरम्} %||4-49-16||

\twolineshloka
{ते नष्टसंज्ञास्तृषिताः सम्भ्रान्ताः सलिलार्थिनः}
{परिपेतुर्बिले तस्मिन्कञ्चित्कालमतन्द्रिताः} %||4-49-17||

\twolineshloka
{ते कृशा दीनवदनाः परिश्रान्ताः प्लवङ्गमाः}
{आलोकं ददृशुर्वीरा निराशा जीविते तदा} %||4-49-18||

\twolineshloka
{ततस्तं देशमागम्य सौम्यं वितिमिरं वनम्}
{ददृशुः काञ्चनान्वृक्षान्दीप्तवैश्वानरप्रभान्} %||4-49-19||

\twolineshloka
{सालांस्तालांश्च पुंनागान्ककुभान्वञ्जुलान्धवान्}
{चम्पकान्नागवृक्षांश्च कर्णिकारांश्च पुष्पितान्} %||4-49-20||

\twolineshloka
{तरुणादित्यसङ्काशान्वैदूर्यमयवेदिकान्}
{नीलवैदूर्यवर्णाश्च पद्मिनीः पतगावृताः} %||4-49-21||

\twolineshloka
{महद्भिः काञ्चनैर्वृक्षैर्वृतं बालार्क संनिभैः}
{जातरूपमयैर्मत्स्यैर्महद्भिश्च सकच्छपैः} %||4-49-22||

\twolineshloka
{नलिनीस्तत्र ददृशुः प्रसन्नसलिलायुताः}
{काञ्चनानि विमानानि राजतानि तथैव च} %||4-49-23||

\twolineshloka
{तपनीयगवाक्षाणि मुक्ताजालावृतानि च}
{हैमराजतभौमानि वैदूर्यमणिमन्ति च} %||4-49-24||

\twolineshloka
{ददृशुस्तत्र हरयो गृहमुख्यानि सर्वशः}
{पुष्पितान्फलिनो वृक्षान्प्रवालमणिसंनिभान्} %||4-49-25||

\twolineshloka
{काञ्चनभ्रमरांश्चैव मधूनि च समन्ततः}
{मणिकाञ्चनचित्राणि शयनान्यासनानि च} %||4-49-26||

\twolineshloka
{महार्हाणि च यानानि ददृशुस्ते समन्ततः}
{हैमराजतकांस्यानां भाजनानां च सञ्चयान्} %||4-49-27||

\twolineshloka
{अगरूणां च दिव्यानां चन्दनानां च सञ्चयान्}
{शुचीन्यभ्यवहार्याणि मूलानि च फलानि च} %||4-49-28||

\threelineshloka
{महार्हाणि च पानानि मधूनि रसवन्ति च}
{दिव्यानामम्बराणां च महार्हाणां च सञ्चयान्}
{कम्बलानां च चित्राणामजिनानां च सञ्चयान्} %||4-49-29||

\twolineshloka
{तत्र तत्र विचिन्वन्तो बिले तत्र महाप्रभाः}
{ददृशुर्वानराः शूराः स्त्रियं काञ्चिददूरतः} %||4-49-30||

\twolineshloka
{तां दृष्ट्वा भृशसन्त्रस्ताश्चीरकृष्णाजिनाम्बराम्}
{तापसीं नियताहारां ज्वलन्तीमिव तेजसा} %||4-49-31||

\fourlineindentedshloka
{ततो हनूमान्गिरिसंनिकाशः}
{ कृताञ्जलिस्तामभिवाद्य वृद्धाम्}
{पप्रच्छ का त्वं भवनं बिलं च}
{ रत्नानि चेमानि वदस्व कस्य} %||4-50-32||


{॥इत्यार्षे श्रीमद्रामायणे वाल्मीकीये आदिकाव्ये किष्किन्धाकाण्डे एकोनपञ्चाशत्तमः सर्गः॥}

\sect{पञ्चाशत्तमः सर्गः}

\twolineshloka
{इत्युक्त्वा हनुमांस्तत्र पुनः कृष्णाजिनाम्बराम्}
{अब्रवीत्तां महाभागां तापसीं धर्मचारिणीम्} %||4-50-1||

\twolineshloka
{इदं प्रविष्टाः सहसा बिलं तिमिरसंवृतम्}
{क्षुत्पिपासा परिश्रान्ताः परिखिन्नाश्च सर्वशः} %||4-50-2||

\threelineshloka
{महद्धिरण्या विवरं प्रविष्टाः स्म पिपासिताः}
{इमांस्त्वेवं विधान्भावान्विविधानद्भुतोपमान्}
{दृष्ट्वा वयं प्रव्यथिताः सम्भ्रान्ता नष्टचेतसः} %||4-50-3||

\twolineshloka
{कस्येमे काञ्चना वृक्षास्तरुणादित्यसंनिभाः}
{शुचीन्यभ्यवहार्याणि मूलानि च फलानि च} %||4-50-4||

\twolineshloka
{काञ्चनानि विमानानि राजतानि गृहाणि च}
{तपनीय गवाक्षाणि मणिजालावृतानि च} %||4-50-5||

\twolineshloka
{पुष्पिताः फालवन्तश्च पुण्याः सुरभिगन्धिनः}
{इमे जाम्बूनदमयाः पादपाः कस्य तेजसा} %||4-50-6||

\twolineshloka
{काञ्चनानि च पद्मानि जातानि विमले जले}
{कथं मत्स्याश्च सौवर्णा चरन्ति सह कच्छपैः} %||4-50-7||

\twolineshloka
{आत्मानमनुभावं च कस्य चैतत्तपोबलम्}
{अजानतां नः सर्वेषां सर्वमाख्यातुमर्हसि} %||4-50-8||

\twolineshloka
{एवमुक्ता हनुमता तापसी धर्मचारिणी}
{प्रत्युवाच हनूमन्तं सर्वभूतहिते रता} %||4-50-9||

\twolineshloka
{मयो नाम महातेजा मायावी दानवर्षभः}
{तेनेदं निर्मितं सर्वं मायया काञ्चनं वनम्} %||4-50-10||

\twolineshloka
{पुरा दानवमुख्यानां विश्वकर्मा बभूव ह}
{येनेदं काञ्चनं दिव्यं निर्मितं भवनोत्तमम्} %||4-50-11||

\twolineshloka
{स तु वर्षसहस्राणि तपस्तप्त्वा महावने}
{पितामहाद्वरं लेभे सर्वमौशसनं धनम्} %||4-50-12||

\twolineshloka
{विधाय सर्वं बलवान्सर्वकामेश्वरस्तदा}
{उवास सुखितः कालं कञ्चिदस्मिन्महावने} %||4-50-13||

\twolineshloka
{तमप्सरसि हेमायां सक्तं दानवपुङ्गवम्}
{विक्रम्यैवाशनिं गृह्य जघानेशः पुरन्दरः} %||4-50-14||

\twolineshloka
{इदं च ब्रह्मणा दत्तं हेमायै वनमुत्तमम्}
{शाश्वतः कामभोगश्च गृहं चेदं हिरण्मयम्} %||4-50-15||

\twolineshloka
{दुहिता मेरुसावर्णेरहं तस्याः स्वयं प्रभा}
{इदं रक्षामि भवनं हेमाया वानरोत्तम} %||4-50-16||

\twolineshloka
{मम प्रियसखी हेमा नृत्तगीतविशारदा}
{तया दत्तवरा चास्मि रक्षामि भवनोत्तमम्} %||4-50-17||

\twolineshloka
{किं कार्यं कस्य वा हेतोः कान्ताराणि प्रपद्यथ}
{कथं चेदं वनं दुर्गं युष्माभिरुपलक्षितम्} %||4-50-18||

\twolineshloka
{इमान्यभ्यवहार्याणि मूलानि च फलानि च}
{भुक्त्वा पीत्वा च पानीयं सर्वं मे वक्तुमर्हथ} %||4-51-19||


{॥इत्यार्षे श्रीमद्रामायणे वाल्मीकीये आदिकाव्ये किष्किन्धाकाण्डे पञ्चाशत्तमः सर्गः॥}

\sect{एकपञ्चाशत्तमः सर्गः}

\twolineshloka
{अथ तानब्रवीत्सर्वान्विश्रान्तान्हरियूथपान्}
{इदं वचनमेकाग्रा तापसी धर्मचारिणी} %||4-51-1||

\twolineshloka
{वानरा यदि वः खेदः प्रनष्टः फलभक्षणात्}
{यदि चैतन्मया श्राव्यं श्रोतुमिच्छामि कथ्यताम्} %||4-51-2||

\twolineshloka
{तस्यास्तद्वचनं श्रुत्वा हनुमान्मारुतात्मजः}
{आर्जवेन यथातत्त्वमाख्यातुमुपचक्रमे} %||4-51-3||

\twolineshloka
{राजा सर्वस्य लोकस्य महेन्द्रवरुणोपमः}
{रामो दाशरथिः श्रीमान्प्रविष्टो दण्डकावनम्} %||4-51-4||

\twolineshloka
{लक्ष्मणेन सह भ्रात्रा वैदेह्या चापि भार्यया}
{तस्य भार्या जनस्थानाद्रावणेन हृता बलात्} %||4-51-5||

\twolineshloka
{वीरस्तस्य सखा राज्ञः सुग्रीवो नाम वानरः}
{राजा वानरमुख्यानां येन प्रस्थापिता वयम्} %||4-51-6||

\twolineshloka
{अगस्त्यचरितामाशां दक्षिणां यमरक्षिताम्}
{सहैभिर्वानरैर्मुख्यैरङ्गदप्रमुखैर्वयम्} %||4-51-7||

\twolineshloka
{रावणं सहिताः सर्वे राक्षसं कामरूपिणम्}
{सीतया सह वैदेह्या मार्गध्वमिति चोदिताः} %||4-51-8||

\twolineshloka
{विचित्य तु वयं सर्वे समग्रां दक्षिणां दिशम्}
{बुभुक्षिताः परिश्रान्ता वृक्षमूलमुपाश्रिताः} %||4-51-9||

\twolineshloka
{विवर्णवदनाः सर्वे सर्वे ध्यानपरायणाः}
{नाधिगच्छामहे पारं मग्नाश्चिन्तामहार्णवे} %||4-51-10||

\twolineshloka
{चारयन्तस्ततश्चक्षुर्दृष्टवन्तो महद्बिलम्}
{लतापादपसञ्छन्नं तिमिरेण समावृतम्} %||4-51-11||

\threelineshloka
{अस्माद्धंसा जलक्लिन्नाः पक्षैः सलिलरेणुभिः}
{कुरराः सारसाश्चैव निष्पतन्ति पतत्रिणः}
{साध्वत्र प्रविशामेति मया तूक्ताः प्लवङ्गमाः} %||4-51-12||

\twolineshloka
{तेषामपि हि सर्वेषामनुमानमुपागतम्}
{गच्छामः प्रविशामेति भर्तृकार्यत्वरान्विताः} %||4-51-13||

\twolineshloka
{ततो गाढं निपतिता गृह्य हस्तौ परस्परम्}
{इदं प्रविष्टाः सहसा बिलं तिमिरसंवृतम्} %||4-51-14||

\twolineshloka
{एतन्नः कायमेतेन कृत्येन वयमागताः}
{त्वां चैवोपगताः सर्वे परिद्यूना बुभुक्षिताः} %||4-51-15||

\twolineshloka
{आतिथ्यधर्मदत्तानि मूलानि च फलानि च}
{अस्माभिरुपभुक्तानि बुभुक्षापरिपीडितैः} %||4-51-16||

\twolineshloka
{यत्त्वया रक्षिताः सर्वे म्रियमाणा बुभुक्षया}
{ब्रूहि प्रत्युपकारार्थं किं ते कुर्वन्तु वानराः} %||4-51-17||

\twolineshloka
{एवमुक्ता तु सर्वज्ञा वानरैस्तैः स्वयम्प्रभा}
{प्रत्युवाच ततः सर्वानिदं वानरयूथपम्} %||4-51-18||

\twolineshloka
{सर्वेषां परितुष्टास्मि वानराणां तरस्विनाम्}
{चरन्त्या मम धर्मेण न कार्यमिह केनचित्} %||4-52-19||


{॥इत्यार्षे श्रीमद्रामायणे वाल्मीकीये आदिकाव्ये किष्किन्धाकाण्डे एकपञ्चाशत्तमः सर्गः॥}

\sect{द्विपञ्चाशत्तमः सर्गः}

\twolineshloka
{एवमुक्तः शुभं वाक्यं तापस्या धर्मसंहितम्}
{उवाच हनुमान्वाक्यं तामनिन्दितचेष्टिताम्} %||4-52-1||

\threelineshloka
{शरणं त्वां प्रपन्नाः स्मः सर्वे वै धर्मचारिणि}
{यः कृतः समयोऽस्माकं सुग्रीवेण महात्मना}
{स तु कालो व्यतिक्रान्तो बिले च परिवर्तताम्} %||4-52-2||

{सा त्वमस्माद्बिलाद्घोरादुत्तारयितुमर्हसि॥\devanumber{3}॥} %||4-52-3|| (Check)

\addtocounter{shlokacount}{1}

\twolineshloka
{तस्मात्सुग्रीववचनादतिक्रान्तान्गतायुषः}
{त्रातुमर्हसि नः सर्वान्सुग्रीवभयशङ्कितान्} %||4-52-4||

\twolineshloka
{महच्च कार्यमस्माभिः कर्तव्यं धर्मचारिणि}
{तच्चापि न कृतं कार्यमस्माभिरिह वासिभिः} %||4-52-5||

\twolineshloka
{एवमुक्ता हनुमता तापसी वाक्यमब्रवीत्}
{जीवता दुष्करं मन्ये प्रविष्टेन निवर्तितुम्} %||4-52-6||

\twolineshloka
{तपसस्तु प्रभावेन नियमोपार्जितेन च}
{सर्वानेव बिलादस्मादुद्धरिष्यामि वानरान्} %||4-52-7||

\twolineshloka
{निमीलयत चक्षूंषि सर्वे वानरपुङ्गवाः}
{न हि निष्क्रमितुं शक्यमनिमीलितलोचनैः} %||4-52-8||

\twolineshloka
{ततः सम्मीलिताः सर्वे सुकुमाराङ्गुलैः करैः}
{सहसा पिदधुर्दृष्टिं हृष्टा गमनकाङ्क्षिणः} %||4-52-9||

\twolineshloka
{वानरास्तु महात्मानो हस्तरुद्धमुखास्तदा}
{निमेषान्तरमात्रेण बिलादुत्तारितास्तया} %||4-52-10||

\twolineshloka
{ततस्तान्वानरान्सर्वांस्तापसी धर्मचारिणी}
{निःसृतान्विषमात्तस्मात्समाश्वास्येदमब्रवीत्} %||4-52-11||

\twolineshloka
{एष विन्ध्यो गिरिः श्रीमान्नानाद्रुमलतायुतः}
{एष प्रसवणः शैलः सागरोऽयं महोदधिः} %||4-52-12||

\twolineshloka
{स्वस्ति वोऽस्तु गमिष्यामि भवनं वानरर्षभाः}
{इत्युक्त्वा तद्बिलं श्रीमत्प्रविवेश स्वयम्प्रभा} %||4-52-13||

\twolineshloka
{ततस्ते ददृशुर्घोरं सागरं वरुणालयम्}
{अपारमभिगर्जन्तं घोरैरूर्मिभिराकुलम्} %||4-52-14||

\twolineshloka
{मयस्य माया विहितं गिरिदुर्गं विचिन्वताम्}
{तेषां मासो व्यतिक्रान्तो यो राज्ञा समयः कृतः} %||4-52-15||

\twolineshloka
{विन्ध्यस्य तु गिरेः पादे सम्प्रपुष्पितपादपे}
{उपविश्य महाभागाश्चिन्तामापेदिरे तदा} %||4-52-16||

\twolineshloka
{ततः पुष्पातिभाराग्राँल्लताशतसमावृतान्}
{द्रुमान्वासन्तिकान्दृष्ट्वा बभूवुर्भयशङ्किताः} %||4-52-17||

\twolineshloka
{ते वसन्तमनुप्राप्तं प्रतिवेद्य परस्परम्}
{नष्टसन्देशकालार्था निपेतुर्धरणीतले} %||4-52-18||

\twolineshloka
{स तु सिंहर्षभ स्कन्धः पीनायतभुजः कपिः}
{युवराजो महाप्राज्ञ अङ्गदो वाक्यमब्रवीत्} %||4-52-19||

\twolineshloka
{शासनात्कपिराजस्य वयं सर्वे विनिर्गताः}
{मासः पूर्णो बिलस्थानां हरयः किं न बुध्यते} %||4-52-20||

\twolineshloka
{तस्मिन्नतीते काले तु सुग्रीवेण कृते स्वयम्}
{प्रायोपवेशनं युक्तं सर्वेषां च वनौकसाम्} %||4-52-21||

\twolineshloka
{तीक्ष्णः प्रकृत्या सुग्रीवः स्वामिभावे व्यवस्थितः}
{न क्षमिष्यति नः सर्वानपराधकृतो गतान्} %||4-52-22||

\twolineshloka
{अप्रवृत्तौ च सीतायाः पापमेव करिष्यति}
{तस्मात्क्षममिहाद्यैव प्रायोपविशनं हि नः} %||4-52-23||

\threelineshloka
{त्यक्त्वा पुत्रांश्च दारांश्च धनानि च गृहाणि च}
{यावन्न घातयेद्राजा सर्वान्प्रतिगतानितः}
{वधेनाप्रतिरूपेण श्रेयान्मृत्युरिहैव नः} %||4-52-24||

\twolineshloka
{न चाहं यौवराज्येन सुग्रीवेणाभिषेचितः}
{नरेन्द्रेणाभिषिक्तोऽस्मि रामेणाक्लिष्टकर्मणा} %||4-52-25||

\twolineshloka
{स पूर्वं बद्धवैरो मां राजा दृष्ट्वा व्यतिक्रमम्}
{घातयिष्यति दण्डेन तीक्ष्णेन कृतनिश्चयः} %||4-52-26||

\twolineshloka
{किं मे सुहृद्भिर्व्यसनं पश्यद्भिर्जीवितान्तरे}
{इहैव प्रायमासिष्ये पुण्ये सागररोधसि} %||4-52-27||

\twolineshloka
{एतच्छ्रुत्वा कुमारेण युवराजेन भाषितम्}
{सर्वे ते वानरश्रेष्ठाः करुणं वाक्यमब्रुवन्} %||4-52-28||

\twolineshloka
{तीक्ष्णः प्रकृत्या सुग्रीवः प्रियासक्तश्च राघवः}
{अदृष्टायां च वैदेह्यां दृष्ट्वास्मांश्च समागतान्} %||4-52-29||

\twolineshloka
{राघवप्रियकामार्थं घातयिष्यत्यसंशयम्}
{न क्षमं चापराद्धानां गमनं स्वामिपार्श्वतः} %||4-52-30||

\fourlineindentedshloka
{प्लवङ्गमानां तु भयार्दितानाम्}
{ श्रुत्वा वचस्तार इदं बभाषे}
{अलं विषादेन बिलं प्रविश्य}
{ वसाम सर्वे यदि रोचते वः} %||4-52-31||

\fourlineindentedshloka
{इदं हि माया विहितं सुदुर्गमम्}
{ प्रभूतवृक्षोदकभोज्यपेयम्}
{इहास्ति नो नैव भयं पुरन्दरा}
{न्न राघवाद्वानरराजतोऽपि वा} %||4-52-32||

\fourlineindentedshloka
{श्रुत्वाङ्गदस्यापि वचोऽनुकूलम्}
{ ऊचुश्च सर्वे हरयः प्रतीताः}
{यथा न हन्येम तथाविधानम्}
{ असक्तमद्यैव विधीयतां नः} %||4-53-33||


{॥इत्यार्षे श्रीमद्रामायणे वाल्मीकीये आदिकाव्ये किष्किन्धाकाण्डे द्विपञ्चाशत्तमः सर्गः॥}

\sect{त्रिपञ्चाशत्तमः सर्गः}

\twolineshloka
{तथा ब्रुवति तारे तु ताराधिपतिवर्चसि}
{अथ मेने हृतं राज्यं हनुमानङ्गदेन तत्} %||4-53-1||

\twolineshloka
{बुद्ध्या ह्यष्टाङ्गया युक्तं चतुर्बलसमन्वितम्}
{चतुर्दशगुणं मेने हनुमान्वालिनः सुतम्} %||4-53-2||

\twolineshloka
{आपूर्यमाणं शश्वच्च तेजोबलपराक्रमैः}
{शशिनं शुक्लपक्षादौ वर्धमानमिव श्रिया} %||4-53-3||

\twolineshloka
{बृहस्पतिसमं बुद्ध्या विक्रमे सदृशं पितुः}
{शुश्रूषमाणं तारस्य शुक्रस्येव पुरन्दरम्} %||4-53-4||

\twolineshloka
{भर्तुरर्थे परिश्रान्तं सर्वशास्त्रविशारदम्}
{अभिसन्धातुमारेभे हनुमानङ्गदं ततः} %||4-53-5||

\twolineshloka
{स चतुर्णामुपायानां तृतीयमुपवर्णयन्}
{भेदयामास तान्सर्वान्वानरान्वाक्यसम्पदा} %||4-53-6||

\twolineshloka
{तेषु सर्वेषु भिन्नेषु ततोऽभीषयदङ्गदम्}
{भीषणैर्बहुभिर्वाक्यैः कोपोपायसमन्वितैः} %||4-53-7||

\twolineshloka
{त्वं समर्थतरः पित्रा युद्धे तारेय वै धुरम्}
{दृढं धारयितुं शक्तः कपिराज्यं यथा पिता} %||4-53-8||

\twolineshloka
{नित्यमस्थिरचित्ता हि कपयो हरिपुङ्गव}
{नाज्ञाप्यं विषहिष्यन्ति पुत्रदारान्विना त्वया} %||4-53-9||

\twolineshloka
{त्वां नैते ह्यनुयुञ्जेयुः प्रत्यक्षं प्रवदामि ते}
{यथायं जाम्बवान्नीलः सुहोत्रश्च महाकपिः} %||4-53-10||

\twolineshloka
{न ह्यहं त इमे सर्वे सामदानादिभिर्गुणैः}
{दण्डेन न त्वया शक्याः सुग्रीवादपकर्षितुम्} %||4-53-11||

\twolineshloka
{विगृह्यासनमप्याहुर्दुर्बलेन बलीयसः}
{आत्मरक्षाकरस्तस्मान्न विगृह्णीत दुर्बलः} %||4-53-12||

\twolineshloka
{यां चेमां मन्यसे धात्रीमेतद्बिलमिति श्रुतम्}
{एतल्लक्ष्मणबाणानामीषत्कार्यं विदारणे} %||4-53-13||

\threelineshloka
{स्वल्पं हि कृतमिन्द्रेण क्षिपता ह्यशनिं पुरा}
{लक्ष्मणो निशितैर्बाणैर्भिन्द्यात्पत्रपुटं यथा}
{लक्ष्मणस्य च नाराचा बहवः सन्ति तद्विधाः} %||4-53-14||

\twolineshloka
{अवस्थाने यदैव त्वमासिष्यसि परन्तप}
{तदैव हरयः सर्वे त्यक्ष्यन्ति कृतनिश्चयाः} %||4-53-15||

\twolineshloka
{स्मरन्तः पुत्रदाराणां नित्योद्विग्ना बुभुक्षिताः}
{खेदिता दुःखशय्याभिस्त्वां करिष्यन्ति पृष्ठतः} %||4-53-16||

\twolineshloka
{स त्वं हीनः सुहृद्भिश्च हितकामैश्च बन्धुभिः}
{तृणादपि भृशोद्विग्नः स्पन्दमानाद्भविष्यसि} %||4-53-17||

\twolineshloka
{न च जातु न हिंस्युस्त्वां घोरा लक्ष्मणसायकाः}
{अपवृत्तं जिघांसन्तो महावेगा दुरासदाः} %||4-53-18||

\twolineshloka
{अस्माभिस्तु गतं सार्धं विनीतवदुपस्थितम्}
{आनुपूर्व्यात्तु सुग्रीवो राज्ये त्वां स्थापयिष्यति} %||4-53-19||

\twolineshloka
{धर्मकामः पितृव्यस्ते प्रीतिकामो दृढव्रतः}
{शुचिः सत्यप्रतिज्ञश्च ना त्वां जातु जिघांसति} %||4-53-20||

\twolineshloka
{प्रियकामश्च ते मातुस्तदर्थं चास्य जीवितम्}
{तस्यापत्यं च नास्त्यन्यत्तस्मादङ्गद गम्यताम्} %||4-54-21||


{॥इत्यार्षे श्रीमद्रामायणे वाल्मीकीये आदिकाव्ये किष्किन्धाकाण्डे त्रिपञ्चाशत्तमः सर्गः॥}

\sect{चतुःपञ्चाशत्तमः सर्गः}

\twolineshloka
{श्रुत्वा हनुमतो वाक्यं प्रश्रितं धर्मसंहितम्}
{स्वामिसत्कारसंयुक्तमङ्गदो वाक्यमब्रवीत्} %||4-54-1||

\twolineshloka
{स्थैर्यं सर्वात्मना शौचमानृशंस्यमथार्जवम्}
{विक्रमैश्चैव धैर्यं च सुग्रीवे नोपपद्यते} %||4-54-2||

\twolineshloka
{भ्रातुर्ज्येष्ठस्य यो भार्यां जीवितो महिषीं प्रियाम्}
{धर्मेण मातरं यस्तु स्वीकरोति जुगुप्सितः} %||4-54-3||

\twolineshloka
{कथं स धर्मं जानीते येन भ्रात्रा दुरात्मना}
{युद्धायाभिनियुक्तेन बिलस्य पिहितं मुखम्} %||4-54-4||

\twolineshloka
{सत्यात्पाणिगृहीतश्च कृतकर्मा महायशाः}
{विस्मृतो राघवो येन स कस्य सुकृतं स्मरेत्} %||4-54-5||

\twolineshloka
{लक्ष्मणस्य भयाद्येन नाधर्मभयभीरुणा}
{आदिष्टा मार्गितुं सीतां धर्ममस्मिन्कथं भवेत्} %||4-54-6||

\twolineshloka
{तस्मिन्पापे कृतघ्ने तु स्मृतिहीने चलात्मनि}
{आर्यः को विश्वसेज्जातु तत्कुलीनो जिजीविषुः} %||4-54-7||

\twolineshloka
{राज्ये पुत्रं प्रतिष्ठाप्य सगुणो निर्गुणोऽपि वा}
{कथं शत्रुकुलीनं मां सुग्रीवो जीवयिष्यति} %||4-54-8||

\twolineshloka
{भिन्नमन्त्रोऽपराद्धश्च हीनशक्तिः कथं ह्यहम्}
{किष्किन्धां प्राप्य जीवेयमनाथ इव दुर्बलः} %||4-54-9||

\twolineshloka
{उपांशुदण्डेन हि मां बन्धनेनोपपादयेत्}
{शठः क्रूरो नृशंसश्च सुग्रीवो राज्यकारणात्} %||4-54-10||

\twolineshloka
{बन्धनाच्चावसादान्मे श्रेयः प्रायोपवेशनम्}
{अनुजानीत मां सर्वे गृहान्गच्छन्तु वानराः} %||4-54-11||

\twolineshloka
{अहं वः प्रतिजानामि न गमिष्याम्यहं पुरीम्}
{इहैव प्रायमासिष्ये श्रेयो मरणमेव मे} %||4-54-12||

\twolineshloka
{अभिवादनपूर्वं तु राजा कुशलमेव च}
{वाच्यस्ततो यवीयान्मे सुग्रीवो वानरेश्वरः} %||4-54-13||

\twolineshloka
{आरोग्यपूर्वं कुशलं वाच्या माता रुमा च मे}
{मातरं चैव मे तारामाश्वासयितुमर्हथ} %||4-54-14||

\twolineshloka
{प्रकृत्या प्रियपुत्रा सा सानुक्रोशा तपस्विनी}
{विनष्टं मामिह श्रुत्वा व्यक्तं हास्यति जीवितम्} %||4-54-15||

\twolineshloka
{एतावदुक्त्वा वचनं वृद्धानप्यभिवाद्य च}
{संविवेशाङ्गदो भूमौ रुदन्दर्भेषु दुर्मनाः} %||4-54-16||

\twolineshloka
{तस्य संविशतस्तत्र रुदन्तो वानरर्षभाः}
{नयनेभ्यः प्रमुमुचुरुष्णं वै वारिदुःखिताः} %||4-54-17||

\twolineshloka
{सुग्रीवं चैव निन्दन्तः प्रशंसन्तश्च वालिनम्}
{परिवार्याङ्गदो सर्वे व्यवस्यन्प्रायमासितुम्} %||4-54-18||

\threelineshloka
{मतं तद्वालिपुत्रस्य विज्ञाय प्लवगर्षभाः}
{उपस्पृश्योदकं सर्वे प्राङ्मुखाः समुपाविशन्}
{दक्षिणाग्रेषु दर्भेषु उदक्तीरं समाश्रिताः} %||4-54-19||

{स संविशद्भिर्बहुभिर्महीधरो; महाद्रिकूटप्रमितैः प्लवङ्गमैः॥\devanumber{20}॥} %||4-54-20|| (Check)

\addtocounter{shlokacount}{1}


{॥इत्यार्षे श्रीमद्रामायणे वाल्मीकीये आदिकाव्ये किष्किन्धाकाण्डे चतुःपञ्चाशत्तमः सर्गः॥}

\sect{पञ्चपञ्चाशत्तमः सर्गः}

\twolineshloka
{उपविष्टास्तु ते सर्वे यस्मिन्प्रायं गिरिस्थले}
{हरयो गृध्रराजश्च तं देशमुपचक्रमे} %||4-55-1||

\twolineshloka
{साम्पातिर्नाम नाम्ना तु चिरजीवी विहङ्गमः}
{भ्राता जटायुषः श्रीमान्प्रख्यातबलपौरुषः} %||4-55-2||

\twolineshloka
{कन्दरादभिनिष्क्रम्य स विन्ध्यस्य महागिरेः}
{उपविष्टान्हरीन्दृष्ट्वा हृष्टात्मा गिरमब्रवीत्} %||4-55-3||

\twolineshloka
{विधिः किल नरं लोके विधानेनानुवर्तते}
{यथायं विहितो भक्ष्यश्चिरान्मह्यमुपागतः} %||4-55-4||

\twolineshloka
{परम्पराणां भक्षिष्ये वानराणां मृतं मृतम्}
{उवाचैवं वचः पक्षी तान्निरीक्ष्य प्लवङ्गमान्} %||4-55-5||

\twolineshloka
{तस्य तद्वचनं श्रुत्वा भक्षलुब्धस्य पक्षिणः}
{अङ्गदः परमायस्तो हनूमन्तमथाब्रवीत्} %||4-55-6||

\twolineshloka
{पश्य सीतापदेशेन साक्षाद्वैवस्वतो यमः}
{इमं देशमनुप्राप्तो वानराणां विपत्तये} %||4-55-7||

\twolineshloka
{रामस्य न कृतं कार्यं राज्ञो न च वचः कृतम्}
{हरीणामियमज्ञाता विपत्तिः सहसागता} %||4-55-8||

\twolineshloka
{वैदेह्याः प्रियकामेन कृतं कर्म जटायुषा}
{गृध्रराजेन यत्तत्र श्रुतं वस्तदशेषतः} %||4-55-9||

\twolineshloka
{तथा सर्वाणि भूतानि तिर्यग्योनिगतान्यपि}
{प्रियं कुर्वन्ति रामस्य त्यक्त्वा प्राणान्यथा वयम्} %||4-55-10||

\twolineshloka
{राघवार्थे परिश्रान्ता वयं सन्त्यक्तजीविताः}
{कान्ताराणि प्रपन्नाः स्म न च पश्याम मैथिलीम्} %||4-55-11||

\twolineshloka
{स सुखी गृध्रराजस्तु रावणेन हतो रणे}
{मुक्तश्च सुग्रीवभयाद्गतश्च परमां गतिम्} %||4-55-12||

\twolineshloka
{जटायुषो विनाशेन राज्ञो दशरथस्य च}
{हरणेन च वैदेह्याः संशयं हरयो गताः} %||4-55-13||

\twolineshloka
{रामलक्ष्मणयोर्वासामरण्ये सह सीतया}
{राघवस्य च बाणेन वालिनश्च तथा वधः} %||4-55-14||

\twolineshloka
{रामकोपादशेषाणां राक्षसानां तथा वधः}
{कैकेय्या वरदानेन इदं हि विकृतं कृतम्} %||4-55-15||

\twolineshloka
{तत्तु श्रुत्वा तदा वाक्यमङ्गदस्य मुखोद्गतम्}
{अब्रवीद्वचनं गृध्रस्तीक्ष्णतुण्डो महास्वनः} %||4-55-16||

\twolineshloka
{कोऽयं गिरा घोषयति प्राणैः प्रियतरस्य मे}
{जटायुषो वधं भ्रातुः कम्पयन्निव मे मनः} %||4-55-17||

\twolineshloka
{कथमासीज्जनस्थाने युद्धं राक्षसगृध्रयोः}
{नामधेयमिदं भ्रातुश्चिरस्याद्य मया श्रुतम्} %||4-55-18||

\twolineshloka
{यवीयसो गुणज्ञस्य श्लाघनीयस्य विक्रमैः}
{तदिच्छेयमहं श्रोतुं विनाशं वानरर्षभाः} %||4-55-19||

\threelineshloka
{भ्रातुर्जटायुषस्तस्य जनस्थाननिवासिनः}
{तस्यैव च मम भ्रातुः सखा दशरथः कथम्}
{यस्य रामः प्रियः पुत्रो ज्येष्ठो गुरुजनप्रियः} %||4-55-20||

\twolineshloka
{सूर्यांशुदग्धपक्षत्वान्न शक्नोमि विसर्पितुम्}
{इच्छेयं पर्वतादस्मादवतर्तुमरिन्दमाः} %||4-56-21||


{॥इत्यार्षे श्रीमद्रामायणे वाल्मीकीये आदिकाव्ये किष्किन्धाकाण्डे पञ्चपञ्चाशत्तमः सर्गः॥}

\sect{षट्पञ्चाशत्तमः सर्गः}

\twolineshloka
{शोकाद्भ्रष्टस्वरमपि श्रुत्वा ते हरियूथपाः}
{श्रद्दधुर्नैव तद्वाक्यं कर्मणा तस्य शङ्किताः} %||4-56-1||

\twolineshloka
{ते प्रायमुपविष्टास्तु दृष्ट्वा गृध्रं प्लवङ्गमाः}
{चक्रुर्बुद्धिं तदा रौद्रां सर्वान्नो भक्षयिष्यति} %||4-56-2||

\twolineshloka
{सर्वथा प्रायमासीनान्यदि नो भक्षयिष्यति}
{कृतकृत्या भविष्यामः क्षिप्रं सिद्धिमितो गताः} %||4-56-3||

\twolineshloka
{एतां बुद्धिं ततश्चक्रुः सर्वे ते वानरर्षभाः}
{अवतार्य गिरेः शृङ्गाद्गृध्रमाहाङ्गदस्तदा} %||4-56-4||

\twolineshloka
{बभूवुर्क्षरजो नाम वानरेन्द्रः प्रतापवान्}
{ममार्यः पार्थिवः पक्षिन्धार्मिकौ तस्य चात्मजौ} %||4-56-5||

\twolineshloka
{सुग्रीवश्चैव वली च पुत्रावोघबलावुभौ}
{लोके विश्रुतकर्माभूद्राजा वाली पिता मम} %||4-56-6||

\twolineshloka
{राजा कृत्स्नस्य जगत इक्ष्वाकूणां महारथः}
{रामो दाशरथिः श्रीमान्प्रविष्टो दण्डकावनम्} %||4-56-7||

\threelineshloka
{लक्ष्मणेन सह भ्रात्रा वैदेह्या चापि भार्यया}
{पितुर्निदेशनिरतो धर्म्यं पन्थानमाश्रितः}
{तस्य भार्या जनस्थानाद्रावणेन हृता बलात्} %||4-56-8||

\twolineshloka
{रामस्य च पितुर्मित्रं जटायुर्नाम गृध्रराट्}
{ददर्श सीतां वैदेहीं ह्रियमाणां विहायसा} %||4-56-9||

\twolineshloka
{रावणं विरथं कृत्वा स्थापयित्वा च मैथिलीम्}
{परिश्रान्तश्च वृद्धश्च रावणेन हतो रणे} %||4-56-10||

\twolineshloka
{एवं गृध्रो हतस्तेन रावणेन बहीयसा}
{संस्कृतश्चापि रामेण गतश्च गतिमुत्तमाम्} %||4-56-11||

\twolineshloka
{ततो मम पितृव्येण सुग्रीवेण महात्मना}
{चकार राघवः सख्यं सोऽवधीत्पितरं मम} %||4-56-12||

\twolineshloka
{माम पित्रा विरुद्धो हि सुग्रीवः सचिवैः सह}
{निहत्य वालिनं रामस्ततस्तमभिषेचयत्} %||4-56-13||

\twolineshloka
{स राज्ये स्थापितस्तेन सुग्रीवो वानरेश्वरः}
{राजा वानरमुख्यानां येन प्रस्थापिता वयम्} %||4-56-14||

\twolineshloka
{एवं रामप्रयुक्तास्तु मार्गमाणास्ततस्ततः}
{वैदेहीं नाधिगच्छामो रात्रौ सूर्यप्रभामिव} %||4-56-15||

\twolineshloka
{ते वयं दण्दकारण्यं विचित्य सुसमाहिताः}
{अज्ञानात्तु प्रविष्टाः स्म धरण्या विवृतं बिलम्} %||4-56-16||

\twolineshloka
{मयस्य माया विहितं तद्बिलं च विचिन्वताम्}
{व्यतीतस्तत्र नो मासो यो राज्ञा सामयः कृतः} %||4-56-17||

\twolineshloka
{ते वयं कपिराजस्य सर्वे वचनकारिणः}
{कृतां संस्थामतिक्रान्ता भयात्प्रायमुपास्महे} %||4-56-18||

\twolineshloka
{क्रुद्धे तस्मिंस्तु काकुत्स्थे सुग्रीवे च सलक्ष्मणे}
{गतानामपि सर्वेषां तत्र नो नास्ति जीवितम्} %||4-57-19||


{॥इत्यार्षे श्रीमद्रामायणे वाल्मीकीये आदिकाव्ये किष्किन्धाकाण्डे षट्पञ्चाशत्तमः सर्गः॥}

\sect{सप्तपञ्चाशत्तमः सर्गः}

\twolineshloka
{इत्युक्तः करुणं वाक्यं वानरैस्त्यक्तजीवितैः}
{सबाष्पो वानरान्गृध्रः प्रत्युवाच महास्वनः} %||4-57-1||

\twolineshloka
{यवीयान्मम स भ्राता जटायुर्नाम वानराः}
{यमाख्यात हतं युद्धे रावणेन बलीयसा} %||4-57-2||

\twolineshloka
{वृद्धभावादपक्षत्वाच्छृण्वंस्तदपि मर्षये}
{न हि मे शक्तिरद्यास्ति भ्रातुर्वैरविमोक्षणे} %||4-57-3||

\twolineshloka
{पुरा वृत्रवधे वृत्ते स चाहं च जयैषिणौ}
{आदित्यमुपयातौ स्वो ज्वलन्तं रश्मिमालिनम्} %||4-57-4||

\twolineshloka
{आवृत्याकाशमार्गेण जवेन स्म गतौ भृशम्}
{मध्यं प्राप्ते च सूर्ये च जटायुरवसीदति} %||4-57-5||

\twolineshloka
{तमहं भ्रातरं दृष्ट्वा सूर्यरश्मिभिरर्दितम्}
{पक्षाभ्यं छादयामास स्नेहात्परमविह्वलम्} %||4-57-6||

\twolineshloka
{निर्दग्धपक्षः पतितो विन्ध्येऽहं वानरोत्तमाः}
{अहमस्मिन्वसन्भ्रातुः प्रवृत्तिं नोपलक्षये} %||4-57-7||

\twolineshloka
{जटायुषस्त्वेवमुक्तो भ्रात्रा सम्पातिना तदा}
{युवराजो महाप्राज्ञः प्रत्युवाचाङ्गदस्तदा} %||4-57-8||

\twolineshloka
{जटायुषो यदि भ्राता श्रुतं ते गदितं मया}
{आख्याहि यदि जानासि निलयं तस्य रक्षसः} %||4-57-9||

\twolineshloka
{अदीर्घदर्शिनं तं वा रावणं राक्षसाधिपम्}
{अन्तिके यदि वा दूरे यदि जानासि शंस नः} %||4-57-10||

\twolineshloka
{ततोऽब्रवीन्महातेजा ज्येष्ठो भ्राता जटायुषः}
{आत्मानुरूपं वचनं वानरान्सम्प्रहर्षयन्} %||4-57-11||

\twolineshloka
{निर्दग्धपक्षो गृध्रोऽहं गतवीर्यः प्लवङ्गमाः}
{वाङ्मात्रेण तु रामस्य करिष्ये साह्यमुत्तमम्} %||4-57-12||

\twolineshloka
{जानामि वारुणाल्लोकान्विष्णोस्त्रैविक्रमानपि}
{देवासुरविमर्दांश्च अमृतस्य च मन्थनम्} %||4-57-13||

\twolineshloka
{रामस्य यदिदं कार्यं कर्तव्यं प्रथमं मया}
{जरया च हृतं तेजः प्राणाश्च शिथिला मम} %||4-57-14||

\twolineshloka
{तरुणी रूपसम्पन्ना सर्वाभरणभूषिता}
{ह्रियमाणा मया दृष्टा रावणेन दुरात्मना} %||4-57-15||

\twolineshloka
{क्रोशन्ती राम रामेति लक्ष्मणेति च भामिनी}
{भूषणान्यपविध्यन्ती गात्राणि च विधुन्वती} %||4-57-16||

\twolineshloka
{सूर्यप्रभेव शैलाग्रे तस्याः कौशेयमुत्तमम्}
{असिते राक्षसे भाति यथा वा तडिदम्बुदे} %||4-57-17||

\twolineshloka
{तां तु सीतामहं मन्ये रामस्य परिकीर्तनात्}
{श्रूयतां मे कथयतो निलयं तस्य रक्षसः} %||4-57-18||

\twolineshloka
{पुत्रो विश्रवसः साक्षाद्भ्राता वैश्रवणस्य च}
{अध्यास्ते नगरीं लङ्कां रावणो नाम राकसः} %||4-57-19||

\twolineshloka
{इतो द्वीपे समुद्रस्य सम्पूर्णे शतयोजने}
{तस्मिँल्लङ्का पुरी रम्या निर्मिता विश्वकर्मणा} %||4-57-20||

\twolineshloka
{तस्यां वसति वैदेही दीना कौशेयवासिनी}
{रावणान्तःपुरे रुद्धा राक्षसीभिः सुरक्षिता} %||4-57-21||

\twolineshloka
{जनकस्यात्मजां राज्ञस्तस्यां द्रक्ष्यथ मैथिलीम्}
{लङ्कायामथ गुप्तायां सागरेण समन्ततः} %||4-57-22||

\twolineshloka
{सम्प्राप्य सागरस्यान्तं सम्पूर्णं शतयोजनम्}
{आसाद्य दक्षिणं कूलं ततो द्रक्ष्यथ रावणम्} %||4-57-23||

\twolineshloka
{तत्रैव त्वरिताः क्षिप्रं विक्रमध्वं प्लवङ्गमाः}
{ज्ञानेन खलु पश्यामि दृष्ट्वा प्रत्यागमिष्यथ} %||4-57-24||

\twolineshloka
{आद्यः पन्थाः कुलिङ्गानां ये चान्ये धान्यजीविनः}
{द्वितीयो बलिभोजानां ये च वृक्षफलाशिनः} %||4-57-25||

\twolineshloka
{भासास्तृतीयं गच्छन्ति क्रौञ्चाश्च कुररैः सह}
{श्येनाश्चतुर्थं गच्छन्ति गृध्रा गच्छन्ति पञ्चमम्} %||4-57-26||

\threelineshloka
{बलवीर्योपपन्नानां रूपयौवनशालिनाम्}
{षष्ठस्तु पन्था हंसानां वैनतेयगतिः परा}
{वैनतेयाच्च नो जन्म सर्वेषां वानरर्षभाः} %||4-57-27||

\twolineshloka
{गर्हितं तु कृतं कर्म येन स्म पिशिताशनाः}
{इहस्थोऽहं प्रपश्यामि रावणं जानकीं तथा} %||4-57-28||

\threelineshloka
{अस्माकमपि सौवर्णं दिव्यं चक्षुर्बलं तथा}
{तस्मादाहारवीर्येण निसर्गेण च वानराः}
{आयोजनशतात्साग्राद्वयं पश्याम नित्यशः} %||4-57-29||

\twolineshloka
{अस्माकं विहिता वृत्तिर्निसार्गेण च दूरतः}
{विहिता पादमूले तु वृत्तिश्चरणयोधिनाम्} %||4-57-30||

\twolineshloka
{उपायो दृश्यतां कश्चिल्लङ्घने लवणाम्भसः}
{अभिगम्य तु वैदेहीं समृद्धार्था गमिष्यथ} %||4-57-31||

\twolineshloka
{समुद्रं नेतुमिच्छामि भवद्भिर्वरुणालयम्}
{प्रदास्याम्युदकं भ्रातुः स्वर्गतस्य महात्मनः} %||4-57-32||

\twolineshloka
{ततो नीत्वा तु तं देशं तीरे नदनदीपतेः}
{निर्दग्धपक्षं सम्पातिं वानराः सुमहौजसः} %||4-57-33||

\twolineshloka
{पुनः प्रत्यानयित्वा वै तं देशं पतगेश्वरम्}
{बभूवुर्वानरा हृष्टाः प्रवृत्तिमुपलभ्य ते} %||4-58-34||


{॥इत्यार्षे श्रीमद्रामायणे वाल्मीकीये आदिकाव्ये किष्किन्धाकाण्डे सप्तपञ्चाशत्तमः सर्गः॥}

\sect{अष्टपञ्चाशत्तमः सर्गः}

\twolineshloka
{ततस्तदमृतास्वादं गृध्रराजेन भाषितम्}
{निशम्य वदतो हृष्टास्ते वचः प्लवगर्षभाः} %||4-58-1||

\twolineshloka
{जाम्बवान्वै हरिश्रेष्ठः सह सर्वैः प्लवङ्गमैः}
{भूतलात्सहसोत्थाय गृध्रराजानमब्रवीत्} %||4-58-2||

\twolineshloka
{क्व सीता केन वा दृष्टा को वा हरति मैथिलीम्}
{तदाख्यातु भवान्सर्वं गतिर्भव वनौकसाम्} %||4-58-3||

\twolineshloka
{को दाशरथिबाणानां वज्रवेगनिपातिनाम्}
{स्वयं लक्ष्मणमुक्तानां न चिन्तयति विक्रमम्} %||4-58-4||

\twolineshloka
{स हरीन्प्रीतिसंयुक्तान्सीता श्रुतिसमाहितान्}
{पुनराश्वासयन्प्रीत इदं वचनमब्रवीत्} %||4-58-5||

\twolineshloka
{श्रूयतामिह वैदेह्या यथा मे हरणं श्रुतम्}
{येन चापि ममाख्यातं यत्र चायतलोचना} %||4-58-6||

\twolineshloka
{अहमस्मिन्गिरौ दुर्गे बहुयोजनमायते}
{चिरान्निपतितो वृद्धः क्षीणप्राणपराक्रमः} %||4-58-7||

\twolineshloka
{तं मामेवङ्गतं पुत्रः सुपार्श्वो नाम नामतः}
{आहारेण यथाकालं बिभर्ति पततां वरः} %||4-58-8||

\twolineshloka
{तीक्ष्णकामास्तु गन्धर्वास्तीक्ष्णकोपा भुजङ्गमाः}
{मृगाणां तु भयं तीक्ष्णं ततस्तीक्ष्णक्षुधा वयम्} %||4-58-9||

\twolineshloka
{स कदाचित्क्षुधार्तस्य मम चाहारकाङ्क्षिणः}
{गतसूर्योऽहनि प्राप्तो मम पुत्रो ह्यनामिषः} %||4-58-10||

\twolineshloka
{स मया वृद्धभावाच्च कोपाच्च परिभर्त्सितः}
{क्षुत्पिपासा परीतेन कुमारः पततां वरः} %||4-58-11||

\twolineshloka
{स ममाहारसंरोधात्पीडितः प्रीतिवर्धनः}
{अनुमान्य यथातत्त्वमिदं वचनमब्रवीत्} %||4-58-12||

\twolineshloka
{अहं तात यथाकालमामिषार्थी खमाप्लुतः}
{महेन्द्रस्य गिरेर्द्वारमावृत्य च समास्थितः} %||4-58-13||

\twolineshloka
{तत्र सत्त्वसहस्राणां सागरान्तरचारिणाम्}
{पन्थानमेकोऽध्यवसं संनिरोद्धुमवाङ्मुखः} %||4-58-14||

\twolineshloka
{तत्र कश्चिन्मया दृष्टः सूर्योदयसमप्रभाम्}
{स्त्रियमादाय गच्छन्वै भिन्नाञ्जनचयोपमः} %||4-58-15||

\twolineshloka
{सोऽहमभ्यवहारार्थी तौ दृष्ट्वा कृतनिश्चयः}
{तेन साम्ना विनीतेन पन्थानमभियाचितः} %||4-58-16||

\twolineshloka
{न हि सामोपपन्नानां प्रहर्ता विद्यते क्वचित्}
{नीचेष्वपि जनः कश्चित्किमङ्ग बत मद्विधः} %||4-58-17||

\twolineshloka
{स यातस्तेजसा व्योम सङ्क्षिपन्निव वेगतः}
{अथाहं खे चरैर्भूतैरभिगम्य सभाजितः} %||4-58-18||

\twolineshloka
{दिष्ट्या जीवसि तातेति अब्रुवन्मां महर्षयः}
{कथञ्चित्सकलत्रोऽसौ गतस्ते स्वस्त्यसंशयम्} %||4-58-19||

\twolineshloka
{एवमुक्तस्ततोऽहं तैः सिद्धैः परमशोभनैः}
{स च मे रावणो राजा रक्षसां प्रतिवेदितः} %||4-58-20||

\twolineshloka
{हरन्दाशरथेर्भार्यां रामस्य जनकात्मजाम्}
{भ्रष्टाभरणकौशेयां शोकवेगपराजिताम्} %||4-58-21||

\twolineshloka
{रामलक्ष्मणयोर्नाम क्रोशन्तीं मुक्तमूर्धजाम्}
{एष कालात्ययस्तावदिति वाक्यविदां वरः} %||4-58-22||

\twolineshloka
{एतमर्थं समग्रं मे सुपार्श्वः प्रत्यवेदयत्}
{तच्छ्रुत्वापि हि मे बुद्धिर्नासीत्काचित्पराक्रमे} %||4-58-23||

\twolineshloka
{अपक्षो हि कथं पक्षी कर्म किञ्चिदुपक्रमेत्}
{यत्तु शक्यं मया कर्तुं वाग्बुद्धिगुणवर्तिना} %||4-58-24||

\threelineshloka
{श्रूयतां तत्प्रवक्ष्यामि भवतां पौरुषाश्रयम्}
{वाङ्मतिभ्यां हि सार्वेषां करिष्यामि प्रियं हि वः}
{यद्धि दाशरथेः कार्यं मम तन्नात्र संशयः} %||4-58-25||

\twolineshloka
{ते भवन्तो मतिश्रेष्ठा बलवन्तो मनस्विनः}
{सहिताः कपिराजेन देवैरपि दुरासदाः} %||4-58-26||

\twolineshloka
{रामलक्ष्मणबाणाश्च निशिताः कङ्कपत्रिणः}
{त्रयाणामपि लोकानां पर्याप्तास्त्राणनिग्रहे} %||4-58-27||

\twolineshloka
{कामं खलु दशग्रीवस्तेजोबलसमन्वितः}
{भवतां तु समर्थानां न किञ्चिदपि दुष्करम्} %||4-58-28||

\twolineshloka
{तदलं कालसङ्गेन क्रियतां बुद्धिनिश्चयः}
{न हि कर्मसु सज्जन्ते बुद्धिमन्तो भवद्विधाः} %||4-59-29||


{॥इत्यार्षे श्रीमद्रामायणे वाल्मीकीये आदिकाव्ये किष्किन्धाकाण्डे अष्टपञ्चाशत्तमः सर्गः॥}

\sect{एकोनषष्टितमः सर्गः}

\twolineshloka
{ततः कृतोदकं स्नातं तं गृध्रं हरियूथपाः}
{उपविष्टा गिरौ दुर्गे परिवार्य समन्ततः} %||4-59-1||

\twolineshloka
{तमङ्गदमुपासीनं तैः सर्वैर्हरिभिर्वृतम्}
{जनितप्रत्ययो हर्षात्सम्पातिः पुनरब्रवीत्} %||4-59-2||

\twolineshloka
{कृत्वा निःशब्दमेकाग्राः शृण्वन्तु हरयो मम}
{तत्त्वं सङ्कीर्तयिष्यामि यथा जानामि मैथिलीम्} %||4-59-3||

\twolineshloka
{अस्य विन्ध्यस्य शिखरे पतितोऽस्मि पुरा वने}
{सूर्यातपपरीताङ्गो निर्दग्धः सूर्यरश्मिभिः} %||4-59-4||

\twolineshloka
{लब्धसंज्ञस्तु षड्रात्राद्विवशो विह्वलन्निव}
{वीक्षमाणो दिशः सर्वा नाभिजानामि किञ्चन} %||4-59-5||

\twolineshloka
{ततस्तु सागराञ्शैलान्नदीः सर्वाः सरांसि च}
{वनान्यटविदेशांश्च समीक्ष्य मतिरागमत्} %||4-59-6||

\twolineshloka
{हृष्टपक्षिगणाकीर्णः कन्दरान्तरकूटवान्}
{दक्षिणस्योदधेस्तीरे विन्ध्योऽयमिति निश्चितः} %||4-59-7||

\twolineshloka
{आसीच्चात्राश्रमं पुण्यं सुरैरपि सुपूजितम्}
{ऋषिर्निशाकरो नाम यस्मिन्नुग्रतपाभवत्} %||4-59-8||

\twolineshloka
{अष्टौ वर्षसहस्राणि तेनास्मिन्नृषिणा विना}
{वसतो मम धर्मज्ञाः स्वर्गते तु निशाकरे} %||4-59-9||

\twolineshloka
{अवतीर्य च विन्ध्याग्रात्कृच्छ्रेण विषमाच्छनैः}
{तीक्ष्णदर्भां वसुमतीं दुःखेन पुनरागतः} %||4-59-10||

\twolineshloka
{तमृषिं द्रष्टु कामोऽस्मि दुःखेनाभ्यागतो भृशम्}
{जटायुषा मया चैव बहुशोऽभिगतो हि सः} %||4-59-11||

\twolineshloka
{तस्याश्रमपदाभ्याशे ववुर्वाताः सुगन्धिनः}
{वृक्षो नापुष्पितः कश्चिदफलो वा न दृश्यते} %||4-59-12||

\twolineshloka
{उपेत्य चाश्रमं पुण्यं वृक्षमूलमुपाश्रितः}
{द्रष्टुकामः प्रतीक्षे च भगवन्तं निशाकरम्} %||4-59-13||

\twolineshloka
{अथापश्यमदूरस्थमृषिं ज्वलिततेजसम्}
{कृताभिषेकं दुर्धर्षमुपावृत्तमुदङ्मुखम्} %||4-59-14||

\twolineshloka
{तमृक्षाः सृमरा व्याघ्राः सिंहा नागाः सरीसृपाः}
{परिवार्योपगच्छन्ति दातारं प्राणिनो यथा} %||4-59-15||

\twolineshloka
{ततः प्राप्तमृषिं ज्ञात्वा तानि सत्त्वानि वै ययुः}
{प्रविष्टे राजनि यथा सर्वं सामात्यकं बलम्} %||4-59-16||

\twolineshloka
{ऋषिस्तु दृष्ट्वा मां तुष्टः प्रविष्टश्चाश्रमं पुनः}
{मुहूर्तमात्रान्निष्क्रम्य ततः कार्यमपृच्छत} %||4-59-17||

\twolineshloka
{सौम्य वैकल्यतां दृष्ट्वा रोंणां ते नावगम्यते}
{अग्निदग्धाविमौ पक्षौ त्वक्चैव व्रणिता तव} %||4-59-18||

\twolineshloka
{द्वौ गृध्रौ दृष्टपूर्वौ मे मातरिश्वसमौ जवे}
{गृध्राणां चैव राजानौ भ्रातरौ कामरूपिणौ} %||4-59-19||

\twolineshloka
{ज्येष्ठस्त्वं तु च सम्पातिर्जटायुरनुजस्तव}
{मानुषं रूपमास्थाय गृह्णीतां चरणौ मम} %||4-59-20||

\twolineshloka
{किं ते व्याधिसमुत्थानं पक्षयोः पतनं कथम्}
{दण्डो वायं धृतः केन सर्वमाख्याहि पृच्छतः} %||4-60-21||


{॥इत्यार्षे श्रीमद्रामायणे वाल्मीकीये आदिकाव्ये किष्किन्धाकाण्डे एकोनषष्टितमः सर्गः॥}

\sect{षष्टितमः सर्गः}

\twolineshloka
{ततस्तद्दारुणं कर्म दुष्करं साहसात्कृतम्}
{आचचक्षे मुनेः सर्वं सूर्यानुगमनं तथा} %||4-60-1||

\twolineshloka
{भगवन्व्रणयुक्तत्वाल्लज्जया चाकुलेन्द्रियः}
{परिश्रान्तो न शक्नोमि वचनं परिभाषितुम्} %||4-60-2||

\twolineshloka
{अहं चैव जटायुश्च सङ्घर्षाद्दर्पमोहितौ}
{आकाशं पतितौ वीरौ जिघासन्तौ पराक्रमम्} %||4-60-3||

\twolineshloka
{कैलासशिखरे बद्ध्वा मुनीनामग्रतः पणम्}
{रविः स्यादनुयातव्यो यावदस्तं महागिरिम्} %||4-60-4||

\twolineshloka
{अथावां युगपत्प्राप्तावपश्याव महीतले}
{रथचक्रप्रमाणानि नगराणि पृथक्पृथक्} %||4-60-5||

\twolineshloka
{क्वचिद्वादित्रघोषांश्च ब्रह्मघोषांश्च शुश्रुव}
{गायन्तीश्चाङ्गना बह्वीः पश्यावो रक्तवाससः} %||4-60-6||

\twolineshloka
{तूर्णमुत्पत्य चाकाशमादित्यपथमास्थितौ}
{आवामालोकयावस्तद्वनं शाद्वलसंस्थितम्} %||4-60-7||

\twolineshloka
{उपलैरिव सञ्छन्ना दृश्यते भूः शिलोच्चयैः}
{आपगाभिश्च संवीता सूत्रैरिव वसुन्धरा} %||4-60-8||

\twolineshloka
{हिमवांश्चैव विन्ध्यश्च मेरुश्च सुमहान्नगः}
{भूतले सम्प्रकाशन्ते नागा इव जलाशये} %||4-60-9||

\twolineshloka
{तीव्रस्वेदश्च खेदश्च भयं चासीत्तदावयोः}
{समाविशत मोहश्च मोहान्मूर्छा च दारुणा} %||4-60-10||

\twolineshloka
{न दिग्विज्ञायते याम्या नागेन्या न च वारुणी}
{युगान्ते नियतो लोको हतो दग्ध इवाग्निना} %||4-60-11||

\twolineshloka
{यत्नेन महता भूयो रविः समवलोकितः}
{तुल्यः पृथ्वीप्रमाणेन भास्करः प्रतिभाति नौ} %||4-60-12||

\twolineshloka
{जटायुर्मामनापृच्छ्य निपपात महीं ततः}
{तं दृष्ट्वा तूर्णमाकाशादात्मानं मुक्तवानहम्} %||4-60-13||

\twolineshloka
{पक्षिभ्यां च मया गुप्तो जटायुर्न प्रदह्यत}
{प्रमादात्तत्र निर्दग्धः पतन्वायुपथादहम्} %||4-60-14||

\twolineshloka
{आशङ्के तं निपतितं जनस्थाने जटायुषम्}
{अहं तु पतितो विन्ध्ये दग्धपक्षो जडीकृतः} %||4-60-15||

\twolineshloka
{राज्येन हीनो भ्रात्रा च पक्षाभ्यां विक्रमेण च}
{सर्वथा मर्तुमेवेच्छन्पतिष्ये शिखराद्गिरेः} %||4-61-16||


{॥इत्यार्षे श्रीमद्रामायणे वाल्मीकीये आदिकाव्ये किष्किन्धाकाण्डे षष्टितमः सर्गः॥}

\sect{एकषष्टितमः सर्गः}

\twolineshloka
{एवमुक्त्वा मुनिश्रेष्ठमरुदं दुःखितो भृशम्}
{अथ ध्यात्वा मुहूर्तं तु भगवानिदमब्रवीत्} %||4-61-1||

\twolineshloka
{पक्षौ च ते प्रपक्षौ च पुनरन्यौ भविष्यतः}
{चक्षुषी चैव प्राणाश्च विक्रमश्च बलं च ते} %||4-61-2||

\twolineshloka
{पुराणे सुमहत्कार्यं भविष्यं हि मया श्रुतम्}
{दृष्टं मे तपसा चैव श्रुत्वा च विदितं मम} %||4-61-3||

\twolineshloka
{राजा दशरथो नाम कश्चिदिक्ष्वाकुनन्दनः}
{तस्य पुत्रो महातेजा रामो नाम भविष्यति} %||4-61-4||

\twolineshloka
{अरण्यं च सह भ्रात्रा लक्ष्मणेन गमिष्यति}
{तस्मिन्नर्थे नियुक्तः सन्पित्रा सत्यपराक्रमः} %||4-61-5||

\twolineshloka
{नैरृतो रावणो नाम तस्या भार्यां हरिष्यति}
{राक्षसेन्द्रो जनस्थानादवध्यः सुरदानवैः} %||4-61-6||

\twolineshloka
{सा च कामैः प्रलोभ्यन्ती भक्ष्यैर्भोज्यैश्च मैथिली}
{न भोक्ष्यति महाभागा दुःखमग्ना यशस्विनी} %||4-61-7||

\twolineshloka
{परमान्नं तु वैदेह्या ज्ञात्वा दास्यति वासवः}
{यदन्नममृतप्रख्यं सुराणामपि दुर्लभम्} %||4-61-8||

\twolineshloka
{तदन्नं मैथिली प्राप्य विज्ञायेन्द्रादिदं त्विति}
{अग्रमुद्धृत्य रामाय भूतले निर्वपिष्यति} %||4-61-9||

\twolineshloka
{यदि जीवति मे भर्ता लक्ष्मणेन सह प्रभुः}
{देवत्वं गतयोर्वापि तयोरन्नमिदं त्विति} %||4-61-10||

\twolineshloka
{एष्यन्त्यन्वेषकास्तस्या रामदूताः प्लवङ्गमाः}
{आख्येया राममहिषी त्वया तेभ्यो विहङ्गम} %||4-61-11||

\twolineshloka
{सर्वथा तु न गन्तव्यमीदृशः क्व गमिष्यसि}
{देशकालौ प्रतीक्षस्व पक्षौ त्वं प्रतिपत्स्यसे} %||4-61-12||

\twolineshloka
{उत्सहेयमहं कर्तुमद्यैव त्वां सपक्षकम्}
{इहस्थस्त्वं तु लोकानां हितं कार्यं करिष्यसि} %||4-61-13||

\twolineshloka
{त्वयापि खलु तत्कार्यं तयोश्च नृपपुत्रयोः}
{ब्राह्मणानां सुराणां च मुनीनां वासवस्य च} %||4-61-14||

\twolineshloka
{इच्छाम्यहमपि द्रष्टुं भ्रातरु रामलक्ष्मणौ}
{नेच्छे चिरं धारयितुं प्राणांस्त्यक्ष्ये कलेवरम्} %||4-62-15||


{॥इत्यार्षे श्रीमद्रामायणे वाल्मीकीये आदिकाव्ये किष्किन्धाकाण्डे एकषष्टितमः सर्गः॥}

\sect{द्विषष्टितमः सर्गः}

\twolineshloka
{एतैरन्यैश्च बहुभिर्वाक्यैर्वाक्यविशारदः}
{मां प्रशस्याभ्यनुज्ञाप्य प्रविष्टः स स्वमाश्रमम्} %||4-62-1||

\twolineshloka
{कन्दरात्तु विसर्पित्वा पर्वतस्य शनैः शनैः}
{अहं विन्ध्यं समारुह्य भवतः प्रतिपालये} %||4-62-2||

\twolineshloka
{अद्य त्वेतस्य कालस्य साग्रं वर्षशतं गतम्}
{देशकालप्रतीक्षोऽस्मि हृदि कृत्वा मुनेर्वचः} %||4-62-3||

\twolineshloka
{महाप्रस्थानमासाद्य स्वर्गते तु निशाकरे}
{मां निर्दहति सन्तापो वितर्कैर्बहुभिर्वृतम्} %||4-62-4||

\threelineshloka
{उत्थितां मरणे बुद्धिं मुनि वाक्यैर्निवर्तये}
{बुद्धिर्या तेन मे दत्ता प्राणसंरक्षणाय तु}
{सा मेऽपनयते दुःखं दीप्तेवाग्निशिखा तमः} %||4-62-5||

\twolineshloka
{बुध्यता च मया वीर्यं रावणस्य दुरात्मनः}
{पुत्रः सन्तर्जितो वाग्भिर्न त्राता मैथिली कथम्} %||4-62-6||

\twolineshloka
{तस्या विलपितं श्रुत्वा तौ च सीता विनाकृतौ}
{न मे दशरथस्नेहात्पुत्रेणोत्पादितं प्रियम्} %||4-62-7||

\twolineshloka
{तस्य त्वेवं ब्रुवाणस्य सम्पातेर्वानरैः सह}
{उत्पेततुस्तदा पक्षौ समक्षं वनचारिणाम्} %||4-62-8||

\twolineshloka
{स दृष्ट्वा स्वां तनुं पक्षैरुद्गतैररुणच्छदैः}
{प्रहर्षमतुलं लेभे वानरांश्चेदमब्रवीत्} %||4-62-9||

\twolineshloka
{निशाकरस्य महर्षेः प्रभावादमितात्मनः}
{आदित्यरश्मिनिर्दग्धौ पक्षौ मे पुनरुत्थितौ} %||4-62-10||

\twolineshloka
{यौवने वर्तमानस्य ममासीद्यः पराक्रमः}
{तमेवाद्यावगच्छामि बलं पौरुषमेव च} %||4-62-11||

\twolineshloka
{सर्वथा क्रियतां यत्नः सीतामधिगमिष्यथ}
{पक्षलाभो ममायं वः सिद्धिप्रत्यय कारकः} %||4-62-12||

\twolineshloka
{इत्युक्त्वा तान्हरीन्सर्वान्सम्पातिः पततां वरः}
{उत्पपात गिरेः शृङ्गाज्जिज्ञासुः खगमो गतिम्} %||4-62-13||

\twolineshloka
{तस्य तद्वचनं श्रुत्वा प्रीतिसंहृष्टमानसाः}
{बभूवुर्हरिशार्दूला विक्रमाभ्युदयोन्मुखाः} %||4-62-14||

\fourlineindentedshloka
{अथ पवनसमानविक्रमाः}
{ प्लवगवराः प्रतिलब्ध पौरुषाः}
{अभिजिदभिमुखां दिशं ययुर्-}
{ जनकसुता परिमार्गणोन्मुखाः} %||4-63-15||


{॥इत्यार्षे श्रीमद्रामायणे वाल्मीकीये आदिकाव्ये किष्किन्धाकाण्डे द्विषष्टितमः सर्गः॥}

\sect{त्रिषष्टितमः सर्गः}

\twolineshloka
{आख्याता गृध्रराजेन समुत्पत्य प्लवङ्गमाः}
{सङ्गताः प्रीतिसंयुक्ता विनेदुः सिंहविक्रमाः} %||4-63-1||

\twolineshloka
{सम्पातेर्वचनं श्रुत्वा हरयो रावणक्षयम्}
{हृष्टाः सागरमाजग्मुः सीतादर्शनकाङ्क्षिणः} %||4-63-2||

\twolineshloka
{अभिक्रम्य तु तं देशं ददृशुर्भीमविक्रमाः}
{कृत्स्नं लोकस्य महतः प्रतिबिम्बमिव स्थितम्} %||4-63-3||

\twolineshloka
{दक्षिणस्य समुद्रस्य समासाद्योत्तरां दिशम्}
{संनिवेशं ततश्चक्रुः सहिता वानरोत्तमाः} %||4-63-4||

\twolineshloka
{सत्त्वैर्महद्भिर्विकृतैः क्रीडद्भिर्विविधैर्जले}
{व्यात्तास्यैः सुमहाकायैरूर्मिभिश्च समाकुलम्} %||4-63-5||

\twolineshloka
{प्रसुप्तमिव चान्यत्र क्रीडन्तमिव चान्यतः}
{क्वचित्पर्वतमात्रैश्च जलराशिभिरावृतम्} %||4-63-6||

\twolineshloka
{सङ्कुलं दानवेन्द्रैश्च पातालतलवासिभिः}
{रोमहर्षकरं दृष्ट्वा विषेदुः कपिकुञ्जराः} %||4-63-7||

\twolineshloka
{आकाशमिव दुष्पारं सागरं प्रेक्ष्य वानराः}
{विषेदुः सहसा सर्वे कथं कार्यमिति ब्रुवन्} %||4-63-8||

\twolineshloka
{विषण्णां वाहिनीं दृष्ट्वा सागरस्य निरीक्षणात्}
{आश्वासयामास हरीन्भयार्तान्हरिसत्तमः} %||4-63-9||

\twolineshloka
{न निषादेन नः कार्यं विषादो दोषवत्तरः}
{विषादो हन्ति पुरुषं बालं क्रुद्ध इवोरगः} %||4-63-10||

\twolineshloka
{विषादोऽयं प्रसहते विक्रमे पर्युपस्थिते}
{तेजसा तस्य हीनस्य पुरुषार्थो न सिध्यति} %||4-63-11||

\twolineshloka
{तस्यां रात्र्यां व्यतीतायामङ्गदो वानरैः सह}
{हरिवृद्धैः समागम्य पुनर्मन्त्रममन्त्रयत्} %||4-63-12||

\twolineshloka
{सा वानराणां ध्वजिनी परिवार्याङ्गदं बभौ}
{वासवं परिवार्येव मरुतां वाहिनी स्थिता} %||4-63-13||

\twolineshloka
{कोऽन्यस्तां वानरीं सेनां शक्तः स्तम्भयितुं भवेत्}
{अन्यत्र वालितनयादन्यत्र च हनूमतः} %||4-63-14||

\twolineshloka
{ततस्तान्हरिवृद्धांश्च तच्च सैन्यमरिन्दमः}
{अनुमान्याङ्गदः श्रीमान्वाक्यमर्थवदब्रवीत्} %||4-63-15||

\twolineshloka
{क इदानीं महातेजा लङ्घयिष्यति सागरम्}
{कः करिष्यति सुग्रीवं सत्यसन्धमरिन्दमम्} %||4-63-16||

\twolineshloka
{को वीरो योजनशतं लङ्घयेत प्लवङ्गमाः}
{इमांश्च यूथपान्सर्वान्मोचयेत्को महाभयात्} %||4-63-17||

\twolineshloka
{कस्य प्रसादाद्दारांश्च पुत्रांश्चैव गृहाणि च}
{इतो निवृत्ताः पश्येम सिद्धार्थाः सुखिनो वयम्} %||4-63-18||

\twolineshloka
{कस्य प्रसादाद्रामं च लक्ष्मणं च महाबलम्}
{अभिगच्छेम संहृष्टाः सुग्रीवं च महाबलम्} %||4-63-19||

\twolineshloka
{यदि कश्चित्समर्थो वः सागरप्लवने हरिः}
{स ददात्विह नः शीघ्रं पुण्यामभयदक्षिणाम्} %||4-63-20||

\twolineshloka
{अङ्गदस्य वचः श्रुत्वा न कश्चित्किञ्चिदब्रवीत्}
{स्तिमितेवाभवत्सर्वा सा तत्र हरिवाहिनी} %||4-63-21||

\twolineshloka
{पुनरेवाङ्गदः प्राह तान्हरीन्हरिसत्तमः}
{सर्वे बलवतां श्रेष्ठा भवन्तो दृढविक्रमाः} %||4-63-22||

\twolineshloka
{व्यपदेश्य कुले जाताः पूजिताश्चाप्यभीक्ष्णशः}
{न हि वो गमने सङ्गः कदाचिदपि कस्यचित्} %||4-63-23||

{ब्रुवध्वं यस्य या शक्तिर्गमने प्लवगर्षभाः॥\devanumber{24}॥} %||4-64-24|| (Check)

\addtocounter{shlokacount}{1}


{॥इत्यार्षे श्रीमद्रामायणे वाल्मीकीये आदिकाव्ये किष्किन्धाकाण्डे त्रिषष्टितमः सर्गः॥}

\sect{चतुःषष्टितमः सर्गः}

\twolineshloka
{ततोऽङ्गदवचः श्रुत्वा सर्वे ते वानरोत्तमाः}
{स्वं स्वं गतौ समुत्साहमाहुस्तत्र यथाक्रमम्} %||4-64-1||

\twolineshloka
{गजो गवाक्षो गवयः शरभो गन्धमादनः}
{मैन्दश्च द्विविदश्चैव सुषेणो जाम्बवांस्तथा} %||4-64-2||

\twolineshloka
{आबभाषे गजस्तत्र प्लवेयं दशयोजनम्}
{गवाक्षो योजनान्याह गमिष्यामीति विंशतिम्} %||4-64-3||

\twolineshloka
{गवयो वानरस्तत्र वानरांस्तानुवाच ह}
{त्रिंशतं तु गमिष्यामि योजनानां प्लवङ्गमाः} %||4-64-4||

\twolineshloka
{शरभो वानरस्तत्र वानरांस्तानुवाच ह}
{चत्वारिंशद्गमिष्यामि योजनानां न संशयः} %||4-64-5||

\twolineshloka
{वानरांस्तु महातेजा अब्रवीद्गन्धमादनः}
{योजनानां गमिष्यामि पञ्चाशत्तु न संशयः} %||4-64-6||

\twolineshloka
{मैन्दस्तु वानरस्तत्र वानरांस्तानुवाच ह}
{योजनानां परं षष्टिमहं प्लवितुमुत्सहे} %||4-64-7||

\twolineshloka
{ततस्तत्र महातेजा द्विविदः प्रत्यभाषत}
{गमिष्यामि न सन्देहः सप्ततिं योजनान्यहम्} %||4-64-8||

\twolineshloka
{सुषेणस्तु हरिश्रेष्ठः प्रोक्तवान्कपिसत्तमान्}
{अशीतिं योजनानां तु प्लवेयं प्लवगर्षभाः} %||4-64-9||

\twolineshloka
{तेषां कथयतां तत्र सर्वांस्ताननुमान्य च}
{ततो वृद्धतमस्तेषां जाम्बवान्प्रत्यभाषत} %||4-64-10||

\twolineshloka
{पूर्वमस्माकमप्यासीत्कश्चिद्गतिपराक्रमः}
{ते वयं वयसः पारमनुप्राप्ताः स्म साम्प्रतम्} %||4-64-11||

\twolineshloka
{किं तु नैवं गते शक्यमिदं कार्यमुपेक्षितुम्}
{यदर्थं कपिराजश्च रामश्च कृतनिश्चयौ} %||4-64-12||

\twolineshloka
{साम्प्रतं कालभेदेन या गतिस्तां निबोधत}
{नवतिं योजनानां तु गमिष्यामि न संशयः} %||4-64-13||

\twolineshloka
{तांश्च सर्वान्हरिश्रेष्ठाञ्जाम्बवान्पुनरब्रवीत्}
{न खल्वेतावदेवासीद्गमने मे पराक्रमः} %||4-64-14||

\twolineshloka
{मया महाबलैश्चैव यज्ञे विष्णुः सनातनः}
{प्रदक्षिणीकृतः पूर्वं क्रममाणस्त्रिविक्रमः} %||4-64-15||

\twolineshloka
{स इदानीमहं वृद्धः प्लवने मन्दविक्रमः}
{यौवने च तदासीन्मे बलमप्रतिमं परैः} %||4-64-16||

\twolineshloka
{सम्प्रत्येतावतीं शक्तिं गमने तर्कयाम्यहम्}
{नैतावता च संसिद्धिः कार्यस्यास्य भविष्यति} %||4-64-17||

\twolineshloka
{अथोत्तरमुदारार्थमब्रवीदङ्गदस्तदा}
{अनुमान्य महाप्राज्ञो जाम्बवन्तं महाकपिम्} %||4-64-18||

\twolineshloka
{अहमेतद्गमिष्यामि योजनानां शतं महत्}
{निवर्तने तु मे शक्तिः स्यान्न वेति न निश्चितम्} %||4-64-19||

\twolineshloka
{तमुवाच हरिश्रेष्ठो जाम्बवान्वाक्यकोविदः}
{ज्ञायते गमने शक्तिस्तव हर्यृक्षसत्तम} %||4-64-20||

\twolineshloka
{कामं शतसहस्रं वा न ह्येष विधिरुच्यते}
{योजनानां भवाञ्शक्तो गन्तुं प्रतिनिवर्तितुम्} %||4-64-21||

\twolineshloka
{न हि प्रेषयिता तत स्वामी प्रेष्यः कथञ्चन}
{भवतायं जनः सर्वः प्रेष्यः प्लवगसत्तम} %||4-64-22||

\twolineshloka
{भवान्कलत्रमस्माकं स्वामिभावे व्यवस्थितः}
{स्वामी कलत्रं सैन्यस्य गतिरेषा परन्तप} %||4-64-23||

\twolineshloka
{तस्मात्कलत्रवत्तात प्रतिपाल्यः सदा भवान्}
{अपि चैतस्य कार्यस्य भवान्मूलमरिन्दम} %||4-64-24||

\twolineshloka
{मूलमर्थस्य संरक्ष्यमेष कार्यविदां नयः}
{मूले हि सति सिध्यन्ति गुणाः पुष्पफलादयः} %||4-64-25||

\twolineshloka
{तद्भवानस्या कार्यस्य साधने सत्यविक्रमः}
{बुद्धिविक्रमसम्पन्नो हेतुरत्र परन्तपः} %||4-64-26||

\twolineshloka
{गुरुश्च गुरुपुत्रश्च त्वं हि नः कपिसत्तम}
{भवन्तमाश्रित्य वयं समर्था ह्यर्थसाधने} %||4-64-27||

\twolineshloka
{उक्तवाक्यं महाप्राज्ञं जाम्बवन्तं महाकपिः}
{प्रत्युवाचोत्तरं वाक्यं वालिसूनुरथाङ्गदः} %||4-64-28||

\twolineshloka
{यदि नाहं गमिष्यामि नान्यो वानरपुङ्गवः}
{पुनः खल्विदमस्माभिः कार्यं प्रायोपवेशनम्} %||4-64-29||

\twolineshloka
{न ह्यकृत्वा हरिपतेः सन्देशं तस्य धीमतः}
{तत्रापि गत्वा प्राणानां पश्यामि परिरक्षणम्} %||4-64-30||

\twolineshloka
{स हि प्रसादे चात्यर्थं कोपे च हरिरीश्वरः}
{अतीत्य तस्य सन्देशं विनाशो गमने भवेत्} %||4-64-31||

\twolineshloka
{तद्यथा ह्यस्य कार्यस्य न भवत्यन्यथा गतिः}
{तद्भवानेव दृष्टार्थः सञ्चिन्तयितुमर्हति} %||4-64-32||

\twolineshloka
{सोऽङ्गदेन तदा वीरः प्रत्युक्तः प्लवगर्षभः}
{जाम्बवानुत्तरं वाक्यं प्रोवाचेदं ततोऽङ्गदम्} %||4-64-33||

\twolineshloka
{अस्य ते वीर कार्यस्य न किञ्चित्परिहीयते}
{एष सञ्चोदयाम्येनं यः कार्यं साधयिष्यति} %||4-64-34||

\fourlineindentedshloka
{ततः प्रतीतं प्लवतां वरिष्ठम्}
{ एकान्तमाश्रित्य सुखोपविष्टम्}
{सञ्चोदयामास हरिप्रवीरो}
{ हरिप्रवीरं हनुमन्तमेव} %||4-65-35||


{॥इत्यार्षे श्रीमद्रामायणे वाल्मीकीये आदिकाव्ये किष्किन्धाकाण्डे चतुःषष्टितमः सर्गः॥}

\sect{पञ्चषष्टितमः सर्गः}

\twolineshloka
{अनेकशतसाहस्रीं विषण्णां हरिवाहिनीम्}
{जाम्बवान्समुदीक्ष्यैवं हनुमन्तमथाब्रवीत्} %||4-65-1||

\twolineshloka
{वीर वानरलोकस्य सर्वशास्त्रविदां वर}
{तूष्णीमेकान्तमाश्रित्य हनुमन्किं न जल्पसि} %||4-65-2||

\twolineshloka
{हनुमन्हरिराजस्य सुग्रीवस्य समो ह्यसि}
{रामलक्ष्मणयोश्चापि तेजसा च बलेन च} %||4-65-3||

\twolineshloka
{अरिष्टनेमिनः पुत्रौ वैनतेयो महाबलः}
{गरुत्मानिव विख्यात उत्तमः सर्वपक्षिणाम्} %||4-65-4||

\twolineshloka
{बहुशो हि मया दृष्टः सागरे स महाबलः}
{भुजगानुद्धरन्पक्षी महावेगो महायशाः} %||4-65-5||

\twolineshloka
{पक्षयोर्यद्बलं तस्य तावद्भुजबलं तव}
{विक्रमश्चापि वेगश्च न ते तेनापहीयते} %||4-65-6||

\twolineshloka
{बलं बुद्धिश्च तेजश्च सत्त्वं च हरिसत्तम}
{विशिष्टं सर्वभूतेषु किमात्मानं न बुध्यसे} %||4-65-7||

\twolineshloka
{अप्सराप्सरसां श्रेष्ठा विख्याता पुञ्जिकस्थला}
{अज्ञनेति परिख्याता पत्नी केसरिणो हरेः} %||4-65-8||

\twolineshloka
{अभिशापादभूत्तात वानरी कामरूपिणी}
{दुहिता वानरेन्द्रस्य कुञ्जरस्य महात्मनः} %||4-65-9||

\twolineshloka
{कपित्वे चारुसर्वाङ्गी कदाचित्कामरूपिणी}
{मानुषं विग्रहं कृत्वा यौवनोत्तमशालिनी} %||4-65-10||

\twolineshloka
{अचरत्पर्वतस्याग्रे प्रावृडम्बुदसंनिभे}
{विचित्रमाल्याभरणा महार्हक्षौमवासिनी} %||4-65-11||

\twolineshloka
{तस्या वस्त्रं विशालाक्ष्याः पीतं रक्तदशं शुभम्}
{स्थितायाः पर्वतस्याग्रे मारुतोऽपहरच्छनैः} %||4-65-12||

\twolineshloka
{स ददर्श ततस्तस्या वृत्तावूरू सुसंहतौ}
{स्तनौ च पीनौ सहितौ सुजातं चारु चाननम्} %||4-65-13||

\twolineshloka
{तां विशालायतश्रोणीं तनुमध्यां यशस्विनीम्}
{दृष्ट्वैव शुभसर्वाग्नीं पवनः काममोहितः} %||4-65-14||

\twolineshloka
{स तां भुजाभ्यां पीनाभ्यां पर्यष्वजत मारुतः}
{मन्मथाविष्टसर्वाङ्गो गतात्मा तामनिन्दिताम्} %||4-65-15||

\twolineshloka
{सा तु तत्रैव सम्भ्रान्ता सुवृत्ता वाक्यमब्रवीत्}
{एकपत्नीव्रतमिदं को नाशयितुमिच्छति} %||4-65-16||

\twolineshloka
{अञ्जनाया वचः श्रुत्वा मारुतः प्रत्यभाषत}
{न त्वां हिंसामि सुश्रोणि मा भूत्ते सुभगे भयम्} %||4-65-17||

\twolineshloka
{मनसास्मि गतो यत्त्वां परिष्वज्य यशस्विनि}
{वीर्यवान्बुद्धिसम्पन्नः पुत्रस्तव भविष्यति} %||4-65-18||

\twolineshloka
{अभ्युत्थितं ततः सूर्यं बालो दृष्ट्वा महावने}
{फलं चेति जिघृक्षुस्त्वमुत्प्लुत्याभ्यपतो दिवम्} %||4-65-19||

\twolineshloka
{शतानि त्रीणि गत्वाथ योजनानां महाकपे}
{तेजसा तस्य निर्धूतो न विषादं ततो गतः} %||4-65-20||

\twolineshloka
{तावदापततस्तूर्णमन्तरिक्षं महाकपे}
{क्षिप्तमिन्द्रेण ते वज्रं क्रोधाविष्टेन धीमता} %||4-65-21||

\twolineshloka
{ततः शैलाग्रशिखरे वामो हनुरभज्यत}
{ततो हि नामधेयं ते हनुमानिति कीर्त्यते} %||4-65-22||

\twolineshloka
{ततस्त्वां निहतं दृष्ट्वा वायुर्गन्धवहः स्वयम्}
{त्रैलोक्ये भृशसङ्क्रुद्धो न ववौ वै प्रभञ्जनः} %||4-65-23||

\twolineshloka
{सम्भ्रान्ताश्च सुराः सर्वे त्रैलोक्ये क्षुभिते सति}
{प्रसादयन्ति सङ्क्रुद्धं मारुतं भुवनेश्वराः} %||4-65-24||

\twolineshloka
{प्रसादिते च पवने ब्रह्मा तुभ्यं वरं ददौ}
{अशस्त्रवध्यतां तात समरे सत्यविक्रम} %||4-65-25||

\twolineshloka
{वज्रस्य च निपातेन विरुजं त्वां समीक्ष्य च}
{सहस्रनेत्रः प्रीतात्मा ददौ ते वरमुत्तमम्} %||4-65-26||

\twolineshloka
{स्वच्छन्दतश्च मरणं ते भूयादिति वै प्रभो}
{स त्वं केसरिणः पुत्रः क्षेत्रजो भीमविक्रमः} %||4-65-27||

\twolineshloka
{मारुतस्यौरसः पुत्रस्तेजसा चापि तत्समः}
{त्वं हि वायुसुतो वत्स प्लवने चापि तत्समः} %||4-65-28||

\twolineshloka
{वयमद्य गतप्राणा भवानस्मासु साम्प्रतम्}
{दाक्ष्यविक्रमसम्पन्नः पक्षिराज इवापरः} %||4-65-29||

\twolineshloka
{त्रिविक्रमे मया तात सशैलवनकानना}
{त्रिः सप्तकृत्वः पृथिवी परिक्रान्ता प्रदक्षिणम्} %||4-65-30||

\twolineshloka
{तदा चौषधयोऽस्माभिः सञ्चिता देवशासनात्}
{निष्पन्नममृतं याभिस्तदासीन्नो महद्बलम्} %||4-65-31||

\twolineshloka
{स इदानीमहं वृद्धः परिहीनपराक्रमः}
{साम्प्रतं कालमस्माकं भवान्सर्वगुणान्वितः} %||4-65-32||

\twolineshloka
{तद्विजृम्भस्व विक्रान्तः प्लवतामुत्तमो ह्यसि}
{त्वद्वीर्यं द्रष्टुकामेयं सर्वा वानरवाहिनी} %||4-65-33||

\twolineshloka
{उत्तिष्ठ हरिशार्दूल लङ्घयस्व महार्णवम्}
{परा हि सर्वभूतानां हनुमन्या गतिस्तव} %||4-65-34||

\twolineshloka
{विषाण्णा हरयः सर्वे हनुमन्किमुपेक्षसे}
{विक्रमस्व महावेगो विष्णुस्त्रीन्विक्रमानिव} %||4-65-35||

\fourlineindentedshloka
{ततस्तु वै जाम्बवताभिचोदितः}
{ प्रतीतवेगः पवनात्मजः कपिः}
{प्रहर्षयंस्तां हरिवीर वाहिनीम्}
{ चकार रूपं महदात्मनस्तदा} %||4-66-36||


{॥इत्यार्षे श्रीमद्रामायणे वाल्मीकीये आदिकाव्ये किष्किन्धाकाण्डे पञ्चषष्टितमः सर्गः॥}

\sect{षट्षष्टितमः सर्गः}

\twolineshloka
{संस्तूयमानो हनुमान्व्यवर्धत महाबलः}
{समाविध्य च लाङ्गूलं हर्षाच्च बलमेयिवान्} %||4-66-1||

\twolineshloka
{तस्य संस्तूयमानस्य सर्वैर्वानरपुङ्गवैः}
{तेजसापूर्यमाणस्य रूपमासीदनुत्तमम्} %||4-66-2||

\twolineshloka
{यथा विजृम्भते सिंहो विवृद्धो गिरिगह्वरे}
{मारुतस्यौरसः पुत्रस्तथा सम्प्रति जृम्भते} %||4-66-3||

\twolineshloka
{अशोभत मुखं तस्य जृम्भमाणस्य धीमतः}
{अम्बरीषोपमं दीप्तं विधूम इव पावकः} %||4-66-4||

\twolineshloka
{हरीणामुत्थितो मध्यात्सम्प्रहृष्टतनूरुहः}
{अभिवाद्य हरीन्वृद्धान्हनुमानिदमब्रवीत्} %||4-66-5||

\twolineshloka
{अरुजन्पर्वताग्राणि हुताशनसखोऽनिलः}
{बलवानप्रमेयश्च वायुराकाशगोचरः} %||4-66-6||

\twolineshloka
{तस्याहं शीघ्रवेगस्य शीघ्रगस्य महात्मनः}
{मारुतस्यौरसः पुत्रः प्लवने नास्ति मे समः} %||4-66-7||

\twolineshloka
{उत्सहेयं हि विस्तीर्णमालिखन्तमिवाम्बरम्}
{मेरुं गिरिमसङ्गेन परिगन्तुं सहस्रशः} %||4-66-8||

\twolineshloka
{बाहुवेगप्रणुन्नेन सागरेणाहमुत्सहे}
{समाप्लावयितुं लोकं सपर्वतनदीह्रदम्} %||4-66-9||

\twolineshloka
{ममोरुजङ्घावेगेन भविष्यति समुत्थितः}
{सम्मूर्छितमहाग्राहः समुद्रो वरुणालयः} %||4-66-10||

\twolineshloka
{पन्नगाशनमाकाशे पतन्तं पक्षिसेवितम्}
{वैनतेयमहं शक्तः परिगन्तुं सहस्रशः} %||4-66-11||

\twolineshloka
{उदयात्प्रस्थितं वापि ज्वलन्तं रश्मिमालिनम्}
{अनस्तमितमादित्यमभिगन्तुं समुत्सहे} %||4-66-12||

\twolineshloka
{ततो भूमिमसंस्पृश्य पुनरागन्तुमुत्सहे}
{प्रवेगेनैव महता भीमेन प्लवगर्षभाः} %||4-66-13||

\twolineshloka
{उत्सहेयमतिक्रान्तुं सर्वानाकाशगोचरान्}
{सागरं क्षोभयिष्यामि दारयिष्यामि मेदिनीम्} %||4-66-14||

\twolineshloka
{पर्वतान्कम्पयिष्यामि प्लवमानः प्लवङ्गमाः}
{हरिष्ये चोरुवेगेन प्लवमानो महार्णवम्} %||4-66-15||

\threelineshloka
{लतानां वीरुधां पुष्पं पादपानां च सर्वशः}
{अनुयास्यति मामद्य प्लवमानं विहायसा}
{भविष्यति हि मे पन्थाः स्वातेः पन्था इवाम्बरे} %||4-66-16||

\twolineshloka
{चरन्तं घोरमाकाशमुत्पतिष्यन्तमेव च}
{द्रक्ष्यन्ति निपतन्तं च सर्वभूतानि वानराः} %||4-66-17||

\twolineshloka
{महामेरुप्रतीकाशं मां द्रक्ष्यध्वं प्लवङ्गमाः}
{दिवमावृत्य गच्छन्तं ग्रसमानमिवाम्बरम्} %||4-66-18||

\twolineshloka
{विधमिष्यामि जीमूतान्कम्पयिष्यामि पर्वतान्}
{सागरं क्षोभयिष्यामि प्लवमानः समाहितः} %||4-66-19||

\threelineshloka
{वैनतेयस्य वा शक्तिर्मम वा मारुतस्य वा}
{ऋते सुपर्णराजानं मारुतं वा महाबलम्}
{न हि भूतं प्रपश्यामि यो मां प्लुतमनुव्रजेत्} %||4-66-20||

\twolineshloka
{निमेषान्तरमात्रेण निरालम्भनमम्बरम्}
{सहसा निपतिष्यामि घनाद्विद्युदिवोत्थिता} %||4-66-21||

\twolineshloka
{भविष्यति हि मे रूपं प्लवमानस्य सागरम्}
{विष्णोः प्रक्रममाणस्य तदा त्रीन्विक्रमानिव} %||4-66-22||

\twolineshloka
{बुद्ध्या चाहं प्रपश्यामि मनश्चेष्टा च मे तथा}
{अहं द्रक्ष्यामि वैदेहीं प्रमोदध्वं प्लवङ्गमाः} %||4-66-23||

\twolineshloka
{मारुतस्य समो वेगे गरुडस्य समो जवे}
{अयुतं योजनानां तु गमिष्यामीति मे मतिः} %||4-66-24||

\threelineshloka
{वासवस्य सवज्रस्य ब्रह्मणो वा स्वयम्भुवः}
{विक्रम्य सहसा हस्तादमृतं तदिहानये}
{लङ्कां वापि समुत्क्षिप्य गच्छेयमिति मे मतिः} %||4-66-25||

\twolineshloka
{तमेवं वानरश्रेष्ठं गर्जन्तममितौजसम्}
{उवाच परिसंहृष्टो जाम्बवान्हरिसत्तमः} %||4-66-26||

\twolineshloka
{वीर केसरिणः पुत्र वेगवन्मारुतात्मज}
{ज्ञातीनां विपुलं शोकस्त्वया तात प्रणाशितः} %||4-66-27||

\twolineshloka
{तव कल्याणरुचयः कपिमुख्याः समागताः}
{मङ्गलं कार्यसिद्ध्यर्थं करिष्यन्ति समाहिताः} %||4-66-28||

\twolineshloka
{ऋषीणां च प्रसादेन कपिवृद्धमतेन च}
{गुरूणां च प्रसादेन प्लवस्व त्वं महार्णवम्} %||4-66-29||

\twolineshloka
{स्थास्यामश्चैकपादेन यावदागमनं तव}
{त्वद्गतानि च सर्वेषां जीवितानि वनौकसाम्} %||4-66-30||

\twolineshloka
{ततस्तु हरिशार्दूलस्तानुवाच वनौकसः}
{नेयं मम मही वेगं प्लवने धारयिष्यति} %||4-66-31||

\twolineshloka
{एतानि हि नगस्यास्य शिलासङ्कटशालिनः}
{शिखराणि महेन्द्रस्य स्थिराणि च महान्ति च} %||4-66-32||

\twolineshloka
{एतानि मम निष्पेषं पादयोः पततां वराः}
{प्लवतो धारयिष्यन्ति योजनानामितः शतम्} %||4-66-33||

\twolineshloka
{ततस्तु मारुतप्रख्यः स हरिर्मारुतात्मजः}
{आरुरोह नगश्रेष्ठं महेन्द्रमरिमर्दनः} %||4-66-34||

\twolineshloka
{वृतं नानाविधैर्वृक्षैर्मृगसेवितशाद्वलम्}
{लताकुसुमसम्बाधं नित्यपुष्पफलद्रुमम्} %||4-66-35||

\twolineshloka
{सिंहशार्दूलचरितं मत्तमातङ्गसेवितम्}
{मत्तद्विजगणोद्घुष्टं सलिलोत्पीडसङ्कुलम्} %||4-66-36||

\twolineshloka
{महद्भिरुच्छ्रितं शृङ्गैर्महेन्द्रं स महाबलः}
{विचचार हरिश्रेष्ठो महेन्द्रसमविक्रमः} %||4-66-37||

\twolineshloka
{पादाभ्यां पीडितस्तेन महाशैलो महात्मना}
{ररास सिंहाभिहतो महान्मत्त इव द्विपः} %||4-66-38||

\twolineshloka
{मुमोच सलिलोत्पीडान्विप्रकीर्णशिलोच्चयः}
{वित्रस्तमृगमातङ्गः प्रकम्पितमहाद्रुमः} %||4-66-39||

\twolineshloka
{नानागन्धर्वमिथुनैः पानसंसर्गकर्कशैः}
{उत्पतद्भिर्विहङ्गैश्च विद्याधरगणैरपि} %||4-66-40||

\twolineshloka
{त्यज्यमानमहासानुः संनिलीनमहोरगः}
{शैलशृङ्गशिलोद्घातस्तदाभूत्स महागिरिः} %||4-66-41||

\twolineshloka
{निःश्वसद्भिस्तदा तैस्तु भुजगैरर्धनिःसृतैः}
{सपताक इवाभाति स तदा धरणीधरः} %||4-66-42||

\twolineshloka
{ऋषिभिस्त्रास सम्भ्रान्तैस्त्यज्यमानः शिलोच्चयः}
{सीदन्महति कान्तारे सार्थहीन इवाध्वगः} %||4-66-43||

\fourlineindentedshloka
{स वेगवान्वेगसमाहितात्मा}
{ हरिप्रवीरः परवीरहन्ता}
{मनः समाधाय महानुभावो}
{ जगाम लङ्कां मनसा मनस्वी} %||4-67-44||


{॥इत्यार्षे श्रीमद्रामायणे वाल्मीकीये आदिकाव्ये किष्किन्धाकाण्डे षट्षष्टितमः सर्गः॥}

